\documentclass[book]{magnolia}

\magtex{tex_driver={pdftex},
        tex_packages={float,caption,titling,epigraph,minitoc,slashbox,tabularx,cancel,pgfplots,xypic,nicefrac},
        tex_pstricks={pstricks,pst-plot}}
\magfiche{document_nom={Cours de Sup},
          auteur_nom={François Fayard},
          auteur_mail={fayard.prof@gmail.com}}
\magcours{cours_matiere={maths},
          cours_niveau={mpsi},
          cours_chapitre_numero={11},
          cours_chapitre={Dérivation}}
% \magmisenpage{misenpage_presentation={tikzvelvia},
%           misenpage_format={presentation},
%           misenpage_nbcolonnes={1},
%           misenpage_preuve={non},
%           misenpage_sol={non}}
\magmisenpage{}
\maglieudiff{}
\magprocess

\usepackage[scale=3]{ccicons}
\title{{\Huge\bf Cours de Mathématiques}\\\vspace{1cm}
%       \textbf{\Huge 2022--2023}\\
       \textbf{\Huge Maths Sup -- Tome V}\\\vspace{1cm}
      %  \textsc{L. Bouyge-Bono, F. Fayard, V. Lambert}\\\vspace{1cm}
       \textsc{F. Fayard}\\\vspace{1cm}
       \includegraphics[width=8cm]{../../Commun/Images/lazos-bde-2024.png}\\%\vspace{1cm}
       \includegraphics[width=8cm]{../../Commun/Images/lazos.png}}
% \author{{\sc François Fayard}, {\sc Victor Lambert}}
% \date{3 Août 2020 -- 1a20aa}

% \pretitle{%
%   \begin{center}
%   \includegraphics[width=12cm]{/Users/fayard/Desktop/lazos.png}\\[\bigskipamount]
% }
% \posttitle{\end{center}}

\mtcsettitle{minitoc}{}
\mtcsettitle{secttoc}{Table des matières}
\mtcsettitle{parttoc}{Table des matières}

\usetikzlibrary{positioning}
% \lstset{%
%   frame            = tb,    % draw frame at top and bottom of code block
%   tabsize          = 1,     % tab space width
%   numbers          = none,  % display line numbers on the left
%   framesep         = 3pt,   % expand outward
%   framerule        = 0.4pt, % expand outward 
%   commentstyle     = \color{green},      % comment color
%   keywordstyle     = \color{blue},       % keyword color
%   stringstyle      = \color{blue},    % string color
%   backgroundcolor  = \color{colorLazoBlue1Light}, % backgroundcolor color
%   showstringspaces = false,              % do not mark spaces in strings
%   basicstyle       = \ttfamily,
%   breaklines       = true
% }

\input{vc.tex}

\begin{document}

\maketitle

La version de ce document est la \textsc{\GITAbrHash}.\\

Merci à tous les élèves des lycées Janson de Sailly, du Parc et des Lazaristes pour leurs remarques et corrections. Je remercie
particulièrement \nom{Samuel Auroy}, \nom{Antonin Barbier}, \nom{Amin Belfkira}, \nom{Martin Bot}, \nom{Alexandre Brousse}, \nom{Élodie Brun}, \nom{Damien Callendrier}, \nom{Lauren Calvosa}, \nom{Sylvain Crosnier}, \nom{Enguerrand De Jaegere}, \nom{Thibaud De Valicourt}, \nom{Victor Déru}, \nom{Raphaël Des Boscs}, \nom{Grégoire Dhimoïla}, \nom{Léo Duhamel-Callot}, \nom{Sacha Evrard}, \nom{Axel Faou}, \nom{Titouan Francheteau}, \nom{Hélène Ghaleb}, \nom{Cédric Holocher},
\nom{Maxime Joubert}, \nom{Maxime Lombard}, \nom{Mira Maamri}, \nom{Gauthier Malandrin}, \nom{Pierre-Antoine Nguyen},
\nom{Hilaire Oudinot}, \nom{Eliott Pradeleix}, \nom{Yann-Ellie Ravon}, \nom{Sixtine Reynaud}, \nom{Vivien Thienot},
\nom{Carole Vacherand}, \nom{Camille Vialet}, \nom{Paul Vilars} et \nom{Antonin Villepontoux}.\\

Je tiens enfin à remercier mes anciens professeurs et
collègues qui ont eu une influence sur la rédaction de ce document~: \nom{Walter Appel}, \nom{Bruno Arsac}, \nom{Jean-Pierre Barani}, \nom{Vincent Bayle}, \nom{Christophe Bertault}, \nom{Laurence Bouyge}, \nom{Gilles Chaffard}, \nom{Alain Chillès}, \nom{Denis Choimet}, \nom{Vincent Clapiès}, \nom{Gérard Esposito}, \nom{Stéphane Gonnord}, \nom{Victor Lambert}, %
%\nom{Pierre-Louis Lions},
\nom{Frédéric Morlot}, \nom{Franz Ridde}, \nom{Emmanuel Roblet} et \nom{Alain Troesch}.\\
%, \nom{Cédric Villani}.\\
\vfill


\begin{center}
  \ccbysa\\
  \vspace{2ex}
  This work is licensed under a Creative Commons\\
  Attribution-ShareAlike 4.0 International License.\\
  \url{https://creativecommons.org/licenses/by-sa/4.0/legalcode.fr}\\
  \vspace{2ex}
  La dernière version de ce document ainsi que\\
  les sources \LaTeX{} sont disponibles à l'adresse\\
  \url{https://github.com/FayardProf/Maths-MPSI-MP2I}
  \end{center}
  \vspace{2ex}
  \begin{center}
  \textbf{Vous êtes autorisés à~:}
  \end{center}
  \vspace{2ex}
  \begin{itemize}
  \item \textbf{Partager}~: copier, distribuer et communiquer le matériel par tous les moyens et sous tous formats.
  \item \textbf{Adapter}~: remixer, transformer et créer à partir du matériel
  pour toute utilisation, y compris commerciale.
  \end{itemize}
  \vspace{2ex}
  \begin{center}
  \textbf{Selon les conditions suivantes~:}
  \end{center}
  \vspace{2ex}
  \begin{itemize}
    \item \textbf{Attribution}~: Vous devez créditer l'œuvre, intégrer un lien vers la licence et indiquer si des modifications ont été effectuées à l'œuvre. Vous devez indiquer ces informations par tous les moyens raisonnables, sans toutefois suggérer que l'offrant vous soutient ou soutient la façon dont vous avez utilisé son œuvre.
    \item \textbf{Partage dans les mêmes conditions}~: Dans le cas où vous effectuez un remix, que vous transformez, ou créez à partir du matériel composant l'œuvre originale, vous devez diffuser l'œuvre modifiée dans les même conditions, c'est à dire avec la même licence avec laquelle l'œuvre originale a été diffusée.
    \item \textbf{Pas de restrictions complémentaires}~: Vous n'êtes pas autorisé à appliquer des conditions légales ou des mesures techniques qui restreindraient légalement autrui à utiliser l'œuvre dans les conditions décrites par la licence.
    \end{itemize}

\setcounter{chapter}{21}

\tableofcontents

% \part{Sup}

\chapter{Fractions rationnelles}
\setcounter{numeroexercicecours}{1}
\documentclass{magnolia}

\magtex{tex_driver={pdftex},
        tex_packages={xypic}}
\magfiche{document_nom={Cours sur fractions rationnelles},
          auteur_nom={François Fayard},
          auteur_mail={fayard.prof@gmail.com}}
\magcours{cours_matiere={maths},
          cours_niveau={mpsi},
          cours_chapitre_numero={16},
          cours_chapitre={Fractions rationnelles}}
\magmisenpage{misenpage_presentation={tikzvelvia},
              misenpage_format={a4},
              misenpage_nbcolonnes={1},
              misenpage_preuve={non},
              misenpage_sol={non}}
\maglieudiff{lieu_lycee={Aux Lazaristes},
             lieu_classe={MPSI 1},
             lieu_annee={2019--2020}}
\magprocess

\begin{document}

%BEGIN_BOOK
\magtoc

\section{Fraction rationnelle}
\subsection{Représentants d'une fraction rationnelle}
\begin{definition}
Soit $\K$ un corps. L'anneau $\polyK$ étant intègre, on admet qu'il existe
un unique corps noté $\fracK$ et appelé \emph{corps des fractions
rationnelles}, possédant les propriétés suivantes~:
\begin{itemize}
\item $\polyK$ est un sous-anneau du corps $\fracK$.
\item Pour tout élément $F$ de $\fracK$, il existe un couple de polynômes
  $(A,B)$ avec $B\not=0$ tel que
  \[F=\frac{A}{B}\]
\end{itemize}
Les éléments de $\fracK$ sont appelés \emph{fractions rationnelles} à
coefficients sur le corps $\K$.
\end{definition}

\begin{remarqueUnique}
\remarque Plus généralement, si $A$ est un anneau intègre, il existe un unique
  corps $\K$ tel que~:
  \begin{itemize}
  \item $A$ est un sous-anneau de $\K$.
  \item Pour tout $x\in\K$, il existe un couple $(a,b)\in A^2$ tel que $b\neq 0$ et $x=a/b$.
  \end{itemize}
  Ce corps est appelé \emph{corps des fractions} de $A$. Le corps des fractions de $\Z$ est $\Q$, celui
  de $\polyK$ est $\fracK$.
\end{remarqueUnique}

\begin{definition}
Soit $F$ une fraction rationnelle.
\begin{itemize}
\item On dit qu'un couple de polynômes $(A,B)$ avec $B\not=0$ est un
  \emph{représentant} de $F$ lorsque
  \[F=\frac{A}{B}\]
\item On dit que ce représentant est \emph{irréductible} lorsque $A$ et $B$
  sont premiers entre eux et qu'il est \emph{unitaire} lorsque $B$ est unitaire.
\end{itemize}
\end{definition}

\begin{exoUnique}
\exo Mettre sous forme irréductible les fractions rationnelles
  \[\frac{X^2-1}{X^3-1} \et \frac{X^2+X-2}{X^3-5X^2+8X-4}.\]
  \begin{sol}
  On trouve
  \[\frac{X+1}{X^2+X+1} \et \frac{X+2}{\p{X-2}^2}\]
  \end{sol}
\end{exoUnique}

\begin{proposition}
Soit $F$ une fraction rationnelle. Alors $F$ admet un unique représentant
irréductible unitaire $(A,B)$. De plus~:
\begin{itemize}
\item Les représentants de $F$ sont les couples $(QA,QB)$ où $Q$ est
  un polynôme non nul.
\item Les représentants irréductibles de $F$ sont les couples
  $(\lambda A,\lambda B)$ où $\lambda$ est un scalaire non nul.
\end{itemize}
\end{proposition}

\begin{preuve}
On simplifie par leur pgcd et on a unicité grâce au dénominateur unitaire.

$\frac{A}{B}=\frac{C}{D}$ devient $AD=BC$. Or $A\wedge B=1$ donc $A\mid C$, i.e $C=QA$. De même $D=RB$, donc $ABR=ABQ$ et donc $R=Q$ par intégrité de $\polyK$.
\end{preuve}

\subsection{Degré}
\begin{definition}
Soit $F$ une fraction rationnelle. L'entier relatif $\deg A-\deg B$ ne
dépend pas de la représentation de $F$ choisie; on l'appelle \emph{degré} de $F$.
\end{definition}

\begin{preuve}
On en prend $2$ :
$$F=\dfrac{A_1}{B_1}=\dfrac{A_2}{B_2} \text{ avec } B_1\neq 0 \et B_2 \neq 0.$$
Donc $A_1B_2=B_1A_2$ puis on écrit l'égalité des degrés d'où comme $\deg(B_1),\deg(B_2)\geq 0$, $\deg(A_1)-\deg(B_1)=\deg(A_2)-\deg(B_2)$.
\end{preuve}

\begin{exempleUnique}
\exemple Par exemple
  \[\deg\p{\frac{X+1}{X+2}}=0 \et \deg\p{\frac{X^2+1}{X^3+1}}=-1.\]
\end{exempleUnique}

\begin{proposition}
Soit $F$ et $G$ deux fractions rationnelles.
\begin{itemize}
\item Soit $n\in\Z$. Si $\deg F\leq n$ et $\deg G\leq n$, alors
  \[\forall \lambda,\mu\in\K \qsep \deg(\lambda F+\mu G) \leq n.\]
\item $\deg F G = \deg F + \deg G$
\end{itemize}
\end{proposition}

\begin{preuve}
Si $F_1=\dfrac{A_1}{B_1}$ et $F_2=\dfrac{A_2}{B_2}$, $$\lambda F_1+\mu F_2=\frac{\lambda A_1B_2+\mu A_2B_1}{B_1B_2}$$ donc 
\begin{eqnarray*}
\deg(\lambda F_1+\mu F_2)&=&\deg(\lambda A_1B_2+\mu A_2B_1)-\deg(B_1B_2)\\
&\leq& \max{\set{\deg A_1+\deg B_2,\deg A_2+\deg B_1}}-\deg B_1-\deg B_2\\
&=&\max{\set{\deg A_1-\deg B_1,\deg A_2-\deg B_2}}\\
&=&\max{\set{\deg F_1,\deg F_2}}\leq n.
\end{eqnarray*}
\end{preuve}

\subsection{Racines, pôles et substitution}
\begin{definition}
Soit $F$ une fraction rationnelle et $(A,B)$ un représentant irréductible de
$F$. On dit qu'un scalaire $\alpha$ est
\begin{itemize}
\item une \emph{racine} de $F$ lorsque $\alpha$ est racine de $A$. Dans ce cas
  on définit l'\emph{ordre de multiplicité de la racine $\alpha$} comme
  son ordre de multiplicité en tant que racine de $A$. 
\item un \emph{pôle} de $F$ lorsque $\alpha$ est racine de $B$. Dans ce cas
  on définit l'\emph{ordre de multiplicité du pôle $\alpha$} comme
  son ordre de multiplicité en tant que racine de $B$. 
\end{itemize}
\end{definition}

\begin{exoUnique}
\exo Donner les racines et les pôles de la fraction rationnelle
  \[F\defeq\frac{X^2+X-2}{X^3-5X^2+8X-4}.\]
\end{exoUnique}

\begin{sol}
$$\frac{X^2+X-2}{X^3-5X^2+8X-4}=\frac{X+2}{(X-2)^2}$$ donc $-2$ est une racine d'ordre $1$ et $2$ est un pôle d'ordre $2$.
\end{sol}

\begin{definition}
Soit $F$ une fraction rationnelle sur le corps $\K$ et $\alpha$ un
scalaire. Si $\alpha$ n'est pas un pôle de $F$, on définit $F\p{\alpha}$ par
\[F(\alpha)\defeq\frac{A(\alpha)}{B(\alpha)}\]
où $(A,B)$ est un représentant de $F$ tel que $\alpha$ ne soit pas
racine de $B$.
\end{definition}

\begin{proposition}
Soit $\alpha\in\K$, $F_1$ et $F_2$ deux fractions rationnelles
n'admettant pas $\alpha$ pour pôle et $\lambda$, $\mu\in\K$.
Alors $\lambda F_1+\mu F_2$ et $F_1 F_2$ n'admettent pas
$\alpha$ pour pôle et
\begin{eqnarray*}
(\lambda F_1+\mu F_2)(\alpha) &=& \lambda F_1(\alpha)+\mu F_2(\alpha)\\
(F_1 F_2)(\alpha) &=& F_1(\alpha) F_2(\alpha)
\end{eqnarray*}
\end{proposition}

\begin{definition}
Soit $F$ une fraction rationnelle. Si $\mathcal{P}$ est l'ensemble des pôles de $F$,
on définit la fonction rationnelle $\tilde{F}:\K\setminus\mathcal{P}\to\K$ par
\[\forall x\in\K\setminus\mathcal{P} \qsep \tilde{F}(x)\defeq F(x)\] 
\end{definition}

% \begin{exos}
% \exo Soit $F\in\polyC$. Montrer que $\im\tilde{F}$ est soit un singleton,
%   soit $\C$ privé d'un nombre complexe, soit $\C$ tout entier.
% \end{exos}

\subsection{Conjugaison sur $\fracC$}

\begin{definition}
Soit $F\in\fracC$. On définit la fraction rationnelle $\conj{F}$ par
\[\conj{F}\defeq\frac{\conj{A}}{\conj{B}}\]
\end{definition}

\begin{proposition}
Soit $F,G\in\fracC$.
\begin{itemize}
\item Si $\lambda$, $\mu\in\C$, alors
\begin{eqnarray*}
\conj{\lambda F+\mu G}&=&\conj{\lambda}\, \conj{F}+\conj{\rule{0pt}{6.75pt}\mu}\, \conj{G}\\
\conj{F G}&=&\conj{F}\,\conj{G}
\end{eqnarray*}
\item $\deg\conj{F}=\deg F$.
\end{itemize}
\end{proposition}

\begin{proposition}
Soit $F\in\fracC$. Alors
\[\conj{\conj{F}}=F \quad\et\quad \cro{F\in\fracR \quad\ssi\quad \conj{F}=F}.\]
\end{proposition}

\begin{proposition}
Soit $F\in\fracC$ et $\alpha\in\C$. Alors
\begin{itemize}
\item $\alpha$ est racine de $F$ d'ordre $\omega\in\N$ si et seulement si $\conj{\alpha}$ est racine de $\conj{F}$
   d'ordre $\omega$.
\item $\alpha$ est pôle de $F$ d'ordre $\omega\in\N$ si et seulement si $\conj{\alpha}$ est pôle de $\conj{F}$
   d'ordre $\omega$.
\end{itemize}
\end{proposition}

\begin{remarqueUnique}
\remarque En particulier, si $\alpha\in\C$ est une racine de $F\in\fracR$, alors $\conj{\alpha}$ est
  une racine de $F$ et son ordre relativement à $F$ est le même que celui de $\alpha$. De même, si $\alpha$
  est un pôle de $F$, alors $\conj{\alpha}$ est un pôle de $F$ et son ordre relativement à $F$ est le même
  que celui de $\alpha$.
\end{remarqueUnique}

\section{Décomposition en éléments simples}

\subsection{Décomposition en éléments simples sur $\fracK$}

\begin{definition}
Soit $F\in\fracK$. Alors il existe un unique couple
$(E,G)$ constitué d'un polynôme $E$ appelé \emph{partie entière} de $F$, et d'une
fraction rationnelle $G$ de degré strictement négatif tel que
\[F=E+G\]
En pratique $E$ s'obtient comme le quotient de la division euclidienne
de $A$ par $B$ où $(A,B)$ est un représentant de $F$. 
\end{definition}

\begin{exoUnique}
\exo Calculer la partie entière de
  \[F\defeq\frac{X^3+2X^2+1}{X^2+1}.\]
  \begin{sol}
  On a
  \[F=X+2-\frac{X+1}{X^2+1}\]
  \end{sol}
\end{exoUnique}

\begin{proposition}[nom={Décomposition en éléments simples}]
Soit $F\in\fracK$. On écrit $F$ sous forme irréductible et on
décompose son dénominateur en produit de polynômes irréductibles $P_1,\ldots,P_r$ deux
à deux distincts
\[F=\frac{A}{B}=\frac{A}{\lambda \prod_{k=1}^r P_k^{\alpha_k}}.\]
Alors, il existe un unique polynôme $E\in\polyK$ ainsi qu'une unique famille
de polynômes $R_{k,l}\in\polyK$ avec $k\in\intere{1}{r}$ et
$l\in\intere{1}{\alpha_k}$ tels que
\[F=E+\sum_{k=1}^r \sum_{l=1}^{\alpha_k} \frac{R_{k,l}}{P_k^l}\]
et $\deg R_{k,l}<\deg P_k$.
\end{proposition}

\subsection{Cas où le dénominateur est scindé}

\begin{proposition}
Soit $F\in\fracK$.  On écrit $F$ sous forme irréductible et on suppose
que son dénominateur est scindé
\[F=\frac{A}{B}=\frac{A}{\lambda \prod_{k=1}^r \p{X-\alpha_k}^{\omega_k}}.\]
Alors, il existe un unique polynôme $E\in\polyK$ ainsi qu'une unique famille
de scalaires $a_{k,l}\in\K$ avec $k\in\intere{1}{r}$ et $l\in\intere{1}{\omega_k}$
tels que
\[F=E+\sum_{k=1}^r \sum_{l=1}^{\omega_k} \frac{a_{k,l}}{\p{X-\alpha_k}^l}.\]
\end{proposition}

\begin{remarqueUnique}
\remarque En reprenant les notations de la proposition ci-dessus, on dit que
  \[\sum_{l=1}^{\omega_k} \frac{a_{k,l}}{\p{X-\alpha_k}^l}\]
  est la \emph{partie polaire} de $F$ relative au pôle $\alpha_k$.
\end{remarqueUnique}

% On cherche ici à déterminer la partie polaire relative à un pôle simple. On se
% donne donc une fraction rationnelle $F$ admettant $\alpha$ pour pôle simple
% ainsi qu'un représentant irréductible $(A,B)$ de $F$~:
% $$F=\frac{A}{B}$$
% Puisque $\alpha$ est un pôle simple de $F$, il existe un polynôme $C$
% n'admettant pas $\alpha$ pour racine tel que $B=(X-\alpha)C$.
% Il existe donc $a\in\C$ tel que~:
% $$\frac{A}{(X-\alpha)C}=\frac{a}{X-\alpha}+G$$
% où $G$ est une fraction rationnelle n'admettant pas $\alpha$ pour pôle. Donc~:
% $$\frac{A}{C}=a+(X-\alpha)G$$
% Puisque $\alpha$ n'est pôle ni de $A/C$, ni de $G$, on peut substituer $\alpha$
% à $X$ et on obtient~:
% $$a=\frac{A(\alpha)}{C(\alpha)}$$
% Remarquons que si on dérive $B=(X-\alpha)C$ avant de substituer $\alpha$ à $X$,
% on obtient $C(\alpha)=B'(\alpha)$. On obtient donc une nouvelle expression de
% $a$~:
% $$a=\frac{A(\alpha)}{B'(\alpha)}$$
% On a donc démontré~:

% \begin{proposition}
%   Soit $F$ une fraction rationnelle admettant $\alpha$ pour pôle simple~:
%   $$F=\frac{A}{B}=\frac{A}{(X-\alpha)C}$$
%   où $A$ et $C$ sont des polynômes n'admettant pas $\alpha$ pour racine.
%   Alors la partie polaire relative au pôle $\alpha$ est~:
%   $$\frac{a}{X-\alpha}$$
%   où $a$ est donnée par les deux formule équivalentes~:
%   $$a=\frac{A(\alpha)}{C(\alpha)} \qquad a=\frac{A(\alpha)}{B'(\alpha)}$$
% \end{proposition}

% Dans le cas où $F$ n'admet que des pôles simples, on a donc~:


\begin{proposition}
Soit $F\in\fracK$ une fraction rationnelle admettant $\alpha\in\K$ pour pôle simple
\[F=\frac{A}{B}=\frac{A}{(X-\alpha) C} \quad \text{ avec } C(\alpha)\neq 0. \]
Alors, la partie polaire relative au pôle $\alpha$ est
\[\frac{a}{X-\alpha}\]
où $a$ est donné par les deux formules équivalentes
\[a=\frac{A(\alpha)}{C(\alpha)}, \qquad a=\frac{A(\alpha)}{B'(\alpha)}.\]
\end{proposition}

\begin{preuve}
Soit $k\in \intere{1}{n}$. On pose $C=\dsp\prod_{\substack{j=1\\j\neq k}}^n
  \p{X-\alpha_j}$. La multiplication/évaluation donne $a_k=\dfrac{A(\alpha_k)}{C(\alpha_k)}=\dsp\prod_{\substack{j=1\\j\neq k}}^n
  \p{\alpha_k-\alpha_j}$.
  
  De plus, $B=(X-\alpha_k)C$ donc $B'=C+(X-\alpha_k)C'$ d'où $B'(\alpha_k)=C(\alpha_k)$. Cette dernière version sera utile si $B$ est sous forme développée !!!
\end{preuve}


\begin{exos}
\exo Décomposer les fractions rationnelles suivantes en éléments simples
  sur $\C$.
  \[\frac{X+3}{(X-1)(X+2)}, \qquad \frac{1}{X^2+1}, \qquad \frac{1}{X^n-1}.\]
  \begin{sol}
  On a~:
  \begin{itemize}
  \item Premier : rien de spécial.
    \[F=-\frac{1}{3}\cdot\frac{1}{X+2}+\frac{4}{3}\cdot\frac{1}{X-1}\]
  \item Deuxième : utiliser le fait que F est réel. Ou on peut utiliser la parité.
    \[F=-\frac{1}{2}\cdot\frac{i}{X-i}+\frac{1}{2}\cdot\frac{i}{x+i}\]
  \item Troisième : rien de spécial (on utilise la multiplication/évaluation avec la deuxième forme $a_k=\frac{A(\alpha_k)}{B'(\alpha_k)}$).
    \[F=\frac{1}{n}\sum_{k=0}^{n-1} \frac{\omega^k}{X-\omega^k}\]
  \end{itemize}
  \end{sol}
\exo Calculer la dérivée $n$-ième de la fonction définie sur
  $\R\setminus\ens{-2,1}$ par
  \[\forall x\in\R\setminus\ens{-2,1} \qsep f(x)\defeq\frac{x+3}{(x-1)(x+2)}.\]
\end{exos}


% \begin{itemize}
%   \item Soit $F$ la fraction rationnelle~:
%     $$F=\frac{X+3}{(X-1)(X+2)}$$
%     Alors $\deg F=-1$. De plus $F$ est sous forme irréductible. Donc les pôles
%     de $F$ sont $-1$ et $2$. Comme ce sont des pôles simples, on a~:
%     $$F=\frac{a}{X-1}+\frac{b}{X+2}$$
%     Grâce à la première formule, on obtient facilement~:
%     $$a=\frac{4}{3} \quad \text{et} \quad b=-\frac{1}{3}$$
%     Donc~:
%     $$\frac{X+3}{(X-1)(X+2)}=\frac{4}{3}\frac{1}{X-1}-
%       \frac{1}{3}\frac{1}{X+2}$$
%   \item Soit $F$ la fraction rationnelle~:
%     $$F=\frac{1}{X^n-1}$$
%     Alors $\deg F=-n$. De plus $F$ est sous forme irréductible. Donc les pôles
%     de $F$ sont les racines $n$-ièmes de l'unité. Comme il y en a $n$, ce sont
%     des pôles simples. Il existe donc $a_1,\ldots,a_n\in\C$ tels que~:
%     $$F=\sum_{k=0}^{n-1} \frac{a_k}{X-\omega^k}$$
%     où $\omega=\exp\p{i\frac{2\pi}{n}}$. En posant $B=X^n-1$, on a $B'=nX^{n-1}$
%     et la seconde formule nous donne~:
%     $$\forall k\in\intere{0}{n-1} \quad a_k=\frac{1}{n \omega^{k(n-1)}}
%       =\frac{1}{n}\omega^k$$
%     Donc~:
%     $$\frac{1}{X^n-1}=\frac{1}{n}\sum_{k=0}^{n-1} \frac{\omega^k}{X-\omega^k}$$
% \end{itemize}

% \subsection{Pôles multiples}
% \subsubsection{Utilisation des symétries}
% Dès qu'une fraction rationnelle a au moins deux pôles d'ordre supérieur à $2$,
% il est essentiel de réduire le nombre d'inconnues en utilisant au mieux les
% symétries. Pour cela on, cherche des transformations du plan complexe
% conservant les pôles de $F$ ainsi que leur ordre de multiplicité. On observe
% ensuite quels sont les effets de ces transformations sur la fraction $F$.
% Voyons cette méthode de réduction du nombre d'inconnues sur deux exemples~:
% \begin{itemize}
%   \item Soit $F$ la fraction rationnelle~:
%     $$F=\frac{X^2+2}{X^2(X^2+1)}$$
%     Alors $\deg F=-2$ dont la partie entière de $F$ est nulle. De plus $F$ est
%     sous forme irréductible. En effet $X^2+2$ et $X^2(X^2+1)$ n'ont aucune
%     racine commune dans $\C$. $F$ admet donc $0$ comme pôle double et $i,-i$
%     comme pôles simples. Il existe donc $a,b,c,d\in\C$ tels que~:
%     $$F(X)=\frac{a}{X}+\frac{b}{X^2}+\frac{c}{X-i}+\frac{d}{X+i}$$
%     Puisque $\alpha$ est pôle de $F$ si et seulement si $-\alpha$ est pôle de
%     $F$, comparons $F(-X)$ et $F(X)$. On remarque que $F(-X)=F(X)$. Or~:
%     $$F(-X)=-\frac{a}{X}+\frac{b}{X^2}-\frac{c}{X+i}-\frac{d}{X-i}$$
%     Donc par unicité de la décomposition en éléments simples~:
%     $$a=-a \quad \text{et} \quad c=-d$$
%     Donc $a=0$ et~:
%     $$F(X)=\frac{b}{X^2}+\frac{c}{X-i}-\frac{c}{X+i}$$
%   \item Soit $F$ la fraction rationnelle~:
%     $$F=\frac{1}{(X^2+X+1)^2}$$
%     Alors $\deg F=-4$ donc la partie entière de $F$ est nulle. De plus $F$ est
%     sous forme irréductible. $F$ admet donc $j$ et $j^2$ comme pôles doubles.
%     Il existe donc $a,b,c,d\in\C$ tels que~:
%     $$F(X)=\frac{a}{X-j}+\frac{b}{(X-j)^2}+\frac{c}{X-j^2}+\frac{d}{(X-j^2)^2}$$
%     puisque $F$ est à coefficients réels, $\conj{F}=F$. Or~:
%     $$\conj{F}(X)=\frac{\conj{a}}{X-j^2}+\frac{\conj{b}}{(X-j^2)^2}+
%                     \frac{\conj{c}}{X-j}  +\frac{\conj{d}}{(X-j)^2}$$
%     Donc par unicité de la décomposition en éléments simples~:
%     $$\conj{c}=a \quad \text{et} \quad \conj{d}=b$$
%     Donc~:
%     $$F(X)=\frac{a}{X-j}+\frac{b}{(X-j)^2}+\frac{\conj{a}}{X-j^2}
%            +\frac{\conj{d}}{(X-j^2)^2}$$
% \end{itemize}
% D'autres transformations sont aussi utilisées. Par exemple si les pôles de $F$
% sont les racines $n$-ièmes de l'unité, et que ces pôles sont tous de même
% ordre, il peut être utile de substituer $\omega X$ à $X$ où $\omega=\exp
% \p{i\frac{2\pi}{n}}$.

% \subsubsection{Pôles doubles}
% On cherche ici à déterminer la partie polaire relative à un pôle double. On se
% donne donc une fraction rationnelle $F$ admettant $\alpha$ pour pôle double
% ainsi qu'un représentant irréductible $(A,B)$ de $F$~:
% $$F=\frac{A}{B}$$
% Puisque $\alpha$ est un pôle double de $F$, il existe un polynôme $C$
% n'admettant pas $\alpha$ pour racine tel que $B=(X-\alpha)^2 C$.
% Il existe donc $a_1,a_2\in\C$ tel que~:
% $$\frac{A}{(X-\alpha)^2C}=\frac{a_1}{X-\alpha}+\frac{a_2}{(X-\alpha)^2}+G$$
% où $G$ est une fraction rationnelle n'admettant pas $\alpha$ pour pôle. Donc~:
% $$\frac{A}{C}=a_1(X-\alpha)+a_2+(X-\alpha)^2 G$$
% Puisque $\alpha$ n'est pôle ni de $A/C$, ni de $G$, on peut substituer $\alpha$
% à $X$ et on obtient~:
% $$a_2=\frac{A(\alpha)}{C(\alpha)}$$
% Remarquons que si on dérive deux fois $B=(X-\alpha)^2 C$ avant de substituer
% $\alpha$ à $X$, on obtient $C(\alpha)=B''(\alpha)/2$. On obtient donc
% une nouvelle expression de $a_2$~:
% $$a_2=\frac{2A(\alpha)}{B''(\alpha)}$$
% On a donc démontré~:

\begin{proposition}
Soit $F\in\fracK$ une fraction rationnelle admettant $\alpha\in\K$ pour pôle double
\[F=\frac{A}{B}=\frac{A}{(X-\alpha)^2C} \quad \text{ avec } C(\alpha)\neq 0. \]
Alors, la partie polaire relative au pôle $\alpha$ est
\[\frac{a_1}{X-\alpha}+\frac{a_2}{(X-\alpha)^2}\]
où $a_2$ est donné par 
\[a_2=\frac{A(\alpha)}{C(\alpha)}.\]
\end{proposition}

\begin{preuve}
Puisque $\alpha$ est un pôle double de $F$, il existe un polynôme $C$
 n'admettant pas $\alpha$ pour racine tel que $B=(X-\alpha)^2 C$.
 Il existe donc $a_1,a_2\in\C$ tel que~:
 $$\frac{A}{(X-\alpha)^2C}=\frac{a_1}{X-\alpha}+\frac{a_2}{(X-\alpha)^2}+G$$
 où $G$ est une fraction rationnelle n'admettant pas $\alpha$ pour pôle. Donc~:
 $$\frac{A}{C}=a_1(X-\alpha)+a_2+(X-\alpha)^2 G$$
 Puisque $\alpha$ n'est pôle ni de $A/C$, ni de $G$, on peut substituer $\alpha$
 à $X$ et on obtient~:
 $$a_2=\frac{A(\alpha)}{C(\alpha)}$$
 \end{preuve}
 
 \begin{sol}
 Notons qu'en théorie, de $$\frac{A}{C}=a_1(X-\alpha)+a_2+(X-\alpha)^2 G$$ on tire $$a_1=\p{\frac{A}{C}}'(\alpha)$$ ce qui n'est pas forcément pratique à calculer en général, mais cela arrive !
 \end{sol}

\begin{exos}
\exo Décomposer les fractions rationnelles suivantes en éléments simples
  \[\frac{X+1}{X^2\p{X-1}}, \qquad \frac{X^2+X+3}{X\p{X-1}^3}.\]
%     F=\frac{1}{\p{X-1}\p{X^n-1}}\]
  \begin{sol}
  On a~:
  \begin{itemize}
  \item Premier~: Trouver les termes en $X^2$, $X-1$ puis multiplier par $X$ et
    faire tendre vers $+\infty$.
    \[F=-\frac{2}{X}-\frac{1}{X^2}+\frac{2}{X-1}\]
  \item Second~: Effectuer un développement limité.
    \[F=-\frac{3}{X}+\frac{3}{X-1}-\frac{2}{\p{X-1}^2}+\frac{5}{\p{X-1}^3}\]
    
    ou 
    $$F=\frac{X^2+X+3}{X\p{X-1}^3}=\frac{a}{X}+\frac{b}{X-1}+\frac{c}{(X-1)^2}+\frac{d}{(X-1)^3}$$ la multiplication/évaluation donne $a=-3$ et $d=5$. Ensuite, on multiplie par $X$, on substitue $x$ à $X$ et on le fait tendre vers $+\infty$, cela donne $a+b=0$ donc $b=3$. Dernière étape, on réunit les deux premiers éléments simples, on multiplie par $X^2$ on substitue $x$ à $X$ et on fait tendre $x$ vers $+\infty$. Cela donne $c=1+a=-2$. 
  \end{itemize}
  \end{sol}
\exo Calculer la limite, si elle existe, de la suite de terme général
  \[\sum_{k=1}^n \frac{2k+1}{k^2(k+1)^2}.\]
  \begin{sol}
  On peut remarquer que $F(-1-X)=-F(X)$ pour réduire le problème, puis... 
  La décomposition en éléments simples donne
  \[\frac{2X+1}{X^2(X+1)^2}=\frac{1}{X^2}-\frac{1}{(X+1)^2}\]
  La limite cherchée est donc 1.
  \end{sol}
\end{exos}

\begin{proposition}
Soit $P$ un polynôme scindé de $\polyK$. On écrit
  \[P=\lambda\prod_{k=1}^n \p{X-\alpha_k}\]
  où $\lambda\in\Ks$ et $\alpha_1,\ldots,\alpha_n\in\K$ . Alors
  \[\frac{P'}{P}=\sum_{k=1}^n \frac{1}{X-\alpha_k}.\] 
\end{proposition}

\begin{remarqueUnique}
\remarque
Si $P$ un polynôme scindé de $\polyK$. On écrit
  \[P=\lambda\prod_{k=1}^r \p{X-\alpha_k}^{\omega_k}\]
  où $\lambda\in\Ks$, $\alpha_1,\ldots,\alpha_r\in\K$ sont deux à deux distincts et
  $\omega_1,\ldots,\omega_r\in\Ns$. Alors
  \[\frac{P'}{P}=\sum_{k=1}^r \frac{\omega_k}{X-\alpha_k}.\] 
\end{remarqueUnique}


% Pour déterminer $a_1$, on utilisera l'une des méthode suivante~:
% \begin{itemize}
%   \item Dans le cas où $F$ est de degré négatif (cas auquel il est toujours
%     possible de se ramener quitte à remplacer $F$ par sa partie fractionnaire),
%     on peut multiplier la décomposition de $F$ en éléments simples par $X$ et
%     substituer le réel $x$ à $X$ avant de faire tendre $x$ vers $+\infty$.
%     On obtient une relation entre $a_2$ et les coefficients $a_{i,1}$ des autres
%     pôles de $F$.
%   \item Il est toujours possible d'étudier la fraction rationnelle~:
%     $$F-\frac{a_2}{(X-\alpha)^2}=\frac{a_1}{X-\alpha}+G$$
%     pour qui $\alpha$ est un pôle simple.
%   \item Enfin, il est aussi possible d'évaluer $F$ en un point bien choisit qui
%     n'est pas un pôle de $F$.
% \end{itemize}

% \subsubsection{Pôles d'ordre supérieur}
% Dans le cas où $F$ admet un pôle d'ordre $3$, et que les $3$ inconnues
% subsistent après l'étude des symétries de $F$, alors~:
% \begin{itemize}
%   \item Si $\alpha=0$, on effectue une division de $A$ par $C$ selon les
%     puissances croissantes. Cette méthode permet de donner immédiatement la
%     partie polaire relative au pôle $0$.
%   \item D'une manière équivalente, toujours si $\alpha=0$, on peut effectuer
%     un développement asymptotique de $F(x)$ au voisinage de $0$ à la
%     précision $1$. Cette méthode donnera immédiatement la partie polaire
%     relative au pôle $0$.
%   \item Si $\alpha\not=0$ et si les symétries ne permettent pas de se ramener
%     à un pôle nul, on substitue $Y+\alpha$ à $X$ et pour se ramener au cas où
%     $\alpha=0$.   
% \end{itemize}


\subsection{Cas où $\K=\R$ et le dénominateur n'est pas scindé}

\begin{proposition}
Soit $F\in\fracR$. On écrit $F$ sous forme irréductible et on
décompose son dénominateur en produit de polynômes irréductibles
\[F=\frac{A}{\lambda \prod_{k=1}^r \p{X-\alpha_k}^{\omega_k}
    \prod_{l=1}^s \p{X^2+b_lX+c_l}^{\omega_l'}}.\]
Alors, il existe un unique polynôme $E\in\polyR$, ainsi que des uniques familles
$\p{a_{k,i}}$, $\p{b_{l,j}}$ et $\p{c_{l,j}}$ de réels tels que
\[F=E+\sum_{k=1}^r \p{\sum_{i=1}^{\omega_k} \frac{a_{k,i}}{(X-\alpha_k)^i}}
 +\sum_{l=1}^s \p{\sum_{j=1}^{\omega_l'} \frac{b_{l,j}X+c_{l,j}}{\p{X^2+b_lX+c_l}^j}}.\]
\end{proposition}

\begin{remarqueUnique}
\remarque Pour décomposer une fraction rationnelle $F\in\fracR$, il est possible
  d'effectuer sa décomposition dans $\fracC$ puis de regrouper les parties polaires
  ayant des pôles conjugués. Le plus souvent, on préfèrera cependant une décomposition
  directe dans $\fracR$.
\end{remarqueUnique}

\begin{exoUnique}
\exo Décomposer les fractions rationnelles suivantes en éléments simples sur
  $\R$.
  \[\frac{1}{X\p{X^2+1}}, \qquad \frac{1}{X\p{X^2+X+1}}, \qquad
    \frac{X^2+2}{X^2\p{X^2+1}}\]
  \[\frac{1}{X^4+1}, \qquad \frac{X^7+2}{\p{X^2+X+1}^3}, \qquad
    \frac{1}{X^{2n}-1}.\]
  \begin{sol}
  On a~:
  \begin{itemize}
  \item Premier~: On peut commencer par faire la décomposition dans $\fracC$
    puis regrouper les pôles conjugués.
    Mieux : On commence par trouver le coefficient en $X$. Ensuite on
    trouve le coefficient en $X^2+1$ en multipliant et en évaluant en $i$.
    \[F=\frac{1}{X}-\frac{X}{X^2+1}\]
  \item Second~: Même méthode. On évalue en $j$ et on utilise la liberté de la famille $\set{1,j}$, cf. plus tard.
    \[F=\frac{1}{X}-\frac{X+1}{X^2+X+1}\]
  \item Troisième~: On utilise la symétrie pour prouver qu'il n'y a pas de terme
    en $X$. On trouve ensuite le terme en $X^2$. On multiplie par $X^2$ et on fait
    tendre vers $+\infty$ (ou on évalue en $i$ après multiplication par $X^2+1$).
    \[F=\frac{2}{X^2}-\frac{1}{X^2+1}\]
  \item Quatrième~: On remarque que $X^4+1=\p{X^2+1}-2X^2$ pour factoriser le
    dénominateur. On utilise la parité, puis on évalue en $0$. Ensuite, on évalue
    en une racine $\omega$ de $X^2-\sqrt{2}X+1$ après multiplication par $X^2-\sqrt{2}X+1$, on a :
    $$\frac{1}{\omega^2+\sqrt{2}\omega+1}=a\omega+b.$$ Or, $w^2+1=\sqrt{2}\omega$ donc $$\frac{1}{2\sqrt{2}}=a\omega^2+b\omega=a\p{\sqrt{2}\omega-1}+b\omega,$$ finalement $$\underbrace{\p{\sqrt{2}a+b}}_{\in \R}\underbrace{\omega}_{\in \C\setminus\R}=\underbrace{a+\frac{1}{2\sqrt{2}}}_{\in \R}$$ donc les deux réels sont nuls, ce qui permet de conclure. 
    \[F=\frac{1}{4}\cdot\frac{\sqrt{2}X+2}{X^2+\sqrt{2}X+1}+\frac{1}{4}\cdot
      \frac{-\sqrt{2}X+2}{X^2-\sqrt{2}X+1}\]
  \item Cinquième~: Il suffit de faire plusieurs divisons euclidiennes
    successives par $X^2+X+1$. On trouve~:
    \[F=\frac{X+2}{\p{X^2+X+1+}^3}-\frac{4X+2}{\p{X^2+X+1}^2}+
        \frac{3X+5}{X^2+X+1}+\p{X-3}\]
  \item Sixième~: On décompose dans $\C$ puis on regroupe les racines complexes
    conjuguées. On trouve~:
    \[\frac{1}{X^{2n}-1}=\frac{1}{2n}\cro{\frac{1}{X-1}-\frac{1}{X+1}+\sum_{k=1}^{n-1}
      \frac{2\cos\p{\frac{k\pi}{n}}X-2}{X^2-2\cos\p{\frac{k\pi}{n}}X+1}}\]
  \end{itemize}
  \end{sol}
\end{exoUnique}

% \begin{proposition}
% Soit $F\in\fracR$. On écrit
%  Alors $F$ est
% la somme de sa partie entière et de ses parties polaires~:
% \[F=E+\sum_{i=1}^n F_{\alpha_i}\]
% Autrement dit, si $F$ est une fraction rationnelle de partie entière $E$ et
% dont les pôles $\alpha_1,\ldots,\alpha_n$ sont d'ordres respectifs
% $r_1,\ldots,r_n$, il existe une unique famille de nombres complexes
% $\p{a_{i,j}}$ telle que~:
% \[F=E+\sum_{i=1}^n \p{\sum_{j=1}^{r_i} \frac{a_{i,j}}{(X-\alpha_i)^j}}\]
% Cette décomposition s'appelle la décomposition en éléments simples de
% $F$ sur $\C$.
% \end{proposition}

\section{Primitive d'expression rationnelle}
\subsection{Fractions rationnelles}
% \begin{proposition}
% $\quad$
% \begin{itemize}
% \item Soit $a$ un nombre réel, $n\in\Z$ et $I$ un intervalle ne contenant pas $a$.
%   Alors~:
%   \begin{itemize}
%     \item si $n\not=-1$~:
%       $$\prim{(x-a)^n}{x}=\frac{1}{n+1}(x-a)^{n+1}$$
%     \item pour $n=-1$~:
%       $$\priminv{x-a}{x}=\ln\abs{x-a}$$
%   \end{itemize}
% \item Soit $a=\alpha+i\beta$ un nombre complexe n'appartenant pas à $\R$, $n\in\Z$
%   et $I$ un intervalle de $\R$. Alors~:
%   \begin{itemize}
%     \item si $n\not=-1$~:
%       $$\prim{(x-a)^n}{x}=\frac{1}{n+1}(x-a)^{n+1}$$
%     \item pour $n=-1$~:
%       $$\priminv{x-a}{x}=\frac{1}{2}\ln\p{(x-\alpha)^2+\beta^2}
%                          +i\arctan\p{\frac{x-\alpha}{\beta}}$$
%   \end{itemize}
% \end{itemize}
% \end{proposition}

\begin{exoUnique}
\exo Calculer
  \[\priminv{x^2-1}{x}, \qquad \prim{\frac{2x+1}{x^2+x-2}}{x},  \qquad
    \priminv{x^3-1}{x}, \qquad \priminv{(x^2+1)^2}{x}.\]
  \begin{sol}
  On trouve~:
  \begin{itemize}
  \item Premier~:
     \[\frac{1}{2}\ln\abs{\frac{1+x}{1-x}}\]
  \item Réfléchir deux secondes...
  \item Troisième~:
     La décomposition en éléments simples donne
     \[\frac{1}{X^3-1}=\frac{1}{3}\cdot\frac{1}{X-1}-\frac{1}{3}\cdot
       \frac{X+2}{X^2+X+1}\]
     Et le calcul de primitive
     \[\frac{1}{3}\ln\abs{x-1}-\frac{1}{6}\ln\p{x^2+x+1}-\frac{1}{\sqrt{3}}
       \arctan\p{\frac{2x+1}{\sqrt{3}}}\]
       avec le classique :
       \begin{eqnarray*}\priminv{x^2+x+1}{x}&=&\priminv{(x^2+1/2)^2+3/4}{x}\\
       &=&\priminv{3/4\p{4/3(x+1/2)^2+1}}{x}=4/3\priminv{(\frac{2x+1}{\sqrt{3}})^2+1}{x}\\&=&4/3\prim{\frac{\sqrt{3}/2}{1+u^2}}{u}=2/\sqrt{3}\arctan(u).
       \end{eqnarray*}
    \item Quatrième~:
    Technique : On fait une IPP à partir de $\priminv{1+x^2}{x}$ qui va faire apparaitre en utilisant $x^2=x^2+1-1$ la primitive demandée. On obtient finalement :
    $$\priminv{(x^2+1)^2}{x}=\frac{1}{2}\p{\frac{x}{1+x^2}+\arctan x}.$$
  \end{itemize}
  \end{sol}
\end{exoUnique}

% \begin{remarqueUnique}
% \remarque Lorsqu'on maîtrise bien le calcul de primitives des fractions
%   rationnelles $F$ par les méthodes données sur ces exemples, on peut accélérer
%   le temps de calcul en cherchant d'abord l'allure de la primitive ainsi
%   calculée.
%   \begin{itemize}
%   \item La partie entière de $F$ donnera un terme polynomial.
%   \item Les pôles $\alpha$ d'ordre $\omega$ de $F$ feront apparaître des
%     termes en
%     \[a_0\ln\abs{x-\alpha}+\sum_{k=1}^{\omega-1} \frac{a_k}{\p{x-\alpha}^k}\]
%   \item Si le trinôme $aX^2+bX+c$ apparaît avec un exposant 1 dans la
%     factorisation du dénominateur de $F$, on aura l'apparition de termes en
%     \[d\ln(ax^2+bx+c)+e\arctan(fx+g)\]
%     où $f$ et $g$ sont des réels tels que les racines complexes du trinôme
%     $aX^2+bX+c$ sont données par l'équation $fx+g=\pm i$.
%   \item Si le trinôme $aX^2+bX+c$ apparaît avec un exposant 2 dans la
%     factorisation du dénominateur de $F$, on aura l'apparition de termes en
%     \[d\ln(ax^2+bx+c)+e\arctan(fx+g)+\frac{hx+p}{ax^2+bx+c}\]
%     où $f$ et $g$ sont calculés comme plus haut.
%   \end{itemize}
%   Une fois que l'on a la forme de la primitive, on peut exploiter la parité
%   ou l'imparité de $F$ pour réduire le nombre d'inconnues (rappelons qu'une
%   fonction paire admet \textit{une} primitive impaire et qu'une fonction impaire
%   admet \textit{une} primitive paire). Il suffit alors de dériver l'expression
%   puis de regrouper les termes pour obtenir la forme d'une décomposition en
%   éléments simples avant de calculer les coefficients recherchés par les
%   méthodes classiques vues plus haut.
% \end{remarqueUnique}

\subsection{Fractions rationnelles en $\e^x$}
% \begin{proposition}
%   Soit $F$ une fraction rationnelle et $I$ un intervalle tel que son image
%   par l'exponentielle ne contienne pas de pôle de $F$. Alors, après
%   avoir fait le changement de variable $u=e^x$, le calcul de~:
%   $$\prim{F(e^x)}{x}$$
%   se ramène au calcul d'une primitive de la fonction d'expression $F(u)/u$.
% \end{proposition}
Pour ces fonctions, il suffit d'effectuer le changement de variable
$u\defeq\e^x$ pour se ramener au calcul d'une primitive d'une fraction rationnelle.

\begin{exoUnique}
\exo Calculer
  \[\priminv{\e^{2x}+1}{x}.\]
  \begin{sol}
  On trouve
  \[\frac{1}{X\p{X^2+1}}=\frac{1}{X}-\frac{X}{X^2+1}\]
  donc
  \[x-\frac{1}{2}\ln\p{e^{2x}+1}\]
  \end{sol}
\end{exoUnique}

\subsection{Fractions rationnelles en $\cos x$, $\sin x$}

% \begin{proposition}
%   Soit $F$ une fraction rationnelle et $I$ un intervalle tel que son image
%   par $\cos$ ne contienne pas de pôle de $F$. Alors, après
%   avoir fait le changement de variable $u=\cos x$, le calcul de~:
%   $$\prim{F(\cos x)\sin x}{x}$$
%   se ramène au calcul d'une primitive de la fonction d'expression $-F(u)$.
% \end{proposition}

% \begin{proposition}
%   Soit $F$ une fraction rationnelle et $I$ un intervalle tel que son image
%   par $\sin$ ne contienne pas de pôle de $F$. Alors, après
%   avoir fait le changement de variable $u=\sin x$, le calcul de~:
%   $$\prim{F(\sin x)\cos x}{x}$$
%   se ramène au calcul d'une primitive de la fonction d'expression $F(u)$.
% \end{proposition}

% \begin{proposition}
%   Soit $F$ une fraction rationnelle et $I$ un intervalle tel que son image
%   par $\tan$ ne contienne pas de pôle de $F$. Alors, après
%   avoir fait le changement de variable $u=\tan x$, le calcul de~:
%   $$\prim{F(\tan x)(1+\tan^2 x)}{x}$$
%   se ramène au calcul d'une primitive de la fonction d'expression $F(u)$.
% \end{proposition}

\begin{proposition}[nom={Règles de \nom{Bioche}}]
Soit $G$ une fraction rationnelle en $\cos x$, $\sin x$.
\begin{itemize}
\item Si $G\p{-x}=-G(x)$, il existe une fraction rationnelle $F$ telle que
  \[G(x)=F\p{\cos x}\sin x.\]
  On est donc ramené, après le changement de variable $u=\cos x$, à un calcul
  de primitive d'une fraction rationnelle.
\item Si $G\p{\pi-x}=-G(x)$, il existe une fraction rationnelle $F$ telle que
  \[G(x)=F\p{\sin x}\cos x.\]
  On est donc ramené, après le changement de variable $u=\sin x$, à un calcul
  de primitive d'une fraction rationnelle.
\item Si $G\p{\pi+x}=G(x)$, il existe une fraction rationnelle $F$ telle que
  \[G(x)=F\p{\tan x}\p{1+\tan^2 x}.\]
  On est donc ramené, après le changement de variable $u=\tan x$, à un calcul
  de primitive d'une fraction rationnelle.
\item Sinon, on effectue le changement de variable $u=\tan\p{x/2}$. En
  remarquant que
  \[\cos x=\frac{1-u^2}{1+u^2} \et \sin x=\frac{2u}{1+u^2}\]
  on est ramené à un calcul de primitive d'une fraction rationnelle.
\end{itemize}
\end{proposition}

\begin{exoUnique}
\exo Calculer
  \[\priminv{\cos x\cos\p{2x}}{x}, \qquad \priminv{\sin x+\sin\p{2x}}{x},
    \qquad \integinv{0}{2\pi}{2+\cos x}{x}.\]
  \begin{sol}
  On trouve~:
  \begin{itemize}
  \item Premier~: Changement de variable $u=\sin x$.
    Décomposition en éléments simples
    \[\frac{1}{\p{1-2X^2}\p{1-X^2}}=\frac{1}{2}\cro{\frac{1}{X-1}-\frac{1}{X+1}}+
      \frac{\sqrt{2}}{2}\cro{\frac{1}{X+\frac{\sqrt{2}}{2}}-
      \frac{1}{X-\frac{\sqrt{2}}{2}}}\]
    Puis regrouper les pôles opposés. On trouve
    \[\sqrt{2}\argth\p{\sqrt{2}\sin x}-\argth\p{\sin x}\]
    
    Version Victor :
    On se place par exemple sur $\intero{-\pi/4}{\pi/4}$ de façon à ce que le dénominateur ne s'annule pas, puis on effectue le changement de variable $u=\sin x$ ($du=\cos x dx$). On a alors :
$$\priminv{\cos x\cos\p{2x}}{x}=\prim{\frac{\cos x}{(1-\sin^2 x)(1-2\sin^2 x)}}{x}=\prim{\frac{1}{(1-u^2)(1-2u^2)}}{u}.$$ On doit donc effectuer la décomposition en éléments simples
    \[\frac{1}{\p{1-2X^2}\p{1-X^2}}=\frac{1}{2}\cro{\frac{1}{X-1}-\frac{1}{X+1}}+
      \frac{\sqrt{2}}{2}\cro{\frac{1}{X+\frac{\sqrt{2}}{2}}-
      \frac{1}{X-\frac{\sqrt{2}}{2}}}\]
    On trouve donc :
    \[\priminv{\cos x\cos\p{2x}}{x}=\frac{1}{2}\ln\p{\abs{\frac{u-1}{u+1}}}+
      \frac{\sqrt{2}}{2}\ln\p{\abs{\frac{u+\frac{\sqrt{2}}{2}}{u-\frac{\sqrt{2}}{2}}}}=\boxed{\frac{1}{2}\ln\p{\abs{\frac{\sin x-1}{\sin x+1}}}+
      \frac{\sqrt{2}}{2}\ln\p{\abs{\frac{\sin x+\frac{\sqrt{2}}{2}}{\sin x-\frac{\sqrt{2}}{2}}}}} .\]
  \item Second~: Changement de variable $u=\cos x$. Décomposition en éléments
    simples
    \[\frac{1}{\p{X^2-1}\p{1+2X}}=\frac{1}{6}\cdot\frac{1}{X-1}+
      \frac{1}{2}\cdot\frac{1}{X+1}-\frac{4}{3}\cdot\frac{1}{1+2X}\]
    On trouve donc
    \[\frac{1}{6}\ln\p{1-\cos x}+\frac{1}{2}\ln\p{1+\cos x}-\frac{2}{3}
      \ln\abs{1+2\cos x}\]
  \item Troisième~: Les règles de Bioche ne marchent pas. Il faut
    faire le changement de variable~: $u=\tan(t/2)$. On trouve
    \[\frac{2\sqrt{3}}{3}\arctan\p{\frac{\sqrt{3}\tan\p{x/2}}{3}}\]
    et l'intégrale vaut
    \[\frac{2\sqrt{3}}{3}\pi\]
    
    Version Victor : On se place sur $\R$. Les règles de {\sc Bioche} ne marchent pas. Il faut
    faire le changement de variable~: $t=\tan(x/2)$. 
    $dt/dx=1/2(1+\tan^2(x/2))=1/2(1+t^2)$ donc $dx=2dt/(1+t^2)$. On a alors $\displaystyle \cos(x)=\frac{1-t^2}{1+t^2}$ et $\displaystyle dx=\frac{2dt}{1+t^2}$, d'où :
    \begin{eqnarray*}
    \priminv{2+\cos(x)}{x}&=&\prim{\frac{2}{(1+t^2)\p{2+\frac{1-t^2}{1+t^2}}}}{t}=\prim{\frac{2}{3+t^2}}{t}\\
    &=&\frac{2}{3}\prim{\frac{1}{1+\p{\frac{t}{\sqrt{3}}}^2}}{t}=\frac{2}{3}\sqrt{3}\arctan\p{\frac{t}{\sqrt{3}}}\\
    &=&\boxed{\frac{2}{\sqrt{3}}\arctan\p{\frac{\tan(x/2)}{\sqrt{3}}}}.
    \end{eqnarray*} 
Ca c'est pour la primitive, mais le changement de variable et donc le terme intégré n'a pas de sens en $\pi$. On va donc se ramener à un autre intervalle :
$$\integinv{0}{2\pi}{2+\cos x}{x}=\integinv{0}{\pi}{2+\cos x}{x}+\integinv{\pi}{2\pi}{2+\cos x}{x}=\integinv{0}{2\pi}{2+\cos x}{x}+\integinv{-\pi}{0}{2+\cos x}{x}=\integinv{-\pi}{\pi}{2+\cos x}{x}$$
Or $$\integinv{-\pi+\epsilon}{\pi-\epsilon}{2+\cos x}{x}=\frac{2}{\sqrt{3}}\p{\arctan\p{\frac{\tan\p{\frac{\pi}{2}-\frac{\epsilon}{2}}}{\sqrt{3}}}-\arctan\p{\frac{\tan\p{\frac{-\pi}{2}+\frac{\epsilon}{2}}}{\sqrt{3}}}}\tendvers{\epsilon}{0^+}\frac{2}{\sqrt{3}}(\frac{\pi}{2}-(-\frac{\pi}{2}))=\frac{2\pi}{\sqrt{3}}.$$

  \end{itemize}
  \end{sol}
\end{exoUnique}

% Si $G$ est une expression rationnelle en $\sin x$ et $\cos x$, on peut se
% demander si il est possible de transformer l'expression en l'un des deux cas
% précédent. Remarquons que si il existe une fraction rationnelle $F$ telle que~:
% $$\forall x\in I \quad G(x)=F(\sin x)\cos x$$
% Alors l'expression $G(x)$ change de signe lorsqu'on transforme $x$ en $\pi-x$.
% De même si il existe une fraction rationnelle $F$ telle que~:
% $$\forall x\in I \quad G(x)=F(\cos x)\sin x$$
% Alors l'expression $G(x)$ change de signe lorsqu'on transforme $x$ en $-x$.
% Enfin si il existe une fraction rationnelle $F$ telle que~:
% $$\forall x\in I \quad G(x)=F(\tan x)(1+\tan^2 x)$$
% Alors l'expression $G(x)$ est invariante lorsqu'on transforme $x$ en $x+\pi$.

% Il est possible de démontrer que les réciproques sont vraies. Face à une
% expression rationnelle en $\cos x$ et $\sin x$, on utilisera donc les règles
% suivantes dites règles de Bicohe~:
% \begin{itemize}
%   \item On change $x$ en $-x$. Si l'expression change de signe, on cherchera
%     une fraction rationnelle $F$ telle que~:
%     $$\forall x\in I \quad G(x)=F(\cos x)\sin x$$
%     On effectue ensuite le changement de variable $u=\cos x$.
%   \item On change $x$ en $\pi-x$. Si l'expression change de signe, on
%     cherchera une fraction rationnelle $F$ telle que~:
%     $$\forall x\in I \quad G(x)=F(\sin x)\cos x$$
%     On effectue ensuite le changement de variable $u=\sin x$.
%   \item On change $x$ en $x+\pi$. Si l'expression est invariante, on
%     cherchera une fraction rationnelle $F$ telle que~:
%     $$\forall x\in I \quad G(x)=F(\tan x)(1+\tan^2 x)$$
%     On effectue ensuite le changement de variable $u=\tan x$.
% \end{itemize}

\subsection{Fractions rationnelles en $\ch x$, $\sh x$}

Si $f(x)$ est une fraction rationnelle en $\ch x$ et $\sh x$,
alors c'est une fraction rationnelle en $\e^x$. On peut donc calculer une primitive
d'une telle fonction en posant $u\defeq\e^x$.

% \begin{proposition}
% Soit $G$ une fraction rationnelle en $\ch x$ et $\sh x$.
% \begin{itemize}
% \item Si lorsqu'on remplace $\ch$ par $\cos$ et $\sh$ par $\sin$ les règles de
%   \nom{Bioche} préconisent le changement de variable $u=\cos x$, on effectue le
%   changement de variable $u=\ch x$. 
% \item Si lorsqu'on remplace $\ch$ par $\cos$ et $\sh$ par $\sin$ les règles de
%   \nom{Bioche} préconisent le changement de variable $u=\sin x$, on effectue le
%   changement de variable $u=\sh x$. 
% \item Si lorsqu'on remplace $\ch$ par $\cos$ et $\sh$ par $\sin$ les règles de
%   \nom{Bioche} préconisent le changement de variable $u=\tan x$, on effectue le
%   changement de variable $u=\th x$. 
% \end{itemize}
% \end{proposition}

\begin{exoUnique}
\exo Calculer
  \[\prim{\frac{\ch x}{\sh x+\ch x}}{x}, \qquad
    \priminv{\sh^2 x+2}{x}.\]
  \begin{sol}
  On trouve
  \begin{itemize}
  \item Tout transformer avec des exponentielles. On trouve
    \[\prim{\frac{\ch x}{\sh x+\ch x}}{x}==\prim{\frac{\e^x+\e^{-x}}{2\e^x}}{x}=\prim{\frac{1}{2}\p{1+\e^{-2x}}}{x}=\frac{1}{2}x
      -\frac{1}{4}e^{-2x}\]
%   Effectuer le changement de variable $u=\th x$. On a~:
%   \[\frac{1}{\p{1-X}\p{1+X}^2}=\frac{1}{4}\cdot\frac{1}{1+X}+
%     \frac{1}{2}\cdot\frac{1}{\p{1+X}^2}+\frac{1}{4}\cdot\frac{1}{1-X}\]
%   On trouve
%   \[\prim{\frac{\ch x}{\sh x+\ch x}}{x}=\frac{1}{2}x
%     -\frac{1}{2}\cdot\frac{1}{1+\th x}\]
  \item Effectuer le changement de variable $t=\th x$. $dt=(1-\th^2 x)dx=(1-t^2)dx$. $\ch^2-\sh^2=1$ donc $\th^2=\sh^2/\ch^2=\sh^2/(1+\sh^2)$ d'où $\sh^2=\th^2/(1-\th^2)$. Donc 
  $$\priminv{\sh^2 x+2}{x}=\prim{\frac{1}{\frac{t^2}{1-t^2}+2}\frac{1}{1-t^2}}{t}=\priminv{2-t^2}{t}$$ ensuite on fait la DES, etc...
  $$\priminv{\sh^2 x+2}{x}=\frac{1}{2\sqrt{2}}\ln\p{\abs{\frac{t+\sqrt{2}}{t-\sqrt{2}}}}=\frac{1}{2\sqrt{2}}\ln\p{\frac{\sqrt{2}+\th(x)}{\sqrt{2}-\th(x)}}$$ car $\th(x)\in \intero{-1}{1}.$

  \end{itemize}
  \end{sol}
\end{exoUnique}

\subsection{Fractions rationnelles en $x$ et $\sqrt[n]{(ax+b)/(cx+d)}$, en $x$ et $\sqrt{ax^2+bx+c}$}

% \begin{remarques}
% \remarque 
Si $f$ s'écrit
  \[f(x)=F\p{x,\sqrt[n]{\frac{ax+b}{cx+d}}}\]
  où $F$ est une fraction rationnelle à deux variables, le calcul d'une
  primitive de $f$ se fait en effectuant le changement de variable
  \[u\defeq\sqrt[n]{\frac{ax+b}{cx+d}}\]
  On a alors
  \[x=\frac{du^n-b}{a-cu^n} \et \text{d}x=G(u)\,\text{d}u\]
  où $G$ est une fraction rationnelle. On est donc ramené au calcul d'une
  primitive d'une fraction rationnelle.\\

% \remarque 
On souhaite désormais calculer une primitive d'une fonction $f$ de la forme
  \[f(x)=F\p{x,\sqrt{ax^2+bx+c}}\]
  où $F$ est une fraction rationnelle à deux variables.
  \begin{itemize}
  \item Si $aX^2+bX+c$ admet une racine double réelle $\alpha$ (cas où $\Delta=0$), on
    a $\sqrt{ax^2+bx+c}=\sqrt{a}\abs{x-\alpha}$. L'expression est donc
    une fraction rationnelle sur chaque intervalle où $x-\alpha$ est de signe
    constant.
  \item Si $aX^2+bX+c$ n'admet pas de racine réelle (cas où $\Delta<0$), après
    mise sous forme canonique de $aX^2+bX+c$ et changement de variable, on
    est ramené au calcul d'une primitive d'une fraction rationnelle
    en $u,\sqrt{1+u^2}$. Il suffit alors d'effectuer le changement de variable
    $u\defeq\sh v$ pour se ramener au calcul d'une primitive d'une fraction
    rationnelle en $\ch v$, $\sh v$.
  \item Si $aX^2+bX+c$ admet deux racines réelles (cas où $\Delta>0$), après
    mise sous forme canonique de $aX^2+bX+c$ et changement de variable, on
    est ramené au calcul d'une primitive d'une fraction rationnelle
    en $u,\sqrt{1-u^2}$ ou en $u,\sqrt{u^2-1}$. Dans le premier cas, on
    effectue le changement de variable $u\defeq\cos v$ alors que dans le second
    cas on effectue le changement de variable $u\defeq\pm\ch v$.
  \end{itemize}
% \end{remarques}

\begin{exoUnique}
\exo Calculer
  \[\priminv{x\sqrt{x+1}}{x}, \qquad \prim{\arctan\sqrt{1+x^2}}{x}.\]
  \begin{sol}
  \begin{itemize}
  \item On se place sur $I$ un intervalle de $]-1;0[\cup ]0;+\infty[$ et on pose $u=\sqrt{1+x}$. $u^2=x+1$ donc $2udu=dx$. On trouve alors
    \[\priminv{x\sqrt{x+1}}{x}=\ln\p{\frac{\sqrt{1+x}-1}{\sqrt{1+x}+1}}\]
  \item 
\begin{eqnarray*}
\primppdi{\arctan\sqrt{1+x^2}}{1}{x}&\underset{IPP}{=}&x\arctan\sqrt{1+x^2}-\prim{x\cdot 2x\cdot\frac{1}{2}(1+x^2)^{-1/2}\frac{1}{1+\p{\sqrt{1+x^2}}^2}}{x}\\
&=&x\arctan\sqrt{1+x^2}-\prim{\frac{x^2}{\sqrt{1+x^2}\p{2+x^2}}}{x}.
\end{eqnarray*}  
En posant $x=\sh(t)$, alors $dx=\ch(t)dt$, donc :
\begin{eqnarray*}
\prim{\frac{x^2}{\sqrt{1+x^2}\p{2+x^2}}}{x}&=&\prim{\frac{\sh^2(t)\ch(t)}{\ch(t)\p{2+\sh^2(t)}}}{t}\\
&=&\prim{\frac{\sh^2(t)}{2+\sh^2(t)}}{t}\\
&\underset{u=\th(t)}{=}&\prim{\frac{\frac{u^2}{1-u^2}}{2+\frac{u^2}{1-u^2}}\frac{1}{1-u^2}}{u}\\
&=&\prim{\frac{u^2}{(u^2-1)(u^2-2)}}{u}\\
&=&\ldots\\
&=& \frac{1}{2}\ln\abs{\frac{u+1}{u-1}}+\frac{1}{\sqrt{2}}\ln\abs{\frac{u-\sqrt{2}}{u+\sqrt{2}}}\\
&=&\frac{1}{2}\ln\p{\frac{1+\th(t)}{1-\th(t)}}+\frac{1}{\sqrt{2}}\ln\p{\frac{\sqrt{2}-\th(t)}{\sqrt{2}+\th(t)}}.
\end{eqnarray*}
  On trouve enfin grâce à $\th(t)=\dfrac{\sh(t)}{\ch(t)}=\dfrac{\sh(t)}{\sqrt{1+\sh^2(t)}}=\dfrac{x}{\sqrt{1+x^2}}$ :
$$\prim{\arctan\sqrt{1+x^2}}{x}=x\arctan\sqrt{1+x^2}+\frac{1}{2}\ln\p{\frac{\sqrt{1+x^2}-x}{\sqrt{1+x^2}+x}}+\frac{1}{\sqrt{2}}\ln\p{\frac{\sqrt{2(1+x^2)}+x}{\sqrt{2(1+x^2)}-x}}.$$
  
  
    \[x\arctan\sqrt{1+x^2}-\argsh x+
      \sqrt{2}\argth\p{\frac{x}{\sqrt{2\p{1+x^2}}}}\]
  \end{itemize}
  \end{sol}
\end{exoUnique}

% Pour le calcul de primitives de fonctions dont l'expression est rationnelle
% en $\sh x$ et $\ch x$, on considère l'expression obtenue en remplaçant
% respectivement $\sh x$ par $\sin x$, $\ch x$ par $\cos x$ et $\th x$ par
% $\tan x$. On regarde ensuite si une des règles de Bioche s'applique~:
% \begin{itemize}
%   \item Si la règle de Bioche préconise le changement de variable $u=\cos x$
%     sur la nouvelle expression, effectuer le changement de variable $u=\ch x$
%     sur l'expression d'origine.
%   \item Si la règle de Bioche préconise le changement de variable $u=\sin x$
%     sur la nouvelle expression, effectuer le changement de variable $u=\sh x$
%     sur l'expression d'origine.
%   \item Si la règle de Bioche préconise le changement de variable $u=\tan x$
%     sur la nouvelle expression, effectuer le changement de variable $u=\th x$
%     sur l'expression d'origine.
% \end{itemize}
%END_BOOK

\end{document}


\section{Exercices}
\setcounter{numeroexercice}{1}
\documentclass{magnolia}

\magtex{tex_driver={pdftex}}
\magfiche{document_nom={Exercices sur les fractions rationnelles},
          auteur_nom={François Fayard},
          auteur_mail={fayard.prof@gmail.com}}
\magexos{exos_matiere={maths},
         exos_niveau={mpsi},
         exos_chapitre_numero={11},
         exos_theme={Fractions rationnelles}}
\magmisenpage{}
\maglieudiff{}
\magprocess
 
\begin{document}

%BEGIN_BOOK
\magsection{Fraction rationnelle}
\magsubsection{Représentants d'une fraction rationnelle}
\magsubsection{Degré}

\exercice{nom={Primitive}}
Montrer que la fraction rationnelle $F\defeq 1/X$ n'a pas de primitive dans $\fracC$.

\magsubsection{Racines, pôles et substitution}

\exercice{nom={Valeurs prises par une fraction rationnelle}}
% Gourdon Algèbre, page 75
Soit $F\in\fracC$ une fraction rationnelle. On note $\mathcal{D}_F\defeq\C\setminus\ens{a_1,\ldots,a_n}$,
où $a_1,\ldots,a_n\in\C$ sont les pôles de $F$, et on définit la fonction
\[\dspappli{f}{\mathcal{D}_F}{\C}{z}{F(z).}\]
On suppose que $F$ n'est pas constante. Montrer que $f(\mathcal{D}_F)=\C$ ou qu'il existe $\alpha\in\C$
tel que $f(\mathcal{D}_F)=\C\setminus\ens{\alpha}$.
\magsubsection{Conjugaison sur $\fracC$}



\magsection{Décomposition en éléments simples}
\magsubsection{Décomposition en éléments simples sur $\fracK$}
\magsubsection{Cas où le dénominateur est scindé}

\exercice{nom={Décomposition en éléments simples}}
Pour $n\in\Ns$, décomposer en éléments simples sur $\C$ les fractions suivantes
\[\frac{X^{n-1}}{X^n-1}, \qquad
  \frac{1}{(X-1)(X^n-1)}.\]
\begin{sol}
$\quad$
\begin{questions}
\question 
  \[\frac{1}{X^n-1}=\sum_{k=0}^{n-1} \frac{1}{n}\cdot\frac{\omega^k}{X-\omega^k}\]
\question 
  \[\frac{X^{n-1}}{X^n-1}=\frac{1}{n}\sum_{k=0}^{n-1} \frac{1}{X-\omega^k}\]
\question 
  \[\frac{1}{\p{X-1}\p{X^n-1}}=-\frac{n-1}{2n}\cdot\frac{1}{X-1}+
    \frac{1}{n}\cdot\frac{1}{\p{X-1}^2}+\sum_{k=1}^{n-1}
    \frac{1}{n}\cdot\frac{\omega^k}{\omega^k-1}\cdot\frac{1}{X-\omega^k}\]
\end{questions}
\end{sol}

\exercice{nom={Une somme}}
Soit $n\in\Ns$. Simplifier
\[\sum_{k=1}^n \frac{1}{k(k+1)(k+2)}.\]

\exercice{nom={Une somme de série}}
Calculer la limite lorsque $n$ tend vers $+\infty$ de
\[\sum_{k=2}^n \frac{3k^2-1}{(k-1)^2k^2(k+1)^2}.\]
\begin{sol}
On trouve
\[\frac{3X^2-1}{(X-1)^2X^2(X+1)^2}=\frac{1}{2}\cdot\frac{1}{(X-1)^2}
  -\frac{1}{X^2}+\frac{1}{2}\cdot\frac{1}{(X+1)^2}\]
(Astuce de Anis Marrakchi : Pour connaître le coefficient au dessus de
$1/(X+1)$ et $-1/(X-1)$ on regroupe ces éléments simples, on multiplie par
$X^2$ et on fait tendre $x$ vers $+\infty$ // Autre astuce de Anis Marrakchi :
on peut remarquer que $F=P'/P^2$).
On obtient donc une somme télescopique. La somme de la série est donc
$\frac{3}{8}$.
\end{sol}

\exercice{nom={Décomposition en éléments simples}}
On considère la fraction
\[F\defeq\frac{1}{(X^3-1)^3}\]
que l'on souhaite décomposer en éléments simples sur $\fracC$.
\begin{questions}
\question Calculer la partie polaire de $F$ relative au pôle 1.
\question En étudiant les symétries de $F$, en déduire sa décomposition en
  éléments simples.
\end{questions}
\begin{sol}
$\quad$
\begin{questions}
\question On trouve
  \[\frac{5}{27}\cdot\frac{1}{X-1}-\frac{1}{9}\cdot\frac{1}{(X-1)^2}
    +\frac{1}{27}\cdot\frac{1}{(X-1)^3}\]
\question On a enfin~:
  \begin{eqnarray*}
  \frac{1}{(X^3-1)^3}&=&
    \frac{5}{27}\cdot\frac{1}{X-1}-\frac{1}{9}\cdot\frac{1}{(X-1)^2}
    +\frac{1}{27}\cdot\frac{1}{(X-1)^3}+\\
    &&\sum_{\alpha\in\ens{j,j^2}}\cro{
    \frac{5}{27}\cdot\frac{\alpha}{X-\alpha}+\frac{1}{9}\cdot\frac{1+\alpha}{(X-\alpha)^2}
    +\frac{1}{27}\cdot\frac{1}{(X-\alpha)^3}}
  \end{eqnarray*}
\end{questions}
\end{sol}




\exercice{nom={Autour de \nom{Chebyshev}}}
\begin{questions}
\question Montrer que pour tout entier $n\in\N$, il existe un unique polynôme
  $P_n\in\polyR$ tel que
  \[X^n+\frac{1}{X^n}=P_n\p{X+\frac{1}{X}}.\]
\question Soit $n\in\Ns$.
  \begin{questions}
  \question Factoriser $P_n$ dans $\polyC$.
  \question Décomposer $1/P_n$ en éléments simples dans $\fracC$.
  \end{questions}
\end{questions}
\begin{sol}
$\quad$
\begin{questions}
\question $P_{n+1}=XP_n-P_{n-1}$
\question
  \begin{questions}
  \question 
    \[P_n=\prod_{k=0}^{n-1} \p{X-2\cos\p{\frac{\pi\p{2k+1}}{2n}}}\]
  \question
    \[\frac{1}{P_n}=\sum_{k=0}^{n-1} \frac{\p{-1}^k}{n}\sin\p{\frac{\pi\p{2k+1}}{2n}}
      \frac{1}{X-2\cos\p{\frac{\pi\p{2k+1}}{2n}}}\]
  \end{questions}
\end{questions}
\end{sol}

\exercice{nom={Racines multiples de $P'$}}
Soit $P\in\polyR$ un polynôme non constant.
\begin{questions}
\question On suppose que $P$ est scindé sur $\polyR$.
\begin{questions}
\question Calculer la décomposition en éléments simples de $P'/P$.
\question En déduire que
  \[\forall x\in\R\qsep P'(x)^2\geq P(x)P''(x)\]
  avec égalité si et seulement si $x$ est racine multiple de $P$.
\question Montrer que $P'$ est scindé, puis montrer que toute
  racine multiple de $P'$ est racine de $P$.
\end{questions}
\question
\begin{questions}
\question Est-il vrai en général que toute racine multiple de $P'$ est racine de $P$~?
\question Est-il vrai, sous l'hypothèse que $P$ est scindé sur $\R$, que toute racine de $P'$ est racine de $P$~?
\end{questions}
\question On suppose que
  \[\forall x\in\R\qsep P'(x)^2\geq P(x)P''(x).\]
  Montrer que $P$ possède une racine réelle.
\end{questions}


\exercice{nom={Théorème de \nom{Gauss}-\nom{Lucas}}}
\begin{questions}
\question Soit $P\in\fracC$ non nul. Calculer la décomposition en éléments simples de $F\defeq P'/P$.
\question Soit $n\in\Ns$. Déterminer la valeur de
  \[\sum_{z\in\mathbb{U}_n\setminus\ens{1}} \frac{1}{1-z}.\]
\question Soit $P\in\polyC$ un polynôme non constant dont les racines sont $z_1,\ldots,z_r\in\C$.
\begin{questions}
\question Montrer que pour toute racine $\alpha\in\C$ de $P'$, il existe $\lambda_1,\ldots,\lambda_r \geq 0$
  tels que $\lambda_1+\cdots+\lambda_r=1$ et \[\alpha=\lambda_1 z_1+\cdots+\lambda_r z_r.\]
\question Interpréter géométriquement ce résultat.
\end{questions}
\end{questions}

\exercice{nom={Calcul de dérivées $n$-ièmes}}
Calculer la dérivée $n$-ième de
\[x\mapsto\ln(x^2-x+2).\]

\exercice{nom={Majoration d'une dérivée $n$-ième}}
Montrer que
\[\forall n\in\Ns \qsep \forall x\in\R \qsep \abs{\arctan^{(n)} x}\leq(n-1)!.\]
\begin{sol}
On trouve grâce à $\arctan'(x)=\dfrac{1}{1+x^2}=\dfrac{\ii}{2}\p{\dfrac{1}{x+i}-\dfrac{1}{x-i}}$, 
\[\forall n\in\Ns \quad \forall x\in\R \quad \arctan^{(n)} x=
  (-1)^{n+1}\frac{(n-1)!}{2}\cro{\frac{1}{(x+i)^n}-\frac{1}{(x-i)^n}}i\]
  puis on majore en passant aux valeurs absolues.
\end{sol}

\exercice{nom={Dérivée $n$-ième}}
On pose
\[R\defeq\frac{1}{X^2+1}.\]
\begin{questions}
\question Décomposer $R$ en éléments simples sur $\fracC$,
  puis calculer $R^{(n)}$ pour tout $n\in\N$.
\question Montrer que pour tout $n\in\N$
  \[R^{(n)}=\frac{(-1)^n (n+1)!}{(X^2+1)^{n+1}}\prod_{k=1}^n \p{X-\cotan\frac{k\pi}{n+1}}.\]
\end{questions}

\magsubsection{Cas où $\K=\R$ et le dénominateur n'est pas scindé}


\exercice{nom={Décomposition en éléments simples}}
Décomposer en éléments simples sur $\R$ les fractions suivantes
\[\frac{10X^3}{(X^2+1)(X^2-4)}, \qquad \frac{X^4+1}{X^4+X^2+1},
  \qquad \frac{X^3-1}{(X-1)(X-2)(X-3)},\]
\[\frac{X^2}{(X^2+X+1)^2}, \qquad \frac{(X^2+4)^2}{(X^2+1)(X^2-2)^2}.\]
\begin{sol}
$\quad$
\begin{questions}
\question 
  \[\frac{10X^3}{(X^2+1)(X^2-4)}=
    \frac{4}{X-2}+\frac{4}{X+2}+\frac{2X}{X^2+1}\]
\question 
  \[\frac{X^4+1}{X^4+X^2+1}=
    1-\frac{1}{2}\cdot\frac{X}{X^2-X+1}+\frac{1}{2}\cdot\frac{X}{X^2+X+1}\]
\question 
  \[\frac{X^3-1}{(X-1)(X-2)(X-3)}=
    1-\frac{7}{X-2}+\frac{13}{X-3}\]
\question 
  \[\frac{X^2}{(X^2+X+1)^2}=
    \frac{1}{X^2+X+1}-\frac{X+1}{\p{X^2+X+1}^2}\]\
\question 
  \begin{eqnarray*}
  \frac{(X^2+4)^2}{(X^2+1)(X^2-2)^2}&=&
    \frac{1}{X^2+1}+\frac{3\sqrt{2}}{4}\cdot\frac{1}{X+\sqrt{2}}
    +\frac{3}{2}\cdot\frac{1}{\p{X+\sqrt{2}}^2}\\
   & &
    -\frac{3\sqrt{2}}{4}\cdot\frac{1}{X-\sqrt{2}}
    +\frac{3}{2}\cdot\frac{1}{\p{X-\sqrt{2}}^2}
  \end{eqnarray*}
\end{questions}
\end{sol}






% \exercice{nom={Racines de P'}}
% Soit $P\in\polyC$ de degré $n$.
% \begin{questions}
% \question Décomposer $P'/P$ en éléments simples.
% \question En déduire que les racines de $P'$ sont dans l'enveloppe convexe des
%   racines de $P$ c'est-à-dire que toute racine de $P'$ s'écrit comme barycentre
%   à coefficients positifs des racines de $P$.
% \end{questions}









\magsection{Primitive d'expression rationnelle}

\magsubsection{Fractions rationnelles}

\exercice{nom={Primitives de fractions rationnelles}}
% \begin{questions}
% \question
Calculer, sur un domaine $\mathcal{D}$ que l'on précisera, les primitives des
  fonctions dont les expressions sont
  \[\frac{1}{x^4-1}, \qquad \frac{x^2}{x^3-1}, \qquad \frac{x}{x^4+1}.\]
% \question Calculer
%   \[\integ{0}{1}{\frac{x^3+x}{(x^2+x+1)^2}}{x}\]
% \end{questions}
\begin{sol}
$\quad$
\begin{questions}
\question
  \[\priminv{1-x^4}{x}=\frac{1}{2}\arctan x+\frac{1}{2}\argth x\]
  \[\prim{\frac{x}{1+x^4}}{x}=\frac{1}{2}\arctan\p{x^2}\]
  \[\prim{\frac{x^2}{x^3-1}}{x}=\frac{1}{3}\ln\abs{x^3-1}\]
\question Le mieux est de faire cet exercice en cherchant une primitive sous
  une forme donnée.
  \[\prim{\frac{x^3+x}{(x^2+x+1)^2}}{x}=\frac{1}{2} \ln\p{1+x+x^2}
    -\frac{7\sqrt{3}}{9}
    \arctan\p{\frac{2 x+1}{\sqrt{3}}}+\frac{x-1}{3 \p{1+x+x^2}}\]
  On trouve
  \[\frac{1}{54}\p{18+27\ln 3-7\sqrt{3}\pi}\]
\end{questions}
\end{sol}

\magsubsection{Fractions rationnelles en $\e^x$}
\magsubsection{Fractions rationnelles en $\cos x$, $\sin x$}


\exercice{nom={Primitives de fractions rationnelles en $\cos,\sin$}}
Calculer sur un domaine $\mathcal{D}$ que l'on précisera les primitives des
fonctions dont les expressions sont
\[\frac{1}{\cos x\sin^3 x}, \qquad \frac{1}{\sin x+\cos x}.\]
\begin{sol}
$\quad$
\begin{questions}
  \question On doit décomposer en éléments simples~:
  \[\frac{1}{X \left(X^2-1 \right)^2}=
    \frac{1}{4 (X-1)^2}-\frac{1}{2 (X-1)}+\frac{1}{X}-\frac{1}{4 \
    (X+1)^2}-\frac{1}{2 (X+1)}\]
  Puis~:
  \[\priminv{\cos x\sin^3 x}{x}=\ln\abs{\tan x}-\frac{1}{2\sin^2 x}\]
% \question
%   \[\priminv{1+2\cos x}{x}=
%     \frac{2\sqrt{3}}{3}\arctan\p{\frac{\sqrt{3}\tan\p{x/2}}{3}}\]
\question
  \[\priminv{\sin x+\cos x}{x}=\sqrt{2}
    \argth\p{\frac{\tan\p{\frac{x}{2}}-1}{\sqrt{2}}}\]
\end{questions}
\end{sol}

\magsubsection{Fractions rationnelles en $\ch x$, $\sh x$}
\magsubsection{Fractions rationnelles en $x$ et $\sqrt[n]{(ax+b)/(cx+d)}$, en $x$ et $\sqrt{ax^2+bx+c}$}

\exercice{nom={Calcul de primitives}}
Calculer sur un domaine $\mathcal{D}$ que l'on précisera, les primitives des
fonctions dont les expressions sont
\[\frac{1}{2}x\sqrt{\frac{x-1}{x+1}}, \qquad \frac{x}{\sqrt{-x^2+x+2}}.\]

% \exercice{nom={Primitives de polynôme-exponentielles}}
% Calculer les primitives des fonctions dont les expressions sont :
% \[(1+2x+x^2)e^x \qquad (1+2x)e^x \sin x\]
% \begin{sol}
% $\quad$
% \begin{questions}
% \question
%   \[\prim{(1+2x+x^2)e^x}{x}=\p{1+x^2}e^x\]
% \question
%   \[\prim{(1+2x)e^x \sin x}{x}=
%     \p{\frac{1}{2}-x}e^x\cos x-\p{\frac{1}{2}+x}e^x\sin x\]
% \end{questions}
% \end{sol}

% \exercice{nom={Primitives de polynômes en $\cos,\sin$}}
% Calculer les primitives des fonctions dont les expressions sont :
% \[\sin^2 x\cos^3 x \qquad \sin^3 x\cos^4 x \qquad \sin^3 x\cos^5 x\]






\exercice{nom={Une intégrale impropre}}
Calculer la limite lorsque $x$ tend vers $\pi/2$ par la gauche de
\[\integ{0}{x}{\sqrt{\tan t}}{t}.\]
\begin{sol}
On se ramène au calcul d'une primitive de $u^2/(1+u^4)$. Après avoir
factorisé $u^4+1$ en $(u^2+\sqrt{2}u+1)(u^2-\sqrt{2}u+1)$, on cherche
l'allure d'une primitive. On utilise ensuite l'imparité de la primitive
cherchée pour se ramener au calcul de 2 coefficients. Remarquer qu'un des
coefficients est inutile, puis calculer le dernier par substitution. On trouve
\[\integ{0}{\frac{\pi}{2}}{\sqrt{\tan t}}{t}=\frac{\pi}{\sqrt{2}}\]
\end{sol}
%END_BOOK

\end{document}

\chapter{Séries}
\setcounter{numeroexercicecours}{1}
\documentclass{magnolia}

\magtex{tex_driver={pdftex},
        tex_packages={epigraph,xypic}}
\magfiche{document_nom={Cours sur les séries},
          auteur_nom={François Fayard},
          auteur_mail={fayard.prof@gmail.com}}
\magcours{cours_matiere={maths},
          cours_niveau={mpsi},
          cours_chapitre_numero={21},
          cours_chapitre={Séries}}
\magmisenpage{misenpage_presentation={tikzvelvia},
          misenpage_format={a4},
          misenpage_nbcolonnes={1},
          misenpage_preuve={non},
          misenpage_sol={non}}
\magmisenpage{}
\maglieudiff{}
\magprocess

\begin{document}

%BEGIN_BOOK


\magtoc
\vspace{2ex}
Dans ce chapitre, $\K$ désignera l'un des corps $\R$ ou $\C$.

\section{Série}
\subsection{Série}

\begin{definition}
Soit $(u_n)$ une suite d'éléments de $\K$. On appelle \emph{série de terme général $u_n$}
et on note $\sum u_n$ la suite $(S_n)$ définie par
\[\forall n\in\N\qsep S_n\defeq \sum_{k=0}^n u_k.\]
Le terme $S_n$ est appelé \emph{somme partielle d'indice $n$} de la série.
\end{definition}

% \begin{remarqueUnique}
% \remarque Soit $(v_n)$ une suite à valeurs dans $\RouC$. On définit la suite $(u_n)$ par
%   \[\forall n\in\N\qsep u_n\defeq v_{n+1} - v_n.\]
%   Alors
%   \[\forall n\in\N\qsep v_n=v_0 + \sum_{k=0}^{n-1} u_k.\]
%   L'étude de la suite $(v_n)$ se ramène donc à l'étude de la série $\sum u_n$. On dit que
%   la suite $(u_n)$ est la suite \emph{dérivée} de la suite $(v_n)$. Il sera parfois utile
%   d'exprimer une suite à l'aide d'une série afin d'appliquer des techniques propres à ces
%   dernières.
% \end{remarqueUnique}

\begin{definition}
On dit qu'une série $\sum u_n$ \emph{converge} lorsque la suite de ses sommes partielles
converge. Si c'est le cas, sa limite $l\in\K$ est appelée \emph{somme} de la série. On
la note
\[\sum_{n=0}^{+\infty} u_n.\]
Dans le cas contraire, on dit qu'elle \emph{diverge}.
\end{definition}

\begin{remarques}
\remarque Si on change un nombre fini de termes de la suite $(u_n)$, on ne change pas
  la nature de la série $\sum u_n$. Par contre, si elle converge, cela peut changer sa somme.
\remarque Soit $\sum u_n$ une série convergente. Alors, quel que soit $n\in\N$
  \[\sum_{k=0}^{+\infty} u_k=\sum_{k=0}^n u_k + \sum_{k=n+1}^{+\infty} u_k.\]
\end{remarques}

% \begin{remarques}
% \remarque Lorsque la série $\sum u_n$ converge vers $l$, on définit la suite $(r_n)$ par
%   \[\forall n\in\N\qsep r_n\defeq l-s_n.\]
%   Quel que soit $n\in\N$, on a alors
%   \[\sum_{k=n+1}^m u_k \tendvers{m}{+\infty} r_n\]
%   et on note
%   \[r_n=\sum_{k=n+1}^{+\infty} u_k.\]
% \end{remarques}

\begin{exoUnique}
\exo Montrer que la série
  \[\sum_{n\geq 1} \frac{1}{n(n+1)}\]
  converge et calculer sa somme.
\end{exoUnique}


\begin{definition}
Soit $\sum u_n$ une série convergente. On définit la suite $(R_n)$ par
\begin{eqnarray*}
\forall n\in\N\qsep R_n &\defeq& \sum_{k=0}^{+\infty} u_k - \sum_{k=0}^n u_k\\
&=& \sum_{k=n+1}^{+\infty} u_k.
\end{eqnarray*}
Le terme $R_n$ est appelé \emph{reste d'indice $n$} de la série.
\end{definition}

\begin{remarqueUnique}
\remarque La suite $(R_n)$ des restes converge vers 0.
\end{remarqueUnique}

\begin{proposition}
Soit $\sum u_n$ une série. Si elle est convergente, alors
\[u_n\tendvers{n}{+\infty}0.\]
Par contraposée, si la suite $(u_n)$ ne converge pas vers 0, la série $\sum u_n$ est
divergente. On dit qu'elle diverge \emph{grossièrement}.
\end{proposition}

\begin{remarqueUnique}
\remarque Il est possible qu'une série diverge sans diverger grossièrement.
  Par exemple, si $(u_n)$ est la suite définie par
  \[\forall n\in\N\qsep u_n\defeq\sqrt{n+1}-\sqrt{n}\]
  alors, la série associée diverge alors que la suite $(u_n)$ converge vers 0.
\end{remarqueUnique}

\begin{proposition}
La suite $(u_n)$ et la série $\sum (u_{n+1}-u_n)$ sont de même nature.
\end{proposition}

\begin{remarqueUnique}
\remarque Si $(u_n)$ est une suite, la suite de terme général $u_{n+1}-u_n$ est appelée
  \emph{dérivée} de la suite $(u_n)$. Par sommation télescopique
  \[\forall n\in\N\qsep u_n=u_0 + \sum_{k=0}^{n-1} (u_{k+1}-u_k).\]
  L'étude de la suite $(u_n)$ se ramène donc à l'étude de la série
  $\sum (u_{n+1}-u_n)$.
\end{remarqueUnique}

\begin{proposition}
Soit $\sum u_n$ et $\sum v_n$ deux séries convergentes et
$\lambda,\mu\in\K$. Alors, la série \mbox{$\sum (\lambda u_n+\mu v_n)$} est convergente et
\[\sum_{n=0}^{+\infty} \p{\lambda u_n+\mu v_n}=\lambda\sum_{n=0}^{+\infty} u_n
  +\mu\sum_{n=0}^{+\infty} v_n.\]
\end{proposition}

\begin{remarques}
\remarque Attention, il est possible que la série $\sum (\lambda u_n+\mu v_n)$ soit convergente
  sans que les séries $\sum u_n$ et $\sum v_n$ le soient. Avant d'écrire
  \[\sum_{n=0}^{+\infty} \p{\lambda u_n+\mu v_n}=\lambda\sum_{n=0}^{+\infty} u_n
    +\mu\sum_{n=0}^{+\infty} v_n\]
  il faudra donc toujours vérifier que les séries $\sum u_n$ et $\sum v_n$ soient convergentes.
  Un tel oubli pourrait conduire à écrire des horreurs comme
  \begin{eqnarray*}
  0 = \sum_{n=0}^{+\infty} 0 &=& \sum_{n=0}^{+\infty} (1+(-1))\\
  &=& \sum_{n=0}^{+\infty} 1 + \sum_{n=0}^{+\infty} (-1).
  \end{eqnarray*}
  La dernière expression n'a en effet aucun sens car les deux séries sont grossièrement divergentes.
\remarque Si $\sum u_n$ est convergente et $\sum v_n$ est divergente, alors
  $\sum (u_n+v_n)$ est divergente.
\end{remarques}


% \begin{exoUnique}
% \exo Déterminer les $z\in\C$ pour lesquels la série
%   \[\sum_{n\in\N} z^n\]
%   converge. Lorsque c'est le cas, calculer sa somme.
% \end{exoUnique}

\begin{proposition}
Soit $z\in\C$. Alors la série
\[\sum z^n\]
converge si et seulement si $\abs{z}<1$. Si tel est le cas, sa somme est
\[\sum_{n=0}^{+\infty} z^n = \frac{1}{1-z}.\]
\end{proposition}

\begin{exoUnique}
\exo Soit $(F_n)$ la suite de Fibonacci définie par
  \[F_0\defeq 0, \quad F_1\defeq 1, \et \forall n\in\N \qsep F_{n+2}\defeq F_{n+1}+F_n.\]
  Démontrer l'existence puis calculer
  \[\sum_{n=0}^{+\infty} \frac{F_n}{2^n}.\]
\begin{sol}
On trouve 2.
\end{sol}
\end{exoUnique}

% \begin{exoUnique}
% \exo Pour tout $x\in\intero{-1}{1}$, calculer
%   \[\sum (n+1)x^n.\]
%   On pourra conjecturer le résultat avant de le prouver rigoureusement.
% \end{exoUnique}

\begin{proposition}
Soit $z\in\C$. Alors la série
\[\sum \frac{z^n}{n!}\]
est convergente et
\[\sum_{n=0}^{+\infty} \frac{z^n}{n!}=\e^z.\]
\end{proposition}

\begin{exoUnique}
\exo Établir l'existence et calculer
    \[\sum_{n=0}^{+\infty} \frac{n^2+1}{n!}.\]
\begin{sol}
On trouve $3\e$.
\end{sol}
\end{exoUnique}



\subsection{Série à termes positifs}

\begin{definition}
On dit qu'une série réelle $\sum u_n$ est à termes positifs lorsque
\[\forall n\in\N\qsep u_n\geq 0.\]
\end{definition}

\begin{remarques}
\remarque La suite des sommes partielles d'une série à termes positifs est croissante.
  Réciproquement, si $(u_n)$ est une suite croissante, sa suite dérivée $(u_{n+1}-u_n)$ est
  une suite à termes positifs.
\remarque Puisque la convergence d'une série ne dépend pas de ses premiers termes, les
  théorèmes de convergence sur les séries à termes positifs s'appliquent même si la série
  est à termes positifs à partir d'un certain rang. Bien entendu, des théorèmes similaires
  aux théorèmes que nous allons énoncer existent pour les séries à termes négatifs.
  Les théorèmes que nous allons énoncer dans cette section sont donc utiles pour étudier
  les séries de signe constant à partir d'un certain rang.
\end{remarques}

\begin{proposition}
Une série à termes positifs converge si et seulement si la suite de ses sommes partielles
est majorée.
\end{proposition}

\begin{proposition}[nom={Série de \nom{Riemann}}]
Soit $\alpha\in\R$. Alors, la série
\[\sum \frac{1}{n^\alpha}\]
est convergente si et seulement si $\alpha > 1$.
\end{proposition}

\begin{remarqueUnique}
\remarque En particulier, la série harmonique $(H_n)$ définie par
  \[\forall n\in\N\qsep H_n\defeq\sum_{k=1}^n \frac{1}{k}\]
  diverge. Une comparaison série intégrale permet de montrer que
  $H_n\equi{n}{+\infty} \ln n$.
\end{remarqueUnique}

\begin{exos}
\exo Montrer que la série
  \[\sum \frac{1}{n^2}\]
  converge et donner un équivalent de son reste.
\exo Prouver la divergence et donner un équivalent des sommes partielles de la série
  \[\sum \frac{1}{n \ln n}.\]
\end{exos}

\begin{proposition}
Soit $\sum u_n$ et $\sum v_n$ deux séries à termes positifs telles que
\[\forall n\in\N\qsep u_n\leq v_n.\]
\begin{itemize}
\item Si la série $\sum v_n$ converge, alors il en est de même pour $\sum u_n$.
\item Si la série $\sum u_n$ diverge, alors il en est de même pour $\sum v_n$.
\end{itemize}
\end{proposition}

\begin{exos}
\exo Donner la nature des séries
  \[\sum\frac{\sin^2 n}{n^3},\qquad\sum\frac{\cos^2\p{\frac{n\pi}{3}}}{\sqrt{n}}.\]
\exo Soit $r\in\RP$. Déterminer la nature de la série
  \[\sum \frac{r^n}{n}\]
  en fonction de $r$.
\end{exos}

\begin{proposition}
Soit $\sum u_n$ et $\sum v_n$ deux séries à termes positifs. On suppose que
\[u_n=\grando{n}{+\infty}{v_n}\]
et que la série $\sum v_n$ est convergente. Alors la série $\sum u_n$ est convergente.
\end{proposition}

\begin{remarqueUnique}
\remarque En particulier, pour des séries à termes positifs, si
\[u_n=\petito{n}{+\infty}{v_n}\]
et si $\sum v_n$ est convergente, alors $\sum u_n$ est convergente. 
\end{remarqueUnique}

\begin{exoUnique}
\exo Déterminer la nature de la série
  \[\sum \frac{\ln n}{n\sqrt{n}}.\]
\end{exoUnique}

\begin{proposition}
Soit $\sum u_n$ et $\sum v_n$ deux séries réelles. On suppose que $\sum v_n$ est à termes
positifs et que
\[u_n\equi{n}{+\infty} v_n.\]
Alors la série $\sum u_n$ est à termes positifs à partir d'un certain rang et les deux séries
sont de même nature.
\end{proposition}

\begin{exos}
\exo Établir la nature des séries suivantes
  \[\sum \frac{1}{3n+1}, \qquad\sum \tan\p{\frac{1}{n}}, \qquad \sum \frac{1}{2^n-n}.\]
\exo Donner la nature des séries
  \[\sum \ln\p{\tan\frac{\pi n}{4n+1}},\qquad \sum \cro{\p{\tanh n}^{\frac{1}{n}}-1}.\]
\exo Montrer qu'il existe $\gamma\in\R$ tel que
  \[\sum_{k=1}^n \frac 1k = \ln n+ \gamma + \petito{n}{+\infty}{1}.\]
  La constante $\gamma\approx 0.577$ est appelée constante d'\nom{Euler}.
\end{exos}


\subsection{Série absolument convergente}

\begin{definition}
Soit $\sum u_n$ une série d'éléments de $\K$. Si la série à termes positifs
\[\sum \abs{u_n}\]
converge, alors la série $\sum u_n$ converge. On dit dans ce cas que la série
$\sum u_n$ est \emph{absolument convergente}.
\end{definition}

\begin{remarqueUnique}
\remarque Une série convergente qui n'est pas absolument convergente est appelée série
  \emph{semi-convergente}.
\end{remarqueUnique}

\begin{exoUnique}
\exo Montrer que la série
  \[\sum\frac{\sin n}{n\sqrt{n}}\]
  est convergente.
\end{exoUnique}


\begin{proposition}
Soit $\sum u_n$ une série d'éléments de $\K$ et $\sum v_n$ une série à termes positifs
telle que
\[u_n=\grando{n}{+\infty}{v_n}.\]
Si $\sum v_n$ est convergente, alors $\sum u_n$ est absolument convergente, donc convergente.
\end{proposition}

\begin{remarqueUnique}
\remarque En particulier, si $(v_n)$ est une suite positive, si
  \[u_n=\petito{n}{+\infty}{v_n}\]
  et si $\sum v_n$ est convergente, alors $\sum u_n$ est absolument convergente, donc
  convergente.
\end{remarqueUnique}

\begin{proposition}[nom={Règle de d'\nom{Alembert}}]
Soit $\sum u_n$ une série d'éléments de $\K$ ne s'annulant pas. On suppose que
\[\abs{\frac{u_{n+1}}{u_n}}\tendvers{n}{+\infty} \omega\in\RP\cup\ens{+\infty}.\]
Alors
\begin{itemize}
\item Si $\omega<1$, la série $\sum u_n$ est absolument convergente.
\item Si $\omega>1$, la série $\sum u_n$ est grossièrement divergente.
\end{itemize}
\end{proposition}

\begin{remarqueUnique}
\remarque Si $\omega=1$, la règle de d'\nom{Alembert} ne permet pas de conclure. Dans ce cas, il peut être
  intéressant d'effectuer une comparaison avec une série de \nom{Riemann}.
\end{remarqueUnique}

\begin{exos}
\exo Déterminer la nature de la série
  \[\sum \frac{n^4}{3^n}.\]
\exo Soit $z\in\C$. Retrouver le fait que la série
  \[\sum \frac{z^n}{n!}\]
  est convergente.
  \exo Soit $a,b\in\R$. Donner une condition nécessaire et suffisante sur $a$ et $b$ pour que la série
  \[\sum \cro{\frac{n^2+2}{n^2+2n+1}-\p{a+\frac{b}{n}}}\]
  soit convergente.
\end{exos}

\subsection{Série semi-convergente}

\begin{theoreme}[nom={Théorème des séries alternées}]
Soit $(u_n)$ une suite à termes positifs, décroissante et convergeant vers 0.
Alors la série
\[\sum (-1)^n u_n\]
converge. De plus, si $(R_n)$ est la suite des restes définie par
\[\forall n\in\N\qsep R_n\defeq\sum_{k=n+1}^{+\infty} (-1)^k u_k\]
alors, pour tout $n\in\N$, $R_n$ est du signe de $(-1)^{n+1}$ et $\abs{R_n}\leq u_{n+1}$.
\end{theoreme}

\begin{remarques}
\remarque La série
  \[\sum \frac{(-1)^n}{\sqrt{n}}\]
  est convergente, mais n'est pas absolument convergente.
\remarque Les séries alternées permettent de construire des suites $(u_n)$ et $(v_n)$
  qui sont équivalentes en $+\infty$ mais qui ne sont pas de même nature. Par exemple,
  si on définit les suites $(u_n)$ et $(v_n)$ par
  \[\forall n\in\Ns\qsep u_n\defeq\frac{(-1)^n}{\sqrt{n}} \et
    v_n\defeq\frac{(-1)^n}{\sqrt{n}}+\frac{1}{n}\]
  alors $(u_n)$ et $(v_n)$ sont équivalentes en $+\infty$ bien que $\sum u_n$ soit convergente
  et que $\sum v_n$ soit divergente.
\end{remarques}

\begin{exos}
\exo Soit $\alpha\in\R$. Donner une condition nécessaire et suffisante sur $\alpha$ pour que
  la série
  \[\sum \frac{(-1)^n}{n^\alpha}\]
  soit convergente.
\exo Soit $\alpha>0$. Discuter, selon $\alpha$, de la nature de la série
  \[\sum \frac{(-1)^n}{n^\alpha + (-1)^n}.\]
\end{exos}

% \begin{remarqueUnique}
% \remarque Il arrive que l'on doive prouver la convergence de séries de la forme
%   $\sum a_n u_n$ où
%   \begin{itemize}
%   \item La suite $(u_n)$ est une suite positive, décroissante et convergeant vers 0.
%   \item La série de terme général $(a_n)$ est bornée.
%   \end{itemize}
%   Pour prouver la convergence d'une telle série on applique ce qu'on appelle une
%   \emph{transformée d'\nom{Abel}} qui est l'équivalent discret d'une intégration par
%   parties. On définit la suite $(A_n)$ par
%   \[\forall n\in\N\qsep A_n\defeq\sum_{k=0}^{n-1} a_k.\]
%   Alors
%   \begin{eqnarray*}
%   \forall n\in\N\qsep
%   \sum_{k=0}^n a_k u_k
%   &=& \sum_{k=0}^n (A_{k+1} - A_k)u_k\\
%   &=& \sum_{k=0}^n A_{k+1} u_k  - \sum_{k=0}^n A_k u_k\\
%   &=& \sum_{k=1}^{n+1} A_k u_{k-1}  - \sum_{k=0}^n A_k u_k\\
%   &=& \p{A_{n+1} u_n - A_0 u_0} - \sum_{k=1}^n A_k (u_{k}-u_{k-1}).
%   \end{eqnarray*}
%   On montre ensuite que $A_{n+1} u_n$ tend vers 0 lorsque $n$ tend vers $+\infty$ et
%   que la série $\sum A_n(u_n - u_{n-1})$ est absolument convergente, ce qui permet de prouver
%   la convergence de $\sum a_n u_n$.
% \end{remarqueUnique}

% \begin{exoUnique}
% \exo Soit $\theta\in\R\setminus 2\pi\Z$ et $\alpha\in\RPs$. Montrer que la série
%   \[\sum \frac{\cos(n\theta)}{n^\alpha}\]
%   est convergente.
% \end{exoUnique}


% \subsection{Développement décimal d'un réel}

% \begin{proposition}[utile=1, nom=Développement décimal propre d'un réel]
% Soit $x$ un réel positif. Alors il existe un unique $a_0\in\N$ et une unique suite
% $(a_n)_{n \geq 1}$ d'éléments de $\intere{0}{9}$, non stationnaire égale à $9$, telle que 
% $$x=\sum\limits_{n=0}^{+\infty} \frac{a_n}{10^n}.$$
% \end{proposition}


% Le but de cette partie est de donner un sens aux écritures suivantes :
% $$\pi = 3,14159265358979\dots \qquad - \frac{7}{11} = -0,6363636363 \dots $$


% \begin{center}
% \begin{tabular}{lr}
% \begin{minipage}{6cm}
% \begin{eqnarray*}
% x & = & 3,14159265358979\dots \\
% 10^3 x & = & 3141,59265358979\dots \\
% \ent{10^3 x} & = & 3141 \\
% \frac{\ent{10^3 x}}{10^3} &= & 3,141
% \end{eqnarray*}
% \end{minipage}

% &
% \begin{minipage}{6cm}
% \begin{eqnarray*}
% x & = & 3,14159265358979\dots \\
% 10^4 x & = & 31415,9265358979\dots \\
% \ent{10^4x} & = & 31415 \\
% \frac{\ent{10^4x}}{10^4} &= & 3,1415
% \end{eqnarray*}
% \end{minipage}
% \end{tabular}
% \end{center}


% On obtient donc une suite $S_n = \frac{\ent{10^n x}}{10^n}$ qui converge vers $x$ et les chiffres $a_n$ qui sont définis par 
% $$a_0 = \ent{x} \text{ et } a_n = \ent{10^n x} - 10 \ent{10^{n-1}x} \text{ pour } n \geq 1$$
% de sorte que 
% $$x = \sum\limits_{n=0}^{\infty} \frac{a_n}{10^n}$$

% Pour que ceci corresponde au développement décimal usuel, il faut prendre $x \geq 0$. Si $x<0$, $\ent{x}$ ne donne pas l'entier avant la virgule.
% La convention, pour $x<0$, est de prendre le développement décimal de $-x$ puis de mettre un signe moins devant.
% $a_0$ n'est pas un chiffre mais un entier quelconque.

% Par contre, si $n \geq 1$ alors $a_n \in \intere{0}{9}$.
% Malheureusement, il n'y a pas unicité d'une telle écriture. Par exemple,
% $$1=1,000 \dots = 0,999\dots$$

% Pour avoir l'unicité, il faut imposer que la suite $(a_n)$ n'est pas stationnaire égale à $9$ i.e. 
% $$\forall N \in \N \qsep \exists n \geq N \qsep a_n \neq 9.$$
% On dit alors que le développement décimal est propre. Montrons que l'on a bien alors l'unicité et l'existence de l'écriture décimale.

% \begin{proposition}[utile=1, nom=Développement décimal propre d'un réel]
% Soit $x$ un réel positif. Alors il existe un unique $a_0 \in \N$ et une unique suite $(a_n)_{n \geq 1}$ d'entiers entre $0$ et $9$ non stationnaire égale à $9$ telle que 
% $$x=\sum\limits_{n=0}^{\infty} \frac{a_n}{10^n}.$$

% \end{proposition}



%END_BOOK
\end{document}
\section{Exercices}
\setcounter{numeroexercice}{1}
\documentclass{magnolia}

\magtex{tex_driver={pdftex}}
\magfiche{document_nom={Exercices sur les développements limités},
          auteur_nom={François Fayard},
          auteur_mail={fayard.prof@gmail.com}}
\magexos{exos_matiere={maths},
         exos_niveau={mpsi},
         exos_chapitre_numero={22},
         exos_theme={Séries}}
\magmisenpage{}
\maglieudiff{}
\magprocess

\begin{document}
%BEGIN_BOOK
\magsection{Série}
\magsubsection{Série}
\exercice{nom={Exercice}}
Établir la convergence et calculer
\[\sum_{n=0}^{+\infty}\frac{2^n}{(n+2)n!},\qquad
  \sum_{n=0}^{+\infty}\frac{2^n(5n+1)}{3^n n!}.\]
\begin{sol}
On trouve $(1/4)(1+\e^2)$, $(13/3)\e^{2/3}$.
\end{sol}

\exercice{nom={Exercice}}
Soit $(u_n)$ une suite positive décroissante tendant vers $0$. Montrer que
si $\sum u_n$ converge, alors
\[u_n = \petito{n}{+\infty}{\frac{1}{n}}.\]
On pourra considérer la suite de terme général $s_{2n} - s_n$. La réciproque est-elle vraie~?

\magsubsection{Série à termes positifs}

\exercice{nom={Exercice}}
Soit $(u_n)$ une suite réelle positive, décroissante.
\begin{questions}
\question Montrer que $\sum u_n$ converge si et seulement si $\sum 2^n u_{2^n}$ converge.
\question Soit $\alpha,\beta\in\R$. Déterminer la nature de
  \[\sum \frac{1}{n^{\alpha} \ln^\beta n}.\]
\end{questions}


\magsubsection{Série absolument convergente}

\exercice{nom={Exercice}}
Déterminer la nature des séries suivantes.
\[\sum \frac{\ln(n)}{n^2},\qquad \sum \frac{n!}{n^n},\qquad \sum \cro{\e - \p{1+\frac 1n}^n},\]
\[\sum \frac{1}{\binom{2n}{n}}, \qquad \sum \frac{1}{\ln(n)^{\ln(n)}}.\]

\exercice{nom={Exercice}}
Montrer que la suite $(u_n)$ définie par
\[\forall n\in\N\qsep u_n \defeq \prod_{k=1}^n \p{1+ \frac{1}{k^2}}\]
converge vers un réel non nul.

\exercice{nom={Exercice}}
Donner la nature de la série de terme général
\[\frac{n^{2 \lambda}}{\lambda^n + \ln n}\]
suivant $\lambda\in\R$.

\exercice{nom={Exercice}}
Soit $a,b\in\R$. Donner la nature de
\[\sum \p{ \sqrt[3]{n^3+n} - \sqrt[2]{n^2+an+b}}.\]

\exercice{nom={Exercice}}
Pour tout $r \in \N$, on pose
\[S_r(x) \defeq \sum_{n=r}^{+\infty} \binom{n}{r} x^{n-r}.\]
\begin{questions}
\question Pour quelles valeurs de $x\in\R$ la somme $S_r(x)$ est-elle définie~?
\question Calculer $(1-x)S_{r+1}(x)$.
\question En déduire la valeur de $S_r(x)$. 
\end{questions}

\exercice{nom={Exercice}}
Montrer qu'il existe $l\in \R$ tel que
\[\sum_{k=1}^n k^{\frac{1}{k}} = n + \frac{1}{2} \ln^2 n+l +\petito{n}{+\infty}{1}.\]

\magsubsection{Série semi-convergente}

\exercice{nom={Exercice}}
On définit les suites $(u_n)$ et $(v_n)$ par
\[\forall n\geq 2\qsep u_n \defeq \frac{(-1)^n}{\sqrt{n}} \et
  v_n \defeq \frac{(-1)^n}{\sqrt{n}+(-1)^n}.\]
\begin{questions}
\question Montrer que \[u_n \equi{n}{+\infty} v_n.\]
\question Montrer que $\sum u_n$ converge et que $\sum v_n$ diverge.
\end{questions}

\exercice{nom={Exercice}}
Déterminer les suites réelles $(v_n)$ telles que pour toute suite réelle positive $(u_n)$
\[\sum u_n \text{ converge } \quad\implique\quad \sum u_n v_n \text{ converge }.\]


%END_BOOK

\end{document}


\chapter{Espaces euclidiens}
\setcounter{numeroexercicecours}{1}
\documentclass{magnolia}

\magtex{tex_driver={pdftex},
        tex_packages={xypic}}
\magfiche{document_nom={Cours sur les espaces euclidiens},
          auteur_nom={François Fayard},
          auteur_mail={fayard.prof@gmail.com}}
\magcours{cours_matiere={maths},
          cours_niveau={mpsi},
          cours_chapitre_numero={22},
          cours_chapitre={Espaces euclidiens}}
\magmisenpage{}
\maglieudiff{}
\magprocess

\begin{document}

%BEGIN_BOOK
\magtoc

\section{Produit scalaire}
\subsection{Produit scalaire}

\begin{definition}[utile=-3]
Soit $E$ un \Rev. On dit qu'une application
\[\dspappli{\ps{.}{.}}{E\times E}{\R}{\p{x,y}}{\ps{x}{y}}\]
est un \emph{produit scalaire} lorsqu'elle est
\begin{itemize}
\item \emph{bilinéaire}
  \begin{eqnarray*}
  \forall x,y,z\in E \qsep \forall \lambda,\mu\in\R,
  &&\ps{x}{\lambda y+\mu z}=\lambda\ps{x}{y}+\mu\ps{x}{z}\\
  &&\ps{\lambda x+\mu y}{z}=\lambda\ps{x}{z}+\mu\ps{y}{z}
  \end{eqnarray*}
\item \emph{symétrique}
  \[\forall x,y\in E \qsep \ps{x}{y}=\ps{y}{x}\]
\item \emph{positive}
  \[\forall x\in E \qsep \ps{x}{x}\geq 0\]
\item \emph{définie}
  \[\forall x\in E \qsep \ps{x}{x}=0 \quad\implique\quad x=0\]
\end{itemize}
On appelle espace \emph{préhilbertien réel} tout \Rev muni d'un produit scalaire.
\end{definition}

\begin{remarques}
\remarque Le produit scalaire $\ps{x}{y}$ est parfois noté $\left\langle x,y\right\rangle$, $(x|y)$,
  $(x,y)$ ou $x\cdot y$.
\remarque Le caractère symétrique du produit scalaire fait qu'il suffit de montrer
  la linéarité par rapport à l'une des variables pour en déduire la bilinéarité.
\remarque Quel que soit $x\in E$, $\ps{x}{0}=0$ et $\ps{0}{x}=0$.
\end{remarques}

\begin{exos}
\exo Soit $n\in\N$. Montrer que l'application $\ps{.}{.}$ définie par
  \[\forall \p{x_1,\ldots,x_n},\p{y_1,\ldots,y_n}\in\R^n \qsep
    \ps{\p{x_1,\ldots,x_n}}{\p{y_1,\ldots,y_n}}\defeq\sum_{k=1}^n x_k y_k.\]
  est un produit scalaire sur $\R^n$.
\exo On pose $E\defeq\polyR$. Montrer que l'application $\ps{.}{.}$ définie par
  \[\forall P,Q\in\polyR \qsep \ps{P}{Q}\defeq\integ{0}{1}{P(x)Q(x)}{x}\]
  est un produit scalaire sur $\polyR$.
\exo On pose $E\defeq\mat{q,p}{\R}$. Montrer que l'application $\ps{.}{.}$ définie par
  \[\forall A,B\in\mat{q,p}{\R} \qsep \ps{A}{B}\defeq \tr\p{\trans{A}B}\]
  est un produit scalaire sur $\mat{q,p}{\R}$.
\end{exos}


\subsection{Norme}

\begin{definition}[utile=-3]
Soit $x\in E$. On définit la \emph{norme} de $x$, notée $\norme{x}$, par
\[\norme{x}=\sqrt{\ps{x}{x}}.\]
On dit que $x$ est \emph{normé} lorsque $\norme{x}=1$.
\end{definition}

\begin{remarqueUnique}
\remarque Le plus souvent, il sera préférable de travailler avec $\norme{x}^2=\ps{x}{x}$
  qu'avec $\norme{x}$.
\end{remarqueUnique}

\begin{proposition}[utile=-3]
\begin{eqnarray*}
\forall x\in E \qsep \forall\lambda\in\R, & & \norme{\lambda x}=
  \abs{\lambda}\norme{x}\\
\forall x\in E, & & \norme{x}=0 \quad\ssi\quad x=0.
\end{eqnarray*}
\end{proposition}

\begin{remarqueUnique}
  \remarque Si $x\in E$ est un vecteur non nul, il y a exactement deux vecteurs 
    normés colinéaires à $x$
    \[\frac{x}{\norme{x}} \qquad\et\qquad -\frac{x}{\norme{x}}.\]
  \end{remarqueUnique}

\begin{proposition}[utile=3, nom=Inégalité de \nom{Cauchy-Schwarz}]
Soit $x,y\in E$. Alors
\[\abs{\ps{x}{y}}\leq\norme{x}\norme{y}.\]
De plus, l'égalité a lieu si et seulement si $x$ et $y$ sont colinéaires.  Plus précisément
\begin{itemize}
\item $\ps{x}{y}=\norme{x}\norme{y}$ si et seulement si il existe $\lambda\geq 0$ tel que
$y=\lambda x$ ou $x=\lambda y$.
\item $\ps{x}{y}=-\norme{x}\norme{y}$ si et seulement si il existe
$\lambda\geq 0$ tel que $y=-\lambda x$  ou $x=-\lambda y$.
\end{itemize}
\end{proposition}

\begin{remarques}
\remarque En particulier, avec le produit scalaire usuel sur $\R^n$, on a
  \[\forall (x_1,\ldots,x_n),(y_1,\ldots,y_n)\in\R^n\qsep
    \abs{\sum_{k=1}^n x_k y_k} \leq \sqrt{\sum_{k=1}^n x_k^2}\sqrt{\sum_{k=1}^n y_k^2}.\]
\remarque L'inégalité de \nom{Cauchy-Schwarz} s'adapte au cas où $\phi:E\times E\to\R$
  est une forme bilinéaire symétrique positive quelconque, même lorsque cette
  dernière n'est pas définie. Si $x,y\in E$, on a alors
  \[\abs{\phi(x,y)}\leq\sqrt{\phi(x,x)}\sqrt{\phi(y,y)}.\]
  Cependant, la condition d'égalité n'est plus valide.
%   De plus, cette inégalité est une égalité si et seulement si il existe $\lambda\in\R$ et
%   $a\in A$ tels que $x=\lambda y+a$ ou $y=\lambda x+a$. Plus précisément
% \begin{itemize}
% \item $\phi(x,y)=\sqrt{\phi(x,x)}\sqrt{\phi(y,y)}$ si et seulement si il existe
%   $\lambda\geq 0$ et $a\in A$ tels que $x=\lambda y+a$ ou
%   $y=\lambda x+a$.
% \item $\phi(x,y)=-\sqrt{\phi(x,x)}\sqrt{\phi(y,y)}$ si et seulement si il existe
%   $\lambda\geq 0$ et $a\in A$ tels que $x=-\lambda y+a$ ou
%   $y=-\lambda x+a$.
% \end{itemize}
\end{remarques}


\begin{preuve}
Définissons $$\begin{array}{ccccc}
f & : & \R & \to & \R \\
 & & \lambda & \mapsto & ||y-\lambda x||^2 
\end{array}.$$

Si $x=0$ $\ldots$ sinon, $\Delta\le 0$, etc $\ldots$
\end{preuve}

\begin{exoUnique}
\exo Montrer que
  \[\forall x,y,z\in\R \qsep x^2+2y^2+3z^2\leq 1 \implique
    \abs{x+y+z}\leq\sqrt{\frac{11}{6}}\]
  Si on suppose que $x^2+2y^2+3z^2\leq 1$, à quelle condition a-t-on
  $\abs{x+y+z}=\sqrt{11/6}$~?
  \begin{sol}
  Considérer le produit scalaire usuel sur $\R^3$ et les vecteurs
  $u=(x,\sqrt{2}y,\sqrt{3}z)$ et $v=\p{1,\frac{1}{\sqrt{2}},\frac{1}{\sqrt{3}}}$.
  \end{sol}
\end{exoUnique}

\begin{proposition}[utile=3, nom=Inégalité de \nom{Cauchy-Schwarz}]
Soit $f,g$ deux fonctions réelles continues sur l'intervalle $I$
et $a,b\in I$. Si $a\leq b$, alors
\[\abs{\integ{a}{b}{f(x)g(x)}{x}}\leq
  \p{\integ{a}{b}{f^2(x)}{x}}^{\frac{1}{2}}
  \p{\integ{a}{b}{g^2(x)}{x}}^{\frac{1}{2}}.\]
De plus, cette inégalité est une égalité si et seulement si il existe
$\lambda\in\R$ tel que
\[\cro{\forall x\in\interf{a}{b} \qsep f(x)=\lambda g(x)}
  \quad\text{ou}\quad
  \cro{\forall x\in\interf{a}{b} \qsep g(x)=\lambda f(x)}.\]
\end{proposition}

\begin{remarqueUnique}
\remarque Cette inégalité reste vraie si $f$ et $g$ sont des fonctions continues
  par morceaux. Cependant, dans ce cas, la condition d'égalité n'est plus valide.
\end{remarqueUnique}

\begin{exos}
\exo Soit $f$ une fonction continue sur $\interf{a}{b}$,
  à valeurs strictement positives. Montrer que
  \[\p{\integ{a}{b}{f(x)}{x}}\p{\integinv{a}{b}{f(x)}{x}}\geq\p{b-a}^2.\]
  Donner une condition nécessaire et suffisante sur $f$ pour que cette inégalité
  soit une égalité.
\exo Soit $a,b\in\RPs$ tels que $a\leq b$. Montrer que
  \[\integinv{a}{b}{x}{x}\leq \frac{b-a}{\sqrt{ab}}.\]
\end{exos}

\begin{proposition}[utile=3, nom=Inégalité triangulaire]
Soit $x,y\in E$.
\begin{itemize}
\item Alors
  \[\norme{x+y}\leq\norme{x}+\norme{y}.\]
  De plus, l'égalité a lieu si et seulement si $x$ et $y$ sont positivement
  liés, c'est-à-dire si et seulement si il existe $\lambda\geq 0$ tel que
  $y=\lambda x$ ou $x=\lambda y$.
\item De plus
  \[\abs{\norme{x}-\norme{y}}\leq\norme{x-y} \et
    \norme{x+y}\geq\norme{x}-\norme{y}.\]
\end{itemize}
\end{proposition}

\begin{preuve}
Différence des deux termes élevés au carré + Cauchy-Schwarz
\end{preuve}

\begin{proposition}[utile=2]
L'identité suivante est appelée identité du parallélogramme
\[\forall x,y\in E, \quad \norme{x+y}^2+\norme{x-y}^2=
  2\p{\norme{x}^2+\norme{y}^2}.\]
\end{proposition}

\begin{remarques}
\remarque Dans un parallélogramme, la somme des carrés des longueurs des
  diagonales est égale à la somme des carrés des côtés.
\remarque Les identités suivantes sont appelées identités de polarisation.
  \begin{eqnarray*}
  \forall x,y\in E \qsep \ps{x}{y}
  &=& \frac{1}{2}\p{\norme{x+y}^2-\p{\norme{x}^2+\norme{y}^2}}\\
  &=& \frac{1}{2}\p{\p{\norme{x}^2+\norme{y}^2}-\norme{x-y}^2}\\
  &=& \frac{1}{4}\p{\norme{x+y}^2-\norme{x-y}^2}.
  \end{eqnarray*}
% \remarque Soit $E$ un \Rev muni de deux produits scalaires $\ps{.}{.}_1$ et
%   $\ps{.}{.}_2$ définissant respectivement les normes $\norme{.}_1$ et
%   $\norme{.}_2$. On suppose que~:
%   \[\forall x\in E \quad \norme{x}_1=\norme{x}_2\]
%   Alors~:
%   \[\forall x,y\in E \quad \ps{x}{y}_1=\ps{x}{y}_2\]
\end{remarques}
\begin{sol} Intérêt, on peut retrouver le produit scalaire à partir de la norme. En particulier si on dispose de deux p.s. sur $E$ qui donne une même norme alors les deux p.s. sont égaux.
\end{sol}

\begin{exos}
\exo Soit $x,y\in E$ deux vecteurs distincts de normes inférieures ou égales
  à 1. Montrer que
  \[\norme{\frac{x+y}{2}}<1.\]
  \begin{sol}
  Faire un dessin dans $\R^2$ de la situation. Puis pour le prouver on va utiliser l'identité du parallélogramme :
  $$\norme{x+y}^2+\norme{x-y}^2=
  2\p{\norme{x}^2+\norme{y}^2}\leq 4$$ donc $$\norme{x+y}^2\leq 4-\norme{x-y}^2<4$$ car $x$ et $y$ sont distincts. D'où le résultat.
  \end{sol}
\exo Soit $u\in\Endo{E}$ tel que, pour tout $x\in E$, $\ps{u(x)}{x}=0$.
  Montrer que
  \[\forall x,y\in E \qsep \ps{u(x)}{y}=-\ps{x}{u(y)}.\]
  \begin{sol} On applique l'hypothèse à $x+y$ et on développe le tout en utilisant à nouveau l'hypothèse on peut conclure.
  \end{sol}
\end{exos}

\subsection{Notion d'orthogonalité}

\begin{definition}[utile=-3]
$\quad$
\begin{itemize}
\item Soit $x,y\in E$. On dit que $x$ et $y$ sont \emph{orthogonaux} lorsque
  $\ps{x}{y}=0$.
\item Soit $A$ et $B$ deux parties de $E$. On dit que $A$ et $B$ sont
  \emph{orthogonales} lorsque
  \[\forall a\in A \qsep \forall b\in B \qsep \ps{a}{b}=0.\]
\end{itemize}
\end{definition}

\begin{remarqueUnique}
\remarque Le vecteur nul est orthogonal à tout vecteur.
\end{remarqueUnique}

\begin{definition}[utile=-3]
Soit $(x_i)_{i\in I}$ une famille de vecteurs de $E$.
\begin{itemize}
\item On dit que cette famille est \emph{orthogonale} lorsque
  \[\forall i,j\in I \qsep i\neq j \quad\implique\quad \ps{x_i}{x_j}=0.\]
\item On dit que cette famille est \emph{orthonormée} lorsque
  \[\forall i,j\in I \qsep \ps{x_i}{x_j}=\delta_{i,j}\defeq
  \begin{cases}
  0 & \text{si $i\neq j$}\\
  1 & \text{si $i=j$.}
  \end{cases}\]
\end{itemize}
\end{definition}

\begin{remarqueUnique}
\remarque Sur $\R^n$ muni du produit scalaire usuel, la base canonique est
  orthonormée.
\end{remarqueUnique}

\begin{proposition}[utile=-3]
Toute famille orthogonale ne contenant aucun vecteur nul est libre.
\end{proposition}

\begin{preuve}
$$\sum_{i=1}^p\lambda_i x_i=0$$ puis produit scalaire avec $x_j$ donne le résultat car $x_j\neq 0$.
\end{preuve}

\begin{sol}
Remarque : toute famille orthonormée est donc libre.
\end{sol}

\begin{remarqueUnique}
\remarque En particulier, toute famille orthonormée est libre.
\end{remarqueUnique}

\begin{exoUnique}
\exo Soit $E$ le \Rev des fonctions réelles continues, $2\pi$-périodiques sur $\R$.
  On définit sur $E$ le produit scalaire
  \[\forall f,g\in E \qsep \ps{f}{g}\defeq\integ{0}{2\pi}{f(x)g(x)}{x}.\]
  Pour tout $k\in\Ns$, on définit la fonction $s_k$ par
  $s_k(x)\defeq\sin\p{kx}$ et pour tout $k\in\N$, la fonction $c_k$ par
  $c_k(x)\defeq\cos\p{kx}$. Montrer que pour tout $n\in\N$, la famille
  $(c_0,s_1,c_1,\ldots,s_n,c_n)$ est orthogonale.
\end{exoUnique}

\begin{sol}
Soit $k_1\neq k_2$ dans $\intere{1}{n}$ :
\begin{eqnarray*}
\ps{s_{k_1}}{s_{k_2}}&=&\int_{0}^{2\pi}\sin(k_1 x)\sin(k_2 x) dx \\
&=& \frac{1}{2}\int_{0}^{2\pi}\p{\cos((k_1-k_2)x)-\cos((k_1+k_2)x)} dx\\
&=&...=0
\end{eqnarray*}
De même, $\ps{c_{k_1}}{c_{k_2}}=0$ pour $k_1\neq k_2$ dans $\intere{0}{n}$.\\
Soit maintenant $k_1\in \intere{1}{n}$, $k_2\in\intere{0}{n}$. Alors : 
\begin{eqnarray*}
\ps{s_{k_1}}{c_{k_2}}&=&\int_{0}^{2\pi}\sin(k_1 x)\cos(k_2 x) dx \\
&=& \frac{1}{2}\int_{0}^{2\pi}\p{\sin((k_1+k_2)x)-\sin((k_2-k_1)x)} dx\\
&=&-\frac{1}{2}\int_{0}^{2\pi}\sin((k_2-k_1)x) dx\\
&=&0 \quad (\text{en séparant tout de même les cas } k_2\neq k_1 \et k_1=k_2).
\end{eqnarray*}

\end{sol}


\begin{proposition}[utile=-3, nom=\nom{Pythagore}]
Soit $(x_1,\ldots,x_n)$ une famille orthogonale. Alors
\[\norme{x_1+\cdots+x_n}^2=\sum_{k=1}^n \norme{x_k}^2.\]
\end{proposition}

\begin{preuve}
On développe la norme au carré par bilinéarité.
\end{preuve}

\begin{remarqueUnique}
\remarque Si $x$ et $y$ sont deux vecteurs de $E$, alors
  $\norme{x+y}^2=\norme{x}^2+\norme{y}^2$ si et seulement si $x$ et $y$ sont
  orthogonaux. C'est le théorème de \nom{Pythagore}.
\end{remarqueUnique}

\begin{sol}
La réciproque est vraie seulement pour deux vecteurs. On peut considérer $(1;2)$, $(0;2)$, $(0;-1)$ dans $\R^2$ pour un contre-exemple. Et dans tout $\R^n$, suffit de rajouter des $0$ à ces trois vecteurs.
\end{sol}

\begin{definition}[utile=-3]
Soit $A$ une partie de $E$. On appelle \emph{orthogonal de $A$} et on note $A^{\perp}$
l'ensemble des vecteurs orthogonaux à tous les éléments de $A$
\[A^{\perp}\defeq\enstq{x\in E}{\forall a\in A \qsep \ps{a}{x}=0}.\]  
En particulier $A$ et $A^\perp$ sont orthogonaux.
\end{definition}

\begin{remarques}
\remarque $\ens{0}^\perp=E$ et $E^\perp=\ens{0}$.
\remarque Si $A$ et $B$ sont deux parties de $E$ telles que $A\subset B$, alors
  $B^\perp\subset A^\perp$.
\end{remarques}

\begin{sol}
Très important de faire ces démos et notamment d'insister sur "quand on est orthogonal à tout le monde, on est en particulier orthogonal à soi-même".
\end{sol}

\begin{proposition}[utile=-3]
Soit $A$ une partie de $E$.
\begin{itemize}
\item $A^{\perp}$ est un sous-espace vectoriel de $E$.
\item $A^{\perp}=\p{\vect A}^{\perp}$.
\end{itemize}
\end{proposition}

\begin{preuve}
\begin{itemize}
\item Par caractérisation des sous-espaces vectoriels.
\item $A\subset \vect(A)$ donc $\p{\vect A}^{\perp}\subset A^{\perp}$. Réciproquement, on le fait à la main.
\end{itemize}
\end{preuve}

\begin{remarqueUnique}
\remarque Si $F$ est un sous-espace vectoriel de $E$, $F$ et $F^\perp$ sont en
  somme directe.
\end{remarqueUnique}

\begin{sol}
Suffit de montrer que l'intersection est réduite à $\set{0}$. Notons qu'en dimension finie, ils seront supplémentaires mais pas nécessairement en dimension infinie où on verra des contre-exemples, à commencer par l'exercice suivant.
\end{sol}

\begin{exoUnique}
\exo Soit $E\defeq\classec{0}\p{\interf{-1}{1},\R}$ muni du produit scalaire
  \[\forall f,g\in E \qsep \ps{f}{g}\defeq\integ{-1}{1}{f(x)g(x)}{x}.\]
  On pose $F\defeq\enstq{f\in E}{\forall x\in\interf{0}{1} \qsep f(x)=0}$.
  Calculer $F^\perp$, puis $F\oplus F^\perp$.
  \begin{sol}
  On pose $G=\enstq{g\in E}{\forall x\in\interf{-1}{0} \quad g(x)=0}$ et on montre que $F^\perp=G$. L'inclusion droite-gauche ne pose aucun soucis. Considérons maintenant $h\in F^\perp$. Montrons que $\forall x \in \interf{-1}{0}, h(x)=0$. Pour cela introduisons $$\dspappli{\phi}{\interf{-1}{1}}{\R}{x}{\begin{cases}|x| \text{ si } x\leq 0\\0 \text{ si } x\geq 0\end{cases}}.$$ $\phi$ est continue et $h\phi \in F$. Le fait que $h$ soit orthogonal à $f\phi$ nous conduit alors au résultat.\\
  On montre ensuite que $F\oplus F^\perp=\set{f\in E, f(0)=0}$ par double inclusion.
  \end{sol}
\end{exoUnique}


\section{Espace euclidien}

\begin{definition}[utile=-3]
On appelle \emph{espace euclidien} tout \Rev de dimension finie muni d'un produit
scalaire.
\end{definition}

\begin{remarqueUnique}
\remarque Dans la suite de ce chapitre, excepté dans la section sur la minimisation de la distance à un
  sous-espace vectoriel, $E$ désignera un espace euclidien.  
\end{remarqueUnique}



\subsection{Supplémentaire orthogonal}

\begin{proposition}[utile=2]
Soit $F$ un sous-espace vectoriel de $E$. Alors
\[E=F\oplus F^{\perp}.\]
De plus, $F^\perp$ est l'unique supplémentaire de $F$ orthogonal à $F$; on
l'appelle le \emph{supplémentaire orthogonal} de $F$.
\end{proposition}

\begin{preuve}
Soit $E$ un espace euclidien et $F$ un sev de $E$. Montrons que $E=F\oplus F^{\perp}$. On note $p=\dim(F)$ et $(e_1,\ldots,e_p)$ une base de $F$. Alors $$F^\perp=\p{\vect(e_1,\ldots,e_p)}^\perp=\set{e_1,\ldots,e_p}^\perp.$$
Ainsi, en définissant $$\dspappli{\phi}{E}{\K^p}{x}{(\ps{e_1}{x},\ldots,\ps{e_p}{x})}$$ on montre que $\phi$ est linéaire et que $\ker(\phi)=F^\perp$. D'après le théorème du rang, on a alors :
$$\dim(F^\perp)=\dim(E)-\dim(\im(\phi))\geq \dim(E)-p$$ donc comme la somme est directe :
$$\dim(F\oplus F^{\perp})=\dim(F)+\dim(F^{\perp})\geq p+\dim(E)-p=\dim(E).$$ Donc on a égalité des dimensions et inclusion d'où $E=F\oplus F^{\perp}$.

Montrons désormais l'unicité d'un tel supplémentaire orthogonal. Si on en considère $G$ un autre, comme tout élément de $G$ est orthogonal à $F$, on a $G\subset F^\perp$ et le fait qu'ils soient tous deux supplémentaires à $F$ dans $E$ nous donne l'égalité des dimensions.
\end{preuve}


\begin{remarqueUnique}
\remarque Lorsque $E$ n'est pas de dimension finie, la somme $F+ F^\perp$ reste
  directe, mais elle n'est pas toujours égale à $E$.
\end{remarqueUnique}


\begin{exoUnique}
  \exo On considère $E\defeq\mat{n}{\R}$ muni du produit scalaire $\ps{.}{.}$ défini par
    \[\forall A,B\in\mat{n}{\R} \qsep \ps{A}{B}\defeq\tr\p{\trans{A}B}.\]
    Montrer que $\mathcal{S}_n\p{\R}^\perp=\mathcal{A}_n\p{\R}$.
    \begin{sol}
    On sait déjà qu'ils sont supplémentaires, reste à voir qu'ils sont orthogonaux pour ce produit scalaire là.
    \end{sol}
  \end{exoUnique}

\begin{proposition}[utile=-3]
Soit $F$ un sous-espace vectoriel de $E$. Alors
\[\dim E=\dim F+\dim F^\perp.\]
\end{proposition}



\begin{proposition}[utile=2]
Soit $F$ un sous-espace vectoriel de $E$. Alors
\[\p{F^\perp}^\perp=F.\]
\end{proposition}

\begin{preuve}
$F$ est un supplémentaire de $F^\perp$ qui lui est orthogonal, c'est donc SON supplémentaire orthogonal.
\end{preuve}



\subsection{Base orthonormée}

\begin{proposition}[utile=-3]
Soit $F$ un sous-espace vectoriel de $E$. Si $(f_1,\ldots,f_p)$ une base
orthonormée de $F$ et $(g_1,\ldots,g_q)$ est une base orthonormée de $F^\perp$,
alors $(f_1,\ldots,f_p,g_1,\ldots,g_q)$ est une base orthonormée de $E$.
\end{proposition}

\begin{proposition}[utile=-3]
Tout espace euclidien admet une base orthonormée.  
\end{proposition}

\begin{preuve}
Par récurrence sur $n=\dim(E)$. On peut commencer à $0$ avec la famille vide. Et sinon, pour l'hérédité, on se donne un élément $a\neq 0$ (possible car $n+1\geq 1$) et on pose $e_{n+1}=\dfrac{1}{\norme{a}}a$. On applique alors l'HR à $\p{\vect(e_{n+1})}^\perp$ et on applique la proposition précédente pour conclure.
\end{preuve}


\begin{proposition}[utile=-3,nom={Théorème de la base incomplète}]
On peut compléter toute famille orthonormée de $E$ en une base orthonormée de $E$.
\end{proposition}

\begin{preuve}
On complète en prenant une base orthogonale du supplémentaire orthogonal de $\vect(e_1,\ldots,e_m)$.
\end{preuve}

\begin{proposition}[utile=-3]
Soit $\mathcal{B}\defeq\p{e_1,\ldots,e_n}$ une base orthonormée de $E$. Alors, pour
tout $x\in E$
\[x=\sum_{k=1}^n \ps{e_k}{x} e_k.\]
% Autrement, dit~:
% \[\forall k\in\intere{1}{n} \quad e_k^*(x)=\ps{e_k}{x}\]
\end{proposition}

\begin{preuve}
Soit $x\in E$ qu'on écrit $x=\sum_{k=1}^n \lambda_k e_k$. On a le résultat en calculant le produit scalaire avec chaque $e_k$.
\end{preuve}

\begin{sol}
Remarque : 
Autrement, dit~:
 \[\forall k\in\intere{1}{n} \quad e_k^*(x)=\ps{e_k}{x}\]
\end{sol}

\begin{remarques}
\remarque Le calcul des coordonnées de $x\in E$ dans une base orthonormée se fait
  donc en effectuant des produits scalaires.
\remarque En termes de base duale, la proposition précédente s'énonce
  \[\forall x\in E\qsep \forall k\in\intere{1}{n}\qsep e_k^{\star}(x)=\ps{e_k}{x}.\]
\end{remarques}

\begin{proposition}[utile=-3]
Soit $\mathcal{B}\defeq\p{e_1,\ldots,e_n}$ une base orthonormée de $E$.
\begin{itemize}
\item Soit $x,y\in E$ dont les coordonnées dans la base $\mathcal{B}$ sont
  respectivement $\p{x_1,\ldots,x_n}$ et $\p{y_1,\ldots,y_n}$. Alors
  \[\ps{x}{y}=\sum_{k=1}^n x_k y_k.\]
  Autrement dit, si $X\defeq\mat{\mathcal{B}}{x}$ et $Y\defeq\mat{\mathcal{B}}{y}$
  \[\ps{x}{y}=\trans{X}Y.\]
  En particulier
  \[\norme{x}^2=\sum_{k=1}^n x_k^2=\trans{X}X.\]
\item Soit $(x_1,\ldots,x_p)$ une famille de vecteurs de $E$ et
  $M\defeq\mat{\mathcal{B}}{x_1,\ldots,x_p}\in\mat{n,p}{\R}$. Alors
  \[\trans{M}M=\p{\ps{x_i}{x_j}}_{\substack{1\leq i\leq p\\ 1\leq j\leq p}}\]
\end{itemize}
\end{proposition}

\begin{preuve}
\begin{itemize}
\item On a $$x=\sum_{i=1}^n x_i e_i \qquad ; \qquad y=\sum_{i=1}^n y_i e_i.$$
On a alors $X=\begin{pmatrix}x_1\\ \vdots \\ x_n \end{pmatrix} \text{ et } Y=\begin{pmatrix}y_1\\ \vdots \\ y_n \end{pmatrix}$.
Or, 
\begin{eqnarray*}\ps{x}{y}&=&<\sum_{i=1}^n x_i e_i,\sum_{j=1}^n y_j e_j > \\
&=& \sum_{i=1}^n x_i< e_i,\sum_{j=1}^n y_j e_j >\\
&=& \sum_{i=1}^n\sum_{j=1}^n  x_i y_j< e_i, e_j >\\
&=& \sum_{i=1}^n  x_i y_i\\
&=& \trans{X}Y.
\end{eqnarray*}
\item $\forall j \in \intere{1}{n}$, $x_j=\sum_{i=1}^n m_{i,j}e_i$. Alors~:
  \[\cro{\trans{M}M}_{i,j}=\sum_{k=1}^n \cro{\trans{M}}_{i,k}\cro{M}_{k,j}=\sum_{k=1}^n m_{k,i}m_{k,j}=\ps{x_i}{x_j}\]
  
  En particulier $x_1,\ldots,x_n$ est une base orthonormée de $E$ si et
  seulement si $\forall i,j$, $\ps{x_i}{x_j}=\delta_{i,j}$, i.e. ssi $\trans{M}M=I_n$.
\end{itemize}

\end{preuve}

\begin{proposition}
Soit $\mathcal{B}$ une base orthonormée de $E$ de dimension $n$ et $(x_1,\ldots,x_n)$
une famille de $n$ vecteurs de $E$. On pose $M\defeq\mat{\mathcal{B}}{x_1,\ldots,x_n}$. 
Alors $(x_1,\ldots,x_n)$
est une base orthonormée de $E$ si et seulement si
\[\trans{M}M=I_n.\]
\end{proposition}

% \begin{remarqueUnique}
% \remarque Puisqu'une matrice est inversible à gauche si et seulement si elle est
%   inversible à droite, cette condition est équivalente à $M\trans{M}=I_n$.
%   De telles matrices sont dites \emph{orthogonales}.
% \end{remarqueUnique}

% \subsection{Projecteurs et symétries orthogonaux}

% \begin{definition}[utile=-3]
% Soit $F$ un sous-espace vectoriel de $E$. On appelle \emph{symétrie orthogonale par
% rapport à $F$}, la symétrie par rapport à $F$ parallèlement à $F^\perp$.
% \end{definition}

% \begin{proposition}[utile=-3]
% Soit $s$ une symétrie orthogonale. Alors
% \[\forall x\in E \qsep \norme{s(x)}=\norme{x}.\]
% \end{proposition}

% \begin{preuve}
% En se munissant de $F$ tel que $s$ est la symétrie orthogonale par rapport à $F$, en décomposant $x$ sur $F\oplus F^\perp$ en $x=f+g$, on a $s(x)=f-g$ et $$\norme{s(x)}^2=\norme{f+(-g)}^2=\norme{f}^2+\norme{-g}^2=\norme{f}^2+\norme{g}^2=\norme{x}^2$$ où on a utilisé Pythagore.

% \end{preuve}

% \begin{definition}[utile=-3]
% On appelle \emph{réflexion} toute symétrie orthogonale par rapport à un hyperplan.
% \end{definition}


% \begin{proposition}[utile=2]
% Soit $H$ un hyperplan de $E$ et $u$ un vecteur non nul orthogonal à $H$.  Alors
% \[\forall x\in E \qsep \rho(x)=x-2\frac{\ps{u}{x}}{\norme{u}^2}u.\]
% \end{proposition}

% \begin{preuve}
% $H^\perp$ est de dimension $1$ et $u$ en est un vecteur unitaire (donc non nul) donc en est une BON. Tout $x$ de $E$ se décompose donc sur $H\oplus H^\perp$ en $x=h+g=h+\ps{u}{x}u$. 
% Donc $$s(x)=h-g=h+g-2g=x-2\ps{u}{x}u.$$
% \end{preuve}

% \begin{remarques}
% \remarque Soit $u$ est un vecteur non nul de $E$, et $H=\ens{u}^\perp$. Soit
%   $p$ la projection orthogonale sur $H$ et $s$ la symétrie orthogonale par
%   rapport à $H$. Alors
%   \[\forall x\in E \quad s(x)=x-2\frac{\ps{x}{u}}{\norme{u}^2}u \et
%                          p(x)=x-\frac{\ps{x}{u}}{\norme{u}^2}u\]
% \end{remarques}

% \begin{exos}
% \exo Déterminer la matrice dans la base canonique de $\R^3$ de la
%   réflexion par rapport au plan $P$ d'équation $ax+by+cz=0$ où $a,b,c\in\R$
%   vérifient $a^2+b^2+c^2=1$.
%   \begin{sol}
%   On pose $u=(a,b,c)$, il s'agit alors de la symétrie par rapport à $\set{u}^\perp$ et on a $s(x)=x-2\ps{u}{x}u$ ce qu'on applique en chaque vecteur de la base canonique et on trouve
%   \[S=\begin{pmatrix}
%     1-2a^2 & -2ab & -2ac\\
%     -2ab & 1-2b^2 & -2bc\\
%     -2ac & -2bc & 1-2c^2
%     \end{pmatrix}\]
%   \end{sol}
% \exo Dans $\R^3$, déterminer la matrice dans la base canonique de la
%   réflexion par rapport au plan $P$ d'équation $x-y+z=0$.
%   \begin{sol}
%   On applique l'exercice précédent avec $a=\dfrac{1}{\sqrt{3}}$, $b=-\dfrac{1}{\sqrt{3}}$ et $c=\dfrac{1}{\sqrt{3}}$.
%   On trouve
%   \[S=\frac{1}{3}\begin{pmatrix}
%     1 & 2 & -2\\
%     2 & 1 & 2\\
%     -2 & 2 & 1
%     \end{pmatrix}\]
%   \end{sol}
% \end{exos}

% \begin{proposition}[utile=2]
% Soit $x$ et $y$ deux vecteurs distincts de même norme. Alors il existe une et
% une seule réflexion $\rho$ telle que $\rho(x)=y$ et $\rho(y)=x$.
% \end{proposition}

% \begin{preuve}FAIRE UN DESSIN.\\
% On procède par analyse-synthèse :
% \begin{itemize}
% \item[$\bullet$] Analyse : Supposons trouvé une telle réflexion. On a alors :
% $$\rho(x-y)=\rho(x)-\rho(y)=y-x=-(x-y)$$ donc $x-y \in \ker(\rho+\id)$. Or $\rho$ étant une réflexion, c'est une symétrie par rapport à $\ker(\rho-\id)$ (de dimension $n-1$) parallèlement à $\ker(\rho+\id)$ qui est donc de dimension $1$. Comme $x\neq y$, $x-y$ est une base de $\ker(\rho+\id)$. Donc $\rho$ est la symétrie orthogonale par rapport à $\set{x-y}^\perp$, d'où l'unicité.
% \item[$\bullet$] Synthèse : Considérons donc $\rho$ la symétrie orthogonale par rapport à $\set{x-y}^\perp$. Il s'agit bien d'une réflexion puisque $H:=\set{x-y}^\perp$ est un hyperplan. \\
% Faisons un dessin des vecteurs $x$, $y$, $\dfrac{x+y}{2}$, et  $\dfrac{x-y}{2}$.\\
% On a $\ps{\dfrac{x+y}{2}}{\dfrac{x-y}{2}}=\dfrac{1}{4}\p{\norme{x}^2-\norme{y}^2}=0$ par hypothèse.
% Ainsi, on a $$x=\underbrace{\dfrac{x+y}{2}}_{\in H}+\underbrace{\dfrac{x-y}{2}}_{\in H^\perp}$$ donc $$\rho(x)=\dfrac{x+y}{2}-\dfrac{x-y}{2}=y.$$ On obtient de même $\rho(y)=x$.
% \end{itemize}
% \end{preuve}

\subsection{Projecteur orthogonal}

Dans cette section, et dans cette section seulement, $E$ désigne un espace préhilbertien.\\

\begin{definition}
Soit $F$ un sous-espace vectoriel de dimension finie de $E$. Alors
\[E=F\oplus F^\perp.\]
De plus, $F^\perp$ est l'unique supplémentaire de $F$ orthogonal à $F$; on l'appelle le \emph{supplémentaire orthogonal}
de $F$.
\end{definition}


\begin{definition}[utile=-3]
Soit $F$ un sous-espace vectoriel de dimension finie de $E$. On appelle \emph{projecteur orthogonal} sur
$F$, le projecteur sur $F$ parallèlement à $F^\perp$.
\end{definition}

% \begin{remarqueUnique}
% \remarque Soit $p\in\Endo{E}$ un projecteur sur $A$ parallèlement à $B$. Pour
%   montrer que $p$ est un projecteur orthogonal, il suffit de vérifier que $A$
%   et $B$ sont orthogonaux.
% \end{remarqueUnique}

\begin{proposition}[utile=-3]
  Soit $p$ un projecteur orthogonal. Alors
  \[\forall x\in E \qsep \norme{p(x)}\leq\norme{x}.\]
  \end{proposition}
  
  \begin{preuve}
  En se munissant de $F$ tel que $p$ est le projecteur orthogonal sur $F$, en décomposant $x$ sur $F\oplus F^\perp$ en $x=p_F(x)+(x-p_F(x))$, on a $$\norme{x}^2=\norme{p_F(x)+(x-p_F(x))}^2=\norme{p_F(x)}^2+\norme{x-p_F(x)}^2\geq \norme{p_F(x)}^2$$ où on a utilisé Pythagore.
  \end{preuve}
  
  \begin{sol}
  Faire un dessin pour expliquer que si $F$ n'est pas orthogonal, ce n'est pas vrai.
  \end{sol}
  
  \begin{proposition}[utile=3]
  Soit $F$ un sous-espace vectoriel de dimension finie de $E$, $(e_1,\ldots,e_r)$ une base orthonormée
  de $F$ et $p$ le projecteur orthogonal sur $F$. Alors
  \[\forall x\in E \qsep p(x)=\sum_{k=1}^r \ps{e_k}{x} e_k.\]
  \end{proposition}
  
  \begin{preuve}
  On décompose $x$ sur une base adaptée au problème.
  \end{preuve}
  
  \begin{remarques}
  \remarque Si la base $(e_1,\ldots,e_r)$ de $F$ est orthogonale, alors
    \[\forall x\in E\qsep p(x)=\sum_{k=1}^r \frac{\ps{e_k}{x}}{\norme{e_k}^2}e_k.\]
    De plus, si $q$ est le projecteur sur $F^\perp$ parallèlement à $F$, alors $p+q=\id$ donc 
    \[\forall x\in E\qsep q(x)=x-\sum_{k=1}^r \frac{\ps{e_k}{x}}{\norme{e_k}^2}e_k.\]
  \remarque En particulier, si $E$ est un espace euclidien, $u$ est un vecteur non nul de $E$ et
    $p$ est le projecteur orthogonal sur $\vect(u)^\perp$, alors
    \[\forall x\in E\qsep p(x)=x-\frac{\ps{u}{x}}{\norme{u}^2}u.\]
  \end{remarques}
  
  \begin{exoUnique}
  \exo Déterminer la matrice dans la base canonique de $\R^4$ de la
    projection orthogonale sur le sous-espace vectoriel $F$ d'équation
    \[\syslin{x&+y&+z&+t&=&0\hfill\cr
              x&-y&+z&-t&=&0.}\]
    \begin{sol}
    En posant $u_1=(1,1,1,1)$ et $u_2=(1,-1,1,-1)$, on peut voir que $F=\set{u_1,u_2}^\perp=\p{\vect(u_1,u_2)}^\perp$. On orthonormalise facilement la famille $u_1, u_2$ (elle est déjà orthogonale) en $(b_1,b_2)$ qui est donc une BON de $\vect(u_1,u_2)=F^\perp$. On peut donc avoir facilement l'expression du projecteur sur $F^\perp$. Mais alors, utilisons qu'avec $p$ le projecteur orthogonal sur $F$ et $q$ le projecteur orthogonal sur $F^\perp$, on a $p+q=\id$. On obtient facilement la matrice de $q$ dans la base canonique et l'identité moins cette-dernière est bien la matrice de $p$ dans la base canonique et vaut
    \[P=\frac{1}{2}
      \begin{pmatrix}
      1 & 0 & -1 & 0\\
      0 & 1 & 0 & -1\\
      -1 & 0 & 1 & 0\\
      0 & -1 & 0 & 1
      \end{pmatrix}\]
    \end{sol}
  \end{exoUnique}


% \begin{remarqueUnique}
% \remarque Si $F$ est un sous-espace vectoriel de dimension finie d'un espace préhilbertien $E$, on peut
%   donc définir le projecteur orthogonal $p$ sur $F$ comme le projecteur sur $F$ parallèlement à $F^\perp$.
% \end{remarqueUnique}


% \begin{proposition}
%   Soit $F$ un sous-espace vectoriel de dimension finie d'un espace préhilbertien $E$, $(e_1,\ldots,e_r)$ une base orthonormée
%   de $F$ et $p$ le projecteur orthogonal sur $F$. Alors
%   \[\forall x\in E \qsep p(x)=\sum_{k=1}^r \ps{e_k}{x} e_k.\]
% \end{proposition}

\begin{definition}[utile=-3]
Soit $A$ une partie non vide de  $E$ et $x\in E$. On définit
\emph{la distance de $x$ à $A$}, que l'on note ${\rm d}(x,A)$, par
\[{\rm d}(x,A)\defeq\inf_{a\in A} \norme{x-a}.\]
\end{definition}

\begin{preuve}
Le caractère non vide de $A$ est primordiale ici et justifie la définition.\\
Parlons avec les élèves du caractère atteint ou non atteint...
\end{preuve}

\begin{proposition}[utile=3]
Soit $F$ un sous-espace vectoriel de dimension finie de $E$ et $x\in E$. Alors, la borne
inférieure
\[{\rm d}(x,F)=\inf_{f\in F} \norme{x-f}\]
est un minimum qui est atteint en un unique point. Ce point est le projeté orthogonal de $x$ sur $F$.  
\end{proposition}

\begin{preuve}
$\forall f \in F$, $$\norme{x-f}^2=\norme{\underbrace{(x-p_F(x))}_{\in F^\perp}+\underbrace{\p{p_F(x)-f}}_{\in F}}^2=\norme{x-p_F(x)}^2+\norme{p_F(x)-f}^2\geq \norme{x-p_F(x)}^2$$ d'où le résultat et on a aussi l'unicité car on a égalité ssi $\norme{p_F(x)-f}^2=0$.
\end{preuve}

\begin{exoUnique}
\exo Montrer que la borne inférieure
  \[\inf_{(a,b)\in\R^2} \integ{-1}{1}{\cro{\e^x-\p{ax+b}}^2}{x}\]
  est atteinte en un unique couple $(a,b)\in\R^2$ que l'on calculera.
  \begin{sol}
  On définit $\dspappli{g}{\interf{0}{1}}{\R}{x}{e^x}$, $\dspappli{f_1}{\interf{0}{1}}{\R}{x}{1}$ et $\dspappli{f_2}{\interf{0}{1}}{\R}{x}{x}$. On considère alors $E=\vect(f_1,f_2,g)$ muni du produit scalaire usuel et $F=\vect(f_1,f_2)$. La question est de calculer la distance de $g$ à $F$ et nous allons pouvoir pour cela déterminer la projection orthogonal de $g$ sur $F$ puisqu'on sait que c'est là où est atteint notre minimum et que celui vaut $\norme{g-p_F(g)}^2$. FAIRE UN DESSIN POUR EXPLIQUER TOUT CELA\\
  Pour chercher $p_F(g)$ on a besoin d'une BON de $F$. Commençons par orthogonaliser $f_1,f_2$. On pose $u_1=f_1$ puis on cherche $u_2$ de la forme $u_2=f_2+\lambda f_1$ de sorte que $\ps{u_2}{u_1}=0$. On trouve $\lambda=-1/2$. Il reste alors à normaliser cette base. On trouve alors $b_1$ et $b_2$ définies par :
  \[b_1(x)=1 \et b_2(x)=2\sqrt{3}\p{x-1/2}\]
  On détermine alors $p_F(g)=\ps{g}{b_1}b_1+\ps{g}{b_2}b_2$ puis $g-p_F(g)$ dont on calcule la norme au carrée. On trouve $p_F(g)(x)=6(3-e)x+4e-10$ et pour la distance ... $\dfrac{1}{2}(e-3)(7e-19)$.
  \end{sol}  
\end{exoUnique}

\subsection{Algorithme d'orthonormalisation de \nom{Gram-Schmidt}}

Dans la suite de ce chapitre, $E$ désigne de nouveau un espace euclidien.\\

\begin{proposition}[nom={Algorithme d'orthogonalisation de \nom{Gram-Schmidt}}]
Soit $\mathcal{B}\defeq(e_1,\ldots,e_n)$ une base de $E$. Alors il existe une base
orthogonale $(f_1,\ldots,f_n)$ de $E$ telle que pour tout
$k\in\intere{0}{n}$,
$(f_1,\ldots,f_k)$ est une base orthogonale de $E_k\defeq\vect(e_1,\ldots,e_k)$.
\end{proposition}

\begin{preuve}
\begin{itemize}
\item[$\bullet$] Existence : On va faire une preuve constructive.\\
Pour $k\in \intere{1}{n}$, posons $F_k=\vect(e_1,\ldots,e_k)$.
\begin{itemize}
\item On commence par poser $f_1=\dfrac{1}{\norme{e_1}}e_1$ qui est bien de norme $1$. De plus, $\vect(f_1)=F_1$.
\item Posons $u_2=p_{F_1^\perp}(e_2)=e_2-\ps{e_2}{f_1}f_1$. On a $u_2\neq 0$ car $(e_1,e_2)$ est libre. On a $u_2$ orthogonal à $f_1$ par définition. Donc $(u_2,f_1)$ est libre. Ainsi, $\vect(u_2,f_1)=F_2$ par inclusion plus égalité des dimensions. Il nous reste à normaliser $u_2$ en un vecteur $f_2$ ce qui ne change rien au vect.
\item Soit $k\in \intere{1}{n-1}$. Supposons construit $(u_1,\ldots,u_k)$ une FON vérifiant $\vect(u_1,\ldots,u_k)=F_k$.\\
Posons $$u_{k+1}=p_{F_k^\perp}(e_{k+1})=e_{k+1}-\sum_{i=1}^k\dfrac{\ps{e_{k+1}}{f_i}}{f_i}.$$
$u_{k+1}$ est orthogonal à $f_1,\ldots,f_k$ par définition. De plus, $u_{k+1}\neq 0$ car sinon $e_{k+1}\in \vect(u_1,\ldots,u_k)=F_k$ par hypothèse, ce qui est faux. Ainsi, $(f_1,\ldots,f_k,u_{k+1})$ est une FOG de $F_{k+1}$ (sans vecteurs nuls) donc une famille libre donc une base (cardinal égal à la dimension) de $E_{k+1}$. Il nous reste à normaliser $u_{k+1}$ en $f_{k+1}$ ce qui ne change rien à toutes les propriétés énoncées. 
\end{itemize}
On a ainsi construit une BON $B'=(f_1,\ldots, f_n)$ et on observe finalement que la matrice de $B'$ dans $B$ est bien comme demandée.
\item[$\bullet$] Pour l'unicité (on peut la zapper), si on a $B_1$, $B_2$ deux BON telles que $P(B,B_1)$ et $P(B,B_2)$ soient triangulaires supérieures à coefficients diagonaux strictement positifs, alors $$A:=P(B_1,B_2)=P(B_1,B)P(B,B_2)=P(B,B_1)^{-1}P(B,B_2)$$ est aussi triangulaire supérieure à coefficients diagonaux strictement positifs. De plus (cf. plus tard), comme on a passe d'une BON à une BON, la famille des colonnes est une BON. Et ce $\ps{c_i}{c_j}=\delta_{i,j}$ combiné au fait que c'est triangulaire supérieure à coefficients positifs nous permet de montrer en commençant par la gauche et en allant de proche en proche que $A=I_n$ ce qui signifie que $B_1=B_2$.
\end{itemize}
\end{preuve}

\begin{remarques}
% \remarque L'\emph{algorithme d'orthonormalisation de \nom{Gram-Schmidt}} permet de construire
%   une famille $(f_1,\ldots,f_n)$ d'éléments de $E$ telle que pour tout
%   $k\in\intere{0}{n}$, $(f_1,\ldots,f_k)$ est une base orthonormée de $E_k\defeq \vect(e_1,\ldots,e_k)$.
%   Sa construction se fait de la manière suivante.
%   \begin{itemize}
%   \item On pose
%     \[f_1\defeq \frac{1}{\norme{e_1}}e_1.\]
%     La famille $(f_1)$ est alors une base orthonormée de $E_1$.
%   \item Soit $k\in\intere{1}{n-1}$. On suppose qu'on a construit une base
%     orthonormée $(f_1,\ldots,f_k)$ de $E_k$. On définit $p_{k+1}$ comme étant le projeté
%     orthogonal de $e_{k+1}$ sur $E_k$. Puisque $(f_1,\ldots,f_k)$ est une base orthonormée
%     de $E_k$, on a
%     \[p_{k+1}= \sum_{i=1}^k \ps{f_i}{e_{k+1}} f_i.\]
%     On note $g_{k+1}\defeq e_{k+1} - p_{k+1}$ le projeté orthogonal de $e_{k+1}$ sur
%     $E_k^\perp$. On a alors
%     \[g_{k+1}=e_{k+1} - \sum_{i=1}^k \ps{f_i}{e_{k+1}} f_i.\]
%     Il ne reste plus qu'à normer ce vecteur en posant
%     \[f_{k+1}\defeq\frac{1}{\norme{g_{k+1}}}g_{k+1}.\]
%     On a ainsi une base orthonormée $(f_1,\ldots,f_{k+1})$ de $E_{k+1}$.
%   \end{itemize}
\remarque Une telle base est donnée par l'algorithme d'\emph{orthogonalisation de \nom{Gram-Schmidt}}.
  \begin{itemize}
  \item On pose $f_1\defeq e_1$. La famille $(f_1)$ est alors une base orthogonale de $E_1$.
  \item Soit $k\in\intere{1}{n-1}$. On suppose qu'on a construit une base
    orthogonale $(f_1,\ldots,f_k)$ de $E_k$. On définit $p_{k+1}$ comme le projeté
    orthogonal de $e_{k+1}$ sur $E_k$. Puisque $(f_1,\ldots,f_k)$ est une base orthogonale
    de $E_k$, on a
    \[p_{k+1}= \sum_{i=1}^k \frac{\ps{f_i}{e_{k+1}}}{\norme{f_i}^2} f_i.\]
    On note $f_{k+1}\defeq e_{k+1} - p_{k+1}$ le projeté orthogonal de $e_{k+1}$ sur
    $E_k^\perp$. On a alors
    \[f_{k+1}=e_{k+1} - \sum_{i=1}^k \frac{\ps{f_i}{e_{k+1}}}{\norme{f_i}^2} f_i.\]
    On a ainsi construit une base orthogonale $(f_1,\ldots,f_{k+1})$ de $E_{k+1}$.
  \end{itemize}
  La base orthogonale $(f_1,\ldots,f_n)$ ainsi construite 
  est solution de notre problème.
\remarque Dans la proposition précédente, il n'y a pas unicité d'une telle base
  $(f_1,\ldots,f_n)$. Cependant, une famille $(g_1,\ldots,g_n)$ est une autre solution
  de notre problème si et seulement si il existe $\lambda_1,\ldots,\lambda_n\in\Rs$ tels
  que
  \[\forall k\in\intere{1}{n}\qsep g_k=\lambda_k f_k.\]
\remarque L'algorithme d'orthogonalisation de Gram-Schmidt peut s'utiliser pour une
  famille $(e_1,\ldots,e_n)$ quelconque d'un espace euclidien de dimension $n$.
  \begin{itemize}
  \item Si l'un des vecteurs $f_k$ est nul, alors $e_k \in\vect(e_1,\ldots,e_{k-1})$
    et l'algorithme s'arrête.
  \item Sinon, l'algorithme construit une base orthogonale de $E$ tout en prouvant
    la liberté de $(e_1,\ldots,e_n)$.
  \end{itemize}
% \remarque Remarquons que si $(f_1,\ldots,f_n)$ est une base orthonormée telle que
%   \[\forall k\in\intere{0}{n}\qsep \vect(f_1,\ldots,f_k)=E_k\]
%   alors on a
%   \[\forall k\in\intere{1}{n}\qsep \exists\lambda\in\Rs\qsep \exists x\in E_{k-1}\qsep f_k = \lambda e_k + x.\]
%   La seconde condition de la proposition est donc uniquement présente pour imposer
%   à $\lambda$ d'être strictement positif et assurer l'unicité de la base $(f_1,\ldots,f_n)$.
%   Remarquons que si $k\in\intere{1}{n}$ et $\lambda\in\Rs, x\in E_{k-1}$
%   sont tels que $f_k=\lambda e_k + x$, en effectuant le produit scalaire par
%   $f_k$, on obtient
%   \[\lambda=\frac{1}{\ps{f_k}{e_k}}.\]
%   C'est la raison pour laquelle la proposition précédente est parfois énoncée avec 
%   la condition $\ps{f_k}{e_k}>0$ comme condition d'unicité.
% \remarque Si $(f_1,\ldots,f_n)$ est la base obtenue par le procédé
%   d'othonormalisation de \nom{Gram-Schmidt}, alors une famille $(f_1',\ldots,f_n')$
%   est une base orthonornée de $E$ vérifiant
%   \[\forall k\in\intere{0}{n}\qsep \vect(f_1',\ldots,f_k')=E_k\]
%   si et seulement si il existe $\epsilon_1,\ldots,\epsilon_n\in\ens{-1,1}$ tels que
%   \[\forall k\in\intere{1}{n}\qsep f_k'=\epsilon_k f_k.\]
  % Dans la proposition précédente, c'est donc le second point qui assure l'unicité.
% \remarque Il est possible de remplacer le second point par la condition
%   \[\forall k\in\intere{1}{n}\qsep \ps{f_k}{e_k}>0.\]
%   On a toujours unicité et la base ainsi construite est la même que dans la proposition
%   d'origine.
% \remarque La matrice de passage de $\mathcal{B}$ à $\mathcal{B}'$ est une matrice
%   triangulaire supérieure dont les coefficients diagonaux sont strictement positifs.
% \remarque Si $\mathcal{B}=(e_1,\ldots,e_n)$ est une base de $E$, on construit la
%   base $\mathcal{B}'=(f_1,\ldots,f_n)$ de la manière suivante.
%   \begin{itemize}
%   \item On pose $f_1\defeq e_1/\norme{e_1}$.
%   \item On suppose que l'on a construit $(f_1,\ldots,f_k)$. Soit $c_{k+1}$ le
%     projeté orthogonal de $e_{k+1}$ sur $\vect\p{f_1,\ldots,f_k}$. Comme
%     $(f_1,\ldots,f_k)$ est une base orthonormée de $\vect\p{f_1,\ldots,f_k}$, on
%     a
%     \[c_{k+1}=\sum_{l=1}^k \ps{f_l}{e_{k+1}}f_l\]
%     On pose ensuite $f_{k+1}'\defeq e_{k+1}-c_{k+1}$, puis
%     $f_{k+1}\defeq f_{k+1}'/\norme{f_{k+1}'}$.
%   \end{itemize}
% \remarque Il arrive que la construction d'une base orthogonale suffise. On
%   utilise alors la variante suivante du théorème précédent~: \og 
%   Soit $\mathcal{B}=(e_1,\ldots,e_n)$ une base de $E$. Alors, il existe une unique
%   base orthogonale $\mathcal{B}''=(g_1,\ldots,g_n)$ de $E$ telle que
%   $P\p{\mathcal{B},\mathcal{B}''}$ soit une matrice triangulaire supérieure
%   dont les coefficients diagonaux sont égaux à 1. On dit que $\mathcal{B}''$ est
%   la base orthogonale de $E$ obtenue à partir de $\mathcal{B}$ par le procédé
%   d'orthogonalisation de Schmidt. \fg
% \remarque
%   Si $\mathcal{B}=(e_1,\ldots,e_n)$ est une base de $E$, on construit la
%   base orthogonale $\mathcal{B}''=(g_1,\ldots,g_n)$ de la manière suivante.
%   \begin{itemize}
%   \item On pose $g_1\defeq e_1$.
%   \item On suppose que l'on a construit $(g_1,\ldots,g_k)$. Soit $c_{k+1}$ le
%     projeté orthogonal de $e_{k+1}$ sur $\vect\p{g_1,\ldots,g_k}$. Comme
%     $(g_1,\ldots,g_k)$ est une base orthogonale de $\vect\p{g_1,\ldots,g_k}$, on
%     a
%     \[c_{k+1}=\sum_{l=1}^k \frac{\ps{g_l}{e_{k+1}}}{\norme{g_l}^2}g_l.\]
%     On pose ensuite $g_{k+1}\defeq e_{k+1}-c_{k+1}$.
%   \end{itemize}
%   Une fois la base orthogonale $\mathcal{B}''=(g_1,\ldots,g_n)$ obtenue, on définit
%   la base $\mathcal{B}'=(f_1,\ldots,f_n)$ par
%   \[\forall k\in\intere{1}{n}\qsep f_k\defeq \frac{g_k}{\norme{g_k}}.\]
\end{remarques}

\begin{exoUnique}
\exo On pose $E\defeq\polyR[n]$, que l'on munit du produit scalaire
  \[\forall P,Q\in\polyR[n]\qsep \ps{P}{Q}\defeq\integ{-1}{1}{P(x)Q(x)}{x}.\]
  Montrer qu'il existe une unique base orthogonale de $E$, formée de
  polynômes $P_0,\ldots,P_n$ de coefficients dominants égaux à 1, tels que
  \[\forall k\in\intere{0}{n}\qsep \deg P_k=k.\]
  Calculer cette base pour $n=3$.
\end{exoUnique}

\begin{proposition}[nom={Algorithme d'orthonormalisation de \nom{Gram-Schmidt}}]
Soit $\mathcal{B}\defeq(e_1,\ldots,e_n)$ une base de $E$. Alors il existe une base
orthonormée $(g_1,\ldots,g_n)$ de $E$ telle que pour tout
$k\in\intere{0}{n}$,
$(g_1,\ldots,g_k)$ est une base orthonormée de $E_k\defeq\vect(e_1,\ldots,e_k)$.
\end{proposition}

\begin{remarques}
\remarque Si l'on souhaite obtenir une telle base,
  il suffit de normer les vecteurs de la famille $(f_1,\ldots,f_n)$
  obtenue par l'algorithme d'orthogonalisation de \nom{Gram-Schmidt} en posant
  \[\forall k\in\intere{1}{n}\qsep g_k=\frac{1}{\norme{f_k}}f_k.\]
  On dit qu'on a obtenu la base $(g_1,\ldots,g_n)$ par l'algorithme
  d'\emph{orthonormalisation de \nom{Gram-Schmidt}}.
\remarque Dans la proposition précédente, il n'y a pas unicité d'une telle base
  $(g_1,\ldots,g_n)$. Cependant, une famille $(h_1,\ldots,h_n)$ est une autre solution
  de notre problème si et seulement si il existe $\epsilon_1,\ldots,\epsilon_n\in\ens{-1,1}$ tels
  que
  \[\forall k\in\intere{1}{n}\qsep h_k=\epsilon_k g_k.\]
\remarque Pour avoir l'unicité, il suffit de rajouter la condition
  \[\forall k\in\intere{1}{n}\qsep \ps{e_k}{g_k}>0.\]
  L'unique famille solution de ce problème est la famille obtenue par le procédé
  d'orthonormalisation de \nom{Gram-Schmidt}.
\end{remarques}


\begin{exoUnique}
\exo Dans $\R^3$ muni du produit scalaire usuel, orthonormaliser la
  famille
  \[e_1\defeq\p{1,-2,2}, \qquad e_2\defeq\p{-1,0,-1}, \qquad e_3\defeq\p{5,-3,7}.\]
  \begin{sol}
  On trouve
  \[b_1=\frac{1}{3}\p{1,-2,2}, \quad b_2=\frac{1}{3}\p{-2,-2,-1}, \quad
    b_3=\frac{1}{3}\p{-2,1,2}.\]
  \end{sol}
\end{exoUnique}

% \begin{remarques}

% \item Si $(f_1,\ldots,f_n)$ est une base orthonormée de $E$ telle que 
% \[\forall k\in\intere{0}{n}\qsep \vect(f_1,\ldots,f_k)=\vect(e_1,\ldots,e_k)\]
%   alors, la famille $(g_1,\ldots,g_n)$ est une base orthonormée de $E$ telle que
%   $\vect(g_1,\ldots,g_k)=\vect(e_1,\ldots,e_k)$ pour tout $k\in\intere{0}{n}$, si et
%   seulement si il existe $\epsilon_1,\ldots,\epsilon_n\in\ens{-1,1}$ tels que
%   \[\forall k\in\intere{1}{n}\qsep g_k=\epsilon_k f_k.\]
% \end{itemize}
% \remarque On peut ainsi démontrer le théorème de Gram-Schmidt qui s'énonce de la
%   manière suivante.
%   Soit $\mathcal{B}\defeq(e_1,\ldots,e_n)$ une base de $E$. Alors, il
%   existe une unique base orthonormée $\mathcal{B}'\defeq(f_1,\ldots,f_n)$ de $E$ telle que~:
%   \begin{itemize}
% \item $\forall k\in\intere{0}{n}\qsep \vect(f_1,\ldots,f_k)=\vect(e_1,\ldots,e_k)$.
% \item $\forall k\in\intere{1}{n}\qsep \ps{e_k}{f_k}>0$.
%   \end{itemize}
%   C'est la base donnée par l'algorithme d'orthonormalisation de \nom{Gram-Schmidt}.
% \end{remarques}

% \begin{exos}
% \exo Soit $E$ le \Rev des fonctions réelles continues sur $\interf{0}{1}$
%   muni du produit scalaire $\ps{.}{.}$ défini par
%   \[\forall f,g\in E \quad \ps{f}{g}=\integ{0}{1}{f(x)g(x)}{x}\]
%   Soit $n\in\Ns$, $A_n$ le sous-espace vectoriel de $E$ des fonctions
%   polynôme de degré inférieur ou égal à $n$ et $f\in E$ un vecteur
%   orthogonal à $A_n$. Monter que $f$ s'annule au moins $n+1$ fois sur
%   $\intero{0}{1}$.
% \end{exos}

% \begin{proposition}
% Soit $\mathcal{B}\defeq(e_1,\ldots,e_n)$ une base de $E$. Alors il existe une unique base
% orthonormée $(f_1,\ldots,f_n)$ de $E$ telle que pour tout $k\in\intere{0}{n}$,
% $(f_1,\ldots,f_k)$ est une base orthonormée de $E_k\defeq\vect(e_1,\ldots,e_k)$ et
% $\ps{e_k}{f_k}>0$. Une telle base est donnée par l'algorithme
% d'\emph{orthonormalisation de \nom{Gram-Schmidt}}.
% \end{proposition}

% \begin{remarqueUnique}
% \remarque 
% \end{remarqueUnique}


\subsection{Dual}

\begin{proposition}[utile=3]
  Pour toute forme linéaire $\phi$ sur $E$, il existe un unique $a\in E$ tel que
  \[\forall x\in E \qsep \phi(x)=\ps{a}{x}.\]
  \end{proposition}

% Soit $\mathcal{B}\defeq\p{e_1,\ldots,e_n}$ une base orthonormée de $E$ et
% $\p{a_1,\ldots,a_n}\in\R^n$ un $n$-uplet non nul. On pose
% \[a\defeq \sum_{k=1}^n a_k e_k.\]
% Alors, si $H$ est une partie de $E$, les trois propositions suivantes sont équivalentes~:
% \begin{itemize}
% \item $H$ est le noyau de la forme linéaire $x\mapsto\ps{a}{x}$.
% \item $H$ admet pour équation dans $\mathcal{B}$
%   \[\sum_{k=1}^n a_k x_k=0.\]
% \item $a$ est un vecteur orthogonal à $H$.
% \end{itemize}
% Si tel est le cas, $H$ est un hyperplan de $E$.

% \begin{proposition}[utile=3]
%   Pour toute forme linéaire $\phi$ sur $E$, il existe un unique $a\in E$ tel que
%   \[\forall x\in E \qsep \phi(x)=\ps{a}{x}.\]
%   \end{proposition}

% \begin{preuve}
% On montre que $$\dspappli{\psi}{E}{E*}{a}{\dspappli{\phi_a}{E}{\R}{x}{\ps{a}{x}}}$$ est bien définie, linéaire, injective. Donc c'est un isomorphisme par égalité des dimensions, ce qui correspond au résultat souhaité.
% \end{preuve}

% \begin{proposition}[utile=-3]
% Soit $\mathcal{B}\defeq\p{e_1,\ldots,e_n}$ une base orthonormée de $E$ et
% $\p{a_1,\ldots,a_n}\in\R^n$ un $n$-uplet non nul. On pose
% \[a\defeq \sum_{k=1}^n a_k e_k.\]
% Alors, si $H$ est une partie de $E$, les trois propositions suivantes sont équivalentes~:
% \begin{itemize}
% \item $H$ est le noyau de la forme linéaire $x\mapsto\ps{a}{x}$.
% \item $H$ admet pour équation dans $\mathcal{B}$
%   \[\sum_{k=1}^n a_k x_k=0.\]
% \item $a$ est un vecteur orthogonal à $H$.
% \end{itemize}
% Si tel est le cas, $H$ est un hyperplan de $E$.
% \end{proposition}

% \begin{preuve}
% On montre $(i)\Rightarrow (ii)\Rightarrow (iii) \Rightarrow (i)$. Pour $(iii) \Rightarrow (i)$, on montre l'inclusion et l'égalité des dimensions, le noyau de la forme linéaire étant de dimension $1$* puisque $a\neq 0$ nous assure que la forme linéaire est non nulle.
% \end{preuve}
%END_BOOK
\end{document}

%%%%%%%%%%%%%%%%%%%%%%%%%%%%%%%%%%%%%%%%%%%%
%%% Démonstration du procédé de Gram-Schmidt
\begin{preuve}
Montrons ce résultat par récurrence sur la dimension $n$ de $E$.
\begin{center}
$\mathcal{P}_n$ : \flqq\ 
\parbox[t]{0.75\linewidth}{%
Soit $E$ un espace euclidien de dimension $n$ et $e_1,\ldots,e_n$ une base de
$E$. Alors, il existe une unique base orthonormée $b_1,\ldots,b_n$ de $E$ telle
que quelque soit $k\in\intere{1}{n}$, $b_1,\ldots,b_k$ est une base de
$\vect\p{e_1,\ldots,e_k}$ et $\ps{e_k}{b_k}\geq 0$\ \frqq}
\end{center}
\begin{itemize}
\item $\mathcal{P}_0$ est vraie. En effet, si $E$ est un espace euclidien de
  dimension nulle, alors $E=\ens{0}$ et l'unique base de $E$ est la famille
  ne contenant aucun élément. Cette base est orthonormée et vérifie bien les
  conditions demandées (puisqu'il n'y a rien à vérifier).
\item $\mathcal{P}_n \implique \mathcal{P}_{n+1}$. On suppose que
  $\mathcal{P}_n$ est vraie et on souhaite montrer que $\mathcal{P}_{n+1}$ est
  vraie. Soit $E$ un espace euclidien de dimension $n+1$ et $e_1,\ldots,e_{n+1}$
  une base de $E$.
  \begin{itemize}
  \item {\it unicité}\\
    Soit $b_1,\ldots,b_{n+1}$ et $b_1',\ldots,b_{n+1}'$ deux bases orthonormées
    de $E$ répondant au problème posé. Montrons que ces deux bases sont
    identiques.\\
    Soit $F=\vect\p{e_1,\ldots,e_n}$. Alors $F$ est un espace euclidien de
    dimension $n$ et $b_1,\ldots,b_n$ et $b_1',\ldots,b_n'$ sont deux bases
    orthonormées de $F$ répondant au problème posé. Puisque $\mathcal{P}_n$
    est vraie, on en déduit que~:
    \[\forall k\in\intere{1}{n} \quad b_k=b_k'\]
    Montrons désormais que $b_{n+1}=b_{n+1}'$. Puisque
    $b_{n+1}\in\vect\p{e_1,\ldots,e_{n+1}}$ et que
    $\vect\p{e_1,\ldots,e_n}=\vect\p{b_1,\ldots,b_n}$, on en déduit que
    $b_{n+1}\in\vect\p{b_1,\ldots,b_n,e_{n+1}}$. Il existe donc
    $\lambda_1,\ldots,\lambda_n\in\R$ et $\lambda\in\R$ tels que~:
    \[b_{n+1}=\sum_{k=1}^n \lambda_k b_k + \lambda e_{n+1}\]
    De plus, $b_{n+1}\not\in\vect\p{b_1,\ldots,b_n}$ car sinon, la famille
    $b_1,\ldots,b_{n+1}$ ne serait pas libre; donc $\lambda\neq 0$. On définit
    alors pour tout $k\in\intere{1}{n}$, $\mu_k=\lambda_k/\lambda$. On a donc~:
    \[b_{n+1}=\lambda\p{\sum_{k=1}^n \mu_k b_k + e_{n+1}}\]
    Puisque $b_{n+1}$ est orthogonal à tous les $b_k$ pour $k\in\intere{1}{n}$,
    on en déduit que~:
    \begin{eqnarray*}
    \forall k\in\intere{1}{n} \quad
               & & \ps{b_k}{b_{n+1}}=0\\
    \text{donc}& & \ps{b_k}{\lambda\p{\sum_{i=1}^n \mu_i b_i + e_{n+1}}}=0\\
    \text{donc}& & \lambda\ps{b_k}{\sum_{i=1}^n \mu_i b_i + e_{n+1}}=0\\
    \text{donc}& & \ps{b_k}{\sum_{i=1}^n \mu_i b_i + e_{n+1}}=0 \quad
                   \text{car $\lambda\neq 0$}\\
    \text{donc}& & \sum_{i=1}^n \mu_i\underbrace{\ps{b_k}{b_i}}_{\delta_{k,i}} +
                   \ps{b_k}{e_{n+1}}=0\\
    \text{donc}& & \mu_k+\ps{b_k}{e_{n+1}}=0\\
    \text{donc}& & \mu_k=-\ps{b_k}{e_{n+1}}
    \end{eqnarray*}
    De plus, $b_{n+1}$ est normé donc~:
    \[\abs{\lambda}\norme{\sum_{k=1}^n \mu_k b_k + e_{n+1}}=1\]
    Donc~:
    \[\abs{\lambda}=\frac{1}{\norme{\sum_{k=1}^n \mu_k b_k + e_{n+1}}}\]
    Enfin, $\ps{e_{n+1}}{b_{n+1}}\geq 0$ donc~:
    \begin{eqnarray*}
    \ps{\frac{1}{\lambda}\p{b_{n+1}-\sum_{k=1}^n \lambda_k b_k}}{b_{n+1}}
    &=& \frac{1}{\lambda}\p{\norme{b_{n+1}}^2-
                    \underbrace{\ps{\sum_{k=1}^n \lambda_k b_k}{b_{n+1}}}_{=0}}\\
    &=& \frac{1}{\lambda}\geq 0
    \end{eqnarray*}
    Donc $\lambda\geq 0$. On en déduit que~:
    \[\lambda=\frac{1}{\norme{\sum_{k=1}^n \mu_k b_k + e_{n+1}}}\]
    On a donc démontré que les $\mu_k$ et $\lambda$ sont entièrement
    déterminés les vecteurs $b_1,\ldots,b_n,e_{n+1}$. Puisque $b_k=b_k'$ pour
    tout $k\in\intere{1}{n}$, on en déduit que $b_{n+1}=b_{n+1}'$. D'où l'unicité
    de la base $b_1,\ldots,b_{n+1}$.
  \item {\it existence}\\
    Soit $F=\vect\p{e_1,\ldots,e_n}$. Puisque $F$ est un espace euclidien de
    dimension $n$, il existe une base orthonormée $b_1,\ldots,b_n$ de
    $F$ telle que pour tout $k\in\intere{1}{n}$, $b_1,\ldots,b_k$ est une base de
    $\vect\p{e_1,\ldots,e_k}$ et $\ps{e_k}{b_k}\geq 0$.\\
    Soit $\mu_1,\ldots,\mu_n\in\R$ et $u$ le vecteur~:
    \[u=\sum_{k=1}^n \mu_kb_k + e_{n+1}\]
    Soit $k\in\intere{1}{n}$. Alors $u$ est orthogonal à $b_k$ si et seulement
    si~:
    \begin{eqnarray*}
    \ps{b_k}{u}=0
    &\ssi& \ps{b_k}{\sum_{i=1}^n \mu_ib_i + e_{n+1}}=0\\
    &\ssi& \sum_{i=1}^n \mu_i \underbrace{\ps{b_k}{b_i}}_{\delta_{k,i}}+
            \ps{b_k}{e_{n+1}}\\
    &\ssi& \mu_k=-\ps{b_k}{e_{n+1}}
    \end{eqnarray*}
    On choisit de tels $\mu_k$ et on pose enfin~:
    \[b_{n+1}=\frac{u}{\norme{u}}\]
    Alors $b_{n+1}$ est normé et est orthogonal à  la famille $b_1,\ldots,b_n$.
    La famille $b_1,\ldots,b_{n+1}$ est donc une famille orthonormée; elle
    est donc libre. Puisqu'elle possède $n+1$ éléments, c'est une base de $E$.
    Enfin~:
    \begin{eqnarray*}
    \ps{e_{n+1}}{b_{n+1}}
    &=& \frac{1}{\norme{u}}\ps{e_{n+1}}{u}\\
    &=& \frac{1}{\norme{u}}\ps{u-\sum_{k=1}^n \mu_k b_k}{u}\\
    &=& \frac{1}{\norme{u}}\p{\norme{u}^2-
           \underbrace{\ps{\sum_{k=1}^n \mu_k b_k}{u}}_{=0}}\\
    &=& \norme{u}\geq 0
    \end{eqnarray*}
    D'où l'existence d'une telle base.
  \end{itemize}
  En conclusion $\mathcal{P}_{n+1}$ est vraie.
\end{itemize}
Par récurrence sur $n$, on en déduit que $\mathcal{P}_n$ est vraie pour tout
$n\in\N$.
\end{preuve}

\section{Exercices}
\setcounter{numeroexercice}{1}
\documentclass{magnolia}

\magtex{tex_driver={pdftex}}
\magfiche{document_nom={Exercices sur les espaces euclidiens},
          auteur_nom={François Fayard},
          auteur_mail={fayard.prof@gmail.com}}
\magexos{exos_matiere={maths},
         exos_niveau={mpsi},
         exos_chapitre_numero={22},
         exos_theme={Espaces euclidiens}}
\magmisenpage{}
\maglieudiff{}
\magprocess

\begin{document}

%BEGIN_BOOK
\magsection{Produit scalaire}
\magsubsection{Produit scalaire}

\exercice{nom={Produit scalaire}}
On pose
\[E\defeq\enstq{P\in\polyR}{P(0)=0 \et P(1)=0}.\]
On définit l'application $\ps{.}{.}$ de $E\times E$ dans $\R$ par
\[\forall P,Q\in E\qsep \ps{P}{Q}\defeq-\integ{0}{1}{P''(x)Q(x)}{x}.\]
Montrer que $\ps{.}{.}$ définit un produit scalaire.

\magsubsection{Norme}

% \exercice{nom={Inégalités}}
% Soit $n\in\Ns$. 
% \begin{questions}
% \question Montrer que
%   \[\forall x_1,\ldots,x_n\in\Rs\qsep \p{\sum_{k=1}^n x_k^2}\p{\sum_{k=1}^n\frac{1}{x_k^2}}\geq n^2.\]
% \question En déduire que
%   \[\sum_{k=1}^n \frac{1}{k^2}\geq \frac{6n}{(n+1)(2n+1)}.\]
% \end{questions}

\exercice{nom={Intégrales}}
Soit $f$ une fonction continue de $[0,1]$ dans $\R$. Montrer que
\[\abs{\integ{0}{1}{\frac{f(t)}{1+t^2}}{t}}\leq\frac{\sqrt{\pi}}{2}
  \p{\integ{0}{1}{\frac{f^2(t)}{1+t^2}}{t}}^{\frac{1}{2}} \et 
  \abs{\integ{0}{1}{\frac{\sqrt{t}f(t)}{1+t^2}}{t}}\leq\frac{\sqrt{2 \ln(2)}}{2}
  \p{\integ{0}{1}{\frac{f^2(t)}{1+t^2}}{t}}^{\frac{1}{2}}.\]
Déterminer dans chacun des deux cas à quelle condition sur $f$ l'inégalité est une
égalité.

\exercice{nom={Exercice}}
Soit $(u_n)$ une suite positive telle que $\sum u_n$ converge. Montrer que
\[\sum \frac{\sqrt{u_n}}{n}\]
converge.

\exercice{nom={Mines}}

On définit, pour tout $A,B\in\mat{n}{\R}$
\[\ps{A}{B}\defeq\tr\p{\trans{A}B}.\]
\begin{questions}
\question Montrer que $\ps{.}{.}$ définit un produit scalaire sur $\mat{n}{\R}$.
\question Montrer que
  \[\forall A,B\in\mat{n}{\R}\qsep \norme{AB}\leq\norme{A}\norme{B}.\]
\end{questions}

\magsubsection{Notion d'orthogonalité}

\exercice{nom={Vecteurs orthogonaux}}
Soit $E$ un espace préhilbertien et $x,y\in E$. Montrer que $x$ et $y$ sont orthogonaux si
et seulement si
\[\forall \lambda\in\R \qsep \norme{x+\lambda y}\geq\norme{x}.\]

\exercice{nom={Une base orthonormale}}
Soit $E$ un espace euclidien, et $(e_{1},...,e_{n})$ une famille de vecteurs unitaires
tels que
\[\forall x\in E\qsep \norme{x}^2=\sum_{i=1}^{n}\ps{e_i}{x}^{2}.\]
Notons qu'on ne suppose pas que $E$ est de dimension $n$.
\begin{questions}
\question Montrer que $(e_{1},...,e_{n})$ est une famille orthogonale.  
\question Montrer que $(e_{1},...,e_{n})$ est une base orthonormale.                                   
\end{questions}

\magsection{Espace euclidien}
\magsubsection{Supplémentaire orthogonal}

\exercice{nom={Orthogonal et somme}}
Soit $E$ un espace euclidien et $F,G$ deux sous-espaces vectoriels de $E$.
\begin{questions}
\question Montrer que l'orthogonal de $F+G$ et l'intersection des orthogonaux
  de $F$ et $G$ sont égaux.
\question Montrer que l'orthogonal de l'intersection de $F$ et $G$ et la somme des
  orthogonaux de $F$ et de $G$ sont égaux.
\end{questions}
\begin{sol}
$\quad$
\begin{questions}
\question
\question Il suffit d'appliquer la première question à l'orthogonal de $F$ et
  l'orthogonal de $G$.
\end{questions}
\end{sol}

\magsubsection{Base orthonormée }
\magsubsection{Projecteur orthogonal}


\exercice{nom={Projecteur orthogonal}}
Soit $E$ un espace euclidien et $p$ un projecteur de $E$. Montrer que $p$ est orthogonal si et seulement si
\[\forall x\in E \qsep \norme{p(x)}\leq\norme{x}.\]


\exercice{nom={Calcul d'une projection orthogonale}}
On se place dans le $\R$-espace euclidien usuel $\R^n$.
\begin{questions}
\question Soit $F$ le
  sous-espace vectoriel de $\R^4$ défini par le système d'équations
  $$\syslin{
  x_1 &+ x_2 &+ x_3 &+ x_4 &=& 0\hfill\cr
  x_1 &- x_2 &+ x_3 &- x_4 &=& 0.\hfill}
  $$
  Déterminer la matrice dans la base canonique de la projection
  orthogonale sur $F$.
\question Soit $F$ le
  sous-espace vectoriel de $\R^4$ défini par le système d'équations
  $$
  \syslin{
  x_1 &+ x_2  &+ x_3  &+ x_4  &=& 0\hfill\cr
  x_1 &+ 2x_2 &+ 3x_3 &+ 4x_4 &=& 0.\hfill}
  $$
  Déterminer la matrice dans la base canonique de la projection
  orthogonale sur à $F$.
\end{questions}
\begin{sol}
$\quad$
\begin{questions}
\question fait en cours
\question On cherche une base. On trouve $\p{1,-2,1,0},\p{2,-3,0,1}$.
  Après orthonormalisation, on a $\p{1,-2,1,0}/\sqrt{6}$ et
  $\p{2,-1,-4,3}/\sqrt{30}$. On a donc
  \[S=\frac{1}{5}\begin{pmatrix}
      -2 & -4 & -1 & 2\\
      -4 & 2 & -2 & -1\\
      -1 & -2 & 2 & -4\\
      2 & -1 & -4 & -2
      \end{pmatrix}\]
  On peut aussi très bien faire cela sans orthonormalisation de Schmidt.
\end{questions}
\end{sol}


\exercice{nom={Rang d'un projecteur orthogonal}}
Soit $p$ un projecteur orthogonal d'un espace euclidien $E$.
\begin{questions}
\question Montrer que pour tout $x\in E$, $\norme{p(x)}^2=\ps{p(x)}{x}$.
\question Montrer que pour toute base orthonormée $\mathcal{B}\defeq(e_1,\ldots,e_n)$ de $E$
  \[\sum_{k=1}^n \norme{p(e_k)}^2=\rg(p).\]
\end{questions}



\magsubsection{Algorithme d'orthonormalisation de \nom{Gram-Schmidt}}


\exercice{nom={Orthonormalisation}}
On munit $\R^3$ du produit scalaire usuel. Montrer que les vecteurs
\[e_1\defeq(1,0,1),\quad e_2\defeq(1,0,2) \et e_3\defeq(1,1,1)\]
forment une base de $\R^3$ et en déterminer l'orthonormalisée de \nom{Gram-Schmidt}.



\exercice{nom={Une distance}}
Calculer
\[\inf_{(a,b)\in \R^2} \integ{0}{1}{\cro{x^2-(ax+b)}^2}{x}.\]
\begin{sol}
Dans $\R_2[X]$ muni du produit scalaire $\ps{P}{Q} = \displaystyle\int_0^1 PQ$, c'est la distance de $X^2$ à $\R_1[X]$. 
Notons $P_0$ le projeté orthogonale de $X^2$ sur $\R_1[X]$, on a $\ps{X^2-P_0}{1} = \ps{X^2-P_0}{X} = 0$ ce qui donne $P_0 = X-\dfrac{1}{6}$. 
La borne inférieure recherchée est alors $\norme{X^2-P_0} = \dfrac{1}{180}$.
\end{sol}


\exercice{nom={Sur les polynômes}}
Soit $n\in\N$ et $a_0,\ldots,a_n\in\R$. On pose $E\defeq\polyR[n]$ et on définit
l'application $\ps{.}{.}$ de $E\times E$ dans $\R$ par
\[\forall P,Q\in E \qsep \ps{P}{Q}\defeq\sum_{k=0}^n P^{(k)}(a_k)Q^{(k)}(a_k).\]
\begin{questions}
\question Montrer que $\ps{.}{.}$ est un produit scalaire sur $E$.
\question On suppose dans cette question que $n=2$, $a_0=1$, $a_1=2$ et $a_2=3$. Déterminer
  une base orthonormée de $E$.
\end{questions}

\exercice{nom={Racines des polynômes orthogonaux}}
Soit $E\defeq\polyR[n]$, muni de la forme
\[\forall P,Q\in\polyR[n]\qsep \ps{P}{Q}\defeq\integ{-1}{1}{P(t)Q(t)}{t}.\]
\begin{questions}
\question Montrer que $\ps{.}{.}$ est un produit scalaire.
\question Montrer que $E$ admet une base orthonormée $(P_k)_{0\leq k\leq n}$ telle que
  \[\forall k\in\intere{0}{n} \qsep \deg P_k=k.\]
\question Montrer que pour tout $k\in\intere{0}{n}$, $P_k$ admet exactement $k$
  racines sur $\intero{-1}{1}$.
\end{questions}


\magsubsection{Dual}




% \exercice{utile=-3, nom={Quelques produits scalaires}}
% Parmi les formes suivantes, lesquelles sont des produits
% scalaires :
% \begin{questions}
% \question Sur $E=\polyR$ :
%   $$\ps{P}{Q}=\integ{-1}{1}{P(t)Q(t)}{t}$$
% \question Sur $E=\classec{0}(\R,\R)$ :
%   $$\ps{f}{g}=\integ{-1}{1}{f(t)g(t)}{t}$$
% \question Sur $\mat{n}{\R}$ :
%   $$\ps{A}{B}=\tr\p{A\trans{B}}$$
%   On rappelle que si $A\in\mat{n}{\K}$, $\tr A=\sum_{k=1}^n a_{k,k}$.
% \end{questions}

% \exercice{nom={Intégrales}}
% On considère l'ensemble $E$ des fonctions continues et strictement
% positives sur un segment $\interf{a}{b}$. Démontrer que la borne
% inférieure
% $$\inf_{f\in E} \p{\integ{a}{b}{f(x)}{x}
%                    \integinv{a}{b}{f(x)}{x}}$$
% existe. Est-elle atteinte ?

% \magsection{Polynômes orthogonaux}
% \exercice{nom={Calcul d'un minimum}}
% Sur $\polyR$, on définit :
% $$\ps{P}{Q}=\integ{0}{1}{P(t)Q(t)}{t}$$
% \begin{questions}
% \question Montrer que $\ps{.}{.}$ est un produit scalaire.
% \question Trouver une base orthonormée de $\polyR[2]$ pour ce produit
%   scalaire.
% \question Calculer le minimum pour $a,b\in\R$ de :
%   $$\integ{0}{1}{\p{x^2-\p{ax+b}}^2}{x}$$
% \end{questions}



% \magsection{Espace euclidien}






% \magsection{Endomorphismes orthogonaux}
% \exercice{nom={Matrices orthogonales}}
% Soit $A\in\mat{n}{\R}$ une matrice orthogonale. Montrer que~:
% \[\abs{\sum_{1\leq i,j\leq n} a_{i,j}}\leq n\]
% \[\sum_{1\leq i,j\leq n} \abs{a_{i,j}} \leq n\sqrt{n}\]

% \exercice{nom={Construction d'un produit scalaire}}
% Soit $E=\R^n$. On se donne $G$ un sous-groupe fini de $GL\p{\R^n}$ et
% $\ps{.}{.}$ le produit scalaire usuel. On définit la forme $\ps{.}{.}_2$
% sur $\R^n$ par :
% $$\forall x,y\in E \quad \ps{x}{y}_2=\frac{1}{\card G}
%   \sum_{u\in G} \ps{u(x)}{u(y)}$$
% \begin{questions}
% \question Montrer que $\ps{.}{.}_2$ est un produit scalaire.
% \question Montrer que pour ce produit scalaire $G$ est un sous-groupe du
%   groupe unitaire.
% \end{questions}


% \magsection{Produit mixte}
% \exercice{nom={Inégalité de Hadamard}}
% Soit $E$ un espace euclidien de dimension $n$ et $x_1,\ldots,x_n$ une
% famille de $n$ vecteurs de $E$. Montrer que :
% $$\abs{\pmixte{x_1,\ldots,x_n}} \leq \norme{x_1}\cdots\norme{x_n}$$
% On pourra montrer le résultat lorsque la famille $x_1,\ldots,x_n$ est
% orthogonale, avant de passer au cas général.

% \exercice{nom={Déterminant de Gram}}
% Soit $E$ un espace euclidien de dimension $n$ et $x_1,\ldots,x_n$ $n$
% vecteurs de $E$. On appelle matrice de Gram de $x_1,\ldots,x_n$ la
% matrice :
% $$\p{\ps{x_i}{x_j}}_{1\leq i,j\leq n}$$
% et déterminant de Gram le déterminant de cette matrice, noté
% $G(x_1,\ldots,x_n)$.
% \begin{questions}
% \question Montrer que la matrice de Gram peut s'écrire comme le produit
%   d'une matrice et de sa transposée. 
% \question En déduire que :
%   $$G(x_1,\ldots,x_n)=\pmixte{x_1,\ldots,x_n}^2$$
% \question On suppose que la famille $x_1,\ldots,x_p$ est libre et on note $F$
%   le sous-espace vectoriel qu'elle engendre. Étant donné $x\in E$, montrer
%   que la distance $d$ du vecteur $x$ au sous-espace $F$ vérifie :
%   $$d^2=\frac{G(x_1,\ldots,x_p,x)}{G(x_1,\ldots,x_p)}$$
% \end{questions}


% \magsection{Espace affine}

% \exercice{nom={Plan passant par une droite et un point}}
% Soit $\mathcal{E}=\R^3$ muni de sa structure d'espace affine usuelle et
% $\mathcal{R}=\p{O,\p{\ve{e_1},\ve{e_2},\ve{e_3}}}$ où $O=\p{0,0,0}$ et
% $\ve{e_1},\ve{e_2},\ve{e_3}$ est la base canonique de $E$. Soit $A$ le point de
% coordonnnées $\p{1,-1,1}$ et $\mathcal{D}$ la droite d'équation~:
% \[\syslin{x &-3y&+2z&=&1\hfill\cr
%           2x& +y&-3z&=&-1\hfill}\]
% Donner l'équation cartesienne du plan $\mathcal{P}$ passant par $A$ et
% $\mathcal{D}$.

% \exercice{nom={Droites coplanaires}}
% Soit $\mathcal{E}=\R^3$ muni de sa structure d'espace affine usuelle et
% $\mathcal{R}=\p{O,\p{\ve{e_1},\ve{e_2},\ve{e_3}}}$ où $O=\p{0,0,0}$ et
% $\ve{e_1},\ve{e_2},\ve{e_3}$ est la base canonique de $E$. Soit $\mathcal{D}$
% et $\mathcal{D}'$ les droites d'équation~:
% \[\mathcal{D} \quad \syslin{x& &-2z&=&1\cr
%                              &y& -z&=&2} \et
%   \mathcal{D}' \quad \syslin{x& +y& +z&=&1\cr
%                              x&-2y&+2z&=&a}\]
% \begin{questions}
% \question Pour quelles valeurs de $a$, $\mathcal{D}$ et $\mathcal{D}'$
%   sont-elles coplanaires~?
% \question Donner alors l'équation du plan contenant $\mathcal{D}$ et
%   $\mathcal{D}'$.
% \end{questions}

% \exercice{nom={Perpendiculaire commune}}
% L'espace est rapporté au repère orthonormé $Oxyz$.
% \begin{questions}
% \question Déterminer la perpendiculaire commune et la distance des droites~:
%   \[\mathcal{D}_1=A_1+\R\ve{u_1} \et \mathcal{D}_2=A_2+\R\ve{u_2}\]
%   où $A_1$ et $A_2$ sont de composantes $\p{1,0,0}$ et $\p{0,0,0}$ et
%   $\ve{u_1}$ et $\ve{u_2}$ de composantes $\p{1,-1,1}$ et $\p{1,1,1}$.
% \question Déterminer la perpendiculaire commune et la distance des droites~:
%   \[\mathcal{D}_1 \quad \syslin{x&+y&+z&=&4\cr
%                                 x&-y&+z&=&6} \et
%     \mathcal{D}_2 \quad \syslin{x&+y&-z&=&4\hfill\cr
%                                 x&+y&+z&=&-2}\]
% \end{questions}

% \exercice{nom={Symétries}}
% Soit $\mathcal{E}=\R^3$ muni de sa structure d'espace affine usuelle et
% $\mathcal{R}=\p{O,\p{\ve{e_1},\ve{e_2},\ve{e_3}}}$ où $O=\p{0,0,0}$ et
% $\ve{e_1},\ve{e_2},\ve{e_3}$ est la base canonique de $E$. Déterminer les
% expressions analytiques des applications suivantes~:
% \begin{questions}
% \question La symétrie par rapport au plan d'équation $x+2y+z=1$ et de direction
%   $\vect\p{\ve{e_1}+\ve{e_2}+\ve{e_3}}$.
% \question La symétrie par rapport à la droite d'équation~:
%   \[\syslin{x&+y&  &+1&=&0\cr
%              &2y&+z&+2&=&0}\]
%   et de direction le plan vectoriel d'équation $3x+3y-2z=0$.
% \end{questions}

% \exercice{nom={Symétries-translation}}
% Soit $f:\mathcal{E}\to\mathcal{E}$ une application affine. On dit que $f$ est
% une symétrie-translation lorsqu'il existe une symétrie $s$ et une translation
% $t$ telles que $f=s\circ t=t\circ s$.
% \begin{questions}
% \question Soit $s$ une symétrie par rapport à $\mathcal{F}$ de direction $G$
%   et $t$ une translation de vecteur $\ve{u}$. Montrer que $s\circ t=t\circ s$
%   si et seulement si $\ve{u}\in F$.
% \question Soit $f$ affine quelconque. Montrer que $f$ est une
%   symétrie-translation si et seulement si $f\circ f$ est une translation.
% \question En déduire que le produit d'une symétrie quelconque par une
%   translation quelconque est une symétrie-translation.
% \question Décomposer l'application $f$ d'expression analytique dans le repère
%   $\mathcal{R}=\p{O,\p{\ve{e_1},\ve{e_2},\ve{e_3}}}$~:
%   \[\syslin{x'&=&\frac{1}{3}(&  x&-2y&-2z&+1)\cr
%             y'&=&\frac{1}{3}(&-2x& +y&-2z&+2)\cr
%             z'&=&\frac{1}{3}(&-2x&-2y& +z&-1)}\]

% \exercice{nom={Isométrie affine de l'espace}}
% Soit $\mathcal{E}=\R^3$ et $f$ l'application de $\mathcal{E}$ dans $\mathcal{E}$
% dont l'expression analytique est~:
% \[\syslin{x'&=&\quad& &z&  \cr
%           y'&=&x&\quad& &+1\cr
%           z'&=& &y&\quad&+2}\]
% \begin{questions}
% \question Montrer que la partie linéaire de $f$ est une rotation. Précisez son
%   axe et son angle.
% \question En déduire la nature et les éléments caractéristiques de $f$.
% \end{questions}

% \exercice{nom={Points alignés}}
% Soit $I,J,K$ trois points du plan affine $\mathcal{P}$ (espace affine de
% dimension 2). Montrer l'équivalence entre les trois propriétés suivantes~:
% \begin{questions}
% \question $I$, $J$ et $K$ sont alignés.
% \question Il existe $M\in\mathcal{P}$ tel que~:
%   \[\det\p{\ve{MI},\ve{MJ}}+\det\p{\ve{MJ},\ve{MK}}+\det\p{\ve{MK},\ve{MI}}=0\]
% \question Pour tout $M\in\mathcal{P}$~:
%   \[\det\p{\ve{MI},\ve{MJ}}+\det\p{\ve{MJ},\ve{MK}}+\det\p{\ve{MK},\ve{MI}}=0\]
% \end{questions}

% \magsection{Barycentres}
% \exercice{nom={Coordonnées barycentriques}}
% Soit $A,B,C$ trois points du plan non alignés. On se donne $M_1,M_2,M_3$ trois
% points du plan dont les coordonnées barycentriques dans le système $(A,B,C)$
% sont respectivement $(\alpha_1,\beta_1,\gamma_1)$, $(\alpha_2,\beta_2,\gamma_2)$
% et $(\alpha_3,\beta_3,\gamma_3)$. Montrer que $M_1,M_2$ et $M_3$ sont alignés si
% et seulement si~:
% \[\begin{vmatrix}
%   \alpha_1 & \alpha_2 & \alpha_3\\
%   \beta_1  & \beta_2 & \beta_3\\
%   \gamma_1 & \gamma_2 & \gamma_3
%   \end{vmatrix}=0\]

% \exercice{nom={Droites concourantes}}
% Dans le plan affine $\mathcal{P}$ rapporté à un repère $\mathcal{R}$, on se
% donne $3$ droites $\mathcal{D}_1,\mathcal{D}_2$ et $\mathcal{D}_3$ d'équations
% respectives~:
% \[(\mathcal{D}_i) \quad \alpha_i x+\beta_i y+\gamma_i=0\]
% On suppose qu'au moins deux des droites ne sont pas parrallèles. Montrer
% que les trois droites sont concourantes si et seulement si :
% \[\begin{vmatrix}
%   \alpha_1 & \alpha_2 & \alpha_3\\
%   \beta_1  & \beta_2  & \beta_3\\
%   \gamma_1 & \gamma_2 & \gamma_3\\
%   \end{vmatrix}=0\]


%END_BOOK
\end{document}


















\chapter{Espérance, variance}
\setcounter{numeroexercicecours}{1}
\documentclass{magnolia}

\magtex{tex_driver={pdftex},
        tex_packages={epigraph,tikz,pgfplots,slashbox,nicefrac,listings,xypic}}
\magfiche{document_nom={Cours sur l'espérance et la variance},
          auteur_nom={François Fayard},
          auteur_mail={fayard.prof@gmail.com}}
\magcours{cours_matiere={maths},
          cours_niveau={mpsi},
          cours_chapitre_numero={24},
          cours_chapitre={Espérance, variance}}
\magmisenpage{}
\maglieudiff{}
\magprocess

\usetikzlibrary{positioning}
\lstset{%
  frame            = tb,    % draw frame at top and bottom of code block
  tabsize          = 1,     % tab space width
  numbers          = none,  % display line numbers on the left
  framesep         = 3pt,   % expand outward
  framerule        = 0.4pt, % expand outward 
  commentstyle     = \color{green},      % comment color
  keywordstyle     = \color{blue},       % keyword color
  stringstyle      = \color{blue},    % string color
  backgroundcolor  = \color{colorLazoBlue1Light}, % backgroundcolor color
  showstringspaces = false,              % do not mark spaces in strings
  basicstyle       = \ttfamily,
  breaklines       = true
}

\begin{document}


% 30.35 de Troesch: Cauchy-Schwarz
% 

%BEGIN_BOOK

\setlength\epigraphwidth{.7\textwidth}
\epigraph{\og Je ne crois aux statistiques que lorsque je les ai moi-même falsifiées. \fg}{--- {\sc Winston Churchill (1874--1965)}}
\hidemesometimes{
\hfill\includegraphics[width=0.25\textwidth]{../../Commun/Images/maths-cours-launch_risk.png}}

\magtoc


% \begin{proposition}
% Soit $P_1,\ldots,P_n$ une famille de lois sur les ensembles finis
% respectifs $E_1,\ldots,E_n$. Alors, il existe une probabilité $\mathbb{P}$ sur
% $\Omega\defeq E_1\times\cdots\times E_n$ telle que les variables aléatoires
% \[\dspappli{X_k}{\Omega}{E_k}{(x_1,\ldots,x_n)}{x_k}\]
% soient mutuellement indépendantes et suivent respectivement les lois $P_1,\ldots,P_n$.
% \end{proposition}

% \begin{remarqueUnique}
% \remarque Si $P_1,\ldots,P_n$ sont des lois uniformes, alors la probabilité $\mathbb{P}$ est
%   uniforme.
% \end{remarqueUnique}

\section{Espérance, variance}

\subsection{Espérance}

\begin{definition}
Soit $X:\Omega\to\R$ une variable aléatoire réelle. On appelle \emph{espérance} de $X$ le
réel
\[\mathbb{E}(X)\defeq\sum_{x\in X(\Omega)} x \mathbb{P}(X=x).\] 
\end{definition}

\begin{remarques}
\remarque Dans cette formule, on peut remplacer $X(\Omega)$ par une partie finie $E'$ de $E$ telle que
  $X(\Omega)\subset E'$. 
\remarque L'espérance de $X$ ne dépend que de la loi suivie par $X$.
\remarque Si $A\in\mathcal{P}(\Omega)$ est un évènement, alors
  $\mathbb{E}(\mathds{1}_A)=\mathbb{P}(A).$
\end{remarques}

\begin{exoUnique}
\exo On lance $n$ fois une pièce équilibrée. On s'intéresse à la première apparition de
  pile. On note $X$ la variable aléatoire égale à $k$ si le premier pile a lieu au $k$-ième
  lancer et on convient que $X=0$ si aucun lancer n'amène pile. Déterminer la loi et
  l'espérance de $X$.
  % \begin{prof}
  % $\pr{X=k} = \frac 1{2^k}$ puis $\esp{X} = \sum\limits_{k=1}^n \frac{k}{2^k} = 2 \p{1 - \frac{n+2}{2^{n+1}}}$.
  % \end{prof}
\end{exoUnique}

\begin{proposition}
Soit $X:\Omega\to\R$ une variable aléatoire réelle. Alors
\[\mathbb{E}(X)\defeq\sum_{\omega\in\Omega} X(\omega) \mathbb{P}(\ens{\omega}).\]
\end{proposition}


\begin{proposition}
Soit $X,Y:\Omega\to\R$ deux variables aléatoires réelles. Alors, quels que soient
$a,b\in\R$
\[\mathbb{E}(aX+bY)=a\mathbb{E}(X)+b\mathbb{E}(Y).\]
\end{proposition}

\begin{exos}
\exo On lance deux dés à 6 faces. Calculer l'espérance de leur somme.
\exo Le chat \nom{Speed} se fait chaque jour les griffes, soit sur le canapé soit sur les
  rideaux. Il ne se fait jamais les griffes deux jours de suite sur le canapé. S'il se fait les
  griffes sur les rideaux un jour donné, alors il choisira le lendemain le canapé avec une
  probabilité \nicefrac{1}{3}. Vous devez vous occuper de \nom{Speed} pendant $m+1$ jours.
  Le premier jour, jour d'indice 0, il attaque les rideaux. Pour tout $n\in\intere{0}{m}$, on
  définit la variable aléatoire $X_{n}$ égale au nombre de jours où le chat a fait ses griffes
  sur les rideaux parmi les jours d'indices 0 à $n$. Calculer $\mathbb{E}(X_n)$.
% \begin{reponse}
% On a un processus de Markov sur un graphe :
% \medskip
% On note $Y_i$ la variable aléatoire qui vaut $1$ si le chat a fait ses griffes sur le rideau le jour $i$ et $0$ sinon. Il s'agit d'une variable de Bernoulli
% de paramètre $p_i$.
% On a $X_n=\sum\limits_{i=1}^n Y_i$ donc $\esp{X}=\sum\limits_{i=1}^n \esp{Y_i}=\sum\limits_{i=1}^n p_i$.
% $$p_i=\pr{Y_i=1}= \prc{Y_i=1}{Y_{i-1}=0} \pr{Y_{i-1}=0} + \prc{Y_i=1}{Y_{i-1}=1} \pr{Y_{i-1}=1}$$
% $$\boxed{p_i}= 1 \times \pr{Y_{i-1}=0} + \frac 23 \pr{Y_{i-1}=1}=1-p_{i-1} + \frac 23 p_{i-1}= \boxed{1 - \frac 13 p_{i-1}}$$
% Avec $p_1=1$ on trouve par la méthode usuelle $p_i=\frac 34 + \frac 14 \p{- \frac 13 }^{i-1}$ puis en sommant
% $$\boxed{\esp{X_n} = \frac 34 n + \frac{3}{16} \p{1 - \p{ - \frac 13 }^{n}} \sim \frac 34 n}$$
% Attention, on peut aussi faire le calcul directement en sommant la relation de récurrence si l'on ne s'intéresse qu'à la limite de l'espérance.
% \end{reponse}
\exo Soit $A_1,\ldots,A_n\in\mathcal{P}(\Omega)$. Montrer la formule dite du crible
  \[\mathbb{P}(A_1\cup\ldots\cup A_n)=\sum_{k=1}^n
    (-1)^{k+1}\sum_{1\leq i_1<\cdots<i_k\leq n}
    \mathbb{P}(A_{i_1}\cap\ldots\cap A_{i_k}).\]
  En déduire que si $A_1,\ldots,A_n$ sont des parties d'un ensemble fini $E$
  \[\card(A_1\cup\ldots\cup A_n)=\sum_{k=1}^n
    (-1)^{k+1}\sum_{1\leq i_1<\cdots<i_k\leq n}
    \card(A_{i_1}\cap\ldots\cap A_{i_k}).\] 
\end{exos}


\begin{proposition}
Soit $X,Y:\Omega\to\R$ deux variables aléatoires réelles.
\begin{itemize}
\item Si $X\geq 0$, alors $\mathbb{E}(X)\geq 0$.
\item Si $X\leq Y$, alors $\mathbb{E}(X)\leq\mathbb{E}(Y)$.
\item On a $\abs{\mathbb{E}(X)}\leq \mathbb{E}(\abs{X})$.
\end{itemize}
\end{proposition}

\begin{proposition}
Soit $X:\Omega\to\R$ une variable aléatoire réelle.
\begin{itemize}
\item Si $X$ est une variable aléatoire constante égale à $a$, alors $\mathbb{E}(X)=a$.
\item S'il existe $a,b\in\Z$ tel que $X\sim\mathcal{U}(\intere{a}{b})$, alors
  $\mathbb{E}(X)=(a+b)/2$.
\item S'il existe $p\in[0,1]$ tel que $X\sim\mathcal{B}(p)$, alors $\mathbb{E}(X)=p$.
\item S'il existe $n\in\N$ et $p\in[0,1]$ tels que $X\sim\mathcal{B}(n, p)$, alors
  $\mathbb{E}(X)=np$. 
\end{itemize}
\end{proposition}

% \begin{remarqueUnique}
% \remarque Dans les formules donnant les espérances d'une loi de \nom{Bernoulli} et d'une loi binomiale,
%   on a posé $q\defeq 1-p$.
% \end{remarqueUnique}

\begin{exos}
\exo On dispose de $n$ fléchettes. À chaque lancer, on a une probabilité de \nicefrac{1}{10} de
  tirer dans le mille. On suppose les lancers indépendants et on note $X$ la variable aléatoire
  égale au nombre de fléchettes placées dans le mille.
\begin{questions}
\question Donner la loi de $X$.
\question Calculer la probabilité de mettre au plus 1 fléchette dans le mille.
\question Calculer l'espérance de $X$.
\end{questions}
% \begin{reponse}
% \begin{enumerate}
% \item \fbox{$X$ suit une loi binomiale de paramètres $n=20$ et $p=0,1$}. 
% On a donc  $\pr{X=k}=\binom{20}{k} p^k (1-p)^{20-k}$.
% \item $\pr{X \leq 1} = \pr{X=0}+\pr{X=1}=q^{20}+20 q^{19} p \simeq \boxed{0,39}$.
% \item $\esp{X}=20 \times 0,1 = \boxed{2}$.
% \end{enumerate}
% \end{reponse}
\exo Soit $n\in\Ns$. On pose $I_n\defeq\intere{1}{n}$ et on munit
  $\Omega_n\defeq\mathcal{F}(I_n,I_n)$ de la
  probabilité uniforme. On note $X_n$ la variable aléatoire sur $\Omega_n$ égale au cardinal de
  l'image d'un élément de $\Omega_n$. Déterminer l'espérance de $X_n$ puis en donner un équivalent quand
  $n$ tend vers $+\infty$.
% \begin{reponse}
% On note $X_k(f) = 1$ si $k$ a un antécédent par $f$ et $0$ sinon. On a donc une variable de Bernoulli de paramètre $p$.
% $1-p=\pr{X_k = 0} = \frac{ \card {\cal F}(I_n, I_k \backslash \{k \})}{\card \Omega} = \frac{(n-1)^n}{n^n}$ donc 
% $p = 1 - \p{1 - \frac 1n}^n$.
% Mais $X = \sum\limits_{k=1}^n X_k$ donc $\boxed{\esp{X} =n \p{1 - \p{1 - \frac 1n}^n} \sim \frac{n}{e}}$.
% \end{reponse}
% \exo On s'intéresse au nombre moyen d'affectations réalisées dans l'algorithme usuel de
%   recherche du maximum dans un tableau. On adopte le modèle suivant~: un tableau d'entiers de
%   longueur $n$ est une permutation des entiers de $1$ à $n$ et ces permutations sont
%   équiprobables. Donner un équivalent du nombre moyen d'affectations réalisées.
% % \begin{lstlisting}[language = Python, numbers = none, escapechar = !,
% %     basicstyle = \ttfamily\bfseries, linewidth = .6\linewidth] 
% \begin{lstlisting}[language=python]
% def find_max(t):
%   """find_max(t: list[int]) -> int"""
%   n = len(t)
%   maximum = t[0]
%   for k in range(n):
%     if t[k] >= maximum:
%       maximum = t[k]
%   return maximum
% \end{lstlisting}
\end{exos}

\begin{definition}
On dit qu'une variable aléatoire réelle $X:\Omega\to\R$ est centrée lorsque
\[\mathbb{E}(X)=0.\]
\end{definition}

\begin{remarqueUnique}
\remarque Si $X:\Omega\to\R$ est une variable aléatoire, alors $X-\mathbb{E}(X)$ est
  centrée.
\end{remarqueUnique}

\begin{proposition}[nom={Formule de transfert}]
Soit $X:\Omega\to E$ une variable aléatoire et $f:E\to\R$. Alors
\[\mathbb{E}(f(X))=\sum_{x\in X(\Omega)} f(x) \mathbb{P}(X=x).\]
\end{proposition}

\begin{remarqueUnique}
\remarque Dans cette formule, on peut remplacer $X(\Omega)$ par une partie finie $E'$ de $E$ telle que
  $X(\Omega)\subset E'$. 
\end{remarqueUnique}

\begin{exoUnique}
\exo Soit $X$ une variable aléatoire suivant une loi uniforme sur $\intere{-2}{2}$. Déterminer
  l'espérance de $X^2$ en utilisant la définition de l'espérance, puis la formule de
  transfert.
\end{exoUnique}

\begin{proposition}
Soit $X,Y:\Omega\to\R$ deux variables aléatoires réelles indépendantes. Alors
\[\mathbb{E}(XY)=\mathbb{E}(X)\mathbb{E}(Y).\]
Plus généralement, soit $X_1,\ldots,X_n:\Omega\to\R$ des variables aléatoires mutuellement
indépendantes. Alors
\[\mathbb{E}\p{\prod_{i=1}^n X_i}=\prod_{i=1}^n \mathbb{E}(X_i).\]
\end{proposition}

\begin{exos}
\exo Dans le quadrillage $\intere{0}{n}\times\intere{0}{p}$, on place uniformément et de manière indépendante
  un point $A$ sur l'axe des abscisses et un point $B$ sur l'axe
  des ordonnées. Quelle est, en moyenne, l'aire du triangle $OAB$~?
% \begin{reponse}
% On note $X$ l'abscisse de $A$ et $Y$ l'ordonnée de $B$. $X$ et $Y$ suivent
% des lois uniformes. L'aire de $OAB$ est $Z=\frac 12 XY$ donc $\exp{Z}=\frac
% 12 \esp{X} \esp{Y}$ puisque $X$ et $Y$ sont indépendantes.
% Mais $\esp{X}=\frac n2$ et $\esp{Y}=\frac p2$  donc $\boxed{\esp{Z}=\frac{np}{8}}$.
% \end{reponse}
\exo On effectue $p\in\Ns$ tirages avec remise dans une urne contenant des boules numérotées
  de 1 à $n$. On note $X_1,\ldots,X_p$ les résultats des tirages successifs et on pose
  $Y\defeq\max(X_1,\ldots,X_p)$.
  \begin{questions}
  \question Pour tout $k\in\intere{0}{n}$, calculer $\mathbb{P}(Y\leq k)$. En déduire $\mathbb{E}(Y)$.
  \question Déterminer un équivalent de $\mathbb{E}(Y)$ lorsqu'à $p$ fixé, $n$ tend vers $+\infty$.
  \question Déterminer un équivalent de $\mathbb{E}(Y)$ lorsqu'à $n$ fixé, $p$ tend vers $+\infty$.
  \end{questions}
\end{exos}



\subsection{Variance}

\begin{definition}
Soit $X$ une variable aléatoire réelle. On appelle respectivement \emph{variance} et
\emph{écart-type} de $X$ les réels
\[\mathbb{V}(X)\defeq\mathbb{E}\p{\cro{X-\mathbb{E}(X)}^2} \quad\et\quad
  \sigma(X)\defeq\sqrt{\mathbb{V}(X)}.\]
\end{definition}

\begin{remarqueUnique}
\remarque La variance et l'écart-type de $X$ ne dépendent que de la loi suivie par $X$.
\end{remarqueUnique}

\begin{proposition}[nom={Formule de \nom{König-Huygens}}]
Soit $X:\Omega\to\R$ une variable aléatoire réelle. Alors
\[\mathbb{V}(X)=\mathbb{E}(X^2)-\mathbb{E}(X)^2.\]
\end{proposition}

\begin{remarqueUnique}
\remarque On calcule souvent la variance de $X$ en utilisant la formule d'\nom{Huygens}
  ou en passant par $\mathbb{E}(X(X-1))$ qui est parfois plus commode à calculer que
  $\mathbb{E}(X^2)$. On utilise alors la relation
  \[\mathbb{V}(X)=\mathbb{E}(X(X-1))+\mathbb{E}(X)-\mathbb{E}(X)^2.\]
\end{remarqueUnique}

\begin{exoUnique}
\exo Soit $X$ une variable aléatoire réelle. Quelle valeur de $\lambda\in\R$ minimise
  $\mathbb{E}((X-\lambda)^2)$~? Que vaut ce minimum~?
\end{exoUnique}

\begin{proposition}
Soit $X:\Omega\to\R$ une variable aléatoire réelle. Alors, quels que soient $a,b\in\R$
\[\mathbb{V}(aX+b)=a^2\mathbb{V}(X) \quad\et\quad
  \sigma(aX+b)=\abs{a}\sigma(X).\]
\end{proposition}

\begin{proposition}
Soit $X:\Omega\to\R$ une variable aléatoire réelle.
\begin{itemize}
\item Si $X$ est une variable aléatoire constante égale à $a\in\R$, alors $\mathbb{V}(X)=0$.
\item Soit $a,b\in\Z$ tels que $a\leq b$. Si $X\sim\mathcal{U}(\intere{a}{b})$,
  alors $\mathbb{V}(X)=(n^2-1)/12$ où $n\defeq\card(\intere{a}{b})=b-a+1$.
\item S'il existe $p\in[0,1]$ tel que $X\sim\mathcal{B}(p)$, alors $\mathbb{V}(X)=pq$.
\item S'il existe $n\in\N$ et $p\in[0,1]$ tels que $X\sim\mathcal{B}(n, p)$, alors
  $\mathbb{V}(X)=npq$. 
\end{itemize}
\end{proposition}

\begin{remarqueUnique}
\remarque Dans les formules donnant les variances d'une loi de \nom{Bernoulli} et d'une loi binomiale,
  on a posé $q\defeq 1-p$.
\end{remarqueUnique}

\begin{proposition}
Soit $X:\Omega\to\R$ une variable aléatoire réelle. Alors $\mathbb{V}(X)=0$ si et seulement
si il existe un évènement $A$ de probabilité 1 tel que
\[\forall \omega\in A\qsep X(\omega)=\mathbb{E}(X).\]
\end{proposition}

\begin{remarqueUnique}
\remarque Autrement dit, la variance d'une variable aléatoire $X:\Omega\to\R$ est nulle si et
  seulement si $X$ est constante sur un évènement de probabilité 1. Un tel évènement est
  dit \og \emph{presque sûr} \fg.
\end{remarqueUnique}

\begin{definition}
On dit qu'une variable aléatoire réelle $X:\Omega\to\R$ est réduite lorsque
\[\mathbb{V}(X)=1.\]
\end{definition}

\begin{remarqueUnique}
\remarque Si $X:\Omega\to\R$ est une variable aléatoire réelle de variance non nulle, alors
  \[\frac{X-\mathbb{E}(X)}{\sigma(X)}\]
  est centrée réduite.
\end{remarqueUnique}

% \begin{definition}
% Soit $X:\Omega\to\R$ une variable aléatoire réelle. Pour tout $k\in\N$, on appelle moment
% d'ordre $k$ le réel $\mathbb{E}\p{X^k}$.
% \end{definition}

\section{Couple de variables aléatoires}

\subsection{Loi conjointe}

\begin{definition}
Soit $X:\Omega\to E$ et $Y:\Omega\to F$ deux variables aléatoires. On définit la variable
aléatoire $Z\defeq (X,Y)$ par
\[\dspappli{Z}{\Omega}{E\times F}{\omega}{(X(\omega),Y(\omega)).}\]
\end{definition}

\begin{remarqueUnique}
\remarque On a $Z(\Omega)\subset X(\Omega)\times Y(\Omega)$. Cependant, l'inclusion peut être
  stricte. Par exemple, Si $X$ est une variable aléatoire suivant la loi uniforme
  sur $\intere{1}{n}$ et $Z\defeq (X,X)$, alors $Z(\Omega)=\ensim{(k,k)}{k\in\intere{1}{n}}$
  tandis que $X(\Omega)\times X(\Omega)=\intere{1}{n}^2$.
\end{remarqueUnique}

% \begin{exoUnique}
% \exo On lance deux dés à quatre faces. On modélise cette expérience aléatoire en
%   prenant $\Omega\defeq\intere{1}{4}^2$, que l'on munit de la probabilité uniforme.
%   On considère les variables aléatoires $A$ et $B$ donnant respectivement les résultats
%   du premier et du second dé. 
%   \begin{center}
%   \includegraphics[width=0.15\textwidth]{../../Commun/Images/maths-cours-de-4_faces.png}
%   \end{center}
%   \begin{questions}
%   \question On pose $C\defeq(A,B)$. Déterminer $A(\Omega)$, $B(\Omega)$ et $C(\Omega)$.
%   \question On pose $X\defeq \min(A,B)$, $Y\defeq\max(A,B)$ et $Z\defeq(X,Y)$.
%     Déterminer $X(\Omega)$, $Y(\Omega)$ et $Z(\Omega)$. 
%   \end{questions}
% \end{exoUnique}


% \begin{proposition}
% Soit $X:\Omega\to E$ et $Y:\Omega\to F$ deux variables aléatoires. $X(\Omega)$ et
% $Y(\Omega)$ sont des ensembles finis que l'on note respectivement $\ens{x_1,\ldots,x_n}$ et
% $\ens{y_1,\ldots,y_m}$. Alors la famille des $((X=x_i)\cap(Y=y_j))$ pour
% $(i,j)\in\intere{1}{n}\times\intere{1}{m}$ forme un système complet d'évènements.
% \end{proposition}

% \begin{remarqueUnique}
% \remarque Cette proposition reste vraie si on remplace respectivement $X(\Omega)$ et $Y(\Omega)$
%   par des parties finies $E'$ de $E$ et $F'$ de $F$ telles que $X(\Omega)\subset E'$
%   et $Y(\Omega)\subset F'$.
% \end{remarqueUnique}

\begin{definition}
Soit $X:\Omega\to E$ et $Y:\Omega\to F$ deux variables aléatoires.
\begin{itemize}
\item  On appelle \emph{loi conjointe} du couple $(X,Y)$ la loi de $Z\defeq (X,Y)$.
\item  On appelle \emph{première loi marginale} de $(X,Y)$ la loi de $X$ et
\emph{seconde loi marginale} de $(X,Y)$ la loi de $Y$.
\end{itemize}
\end{definition}

\begin{remarques}
\remarque La loi conjointe de $(X,Y)$ est caractérisée par sa distribution de
  probabilités $(\mathbb{P}(Z=z))_{z\in Z(\Omega)}$. Étant donné que
  $Z(\Omega)\subset X(\Omega)\times Y(\Omega)$, cette distribution est
  donnée par
  \[\p{\mathbb{P}((X=x)\cap(Y=y))}_{(x,y)\in X(\Omega)\times Y(\Omega)}.\]
\remarque Soit $X:\Omega\to E$ et $Y:\Omega\to F$ deux variables aléatoires. On note
$X(\Omega)\defeq\ens{x_1,\ldots,x_n}$ et $Y(\Omega)\defeq\ens{y_1,\ldots,y_m}$. Alors
\begin{eqnarray*}
\forall i\in\intere{1}{n}\qsep \mathbb{P}(X=x_i)&=&\sum_{j=1}^m \mathbb{P}((X=x_i)\cap(Y=y_j))\\
\forall j\in\intere{1}{m}\qsep \mathbb{P}(Y=y_j)&=&\sum_{i=1}^n \mathbb{P}((X=x_i)\cap(Y=y_j))
\end{eqnarray*}
En particulier, les lois marginales se déduisent de la loi conjointe. Par contre, la
connaissance des lois marginales ne permet pas d'en déduire la loi conjointe. Par exemple,
si on considère l'expérience qui consiste à jouer deux fois de suite à pile ou face,
modélisée par $\Omega\defeq\ens{P,F}^2$ muni de la loi uniforme, et que
$X_1,X_2:\Omega\to\ens{P,F}$ sont les variables aléatoires donnant les résultats respectifs
des deux lancers, alors $(X_1,X_2)$ et $(X_1,X_1)$ ont les mêmes lois marginales, mais n'ont
pas la même loi conjointe.
\end{remarques}

\begin{exoUnique}
\exo On considère l'expérience aléatoire qui consiste à tirer successivement et sans remise
  deux boules dans une urne contenant $2$ boules rouges et $3$ boules noires. On note $X_1$
  (respectivement $X_2$) la variable aléatoire qui vaut $R$ si la première (respectivement
  deuxième) boule tirée est rouge, et $N$ sinon.
  \begin{questions}
  \question Déterminer la loi du couple $Y\defeq(X_1,X_2)$, puis ses lois marginales.
  \question Même question si le tirage se fait avec remise.
  \end{questions}
\end{exoUnique}

\begin{definition}
Soit $X:\Omega\to E$ une variable aléatoire à valeurs dans un ensemble fini $E$ et $B\in\mathcal{P}(\Omega)$
un évènement de probabilité non nulle. On appelle \emph{loi de $X$ conditionnée par l'évènement $B$}
la probabilité
\[\dspappli{\mathbb{P}_{X|B}}{\mathcal{P}(X(\Omega))}{[0,1]}{A}{\mathbb{P}(X\in A|B).}\]
\end{definition}

\begin{remarques}
\remarque En pratique, lorsqu'il nous sera demandé de déterminer la loi d'une variable aléatoire
  $X:\Omega\to E$ conditionnée par l'évènement $B$, on commencera par déterminer un ensemble fini
  $E'$ tel que $X(\Omega)\subset E'$ et on calculera $\mathbb{P}(X=x|B)$ pour tout $x\in E'$.
\remarque Si $X:\Omega\to E$ et $Y:\Omega\to F$ sont deux variables aléatoires et $y\in F$
  est tel que $\mathbb{P}(Y=y)>0$, il est courant de considérer la loi de $X$ conditionnée
  par l'évènement $(Y=y)$.
\end{remarques}

% \begin{exoUnique}
% \exo \nom{Jinko} le chaton a trois passions dans la vie~: Manger, Dormir et Jouer. On peut
%   considérer qu'il pratique ces activités par tranches de 5 minutes.
%   \begin{itemize}
%   \item Après 5 minutes de repas, il continue de manger les 5 minutes suivantes avec une
%     probabilité de \nicefrac{1}{2}. Sinon, il se met à jouer.
%   \item Après 5 minutes de sommeil, il continue de dormir les 5 minutes suivantes avec une
%     probabilité de \nicefrac{3}{4}. Sinon, il a faim au réveil et va manger.
%   \item Après 5 minutes de jeu, soit il est en appétit et mange les 5 minutes suivantes
%     avec une probabilité de \nicefrac{1}{4}, soit il est fatigué et s'endort.
%   \end{itemize}
%   L'expérience que l'on considère est une journée de \nom{Jinko} de $5(m+1)$ minutes. Sur la
%   tranche d'indice 0 des 5 premières minutes de la journée, \nom{Jinko} se lève et va
%   petit-déjeuner. Pour tout $n\in\intere{0}{m}$, on note $A_n:\Omega\to\ens{M,D,J}$ la
%   variable aléatoire donnant l'activité de \nom{Jinko} sur la tranche de 5 minutes
%   d'indice $n$ de la journée. Pour tout
%   $n\in\intere{0}{m}$, on pose
%   \[X_n\defeq
%     \begin{pmatrix}
%     \mathbb{P}(A_n=M)\\
%     \mathbb{P}(A_n=D)\\
%     \mathbb{P}(A_n=J)
%     \end{pmatrix}.\]
%   \begin{questions}
%   \question Déterminer une matrice $B\in\mat{3}{\R}$ telle que
%     \[\forall n\in\intere{0}{m-1}\qsep X_{n+1}=B X_n.\]
%   \question
%     \begin{questions}
%     \question Déterminer les valeurs propres de $B$, c'est-à-dire les $\lambda\in\R$ tels
%       que $B-\lambda I_3$ n'est pas inversible.
%     \enonce On note $\lambda_1,\lambda_2,\lambda_3\in\R$ ces valeurs et on pose
%       $P\defeq(X-\lambda_1)(X-\lambda_2)(X-\lambda_3)$.
%     \question Montrer que $P(B)=0$. En déduire une expression de $B^n$ valable pour tout
%       $n\in\N$.
%     \question En déduire les limites de
%       \[\mathbb{P}(A_n=M),\qquad
%         \mathbb{P}(A_n=D),\qquad
%         \mathbb{P}(A_n=J)\]
%       lorsque $n$ tend vers $+\infty$.
%     \end{questions}
%   \end{questions}
% \end{exoUnique}


\subsection{Covariance}

\begin{definition}
Soit $X,Y:\Omega\to\R$ deux variables aléatoires réelles. On appelle \emph{covariance} de $X$ et
de $Y$ et on note ${\rm Cov}(X,Y)$ le réel
\[{\rm Cov}(X,Y)=\mathbb{E}\p{\cro{X-\mathbb{E}(X)}\cro{Y-\mathbb{E}(Y)}}.\]
Lorsque ${\rm Cov}(X,Y)=0$, on dit que $X$ et $Y$ sont \emph{décorrélées} et on note $X\perp Y$.
\end{definition}

\begin{remarqueUnique}
\remarque Si $X:\Omega\to\R$ est une variable aléatoire, alors
  \[{\rm Cov}(X,X)=\mathbb{V}(X).\]
\end{remarqueUnique}

\begin{proposition}
Soit $X,Y:\Omega\to\R$ deux variables aléatoires réelles. Alors
\[{\rm Cov}(X,Y)=\mathbb{E}(XY)-\mathbb{E}(X)\mathbb{E}(Y).\]
\end{proposition}

\begin{proposition}
Soit $X,Y,Z:\Omega\to\R$ des variables aléatoires réelles.
\begin{itemize}
\item Si $a,b\in\R$, alors
  \begin{eqnarray*}
{\rm Cov}(aX+bY,Z)&=&a{\rm Cov}(X,Z)+b{\rm Cov}(Y,Z)\\
{\rm Cov}(X,aY+bZ)&=&a{\rm Cov}(X,Y)+b{\rm Cov}(X,Z).
  \end{eqnarray*}
\item De plus
  \[{\rm Cov}(X,Y)={\rm Cov}(Y,X).\]
\item Enfin
  \[{\rm Cov}(X,X)=\mathbb{V}(X)\geq 0.\]
\end{itemize}
Autrement dit, la covariance est une forme bilinéaire symétrique positive sur l'espace
vectoriel des variables aléatoires réelles sur $(\Omega,\mathbb{P})$.
\end{proposition}

\begin{remarques}
\remarque Soit $X:\Omega\to\R$ une variable aléatoire réelle et $a\in\R$. Alors
  \[{\rm Cov}(a,X)=0 \quad\et\quad {\rm Cov}(X,a)=0.\]
\remarque La covariance est presque un produit scalaire. La seule propriété qui lui manque
  pour en être un est la propriété de séparation. En effet, si $X:\Omega\to\R$, alors
  ${\rm Cov}(X,X)=0$ si et seulement si $X$ est constante sur un évènement de probabilité 1.
\remarque En particulier, l'inégalité de Cauchy-Schwarz nous dit que  si $X,Y:\Omega\to\R$ sont
  deux variables aléatoires réelles, alors
  \[\abs{{\rm Cov}(X,Y)}\leq\sigma(X)\sigma(Y).\]
  De plus, cette inégalité est une égalité si et seulement si il existe $\alpha,\beta\in\R$
  et un évènement $A$ de probabilité 1 tel que
  \[\cro{\forall \omega\in A\qsep Y(\omega)=\alpha X(\omega)+\beta} \quad\ou\quad
    \cro{\forall \omega\in A\qsep X(\omega)=\alpha Y(\omega)+\beta}.\]
\remarque Si $\sigma(X)>0$ et $\sigma(Y)>0$, le réel
  \[{\rm Cor}(X,Y)\defeq\frac{{\rm Cov}(X,Y)}{\sigma(X)\sigma(Y)},\]
  appelé \emph{coefficient de corrélation}, est dans $[-1,1]$. De plus,
  ${\rm Cor}(X, Y)=1$ si et seulement si il existe $\alpha>0$, $\beta\in\R$ et
  un évènement $A$ de probabilité 1 tel que
  \[\forall \omega\in A\qsep Y(\omega)=\alpha X(\omega)+\beta.\]
  De même ${\rm Cor}(X, Y)=-1$ si et seulement si il existe $\alpha>0$, $\beta\in\R$ et
  un évènement $A$ de probabilité 1 tel que
  \[\forall \omega\in A\qsep Y(\omega)=-\alpha X(\omega)+\beta.\]
\end{remarques}

\begin{proposition}
Soit $X,Y:\Omega\to\R$ deux variables aléatoires réelles. Si $X$ et $Y$ sont
indépendantes, alors elles sont décorrélées.
\end{proposition}

\begin{remarqueUnique}
\remarque La réciproque de la proposition précédente est fausse. Il est possible que
  ${\rm Cov}(X,Y)=0$ sans que $X$ et $Y$ soient indépendantes. Prenons par exemple
  l'exemple d'une variable aléatoire $X$ suivant une loi uniforme sur $\intere{-1}{1}$.
  Alors ${\rm Cov}(X,X^2)=0$. Pourtant, $X$ et $X^2$ ne sont pas indépendantes. 
\end{remarqueUnique}

\begin{proposition}
Soit $X,Y:\Omega\to\R$ deux variables aléatoires. Alors
\[\mathbb{V}(X+Y)=\mathbb{V}(X)+\mathbb{V}(Y)+2{\rm Cov}(X,Y).\]
Plus généralement, si $X_1,\ldots,X_n:\Omega\to\R$ sont des variables aléatoires, alors
\[\mathbb{V}\p{\sum_{i=1}^n X_i}=\sum_{i=1}^n \mathbb{V}(X_i)+
  2\sum_{1\leq i<j\leq n} {\rm Cov}(X_i,X_j).\]
\end{proposition}

\begin{proposition}
Soit $X,Y:\Omega\to\R$ deux variables aléatoires. Si $X$ et $Y$ sont indépendantes, alors
\[\mathbb{V}(X+Y)=\mathbb{V}(X)+\mathbb{V}(Y).\]
Plus généralement, si $X_1,\ldots,X_n:\Omega\to\R$ sont des variables aléatoires deux à deux
indépendantes, alors
\[\mathbb{V}\p{\sum_{i=1}^n X_i}=\sum_{i=1}^n \mathbb{V}(X_i).\]
\end{proposition}

\begin{remarqueUnique}
\remarque Si $X,Y:\Omega\to\R$ sont deux variables aléatoires réelles indépendantes
  \[\sigma(X+Y)=\sqrt{\sigma^2(X)+\sigma^2(Y)}.\]
\end{remarqueUnique}

\begin{exoUnique}
\exo Les \emph{méthodes probabilistes} permettent de prouver certains résultats mathématiques
  dont l'énoncé ne fait pas apparaître des probabilités.
  \begin{questions}
  \question Soit $X:\Omega\to\R$ une variable aléatoire réelle. Montrer qu'il existe
    $\omega\in\Omega$ tel que $X(\omega)\leq\mathbb{E}(X)$.
  \question Soit $v_1,\ldots,v_n\in\R^n$ des vecteurs de norme 1. Montrer qu'il existe
    $\epsilon_1,\ldots,\epsilon_n\in\ens{-1,1}$ tels que
    \[\norme{\epsilon_1 v_1+\cdots+\epsilon_n v_n}\leq\sqrt{n}.\]
  \end{questions}
\end{exoUnique}

\section{Vers les grands nombres}

\begin{proposition}[nom={Inégalité de \nom{Markov}}]
Soit $X:\Omega\to\R$ une variable aléatoire réelle. Alors, quel que soit $a>0$
\[\mathbb{P}(\abs{X}\geq a)\leq\frac{\mathbb{E}(\abs{X})}{a}.\]
\end{proposition}

\begin{remarqueUnique}
\remarque Cette inégalité a le mérite de pouvoir être appliquée sans aucune hypothèse sur la
  loi de la variable aléatoire. Elle est très générale mais assez mauvaise~! Évidemment, cette
  inégalité n'a aucun intérêt si $a\leq \mathbb{E}(\abs{X})$.
% \remarque Il se peut que l'on soit gêné par la valeur absolue dans l'espérance. Dans ce cas,
%   il y a un moyen simple de s'en sortir. Il suffit de constater que les évènements
%   $(\abs{X} \geq a)$ et $(X^2 \geq a^2)$ sont égaux puisque $\abs{X}$ et $a$ sont positifs.
%   On obtient ainsi
%   \[\mathbb{P}(\abs{X} \geq a)=\mathbb{P}(X^2 \geq a^2)\leq\frac{\mathbb{E}(X^2)}{a^2}.\]
\end{remarqueUnique}

\begin{exoUnique}
\exo Soit $X:\Omega\to\R$ une variable aléatoire positive, majorée par $M$. Montrer que
  \[\forall a\geq 0\qsep \mathbb{E}(X)\leq a+M\mathbb{P}(X\geq a).\]
\end{exoUnique}

\begin{proposition}[nom={Inégalité de \nom{Bienaymé-Tchebychev}}]
Soit $X:\Omega\to\R$ une variable aléatoire réelle. Alors, quel que soit $a>0$
\[\mathbb{P}(\abs{X-\mathbb{E}(X)}\geq a)\leq\frac{\mathbb{V}(X)}{a^2}.\]
\end{proposition}

\begin{remarqueUnique}
\remarque On se donne
  $X_1,\ldots,X_n:\Omega\to\R$ des variables aléatoires indépendantes suivant une
  même loi d'espérance $\alpha$ et de variance $\sigma^2$. On pose
  \[M_n\defeq \frac{X_1+\cdots+X_n}{n}.\]
  Alors 
  \[\mathbb{E}(M_n)=\alpha \quad\et\quad \mathbb{V}(M_n)=\frac{\sigma^2}{n}.\]
  En particulier, d'après l'inégalité de \nom{Bienaymé-Tchebychev}, quel que soit
  $\epsilon>0$
  \[\mathbb{P}\p{\abs{M_n-\alpha}\geq\epsilon}\leq\frac{\sigma^2}{n \epsilon^2}
    \tendvers{n}{+\infty}0.\]
  Toute information de ce type est appelée \emph{loi des grands nombres}. Les lois des
  grands nombres sont fondamentales car elles justifient l'approche fréquentiste que l'on a
  des probabilités. 
\end{remarqueUnique}

\begin{exoUnique}
\exo Un parti politique effectue un sondage pour évaluer son image dans l'opinion. La
  proportion de la population qui lui est favorable est $p\in[0,1]$. On interroge un
  échantillon de $n$ personnes. La réponse de la $k$-ième personne sondée est
  représentée par la variable aléatoire $X_k$ qui vaut 1 si la personne est favorable au
  parti et 0 sinon. Chaque variable $X_k$ suit donc une loi de \nom{Bernoulli}
  de paramètre $p$ et on suppose que $X_1,\ldots,X_n$ sont indépendantes.
  La variable aléatoire 
  \[M_n\defeq\frac{X_1+\ldots+X_n}{n}\]
  est utilisée comme un estimateur de $p$.
  \begin{questions}
  \question Montrer que quel que soit $\epsilon>0$
    \[\mathbb{P}(\abs{M_n-p}\geq\epsilon)\leq \frac{1}{4n\epsilon^2}.\]
  \question Soit $\alpha\in]0,1]$. Combien de personnes doit-on interroger pour obtenir une approximation de $p$
    à $\epsilon$ près  avec une probabilité supérieure à $1-\alpha$~? Faites une application
    numérique pour $\epsilon=0.02$ et $\alpha=0.1$.
  \end{questions}
\end{exoUnique}

%END_BOOK

\end{document}



\section{Exercices}
\setcounter{numeroexercice}{1}
\documentclass{magnolia}

\magtex{tex_driver={pdftex}}
\magfiche{document_nom={Exercices sur les variables aléatoires},
          auteur_nom={François Fayard},
          auteur_mail={fayard.prof@gmail.com}}
\magexos{exos_matiere={maths},
         exos_niveau={mpsi},
         exos_theme={Variables aléatoires}}
\magmisenpage{misenpage_presentation={tikzvelvia},
              misenpage_format={a4},
              misenpage_nbcolonnes={1},
              misenpage_sol={non}}
\maglieudiff{lieu_lycee={Aux Lazaristes},
             lieu_classe={MPSI 1},
             lieu_annee={2019--2020}}
\magprocess

\begin{document}
%BEGIN_BOOK
\magsection{Espérance, variance}
\magsubsection{Espérance}

% \exercice{nom={Espérance du maximum}}
% Soit $X$ et $Y$ deux variables aléatoires de loi uniforme sur $\intere{1}{n}$ et indépendantes. On pose $Z \defeq \max (X, Y)$. Déterminer l'espérance de $Z$.

\exercice{nom={Espérance d'une différence}}%colle BCPST
Soit $n \geq 1$. Soit $X$ et $Y$ deux variables aléatoires de loi uniforme sur $\intere{1}{n}$ et indépendantes.
Déterminer la loi et l'espérance de $Z\defeq X-Y$.

\exercice{nom={Nombre de points fixes d'une permutation}}
L'ensemble $S_n$ des permutations de $\intere{1}{n}$ est muni de la probabilité uniforme.
On effectue un tirage aléatoire d'une permutation. Pour $k \in \intere{1}{n}$, on note $X_k$ la variable aléatoire qui vaut $1$ si $k$ est fixe par la permutation et $0$ sinon.
On note $X$ la variable aléatoire égale au nombre de points fixes de la permutation.
Exprimer $X$ en fonction des $X_k$ et en déduire l'espérance de $X$.

\exercice{nom={Exercice}}
Soit $n \in \N ^*$. On dispose de $n$ urnes numérotées de $1$ à $n$. Pour tout
$k\in\intere{1}{n}$, l'urne $k$ contient $k$ boules numérotées de $1$ à $k$. On choisit une
urne au hasard et on tire une boule dans cette urne. On note $X$ le numéro de la boule tirée.
Déterminer la loi de $X$ puis son espérance.

\exercice{nom={Les vaches}}
On considère une population de $2^n$ vaches susceptibles, avec la probabilité $p$, d'être porteuses d'un virus donné.
On dispose d'un test détectant de façon certaine ce virus dans le lait des vaches.
On fixe $0 \leq k \leq n$. On sépare les vaches en $2^{n-k}$ groupes de $2^k$ vaches. On mélange leur lait, on fait un test sur chacun des mélanges, puis on effectue un test sur chacune des vaches des groupes contaminés.
On note $Y$ le nombre de groupes malades et $X$ le nombre total de tests effectués.
\begin{enumerate}
\item Exprimer $X$ en fonction de $Y$, $k$ et $n$.
\item Déterminer la probabilité qu'un groupe donné soit malade.
\item Donner la loi de $Y$ et son espérance.
\item En déduire l'espérance de $X$.
\item On suppose $n=10$ et $p=0.01$. Déterminer la meilleure valeur de $k$.
\end{enumerate}

\exercice{nom={La puce}}
Une puce se déplace (uniquement en avant) sur une bande numérotée par les entiers naturels et commence à la case $0$. À chaque étape, elle fait un bond d'une case avec la probabilité $p$ et de deux cases avec la probabilité $q=1-p$.
On note $X_n$ le numéro de la case où se trouve la puce après $n$ bonds et $Y_n$ le nombre de bonds d'une case effectués après $n$ bonds.
\begin{enumerate}
\item Déterminer la loi de $Y_n$.
\item Exprimer $X_n$ en fonction de $Y_n$.
\item En déduire l'espérance de $X_n$.
\end{enumerate}

\exercice{nom={Exercice}}%\textbf{[Centrale-PSI-2016]} \textbf{[Python]}%RMS
Soit $X_1, ... , X_{n+1}$ des variables aléatoires indépendantes qui suivent
la loi uniforme sur  $[\![1,n]\!]$.
On pose \[T\defeq\min \enstq{j\in [\![2,n+1]\!]}{X_j\in\{X_1,...,X_{j-1}\}}.\]
\begin{enumerate}
\item Vérifier que $T$ est bien défini.
\item %(Question Python facultative, je laisse tel quel l'énoncé de Centrale).
\begin{enumerate}
\item \'Ecrire une fonction Python qui prend $(X_1,...,X_{n+1})$ en arguments 
et qui renvoie $T$.
\item  On choisit $n=1000$. \'Ecrire une fonction qui prend $N$ en argument 
et qui renvoie la moyenne de $T$ au bout de $N$ essais.
\end{enumerate}
\item
\begin{enumerate}
\item Quelles sont les valeurs que peut prendre $T$? 
Déterminer la loi de $T$.
\item Montrer que $\mathbb{E}(T)=\sum_{k=0}^n \mathbb{P}(T>k)$ et en déduire que
\[\mathbb{E}(T)=\frac{n!}{n^n}\sum_{k=0}^n \frac{n^k}{k!}.\]
\item 
En admettant que \[\sum_{k=0}^n \frac{n^k}{k!} \equi{n}{+\infty} \frac{1}{2} \e^n\] donner un équivalent de $\mathbb{E}(T)$.
\end{enumerate}
\end{enumerate}


\exercice{nom={Exercice}}
On considère une urne contenant $n$ boules numérotées de $1$ à $n$. On tire aléatoirement $k$ boules en une seule prise. On note $X$ la variable aléatoire donnant le numéro de la plus petite boule tirée. Calculer $\esp{X}$.

\exercice{nom={Exercice}}
Un questionnaire comporte 20 questions. Pour chaque question, $k$ réponses  sont possibles dont une seule est bonne. Chaque bonne réponse rapporte 1 point. Un candidat répond au hasard à toutes les questions.
\begin{enumerate}
\item Soit $X$ la variable aléatoire donnant le nombre de points obtenus par le candidat à ce questionnaire. Déterminer la loi de $X$.
\item À chaque question, si le candidat s'est trompé, il a droit à une seconde chance et peut choisir une  autre réponse parmi celles qui restent. Il gagne alors 1/2 point en cas de bonne réponse. Soit $Y$ le nombre de 1/2 points obtenus, déterminer la loi de $Y$.
\item Déterminer $k$ pour que le candidat obtienne en moyenne une note de 5 sur 20.
\end{enumerate}

\exercice{nom={Exercice}}
Soit $N$ une variable aléatoire à valeurs dans $\intere{1}{m}$ et $(X_1,\dots,X_m)$ des variables aléatoires à valeurs entières indépendantes entre elles et de $N$ et suivant toutes la même loi.
On pose pour tout $j \in \intere{1}{m}$
\[S_j\defeq\sum\limits_{i=1}^j X_i.\]
On pose enfin
$$Z\defeq\sum\limits_{i=1}^N X_i.$$
\begin{questions}
\question Déterminer l'espérance de $Z$. %et la variance de $Z$.
\question On tire un entier $N$ au hasard et uniformément dans $\intere{1}{n}$
puis on jette $N$ fois une pièce équilibrée. Calculer l'espérance %et la variance
du nombre de piles obtenus.
\end{questions}

\magsubsection{Variance}
\exercice{nom={Exercice}}
On considère $5$ jetons numérotés de $1$ à $5$.
\begin{enumerate}
\item On tire simultanément $2$ jetons parmi les $5$ et on note $X$ la plus
petite valeur.
Déterminer la loi de $X$, son espérance et sa variance.
\item On tire successivement et avec remise $2$ jetons parmi les $5$ et on note $Y$ la plus petite valeur.
Déterminer la loi de $Y$ et son espérance.
\end{enumerate}

\magsection{Couple de variables aléatoires}

\magsubsection{Loi conjointe}
\exercice{nom={Exercice}}
Soit $X$ et $Y$ deux variables aléatoires réelles indépendantes définies sur un même
espace probabilisé fini. On suppose que $X$ et $Y$ sont symétriques, c'est-à-dire que
$X$ et $-X$ (respectivement $Y$ et $-Y$) ont même loi.
\begin{questions}
\question Montrer que $(X,Y)$ et $(X,-Y)$ ont même loi, puis que
  \[\mathbb{P}(X^2=Y^2)=2\mathbb{P}(X=Y)-\mathbb{P}(X=0)\mathbb{P}(Y=0).\]
\question Montrer que $\mathbb{P}(X+Y\geq 0)=\mathbb{P}(X+Y\leq 0)$.
\end{questions}

\magsubsection{Covariance}

\exercice{nom={Exercice}}
On lance deux fois un dé équilibré à 6 faces et on note $D_1$ et $D_2$ les deux résultats. On pose $S\defeq D_1+D_2$ puis $X_1$ (resp. $X_2$) le reste de la division euclidienne de $S$ par $2$ (resp. $5$). 
\begin{questions}
\question Déterminer la loi conjointe du couple $(X_1,X_2)$ (est-elle uniforme ?) et ses lois marginales.
\question  Calculer $\cov {X_1,X_2}$.  $X_1$ et $X_2$ sont-elles indépendantes ?
\end{questions}

\magsection{Vers les grands nombres}

\exercice{nom={Marche aléatoire}}
Soit $X_1,\ldots,X_n$ des variables aléatoires indépendantes et de même loi: $\mathbb{P}(X_k=1)=\mathbb{P}(X_k=-1)=1/2$. On pose $S_n\defeq X_1+\cdots +X_n$. Soit $\epsilon>0$.
\begin{enumerate}
\item Majorer \[\mathbb{P}\left(\left|\frac{S_n}n\right|\geq \epsilon \right).\]
\item Montrer que, pour tout $t\in\R$, $\mathbb{E}\left(\e^{tS_n}\right) = \ch^n t$.
\item Montrer que, pour tout $t>0$, $\ch (t) \leq \e^{t^2/2}$.
\item Montrer que, pour tout $t>0$
\[\mathbb{P}\left(\frac{S_n}n\geq \epsilon\right) \leq \exp\left( \frac{nt^2}2-nt\epsilon\right).\]
\item Montrer que
\[\mathbb{P}\left(\frac{S_n}n\geq \epsilon\right) \leq \exp\left( -\frac{n\epsilon^2}2 \right).\]
\end{enumerate}

% \begin{reponse}
% \begin{enumerate}

% \item $\Omega$ est l'ensemble des parties de cardinal $2$ parmi $5$. Il y
% a donc $10$ tirages possibles qui sont équiprobables.

% $X \p{\Omega}=\{1,2,3,4\}$.

% $(X=1)=\{ \{1,2\}, \{1,3\}, \{1,4\}, \{1,5\}\}$. On a donc $\pr{X=1}=\frac{4}{10}$.

% $(X=2)=\{ \{2,3\}, \{2,4\}, \{2,5\}\}$. On a donc $\pr{X=1}=\frac{3}{10}$.
% Etc.


% \medskip
% \begin{center}
% \begin{tabular}{|c|cccc|}
% \hline
% $k$ &  1 & 2 & 3 & 4  \\

% \hline

% $\pr{X=k}$ & $\frac{4}{10}$ & $\frac{3}{10}$ & $\frac{2}{10}$ & $\frac{1}{10}$ \\


% \hline
% \end{tabular}
% \end{center}

% On a donc $\boxed{\esp{X}=2}$ puis $\boxed{\var{X}}=\esp{X^2}-\esp{X}^2=5-4=\boxed{1}$.



% \item $\Omega$ est l'ensemble des couples d'éléments de $\{1,\dots,5\}$. Il y a donc $25$ tirages possibles qui sont équiprobables.

% $Y \p{\Omega}=\{1,2,3,4,5\}$.

% $(Y=1)=\{ (1,1), (1,2), (2,1), (1,3), (3,1), (1,4), (4,1) , (1,5), (5,1)
% \}$. On a donc $\pr{Y=1}=\frac{9}{25}$.

% $(Y=2)=\{ (2,2), (2,3), (3,2), (2,4), (4,2) , (2,5), (5,2) \}$. On a donc $\pr{Y=2}=\frac{7}{25}$.

% $(Y=3)=\{ (3,3), (3,4), (4,3) , (3,5), (5,3) \}$. On a donc $\pr{Y=3}=\frac{5}{25}$.

% $(Y=4)=\{ (4,4) , (4,5), (5,4) \}$. On a donc $\pr{Y=4}=\frac{3}{25}$.

% $(Y=5)=\{ (5,5) \}$. On a donc $\pr{Y=4}=\frac{1}{25}$.

% \medskip
% \begin{center}
% \begin{tabular}{|c|ccccc|}
% \hline
% $k$ &  1 & 2 & 3 & 4 & 5 \\

% \hline

% $\pr{Y=k}$ & $\frac{9}{25}$ & $\frac{7}{25}$ & $\frac{5}{25}$ & $\frac{3}{25}$
% & $\frac{1}{25}$ \\


% \hline
% \end{tabular}
% \end{center}

% Donc $\boxed{\esp{Y}=\frac{11}{5}}$.

% \end{enumerate}
% \end{reponse}






% Solution
% P(X_1=0,X_2=j) : 3,5,2,5,3 /36
% P(X_1=1,X_2=j) : 4,2,6,2,4 /36
% P(X_1=0) = 1/2
% P(X_1=1) = 1/2
% P(X_2=j) = 7,7,8,7,7 / 36
% Cov(






% \begin{reponse}
% $Z(\Omega)=\intere{-n+1}{n-1}$. Pour $k \in \intere{-n+1}{n-1}$

% $$\pr{Z=k}=\sum\limits_{i=1}^n \pr{X-Y=k,Y=i}=\sum\limits_{i=1}^n \pr{X=k+i,Y=i} $$

% Mais il faut $1 \leq k+i \leq n$ i.e. $-k+1 \leq i \leq n-k$ et en outre $1 \leq i \leq n$.

% \btr Si $k \geq 0$ alors il reste la condition $i \in \intere{1}{n-k}$ donc 

% $$\boxed{\pr{Z=k}}=\sum\limits_{i=1}^{n-k} \pr{X=k+i,Y=i}= \sum\limits_{i=1}^{n-k} \pr{X=k+i} \pr{Y=i} = \frac 1{n^2} \sum\limits_{i=1}^{n-k} 1 = \boxed{\frac{n-k}{n^2}}$$

% \btr Si $k < 0$ alors il reste la condition $i \in \intere{-k+1}{n}$ et on trouve 

% $$\boxed{\pr{Z=k} = \frac{n+k}{n^2}}$$

% ce que l'on pouvait prévoir par un argument de symétrie.

% On a donc une loi en triangle \dots 

% \end{reponse}



% \begin{reponse}

% On $X(\Omega)=\intere{1}{n}$. On note $U$ la variable aléatoire égale au numéro de l'urne tirée.
% On a alors par la formule des probabilités totales
% $$\boxed{\pr{X=k}}=\sum\limits_{i=1}^n \prc{X=k}{U=i} \pr{U=i}=
% \frac 1n \sum\limits_{i=k}^n \prc{X=k}{U=i}=
% \boxed{\frac 1n \sum\limits_{i=k}^n \frac 1i}$$

% On peut alors calculer l'espérance par une interversion
% $$\boxed{\esp{X}}=\sum\limits_{k=1}^n \frac{k}{n} \sum\limits_{i=k}^n \frac 1i
% = \frac 1n \sum\limits_{i=1}^n \frac 1i \sum\limits_{k=1}^{i} k = \boxed{\frac{n+3}{4}}$$

% \end{reponse}





% \begin{reponse}
% On a évidemment $X=\sum X_k$ donc $\esp{X}=\sum \esp{X_k}$. Mais chaque $X_k$ est une variable de Bernoulli de paramètre $p_k=\pr{X_k=1}$.

% Mais $\pr{X_k=1}=\frac 1n$ (en effet, le nombre de permutations qui fixent $k$ est $(n-1)!$ puisqu'il reste à définir une permutation sur les $n-1$ autres éléments de $\intere{1}{n}$).

% On a donc $\esp{X}=\sum\limits_{k=1}^n \frac 1n$ donc $\boxed{\esp{X}=1}$.
% \end{reponse}




%\exercice{nom={Exercice}}
%Soit $(X_n)_{n \in \N^*}$ une suite de variables aléatoires réelles définies sur un même espace probabilisé, indépendantes, de même espérance $\mu$ et de même variance $\sigma^2$. On pose 
%$\overline{X_n} = \frac 1n \p{X_1 + \dots + X_n}$ et $S_n = \frac 1{n-1} \sum\limits_{k=1}^n \p{X_k - \overline{X_n}}^2$.
%\begin{enumerate}
%\item Calculer l'espérance et la variance de $\overline{X_n}$.
%\item Calculer l'espérance de $S_n$. Rien ne vous surprend ?
%\end{enumerate}
%




% \begin{reponse}
% \begin{enumerate}

% \item $X=2^{n-k}+ 2^k Y$.

% \item Un groupe est sain ssi chaque vache du groupe est saine ce qui donne $1-(1-p)^{2^k}$.

% \item $Y$ suit une loi binomiale de paramètres $2^{n-k}, 1-(1-p)^{2^k}$ donc $\esp{Y}=2^{n-k} \left( 1-(1-p)^{2^k} \right)$.

% \item $\esp{X}=2^{n-k}+ 2^k 2^{n-k} \left( 1-(1-p)^{2^k} \right)=2^{n-k}+ 2^n \left( 1-(1-p)^{2^k} \right)$.

% \item Maple donne $k=3$ et $\esp{X} \sim 207$ (au lieu de $1024$).

% \end{enumerate}
% \end{reponse}




% \begin{reponse}
% \begin{enumerate}

% \item $Y_n$ suit une loi binomiale de paramètres $n,p$. 

% \item $X_n=Y_n + 2 (n-Y_n)$.

% \item $\esp{Y_n}=np$ donc $\boxed{\esp{X_n}=n(2-p)}$.

% \end{enumerate}
% \end{reponse}








% \begin{reponse}

% \begin{enumerate}

% \item 

% Il y a $n+1$ valeurs dans un ensemble de cardinal $n$ donc deux valeurs sont égales (principe des tiroirs) donc il existe $i<j$ tels que $X_i = X_j$ donc $X_j \in \{X_1,...,X_{j-1}\}$. $T$ est donc bien défini.

% \item

% \begin{enumerate}

% \item Il faut bien entendu comprendre que la fonction prend en arguments les valeurs et pas les variables aléatoires \dots

% On code de manière naïve pour aller vite (j'utilise \verb+in+ et on a une complexité quadratique) mais on sait que l'on peut procéder par mémoïzation pour obtenir une complexité linéaire. On suppose que les valeurs sont rangées dans une liste $L$ :

% \begin{lstlisting}
% def T(L):
    
%     n = len(L)
    
%     for j in range(1, n):
%         if L[j] in L[:j]:
%             return j+1
% \end{lstlisting}

% Attention au petit décalage d'indice (on renvoie $j+1$) pour respecter l'énoncé où les variables sont numérotées à partir de $1$.

% \item  On utilise

% \begin{lstlisting}
% import numpy.random as rd
% \end{lstlisting}

% puis \verb+rd.randint(a, b, n)+ qui renvoie une liste de $n$ valeurs aléatoires entières de $[a,b[$ (attention, $b$ est exclu) selon la loi uniforme (voir la fiche probas) :

% \begin{lstlisting}
% n = 1000

% def moyenne(N):
    
%     s = 0
    
%     for _ in range(N):
%         L = rd.randint(1, n+1, n+1)
%         s += T(L)
    
%     return s / N
% \end{lstlisting}

% \begin{lstlisting}
% In [1]: moyenne(10**4)
% Out[1]: 40.3642
% \end{lstlisting}

% En moyenne, il faut $40$ tirages d'un nombre entre $1$ et $1000$ pour retrouver un nombre déjà tiré.

% \end{enumerate}

% \item

% \begin{enumerate}

% \item $T$ peut prendre toutes les valeurs entre $2$ et $n+1$ : il suffit de considérer le tirage $(1,2,\dots,k-1, k-1, \dots)$ pour lequel $T=k$.

% Pour déterminer la loi de $T$, il est préférable de suivre l'indication (bien classique) de la question suivante. On calcule donc $\pr{T>k}$ pour $k \in \intere{0}{n}$.

% L'évènement $(T>k)$ signifie que les $k$ premiers tirages ont amené des nombres deux à deux distincts. Puisque le tirage est uniforme, on peut dénombrer. Sur les $k$ premiers nombres, on a $n^k$ tirages possibles et $n(n-1) \dots (n-k+1)$ tirages favorables donc
% $$\boxed{\pr{T>k} = \frac{n(n-1) \dots (n-k+1)}{n^k} = \frac{n!}{(n-k)!n^k}}$$

% Pour répondre à la question (inutile pour le calcul de l'espérance, mais bon \dots) :

% $$\pr{T=k} = \pr{T>k-1} - \pr{T>k} = \frac{n!}{(n-k+1)!n^{k-1}} - \frac{n!}{(n-k)!n^k}$$

% que l'on simplifie \dots

% \item La formule de l'espérance est classique. Elle donne alors

% $$\boxed{{\rm E}(T)} = \sum_{k=0}^n P(T>k) = \sum_{k=0}^n \frac{n!}{(n-k)!n^k} = n! \sum_{k=0}^n \frac{1}{(n-k)!n^k}  = n! \sum_{k=0}^n \frac{1}{k!n^{n-k}}  = \boxed{\frac{n!}{n^n} \sum_{k=0}^n \frac{n^k}{k!}}$$

% où l'on a changé d'indice en posant $k=n-k$.

% \medskip

% Test numérique (calcul brutal) :

% \begin{lstlisting}
% from math import factorial

% def esperance(n):
    
%     return (factorial(n)/n**n)*sum(n**k/factorial(k) for k in range(n+1))
% \end{lstlisting}

% avec $n=100$ (sinon, problème de taille des factorielles) :

% \begin{lstlisting}
% In [2]: esperance(n)
% Out[2]: 13.209960630215981

% In [3]: moyenne(10**4)
% Out[3]: 13.144
% \end{lstlisting}

% \begin{lstlisting}

% \end{lstlisting}

% \end{enumerate}

% \textbf{Remarque} : on peut montrer que $\sum_{k=0}^n \frac{n^k}{k!} \sim \frac 12 e^n$ (voir par exemple l'exercice 970 de la RMS, Centrale-MP-2016-Python). En utilisant ceci et la formule de Stirling, on a donc
% $$\boxed{{\rm E}(T) \sim \sqrt{\frac{\pi n}{2}}}$$

% \end{enumerate}

% \end{reponse}





% \begin{reponse}
% $X\p{\Omega} \subset \intere{1}{n-k+1}$.

% Pour $i \in \intere{1}{n-k+1}$, $\pr{X \geq i} = \frac{\binom{n-i+1}{k}}{\binom{n}{k}}$ donc 

% $$\esp{X} = \sum\limits_{i=1}^{n-k+1} \pr{X \geq i} = \frac{1}{\binom{n}{k}} \sum\limits_{i=1}^{n-k+1} \binom{n-i+1}{k} = \frac{\binom{n+1}{k+1}}{\binom{n}{k}} = \frac{n+1}{k+1}$$
% ce que l'on vérifie pour $k=1$ et $k=n$ \dots
% \end{reponse}




% \begin{reponse}

% $X \hookrightarrow {\cal B}(20, \frac{1}{k})$.

% Quand $X=i$, $Y \hookrightarrow {\cal B}(20-i, \frac{1}{k-1})$.

% On note $Z$ le nombre de points : $Z = X + \frac{1}{2}Y$.

% Sans calculer la loi (trop long)

% $$\esp{Z} = \sum\limits_{n} n \pr{Y=n} = \sum\limits_{n} n \sum\limits_{i} \pr{X=i} \prc{Y=n}{X=i} = \sum\limits_{i}  \pr{X=i} \sum\limits_{n} n \prc{Y=n}{X=i}  $$

% La somme interne est l'espérance de $Y$ sachant $X=i$ donc c'est $\frac{20-i}{k-1}$ donc

% $$\esp{Z} = \sum\limits_{i}  \pr{X=i} \frac{20-i}{k-1} = \frac{1}{k-1} \p{20 - \esp{X}}$$

% Comme $\esp{X} = \frac{20}{k}$ il reste $\boxed{\esp{Z} = \frac{30}{k}}$.
% \end{reponse}






% \begin{reponse}

% \btr Commençons par calculer l'espérance :

% \medskip

% $\esp{Z}=\sum\limits_{k=0}^{\infty} k \pr{Z=k}$ (cette somme est en fait finie). Mais
% $$\pr{Z=k}=\sum\limits_{j=1}^m \prc{Z=k}{N=j} \pr{N=j}$$

% $\prc{Z=k}{N=j}=\prc{S_j=k}{N=j}=\pr{S_j=k}$ puisque $N$ est indépendante des $X_i$. On a donc

% $$\pr{Z=k}=\sum\limits_{j=1}^m \pr{S_j=k} \pr{N=j}$$


% Alors

% $$\esp{Z}= \sum\limits_{k=0}^{\infty} \sum\limits_{j=1}^m k \pr{S_j=k} \pr{N=j} = 
% \sum\limits_{j=1}^m  \p{\sum\limits_{k=0}^{\infty} k \pr{S_j=k} }\pr{N=j} = 
% \sum\limits_{j=1}^m \esp{S_j} \pr{N=j}$$
% Mais $\esp{S_j}=j \esp{X}$ (où l'on note $\esp{X}$ l'espérance commune des $X_i$) puisque les $X_i$ ont toutes la même espérance. On a donc

% $$\boxed{\esp{Z}}=\p{\sum\limits_{j=1}^m j \pr{N=j}} \esp{X} = \boxed{\esp{N} \esp{X}}$$


% \btr Calculons maintenant la variance :

% \medskip


% $\var{Z}=\esp{Z^2}-\esp{Z}^2$. Il reste donc à calculer l'espérance de $Z^2$. On procède comme avant ($k$ est remplacé par $k^2$) :

% $$\esp{Z^2}=\sum\limits_{j=1}^m \esp{S_j^2} \pr{N=j}=\sum\limits_{j=1}^m \p{\var{S_j} + \esp{S_j}^2} \pr{N=j}$$

% Comme les $X_i$ sont indépendantes, $\var{S_j}=j \var{X}$ (toujours en notant $\var{X}$ la variance commune des $X_i$). Comme $\esp{S_j}^2=j^2 \esp{X}^2$, on a 

% $$\esp{Z^2}=\sum\limits_{j=1}^m \p{ j \var{X} + j^2 \esp{X}^2} \pr{N=j} = \var{X} \esp{N} + \esp{X}^2 \esp{N^2}$$

% Mais $\esp{N^2}=\var{N}+\esp{N}^2$ donc 
% $$\esp{Z^2}=\var{X} \esp{N} + \esp{X}^2 \p{\var{N}+\esp{N}^2} = \var{X} \esp{N} + \esp{X}^2 \var{N}+\esp{X}^2 \esp{N}^2$$

% On termine alors avec $\var{Z}=\esp{Z^2}-\esp{Z}^2= \var{X} \esp{N} + \esp{X}^2 \var{N}+\esp{X}^2 \esp{N}^2 - \esp{X}^2 \esp{N}^2$ donc

% $$\boxed{\var{Z}=\var{X} \esp{N} + \esp{X}^2 \var{N}}$$

% \textbf{Remarque :} c'est homogène puisque $N$ est sans unité.

% \medskip

% \textbf{Application :} la variable $N$ suit une loi uniforme sur $\intere{1}{n}$ donc $\esp{N}=\frac{n+1}{2}$ et $\var{N}=\frac{n^2-1}{12}$.

% Si l'on note $X_i=1$ si le $i$-ième lancer de dé apporte un pile et $0$ sinon, alors les $X_i$ sont indépendantes et suivent toutes la même loi de Bernoulli de paramètre $\frac 12$. On a donc $\esp{X}=\frac 12$ et $\var{X}=\frac 12 \p{1 - \frac 12}=\frac 14$.


% Comme le nombre de piles obtenus est égal à $Z=\sum\limits_{i=1}^N X_i$ et que $N$ est indépendante des $X_i$, on a d'après les calculs précédents,

% $$\boxed{\esp{Z}=\frac{n+1}{4} \text{ et } \var{Z}=\frac{(n+1)(n+5)}{48}}$$
% \end{reponse}



%END_BOOK

\end{document}


















\chapter{Déterminants}
\setcounter{numeroexercicecours}{1}
\documentclass{magnolia}

\magtex{tex_driver={pdftex},
        tex_packages={xypic}}
\magfiche{document_nom={Cours sur les déterminants},
          auteur_nom={François Fayard},
          auteur_mail={fayard.prof@gmail.com}}
\magcours{cours_matiere={maths},
          cours_niveau={mpsi},
          cours_chapitre_numero={21},
          cours_chapitre={Déterminants}}
\magmisenpage{misenpage_presentation={tikzvelvia},
          misenpage_format={a4},
          misenpage_nbcolonnes={1},
          misenpage_preuve={non},
          misenpage_sol={non}}
% \magmisenpage{}
\maglieudiff{}
\magprocess
\begin{document}

%BEGIN_BOOK
\magtoc


\section{Déterminant}
\subsection{Forme $n$-linéaire alternée}

\begin{definition}
Soit $E_1,\ldots,E_n$ et $F$ des \Kevs. On dit qu'une application $\phi:E_1\times\cdots\times E_n\to F$ est
\emph{$n$-linéaire} lorsqu'elle est linéaire par rapport à chacune des variables, c'est-à-dire lorsque
quel que soit $i\in\intere{1}{n}$ et $(x_1,\ldots,x_{i-1},x_{i+1},\ldots,x_n)\in E_1\times\cdots\times E_{i-1}\times E_{i+1}\times\cdots\times E_n$, on a
\[\forall x,y\in E_i\qsep \forall \lambda,\mu\in\K\qsep \phi\p{x_1,\ldots,x_{i-1},\lambda x+\mu y,x_{i+1},\ldots,x_n}=\]
\[\lambda \phi\p{x_1,\ldots,x_{i-1},x,x_{i+1},\ldots,x_n}+
\mu \phi\p{x_1,\ldots,x_{i-1},y,x_{i+1},\ldots,x_n}.\]
Si $F=\K$, on dit que \emph{$\phi$ est une forme $n$-lineaire}.
\end{definition}


% \begin{definition}
% On dit qu'une application $\phi$ de $E^n$ dans $\K$ est une \emph{forme $n$-linéaire}
% lorsque $\phi$ est linéaire par rapport à chacune de ses variables
% \[\forall i\in\intere{1}{n} \qsep \forall x_1,\ldots,x_{i-1},x_{i+1},\ldots,x_n\in E
%   \qsep \forall x,y\in E \qsep \forall \lambda,\mu\in\K,\]
% \[\]
% \[\]
% Muni des lois usuelles, l'ensemble des formes $n$-linéaires sur $E$ est un \Kev.
% \end{definition}

\begin{remarques}
\remarque Lorsque $n=2$, on parle d'application $\emph{bilinéaire}$.
\remarque L'ensemble des applications $n$-linéaires de $E_1\times\cdots\times E_n$ dans $F$ est
  un sous-espace vectoriel de l'ensemble des applications de $E_1\times\cdots\times E_n$ dans $F$.
\remarque Si $\phi:E_1\times\cdots\times E_n\to F$ est $n$-linéaire, $i\in\intere{1}{n}$ et
$(x_1,\ldots,x_{i-1},x_{i+1},\ldots,x_n)\in E_1\times\cdots\times E_{i-1}\times E_{i+1}\times\cdots\times E_n$, alors
  $\phi(x_1,\ldots,x_{i-1},0,x_{i+1},\ldots,x_n)=0$.
\remarque On dit qu'une application $\phi$ est une \emph{forme $n$-linéaire sur $E$} lorsque c'est une application $n$-linéaire de
  $E^n$ dans $\K$.
\end{remarques}

\begin{exemples}
\exemple Si $X$ est un ensemble, l'application
  \[\dspappli{\phi}{\mathcal{F}(X,\K)\times\mathcal{F}(X,\K)}{\mathcal{F}(X,\K)}{(f,g)}{fg}\]
  est bilinéaire.
\exemple Si $E,F,G$ sont des \Kevs, l'application
  \[\dspappli{\phi}{\mathcal{L}(E,F)\times\mathcal{L}(F,G)}{\mathcal{L}(E,G)}{(f,g)}{g\circ f}\]
  est bilinéaire. De manière similaire, si $r,q,p\in\N$, l'application
  \[\dspappli{\phi}{\mat{r,q}{\K}\times\mat{q,p}{\K}}{\mat{r,p}{\K}}{(A,B)}{AB}\]
  est bilinéaire.
\exemple Soit $\phi_1,\ldots,\phi_n$ des formes linéaires sur $E$. Alors
  \[\dspappli{\phi}{E^n}{\K}{(x_1,\ldots,x_n)}{\prod_{k=1}^n \phi_k(x_k)}\]
  est une forme $n$-linéaire sur $E$.
\exemple Dans le plan euclidien, le produit scalaire est bilinéaire. Dans l'espace euclidien orienté, le
  produit scalaire et le produit vectoriel sont bilinéaires.
\end{exemples}

\begin{definition}
On dit qu'une forme $\phi$, $n$-linéaire sur $E$, est \emph{alternée} lorsqu'elle est nulle sur toute
famille de vecteurs dont au moins deux sont égaux.
\end{definition}

\begin{remarques}
\remarque Autrement dit, la forme $n$-linéaire $\phi:E^n\to\K$ est alternée lorsque
  \[\forall x_1,\ldots,x_n\in E\qsep \forall i,j\in\intere{1}{n}\qsep \cro{i\neq j \et x_i=x_j}
    \quad\implique\quad \phi\p{x_1,\ldots,x_n}=0.\]
\remarque L'ensemble des formes $n$-linéaires alternées sur $E$ est noté $\Lambda_n(E)$. C'est
  un sous-espace vectoriel de l'espace des formes $n$-linéaires sur $E$.
\end{remarques}

\begin{proposition}
Soit $\phi$ une forme $n$-linéaire alternée sur $E$. Alors, on ne change pas la valeur de
$\phi(x_1,\ldots,x_n)$ en ajoutant à l'une de ses variables une combinaison linéaire des autres.
\end{proposition}

\begin{remarqueUnique}
\remarque En particulier, si $(x_1,\ldots,x_n)$ est liée, alors $\phi(x_1,\ldots,x_n)=0$.
\end{remarqueUnique}

\begin{proposition}
Soit $\phi$ une forme $n$-linéaire alternée sur $E$. Alors
\[\forall x_1,\ldots,x_n\in E \qsep \forall \sigma\in\gsym{n} \qsep
  \phi\p{x_{\sigma\p{1}},\ldots,x_{\sigma(n)}}=\signature(\sigma)
  \phi\p{x_1,\ldots,x_n}.\]
On dit que $\phi$ est \emph{antisymétrique}.
\end{proposition}


\begin{remarqueUnique}
\remarque Une forme $n$-linéaire alternée est donc antisymétrique. Réciproquement, si $\K$ n'est pas de
  caractéristique 2, les formes antisymétriques sont alternées.
\end{remarqueUnique}
  
\begin{preuve}
Commençons par montrer ce résultat pour une transposition
$$\forall x_1,\ldots,x_n\in E \quad \forall \tau\in\gsym{n} \quad
  \phi\p{x_{\tau\p{1}},\ldots,x_{\tau(n)}}=-
  \phi\p{x_1,\ldots,x_n}.$$
  
Considérons $\tau=(i,j)$ une transposition. Supposons par exemple $i<j$.
Comme $\phi$ est alternée, on a :
$$\phi(x_1,\ldots,\underbrace{x_i+x_j}_{i\text{-ième position}},\ldots,\underbrace{x_i+x_j}_{j\text{-ième position}},\ldots,x_n)=0.$$
Mais alors, par linéarité par linéarité par rapport à la $i$-ième variable et par linéarité par rapport à la $j$-ième variable :
$$\underbrace{\phi(x_1,\ldots,x_i,\ldots,x_i,\ldots,x_n)}_{=0 \text{ car } \phi \text{ est alternée }}+\phi(x_1,\ldots,x_i,\ldots,x_j,\ldots,x_n)+\phi(x_1,\ldots,x_j,\ldots,x_i,\ldots,x_n)+\underbrace{\phi(x_1,\ldots,x_j,\ldots,x_j,\ldots,x_n)}_{=0 \text{ car } \phi \text{ est alternée }}=0$$ d'où
$$\phi(x_1,\ldots,x_i,\ldots,x_j,\ldots,x_n)=-\phi(x_1,\ldots,x_j,\ldots,x_i,\ldots,x_n).$$

Pour le cas général, on écrit $\sigma=\tau_1\circ \ldots \circ \tau_m$ en produit de transpositions. On a :

\begin{eqnarray*}
\phi(x_{\sigma(1)},\ldots,x_{\sigma(n)})&=&\phi(x_{\tau_1\circ \ldots \circ \tau_m(1)},\ldots,x_{\tau_1\circ \ldots \circ \tau_m(n)})\\
&=& \epsilon(\tau_1)\phi(x_{\tau_2\circ \ldots \circ \tau_m(1)},\ldots,x_{\tau_2\circ \ldots \circ \tau_m(n)})\\
&=&\ldots \\
&=&\epsilon(\tau_1)\ldots\epsilon(\tau_m)\phi(x_1,\ldots,x_n)\\
&=&\epsilon(\tau_1\circ \ldots \circ \tau_m)\phi(x_1,\ldots,x_n)\\
&=&\signat{\sigma}
  \phi\p{x_1,\ldots,x_n}.
  \end{eqnarray*}



\end{preuve}

\begin{sol}
Une remarque cruciale (voir proposition) à rajouter ici :
Soit $\phi$ une forme $n$-linéaire alternée sur $E$ et $(x_1,\ldots,x_n)\in E^n$. Alors, $\phi(x_1,\ldots,x_n)$ est inchangé si l'on ajoute à l'un des $x_i$ une combinaison linéaire des autres $x_k$.

En effet, fixons $i \in \intere{1}{n}$ et des $\lambda_k$ pour $k\neq i$ dans $\K$. Par linéarité par rapport à la $i$-ième variable :
$$\phi(x_1,\ldots,x_{i-1},x_i+\sum_{k\neq i}\lambda_k x_k,x_{i+1},\ldots,x_n)=\phi(x_1,\ldots,x_n)+\sum_{k\neq i}\lambda_k \underbrace{\phi(x_1,\ldots,x_k,\ldots,x_n)}_{=0 \text{ car } x_k \text{ est présent deux fois}}.$$
\end{sol}

\begin{proposition}
Soit $E$ un \Kev de dimension $n$, $\mathcal{B}\defeq(e_1,\ldots,e_n)$ une base de
$E$ et $\phi$ une forme $n$-linéaire alternée sur $E$. Alors
\[\forall x_1,\ldots,x_n\in E\qsep \phi\p{x_1,\ldots,x_n}=\p{\sum_{\sigma\in\gsym{n}} \signat{\sigma} \prod_{k=1}^n e_{\sigma(k)}^{\star}(x_k)}
  \phi\p{e_1,\ldots,e_n}.\]
\end{proposition}

\begin{remarqueUnique}
\remarque En particulier, si $\phi$ et $\psi$ sont deux formes $n$-linéaires alternées sur un \Kev de dimension
  $n$ qui prennent la même valeur sur une base de $E$, alors $\phi=\psi$.
\end{remarqueUnique}

\begin{preuve}
On considère donc $\mathcal{B}=(e_1,\ldots,e_n)$ une base de $E$. 
Soit $x_1,\ldots,x_n\in E$. Notons $A=\mat{\mathcal{B}}{x_1,\ldots,x_n}$. Cela signifie que :
  \[\forall j\in\intere{1}{n} \quad x_j=\sum_{i=1}^n a_{i,j} e_i.\]
  On a alors :
  \begin{eqnarray*}
  \phi\p{x_1,x_2,\ldots,x_n}
  &=& \phi\p{\sum_{i_1=1}^n a_{i_1,1} e_{i_1},\sum_{i_2=1}^n a_{i_2,2} e_{i_2},\ldots,
             \sum_{i_n=1}^n a_{i_n,n} e_{i_n}}\\
  &=& \sum_{i_1=1}^n a_{i_1,1} \phi\p{e_{i_1},\sum_{i_2=1}^n a_{i_2,2} e_{i_2},\ldots,
             \sum_{i_n=1}^n a_{i_n,n} e_{i_n}}\\
  & & \text{par linéarité de $\phi$ par rapport à la première variable}\\
  &=& \sum_{i_1=1}^n \sum_{i_2=1}^n \cdots \sum_{i_n=1}^n
       a_{i_1,1} a_{i_2,2} \cdots a_{i_n,n} \phi\p{e_{i_1},e_{i_2},\ldots,e_{i_n}}\\
  & & \text{car $\phi$ est $n$-linéaire}\\
  &=& \sum_{\p{i_1,i_2,\ldots,i_n}\in\intere{1}{n}^n}
       a_{i_1,1} a_{i_2,2} \cdots a_{i_n,n} \phi\p{e_{i_1},e_{i_2},\ldots,e_{i_n}}
  \end{eqnarray*}
Or se donner un $n$-uplet $\p{i_1,i_2,\ldots,i_n}\in\intere{1}{n}^n$, c'est exactement se donner une application $\dspappli{f}{\intere{1}{n}}{\intere{1}{n}}{k}{i_k}$. On obtient alors :
  
$$\phi\p{x_1,x_2,\ldots,x_n}= \sum_{f\in\mathcal{F}\p{\intere{1}{n},\intere{1}{n}}}
       a_{f\p{1},1} a_{f\p{2},2} \cdots a_{f(n),n}
       \phi\p{e_{f\p{1}},e_{f\p{2}},\ldots,e_{f(n)}}.$$
       
       Si $f\in\mathcal{F}\p{\intere{1}{n},\intere{1}{n}}$ n'est pas injective, alors il
  existe $i,j\in\intere{1}{n}$ tels que $i\neq j$ et $f(i)=f(j)$. Comme $\phi$
  est alternée, on en déduit que $\phi\p{e_{f\p{1}},e_{f\p{2}},\ldots,e_{f(n)}}=0$.
  On peut donc restreindre la somme aux fonctions injectives de $\intere{1}{n}$
  dans lui-même. L'ensemble $\intere{1}{n}$ étant fini, ces fonctions sont les
  permutations de $\intere{1}{n}$. Donc~:
  \begin{eqnarray*}
  \phi\p{x_1,x_2,\ldots,x_n}
  &=& \sum_{\sigma\in\gsym{n}}
       a_{\sigma\p{1},1} a_{\sigma\p{2},2} \cdots a_{\sigma(n),n}
       \phi\p{e_{\sigma\p{1}},e_{\sigma\p{2}},\ldots,e_{\sigma(n)}}\\
  &=& \sum_{\sigma\in\gsym{n}}
       a_{\sigma\p{1},1} a_{\sigma\p{2},2} \cdots a_{\sigma(n),n} \signat{\sigma} 
       \phi\p{e_1,e_2,\ldots,e_n}\\
       &=& \underbrace{\cro{\sum_{\sigma\in\gsym{n}} \signat{\sigma}
  a_{\sigma\p{1},1}a_{\sigma\p{2},2}\cdots a_{\sigma(n),n}}}_{\text{indépendant de } \phi}
  \phi\p{e_1,\ldots,e_n}
  \end{eqnarray*}


\end{preuve}

\begin{theoreme}
Soit $E$ un \Kev de dimension $n$ et $\mathcal{B}\defeq\p{e_1,\ldots,e_n}$ une base de
$E$. Alors, il existe une unique forme $n$-linéaire alternée $\phi$ sur $E$ telle
que $\phi\p{e_1,\ldots,e_n}=1$.
% De plus, si $x_1,\ldots,x_n\in E$ est une famille de $n$ vecteurs et
% $A=\mat{\mathcal{B}}{x_1,\ldots,x_n}$, on a~:
% \[\phi\p{x_1,\ldots,x_n}=\sum_{\sigma\in\gsym{n}} \signat{\sigma}
%   a_{\sigma\p{1},1}a_{\sigma\p{2},2}\cdots a_{\sigma(n),n}\]
\end{theoreme}

\begin{preuve}
$\quad$
\begin{itemize}
\item Montrons l'unicité d'une telle application $\phi$. Soit $x_1,\ldots,x_n\in E$
  et $A=\mat{\mathcal{B}}{x_1,\ldots,x_n}$. Alors~:
  \[\forall j\in\intere{1}{n} \quad x_j=\sum_{i=1}^n a_{i,j} e_i\]
  Donc~:
  \begin{eqnarray*}
  \phi\p{x_1,x_2,\ldots,x_n}
  &=& \phi\p{\sum_{i_1=1}^n a_{i_1,1} e_{i_1},\sum_{i_2=1}^n a_{i_2,2} e_{i_2},\ldots,
             \sum_{i_n=1}^n a_{i_n,n} e_{i_n}}\\
  &=& \sum_{i_1=1}^n a_{i_1,1} \phi\p{e_{i_1},\sum_{i_2=1}^n a_{i_2,2} e_{i_2},\ldots,
             \sum_{i_n=1}^n a_{i_n,n} e_{i_n}}\\
  & & \text{par linéarité de $\phi$ par rapport à la première variable}\\
  &=& \sum_{i_1=1}^n \sum_{i_2=1}^n \cdots \sum_{i_n=1}^n
       a_{i_1,1} a_{i_2,2} \cdots a_{i_n,n} \phi\p{e_{i_1},e_{i_2},\ldots,e_{i_n}}\\
  & & \text{car $\phi$ est $n$-linéaire}\\
  &=& \sum_{\p{i_1,i_2,\ldots,i_n}\in\intere{1}{n}^n}
       a_{i_1,1} a_{i_2,2} \cdots a_{i_n,n} \phi\p{e_{i_1},e_{i_2},\ldots,e_{i_n}}\\
  &=& \sum_{f\in\mathcal{F}\p{\intere{1}{n},\intere{1}{n}}}
       a_{f\p{1},1} a_{f\p{2},2} \cdots a_{f(n),n}
       \phi\p{e_{f\p{1}},e_{f\p{2}},\ldots,e_{f(n)}}
  \end{eqnarray*}
  Si $f\in\mathcal{F}\p{\intere{1}{n},\intere{1}{n}}$ n'est pas injective, alors il
  existe $i,j\in\intere{1}{n}$ tels que $i\neq j$ et $f(i)=f(j)$. Comme $\phi$
  est alternée, on en déduit que $\phi\p{e_{f\p{1}},e_{f\p{2}},\ldots,e_{f(n)}}=0$.
  On peut donc restreindre la somme aux fonctions injectives de $\intere{1}{n}$
  dans lui-même. L'ensemble $\intere{1}{n}$ étant fini, ces fonctions sont les
  permutations de $\intere{1}{n}$. Donc~:
  \begin{eqnarray*}
  \phi\p{x_1,x_2,\ldots,x_n}
  &=& \sum_{\sigma\in\gsym{n}}
       a_{\sigma\p{1},1} a_{\sigma\p{2},2} \cdots a_{\sigma(n),n}
       \phi\p{e_{\sigma\p{1}},e_{\sigma\p{2}},\ldots,e_{\sigma(n)}}\\
  &=& \sum_{\sigma\in\gsym{n}}
       a_{\sigma\p{1},1} a_{\sigma\p{2},2} \cdots a_{\sigma(n),n} \signat{\sigma} 
       \underbrace{\phi\p{e_1,e_2,\ldots,e_n}}_{=1}\\
  & & \text{car $\phi$ est alternée donc antisymétrique}\\
  &=& \sum_{\sigma\in\gsym{n}}
      \signat{\sigma} a_{\sigma\p{1},1} a_{\sigma\p{2},2} \cdots a_{\sigma(n),n}
  \end{eqnarray*}
\item Montrons désormais l'existence d'une telle application. On définit $\phi$ de
  $E^n$ dans $\K$ par~:
  \[\forall x_1,\ldots,x_n \quad
    \phi\p{x_1,\ldots,x_n}=\sum_{\sigma\in\gsym{n}} \signat{\sigma}
    a_{\sigma\p{1},1}a_{\sigma\p{2},2}\cdots a_{\sigma(n),n}\]
  où $A=\mat{\mathcal{B}}{x_1,\ldots,x_n}$. Montrons que $\phi$ est $n$-linéaire
  alternée et que $\phi\p{e_1,\ldots,e_n}=1$.\\
  Pour cela, on définit, pour toute permutation $\sigma\in\gsym{n}$, l'application
  $\phi_\sigma$ de $E^n$ dans $\K$ par~:
  \[\forall x_1,\ldots,x_n \quad \phi_\sigma \p{x_1,\ldots,x_n}=
    a_{\sigma\p{1},1}a_{\sigma\p{2},2}\cdots a_{\sigma(n),n}\]
  où $A=\mat{\mathcal{B}}{x_1,\ldots,x_n}$. On remarque que~:
  \[\phi=\sum_{\sigma\in\gsym{n}} \signat{\sigma} \phi_{\sigma}\]
  \begin{itemize}
  \item Montrons que $\phi$ est $n$-linéaire. Pour cela, il suffit de montrer que
    $\phi_\sigma$ est $n$-linéaire pour tout $\sigma\in\gsym{n}$. Soit
    $\sigma\in\gsym{n}$, $i\in\intere{1}{n}$,
    $x_1,\ldots,x_{i-1},x_{i+1},\ldots,x_n\in E$, $x,y\in E$ et $\lambda,\mu\in\K$.
    Pour tout $j\in\intere{1}{n}\setminus\ens{i}$, on décompose $x_j$ dans la
    base $\mathcal{B}$~:
    \[x_j=\sum_{i=1}^n a_{i,j} e_i\]
    De même, on décompose $x$ et $y$ dans la base $\mathcal{B}$~:
    \[x=\sum_{i=1}^n x_i e_i \et y=\sum_{i=1}^n y_i e_i\]
    Alors~:
    \[\lambda x+\mu y=\sum_{i=1}^n \p{\lambda x_i+\mu y_i} e_i\]
    Donc~:
    \begin{eqnarray*}
    & & \phi_\sigma \p{x_1,\ldots,x_{i-1},\lambda x+\mu y,x_{i+1},\ldots,x_n}\\
    &=& a_{\sigma\p{1},1}\cdots a_{\sigma\p{i-1},i-1} \p{\lambda x_{\sigma(i)}
           +\mu y_{\sigma(i)}} a_{\sigma\p{i+1},i+1} \cdots a_{\sigma(n),n}\\
    &=& \lambda a_{\sigma\p{1},1}\cdots a_{\sigma\p{i-1},i-1} x_{\sigma(i)}
        a_{\sigma\p{i+1},i+1} \cdots a_{\sigma(n),n} +\\
    & & \mu a_{\sigma\p{1},1}\cdots a_{\sigma\p{i-1},i-1} y_{\sigma(i)}
        a_{\sigma\p{i+1},i+1ß} \cdots a_{\sigma(n),n}\\
    &=& \lambda \phi_\sigma \p{x_1,\ldots,x_{i-1},x,x_{i+1},\ldots,x_n}+\\
    & & \mu     \phi_\sigma \p{x_1,\ldots,x_{i-1},y,x_{i+1},\ldots,x_n}
    \end{eqnarray*}
    Donc $\phi_\sigma$ est $n$-linéaire.
  \item Montrons que $\phi$ est alternée. Soit $i,j\in\intere{1}{n}$ avec
    $i\neq j$ et $x_1,\ldots,x_n\in E$ tels que $x_i=x_j$. Montrons que
    $\phi\p{x_1,\ldots,x_n}=0$. Pour cela, on pose $\tau=\p{i \quad j}$ et on
    considère l'application~:
    \[\dspappli{f}{\galt{n}}{\gsym{n}\setminus\galt{n}}{\sigma}{\sigma \tau}\]
    Alors~:
    \begin{itemize}
    \item $f$ est bien à valeurs dans $\gsym{n}\setminus\galt{n}$, car si
      $\sigma$ est une permutation paire, $\sigma\tau$ est impaire.
    \item De plus, $f$ est bijective. En effet si $g$ est l'application
      \[\dspappli{g}{\gsym{n}\setminus\galt{n}}{\galt{n}}{\sigma}{\sigma \tau}\]
      on vérifie de même que $g$ est bien définie et que~:
      \[\forall \sigma\in \galt{n} \quad
        \p{g\circ f}\p{\sigma}=g\p{\sigma \tau}=\sigma\underbrace{\tau^2}_{=\id}=
        \sigma\]
      Donc $g\circ f=\id$. De même, on montre que $f\circ g=\id$ donc $f$
      est bijective.     
    \end{itemize}
    On en déduit que~:
    \begin{eqnarray*}
    \phi\p{x_1,\ldots,x_n}
    &=& \sum_{\sigma\in\gsym{n}} \signat{\sigma} \phi_{\sigma}\p{x_1,\ldots,x_n}\\
    &=& \sum_{\sigma\in \galt{n}} \signat{\sigma} \phi_{\sigma}\p{x_1,\ldots,x_n}+
        \sum_{\sigma\in \gsym{n}\setminus\galt{n}} \signat{\sigma}
        \phi_{\sigma}\p{x_1,\ldots,x_n}\\
    &=& \sum_{\sigma\in \galt{n}} \signat{\sigma} \phi_{\sigma}\p{x_1,\ldots,x_n}+
        \sum_{\sigma\in \galt{n}} \signat{\sigma \tau}
        \phi_{\p{\sigma\tau}}\p{x_1,\ldots,x_n}\\
    & & \text{car $f$ est bijective}\\
    &=& \sum_{\sigma\in \galt{n}} \signat{\sigma} \phi_{\sigma}\p{x_1,\ldots,x_n}+
        \sum_{\sigma\in \galt{n}} \signat{\sigma} \underbrace{\signat{\tau}}_{=-1}
        \phi_{\p{\sigma\tau}}\p{x_1,\ldots,x_n}\\
    &=& \sum_{\sigma\in \galt{n}} \signat{\sigma}
        \cro{\phi_{\sigma}\p{x_1,\ldots,x_n}-
             \phi_{\p{\sigma\tau}}\p{x_1,\ldots,x_n}}
    \end{eqnarray*}
    Or, si $A=\mat{\mathcal{B}}{x_1,\ldots,x_n}$, il vient~:
    \begin{eqnarray*}
    \phi_{\p{\sigma\tau}}\p{x_1,\ldots,x_n}
    &=& a_{\sigma\p{\tau\p{1}},1}\cdots a_{\sigma\p{\tau(i)},i} \cdots
        a_{\sigma\p{\tau(j)},j} \cdots a_{\sigma\p{\tau(n)},n}\\
    &=& a_{\sigma\p{1},1}\cdots a_{\sigma(j),i} \cdots
        a_{\sigma(i),j} \cdots a_{\sigma(n),n}\\
    &=& a_{\sigma\p{1},1}\cdots a_{\sigma(j),j} \cdots
        a_{\sigma(i),i} \cdots a_{\sigma(n),n}\\
    & & \text{car $x_i=x_j$}\\
    &=& a_{\sigma\p{1},1}\cdots a_{\sigma(i),i} \cdots
        a_{\sigma(j),j} \cdots a_{\sigma(n),n}\\
    &=& \phi_\sigma \p{x_1,\ldots,x_n}
    \end{eqnarray*}
    Donc $\phi\p{x_1,\ldots,x_n}=0$. On a donc prouvé que $\phi$ est alternée.
  \item Montrons que $\phi\p{e_1,\ldots,e_n}=1$. On remarque que~:
    \[A=\mat{\mathcal{B}}{e_1,\ldots,e_n}=I_n\]
    Donc~:
    \[\phi\p{e_1,\ldots,e_n}=\sum_{\sigma\in\gsym{n}} \signat{\sigma}
      a_{\sigma\p{1},1}a_{\sigma\p{2},2}\cdots a_{\sigma(n),n}\]
    Or, si $\sigma\neq \id$, il existe $i$ tel que $\sigma(i)\neq i$, donc
    $a_{\sigma(i),i}=0$. On en déduit que~:
    \[\phi\p{e_1,\ldots,e_n}=\signat{\id}a_{1,1}a_{2,2}\cdots a_{n,n}\]
    Or $\signat{\id}=1$ et $a_{i,i}=1$ pour tout $i\in\intere{1}{n}$, donc
    $\phi\p{e_1,\ldots,e_n}=1$.
  \end{itemize}
  En conclusion $\phi$ est une forme $n$-linéaire alternée telle que
  $\phi\p{e_1,\ldots,e_n}=1$.
\end{itemize}
\end{preuve}

\begin{preuve}
Par Victor :

\begin{itemize}
\item[$\bullet$] \textbf{\underline{Analyse :}} D'après la proposition précédente, si elle existe, cette application est forcément unique car $\phi\p{e_1,\ldots,e_n}=1$ impose $$\phi\p{x_1,\ldots,x_n}=\sum_{\sigma\in\gsym{n}} \signat{\sigma}
  a_{\sigma\p{1},1}a_{\sigma\p{2},2}\cdots a_{\sigma(n),n}.$$

\item[$\bullet$] \textbf{\underline{Synthèse :}}
 Il nous faut montrer que l'application $\phi$ ainsi définie (à la fin de l'analyse) est bien une forme $n$-linéaire alternée et vérifie bien $\phi\p{e_1,\ldots,e_n}=1$.

Pour simplifier les raisonnements, nous allons utiliser les applications coordonnées qu'on définit pour tout $k$ dans $\intere{1}{n}$ par :
$$\dspappli{e_k^*}{E}{\K}{x=\sum_{i=1}^nx_ie_i}{x_k}.$$
Ce sont des formes linéaires de $E$.

\bigskip
L'application $\phi$ peut alors se réécrire :
$$\dspappli{\phi}{E^n}{\K}{(x_1,\ldots,x_n)}{\sum_{\sigma\in \gsym{n}}\epsilon(\sigma)\prod_{k=1}^ne_{\sigma(k)}^*(x_k)}.$$
On comprend bien cette écriture lorsqu'on interprète bien $e_{\sigma(k)}^*(x_k)$ comme la composante de $x_k$ selon $e_{\sigma(k)}$ dans $\mathcal{B}$. Cette écriture a d'ailleurs le mérite de montrer la dépendance en les $x_k$ de $\phi(x_1,\ldots,x_n)$ qui était dissimulée dans l'"ancienne" écriture.

Mais alors :
$$\phi(e_1,\ldots,e_n)=\sum_{\sigma\in \gsym{n}}\epsilon(\sigma)\underbrace{\prod_{k=1}^n\underbrace{e_{\sigma(k)}^*(e_k)}_{=\delta_{\sigma(k),k}}}_{=\begin{cases}0 \text{ si } \sigma\neq id_{\intere{1}{n}}\\
1 \text{ si } \sigma= id_{\intere{1}{n}}\end{cases}}$$
d'où $$\boxed{\phi(e_1,\ldots,e_n)=\epsilon\p{id_{\intere{1}{n}}}.1=1.}$$

Montrons désormais que $\phi$ est $n$-linéaire alternée.
\begin{itemize}
\item Fixons $i\in \intere{1}{n}$ et vérifions que $\phi$ est linéaire par rapport à sa $i$-ième variable. Fixons alors des $x_k$ pour $k\neq i$. Pour tout $x$ dans $E$, on a :
\begin{eqnarray*}
\phi(x_1,\ldots,x_{i-1},x,x_{i+1},\ldots,x_n)&=&\sum_{\sigma\in \gsym{n}}\underbrace{\cro{\epsilon(\sigma)\prod_{k\neq i}e_{\sigma(k)}^*(x_k)}}_{:=\lambda_{\sigma} \text{ indépendant de } x}e_{\sigma(i)}^*(x)\\
&=&\sum_{\sigma\in \gsym{n}}\lambda_{\sigma}e_{\sigma(i)}^*(x).
\end{eqnarray*} ce qui est bien linéaire en $x$ comme combinaison linéaires de telles fonctions.
\item Vérifions maintenant son caractère alterné. Soit alors $(x_1,\ldots,x_n) \in E^n$ tel que il existe $1\leq i<j\leq n$ tels que $x_i=x_j$. Posons $\tau=(i,j) \in \gsym{n}$. On rappelle que $\gsym{n}=\galt{n}\cup \set{\sigma\circ \tau, \sigma \in \galt{n}}$ et que cette union est disjointe. On a :
\begin{eqnarray*}
\phi(x_1,\ldots,x_n)&=&\sum_{\sigma\in \galt{n}}\epsilon(\sigma)\prod_{k=1}^ne_{\sigma(k)}^*(x_k)+\sum_{\sigma\in \galt{n}}\epsilon(\sigma\circ \tau)\prod_{k=1}^ne_{\sigma\circ \tau(k)}^*(x_k)\\
&=&\sum_{\sigma\in \galt{n}}\prod_{k=1}^ne_{\sigma(k)}^*(x_k)-\underbrace{\sum_{\sigma\in \galt{n}}\prod_{k=1}^ne_{\sigma\circ \tau(k)}^*(x_k)}_{:=C}.
\end{eqnarray*}
Intéressons-nous au deuxième terme :
\begin{eqnarray*}
C&=&\sum_{\sigma\in \galt{n}}e_{\sigma\circ \tau(i)}^*(x_i)e_{\sigma\circ \tau(j)}^*(x_j)\prod_{k\notin\set{i,j}}e_{\sigma\circ \tau(k)}^*(x_k)\\
&=&\sum_{\sigma\in \galt{n}}e_{\sigma(j)}^*(x_i)e_{\sigma(i)}^*(x_j)\prod_{k\notin\set{i,j}}e_{\sigma(k)}^*(x_k)\\
&=&\sum_{\sigma\in \galt{n}}e_{\sigma(j)}^*(x_j)e_{\sigma(i)}^*(x_i)\prod_{k\notin\set{i,j}}e_{\sigma(k)}^*(x_k) \text{ car } x_i=x_j\\
&=&\sum_{\sigma\in \galt{n}}\prod_{k=1}^ne_{\sigma(k)}^*(x_k)
\end{eqnarray*}
donc $\phi(x_1,\ldots,x_n)=0$ ce qu'on souhaitait montrer.
\end{itemize}

\end{itemize}

\end{preuve}

\begin{proposition}
Soit $E$ un \Kev de dimension $n$. Alors $\Lambda_n(E)$ est un \Kev de
dimension $1$.
\end{proposition}

\begin{preuve}
En notant $\det_{\mathcal{B}}$ l'unique application du théorème précédent, on a montré dans la proposition 2.2 que toute forme $n$-linéaire alternée sur $E$ $\phi$ vérifiait $\phi=\phi(e_1,\ldots,e_n)\det_\mathcal{B}$. Donc l'ensemble des formes linéaires alternées est égale à $\vect(\det_{\mathcal{B}})$ qui est bien de dimension $1$ car $\det_{\mathcal{B}}$ n'est pas l'application nulle (elle vaut $1$ sur $\mathcal{B}$).\\

Résumé : Tout élément $f$ de $\Lambda_n(E)$ s'écrit $f=\lambda \det_{\mathcal{B}}$ avec $\lambda \in \K$ unique. 
\end{preuve}

% \begin{proposition}
% Soit $E$ un \Kev de dimension $n$, $\phi$ une forme $n$-linéaire alternée
% sur $E$ non nulle et $x_1,\ldots,x_n$ une famille de $n$ vecteurs de $E$. Alors
% $x_1,\ldots,x_n$ est une base de $E$ si et seulement si
% \[\phi\p{x_1,\ldots,x_n}\neq 0.\]
% Autrement dit, $x_1,\ldots,x_n$ est liée si et seulement si
% \[\phi\p{x_1,\ldots,x_n}= 0.\]
% \end{proposition}

% \begin{preuve}
% Je l'enlèverais bien parce-que je vais faire la proposition 6 avant la proposition 5 pour la déduire.

% \end{preuve}

\subsection{Déterminant d'une famille de $n$ vecteurs}

\begin{definition}
Si $\mathcal{B}\defeq\p{e_1,\ldots,e_n}$ est une base de $E$, on appelle \emph{déterminant relativement à la base
$\mathcal{B}$}, et on note $\det{}_{\mathcal{B}}$, l'unique forme $n$-linéaire alternée sur $E$ telle que
$\det{}_{\mathcal{B}}(e_1,\ldots,e_n)=1$. Si $(x_1,\ldots,x_n)$ est une famille de $n$ vecteurs de $E$, on appelle
\emph{déterminant de la famille $(x_1,\ldots,x_n)$ relativement à la base $\mathcal{B}$} le scalaire
$\det{}_{\mathcal{B}}\p{x_1,\ldots,x_n}$.
\end{definition}

\begin{remarqueUnique}
\remarque On en déduit que si $\mathcal{B}$ est une base de $E$, $x_1,\ldots,x_n\in E$ et
  $A\defeq \mat{\mathcal{B}}{x_1,\ldots,x_n}$, alors
  \[\det{}_{\mathcal{B}}(x_1,\ldots,x_n)=\sum_{\sigma\in\gsym{n}} \signat{\sigma}a_{\sigma(1),1}\cdots a_{\sigma(n),n}.\]
\end{remarqueUnique}

\begin{sol}
On a donc 
$$\det{}_{\mathcal{B}}\p{x_1,\ldots,x_n}=\sum_{\sigma\in \gsym{n}}\epsilon(\sigma)\prod_{k=1}^ne_{\sigma(k)}^*(x_k),$$
où $e_{\sigma(k)}^*(x_k)$ est la composante de $x_k$ selon $e_{\sigma(k)}$ dans $\mathcal{B}$.
\end{sol}

\begin{proposition}
Soit $\mathcal{B}\defeq\p{e_1,\ldots,e_n}$ et $\mathcal{B}'\defeq\p{e_1',\ldots,e_n'}$ deux
bases de $E$. Alors
\[\forall x_1,\ldots,x_n \in E \qsep
  \det{}_{\mathcal{B}'}\p{x_1,\ldots,x_n}=\det{}_{\mathcal{B}'}(\mathcal{B})
  \det{}_{\mathcal{B}}\p{x_1,\ldots,x_n}.\]
\end{proposition}

\begin{preuve}
$\det_{\mathcal{B}'}\in \Lambda_n(E)=\vect(\det_{\mathcal{B}})$ donc il existe $\lambda \in \K$ tel que $\det_{\mathcal{B}'}=\lambda \det_{\mathcal{B}}$. En particulier, $\det_{\mathcal{B}'}(\mathcal{B})=\lambda\det_{\mathcal{B}}(\mathcal{B})=\lambda$ d'où le résultat souhaité.
\end{preuve}

\begin{remarqueUnique}
\remarque En particulier, si $\mathcal{B}$, $\mathcal{B}'$ et $\mathcal{B}''$
  sont des bases de $E$
  \[\det{}_{\mathcal{B}}(\mathcal{B}'')=\det{}_{\mathcal{B}}(\mathcal{B}')
    \det{}_{\mathcal{B}'}(\mathcal{B}'').\]
\end{remarqueUnique}

\begin{proposition}
  Soit $\mathcal{B}\defeq\p{e_1,\ldots,e_n}$ une base de $E$ et $(x_1,\ldots,x_n)$ une
  famille de $n$ vecteurs de $E$.
\begin{itemize}
\item Alors $(x_1,\ldots,x_n)$ est une base de $E$
  si et seulement si
  \[\det{}_{\mathcal{B}}\p{x_1,\ldots,x_n}\neq 0.\]
\item Autrement dit, $(x_1,\ldots,x_n)$ est liée si et seulement si
  \[\det{}_{\mathcal{B}}\p{x_1,\ldots,x_n}= 0.\]
\end{itemize}
  \end{proposition}
  
  \begin{preuve}
  A faire après la proposition d'après.
  On raisonne par double implication :
  \begin{itemize}
  \item[$\bullet$] $\Rightarrow$ : Si $\mathcal{B}':=\p{x_1,\ldots,x_n}$ est une base de $E$, d'après la proposition précédente (suivante actuellement dans le poly) $$\underbrace{\det{}_{\mathcal{B}'}\p{x_1,\ldots,x_n}}_{=1}=\det{}_{\mathcal{B}'}\mathcal{B}\cdot
    \det{}_{\mathcal{B}}\p{x_1,\ldots,x_n}$$ donc $\det{}_{\mathcal{B}}\p{x_1,\ldots,x_n}\neq 0$.
  \item[$\bullet$] $\Leftarrow$ : Raisonnons par contraposée. Si $\mathcal{B}':=\p{x_1,\ldots,x_n}$ n'est pas une base de $E$, comme elle est de cardinal $n=\dim E$, c'est une famille liée. Disons que par exemple $x_{i_0}=\sum_{k\neq i_0}\lambda_k x_k$. Alors
  $$\phi(x_1,\ldots,x_{i_0-1},\sum_{k\neq i_0}\lambda_k x_k,x_{i_0+1},\ldots,x_n)=\sum_{k\neq i_0}\lambda_k \underbrace{\phi(x_1,\ldots,x_k,\ldots,x_n)}_{=0 \text{ car } x_k \text{ est présent deux fois}}.$$
  \end{itemize}
  \end{preuve}

\subsection{Déterminant d'un endomorphisme}
\begin{definition}
Si $f\in\Endo{E}$, il existe un unique scalaire, appelé \emph{déterminant de $f$} et noté
$\det f$, tel que pour toute base $\mathcal{B}$ de E, on a
\[\forall x_1,\ldots,x_n\in E \qsep
  \det{}_{\mathcal{B}}\p{f\p{x_1},\ldots,f\p{x_n}}=\det(f)\det{}_{\mathcal{B}}
  \p{x_1,\ldots,x_n}.\]
En particulier, si $\mathcal{B}\defeq\p{e_1,\ldots,e_n}$ est une base de $E$
\[\det f=\det{}_{\mathcal{B}}\p{f\p{e_1},\ldots,f\p{e_n}}.\]
\end{definition}

\begin{preuve}
Soit $f\in \Endo{E}$. Si on fixe $B$ une base de $E$, l'application $$\dspappli{\phi_B}{E^n}{K}{(x_1,\ldots,x_n)}{\det_B(f(x_1),\ldots,f(x_n))}$$ est dans $\Lambda_n(E)$, donc il existe $\lambda_B$ unique dans $\K$ tel que $\phi_B=\lambda_B\det_B$, i.e
$$\forall (x_1,\ldots,x_n)\in E^n, \det{}_{\mathcal{B}}\p{f\p{x_1},\ldots,f\p{x_n}}=\lambda_B\cdot\det{}_{\mathcal{B}}
  \p{x_1,\ldots,x_n}.$$
  Le point crucial est le fait que $\lambda_B$ est indépendant du choix de $B$. Cela découle de la formule de changement de base. La relation précédente donne $$\lambda_B=\det_B(u(B)).$$
  Soit $B'$ une autre base de $E$.
  $$\lambda_{B'}=\det_{B'}(u(B'))=\det_{B'}(B)\cdot \det_B(u(B'))=\det_{B'}(B)\cdot \lambda_B \cdot \det_{B}(B')$$
  Or, la formule de changement de base nous dit que $\det_{B'}=\det_{B'}(B)\det_B$ donc en appliquant cela à $B'$, on obtient $1=\det_{B'}(B)\cdot \det_{B}(B')$ d'où $\lambda_{B'}=\lambda_B$.

\end{preuve}

\begin{exoUnique}
\exo Soit $E$ un \Kev de dimension $n$ et $s$ une symétrie de $E$. On note
  $p$ la dimension de $\ker\p{s+\id}$. Montrer que $\det s=\p{-1}^p$.
\end{exoUnique}

\begin{sol}
$s$ est la symétrie par rapport à $\ker(s-\id)$ parallèlement à $\ker(s+\id)$. Comme $E=\ker(s-\id)\oplus \ker(s+\id)$, on peut se munir d'une base adaptée $B=(e_1,\ldots,e_q,f_1,\ldots,f_p)$. On a alors
\begin{eqnarray*}
\det(s)&=&\det_B(s(e_1),\ldots,s(e_q),s(f_1),\ldots,s(f_p))\\
&=&\det_B(e_1,\ldots,e_q,-f_1,\ldots,-f_p)\\
&=&(-1)^p\det_B(e_1,\ldots,e_q,f_1,\ldots,f_p)\\
&=&(-1)^p\det_B(B)\\
&=&(-1)^p.
\end{eqnarray*}
\end{sol}

\begin{proposition}
$\quad$
\begin{itemize}
\item $\det \id_E=1$
\item Si $f\in\Endo{E}$ et $\lambda\in\K$, alors
  \[\det\p{\lambda f}=\lambda^n \det f.\]
\item Si $f,g\in\Endo{E}$, alors
  \[\det\p{g\circ f}=\det g\cdot \det f.\]
\end{itemize}
\end{proposition}

\begin{preuve}
\begin{itemize}
\item $\det \id_E=\det_B(\id_E(B))=\det_B(B)=1$.
\item Si $f\in\Endo{E}$ et $\lambda\in\K$, alors
  \[\det\p{\lambda f}=\det_B\p{(\lambda f)(B)}=\det_B\p{\lambda f(e_1),\ldots,\lambda f(e_n)}=\lambda^n\det_B\p{f(e_1),\ldots, f(e_n)}=\lambda^n\det_B(f(B))=\lambda^n \det f.\]
\item Si $f,g\in\Endo{E}$, alors~:
  \[\det\p{g\circ f}=\det_B((g\circ f)(B))=\det g \det_B(f(B))=\det g\cdot \det f\]
\end{itemize}
\end{preuve}


\begin{exoUnique}
\exo Soit $E$ un \Rev de dimension $n$. Montrer qu'il existe $f\in\Endo{E}$
  tel que $f^2=-\id$ si et seulement si $n$ est pair.
\end{exoUnique}

\begin{sol}
\begin{itemize}
\item[$\bullet$] $\Rightarrow$ : $\det(f^2)=\det(-\id)=(-1)^n\det(\id)=(-1)^n$. Or $\det(f^2)=\det(f)^2\geq 0$ donc $n$ est pair.
\item[$\bullet$] Réciproquement, si $n$ est pair, on peut construire un endomorphisme, on peut construire une matrice diagonale par blocs $2\times 2$ avec chaque bloc égal à $\begin{pmatrix}0&-1\\1&0\end{pmatrix}$.
\end{itemize}
\end{sol}


\begin{proposition}
Soit $f\in\Endo{E}$. Alors $f$ est un isomorphisme si et seulement si
\[\det f\neq 0.\]
De plus, si tel est le cas
\[\det f^{-1}=\frac{1}{\det f}.\]
\end{proposition}

\begin{preuve}
Soit $B$ une base de $E$.
$f\in \gl{}{E}\Longleftrightarrow f(B) \text{ est une base de } E \Longleftrightarrow \det_B(f(B))\neq 0 \Longleftrightarrow \det(f)\neq 0.$
Dans ce cas, $\det(f)\det(f^{-1})=\det(f\circ f^{-1})=\det(\id)=1$ d'où \[\det f^{-1}=\frac{1}{\det f}.\]
\end{preuve}

\begin{remarqueUnique}
\remarque On en déduit que l'application
  \[\dspappli{\phi}{{\rm GL}(E)}{\Ks}{f}{\det f}\]
  est un morphisme de groupe de $({\rm GL}(E),\circ)$ dans $(\Ks,\times)$.
\end{remarqueUnique}

\subsection{Déterminant d'une matrice carrée}

\begin{definition}
Soit $A\in\mat{n}{\K}$. On appelle \emph{déterminant de $A$} et on note $\det A$
le déterminant des vecteurs colonnes de $A$ relativement à la base canonique de
$\K^n$.
\end{definition}

\begin{remarques}
\remarque Le déterminant est donc une forme $n$-linéaire alternée par rapport aux
  colonnes de la matrice.
\remarque Si $A\in\mat{n}{\K}$, alors
  \[\det A=\sum_{\sigma\in\gsym{n}} \signat{\sigma}
    a_{\sigma\p{1},1}a_{\sigma\p{2},2}\cdots a_{\sigma(n),n}\]
  Cependant, cette formule comporte une somme de $n!$ termes. Il est donc déconseillé de l'utiliser
  pour un calcul effectif de déterminant.
  % Par contre, elle permettra, par exemple, de démontrer que le déterminant d'une matrice à
  % coefficients entiers est un entier.
  \begin{sol}
  Tout de même, pour $n=2$, $S_2=\set{\id,(1,2)}$ et donc $\det(A)=a_{1,1}a_{2,2}-a_{2,1}a_{1,2}$.
  Si $n=3$, règle de Sarrus à éviter.
  \end{sol}
\remarque Si $A\in\mat{n}{\K}$, son déterminant est noté
  \[\begin{vmatrix}
    a_{1,1} & \cdots & a_{1,n}\\
    \vdots & & \vdots\\
    a_{n,1} & \cdots & a_{n,n}
    \end{vmatrix}\]
\end{remarques}

\begin{proposition}
Soit $\mathcal{B}$ une base de $E$.
\begin{itemize}
\item Si $(x_1,\ldots,x_n)$ est une famille de $n$ vecteurs de $E$, alors
  \[\det{}_{\mathcal{B}}\p{x_1,\ldots,x_n}=
    \det\cro{\mat{\mathcal{B}}{x_1,\ldots,x_n}}.\]
\item Si $f\in\Endo{E}$, alors
  \[\det f=\det\cro{\mat{\mathcal{B}}{f}}.\]
\end{itemize}
\end{proposition}

\begin{preuve}
\begin{itemize}
\item Rien à dire, c'est la même définition, la même somme.
\item Si $f\in\Endo{E}$, alors~:
  \[\det f=\det_B(f(e_1),\ldots,f(e_n))=\det(\mat{\mathcal{B}}{f(e_1),\ldots,f(e_n)}=\det\cro{\mat{\mathcal{B}}{f}}.\]
\end{itemize}

\end{preuve}

\begin{proposition}
$\quad$
\begin{itemize}
\item $\det I_n=1$
\item Si $A\in\mat{n}{\K}$ et $\lambda\in\K$, alors
  \[\det\p{\lambda A}=\lambda^n \det A.\]
\item Si $A,B\in\mat{n}{\K}$, alors
  \[\det\p{AB}=\det A\cdot\det B.\]
\end{itemize}
\end{proposition}

\begin{preuve}
On se fixe une base $B$, on passe par les endomorphismes canoniquement associés pour lesquels on a établi ces résultats.
\end{preuve}

\begin{remarqueUnique}
\remarque Il n'existe aucune formule permettant de calculer
  $\det\p{A+B}$ en fonction de $\det A$ et de $\det B$. En particulier, toute
  formule du type $\det\p{A+B}=\det A+\det B$ est fausse.
\end{remarqueUnique}

\begin{proposition}
Soit $A\in\mat{n}{\K}$. Alors $A$ est inversible si et seulement si
\[\det A\neq 0.\]
De plus, si tel est le cas
\[\det\p{A^{-1}}=\frac{1}{\det A}.\]
\end{proposition}

% \begin{proposition}
% Soit $AX=B$ un système de Cramer à $n$ équations et $n$ inconnues. Si
% $C_1,\ldots,C_n\in\R^n$ sont les $n$ colonnes de la matrice $A$, l'unique
% solution $\p{x_1,\ldots,x_n}\in\R^n$ du système est donnée par~:
% \[\forall k\in\intere{1}{n} \quad
%   x_k=\frac{\det\p{C_1,\ldots,C_{k-1},B,C_{k+1},\ldots,C_n}}{\det A}\]
% \end{proposition}

\begin{remarqueUnique}
\remarque On en déduit que l'application
  \[\dspappli{\phi}{{\rm GL}_n(\K)}{\Ks}{A}{\det A}\]
  est un morphisme de groupe de $({\rm GL}_{n}(\K),\times)$ dans $(\Ks,\times)$.
\end{remarqueUnique}

\begin{proposition}
Soit $A\in\mat{n}{\K}$. Alors
\[\det\trans{A}=\det A.\]
\end{proposition}

\begin{preuve}
On note $B=\trans{A}$. Alors~:
\begin{eqnarray*}
\det\trans{A}
&=& \det B\\
&=& \sum_{\sigma\in\gsym{n}} \signat{\sigma}b_{\sigma\p{1},1}\cdots b_{\sigma(n),n}\\
&=& \sum_{\sigma\in\gsym{n}} \signat{\sigma}a_{1,\sigma\p{1}}\cdots a_{n,\sigma(n)}\\
&=& \sum_{\sigma\in\gsym{n}}\signat{\sigma}a_{\sigma^{-1}\p{1},\sigma\p{\sigma^{-1}\p{1}}}
          \cdots a_{\sigma^{-1}(n),\sigma\p{\sigma^{-1}(n)}}\\
& & \text{par permutation des éléments présents dans le produit}\\
&=& \sum_{\sigma\in\gsym{n}}\signat{\sigma}a_{\sigma^{-1}\p{1},1}
          \cdots a_{\sigma^{-1}(n),n}
\end{eqnarray*}
Or $\sigma\sigma^{-1}=\id$ donc $\signat{\sigma}\signat{\sigma^{-1}}=1$. Comme
$\signat{\sigma}\in\ens{-1,1}$, on en déduit que
$\signat{\sigma^{-1}}=\signat{\sigma}$. Donc~:
\[\det\trans{A}=\sum_{\sigma\in\gsym{n}}\signat{\sigma^{-1}}a_{\sigma^{-1}\p{1},1}
          \cdots a_{\sigma^{-1}(n),n}\]
De plus, l'application~:
\[\dspappli{f}{\gsym{n}}{\gsym{n}}{\sigma}{\sigma^{-1}}\]
est une bijection. En effet, on vérifie facilement que $f\circ f=\id$. Donc~:
\begin{eqnarray*}
\det\trans{A}
&=& \sum_{\sigma\in\gsym{n}}\signat{\sigma}a_{\sigma\p{1},1} \cdots a_{\sigma(n),n}\\
&=& \det A
\end{eqnarray*}
\end{preuve}

\begin{remarqueUnique}
\remarque On en déduit que le déterminant est une forme $n$-linéaire alternée par
  rapport aux lignes de la matrice.
\end{remarqueUnique}
\vspace{2ex}
\begin{exoUnique}
\exo Soit $A\in\mat{3}{\R}$ une matrice antisymétrique. Montrer que
  $\det A=0$.
\end{exoUnique}

\begin{sol}
$\det(\trans{A})=\det(-A)=(-1)^3\det(A)=-\det(A)$. Or $\det(\trans{A})=\det(A)$ d'où $2\det(A)=0$.
\end{sol}

\section{Calcul de déterminant}

\subsection{Méthode du pivot}

\begin{proposition}
Soit la matrice
\[A\defeq
  \begin{pmatrix}
  a_{1,1} & a_{1,2}\\
  a_{2,1} & a_{2,2}
  \end{pmatrix}\in\mat{2}{\K}.\]
Alors $\det A=a_{1,1}a_{2,2}-a_{2,1}a_{1,2}$.
\end{proposition}

\begin{remarques}
\remarque Le déterminant de la matrice de taille nulle est 1.
\remarque Si $a \in\K$, le déterminant de la matrice  à une ligne et une colonne $(a)$
  est $a$. 
\end{remarques}
\vspace{2ex}
\begin{exoUnique}
\exo Soit $\theta_1,\ldots,\theta_n\in\R$. Calculer le rang de la matrice
  $A\in\mat{n}{\R}$ définie par
  \[\forall i,j\in\intere{1}{n} \qsep a_{i,j}\defeq\cos\p{\theta_i+\theta_j}.\]   
  \begin{sol}
  On trouve~:
  \begin{itemize}
  \item $\rg A=0$ si et seulement si~:
    \[\cro{\forall k\in\intere{1}{n} \quad \theta_k \equiv \frac{\pi}{4}\ [\pi]}
      \ou
      \cro{\forall k\in\intere{1}{n} \quad \theta_k \equiv -\frac{\pi}{4}\ [\pi]}\]
  \item $\rg A=2$ si et seulement si~:
    \[\exists i_1,i_2,j_1,j_2\in\intere{1}{n} \quad
      \theta_{i_1}\not\equiv \theta_{i_2}\ [\pi] \et
      \theta_{j_1}\not\equiv \theta_{j_2}\ [\pi]\]
  \item $\rg A=1$ sinon.
  \end{itemize}
  \end{sol}
\end{exoUnique}

%\begin{remarques}
% \remarque Le calcul du déterminant d'une matrice $3\times 3$
%   \[A=
%     \begin{pmatrix}
%     a_{1,1} & a_{1,2} & a_{1,3}\\
%     a_{2,1} & a_{2,2} & a_{2,3}\\
%     a_{3,1} & a_{3,2} & a_{3,3}
%     \end{pmatrix}\in\mat{3}{\K}\]
%   peut (bien que ce ne soit une bonne idée que dans très peu de cas) se faire
%   par la règle de Sarrus
%  \begin{eqnarray*}
% $\det A&=&\p{a_{1,1}a_{2,2}a_{3,3}+a_{2,1}a_{3,2}a_{1,3}+a_{3,1}a_{1,2}a_{2,3}}
%         &&-\p{a_{1,1}a_{3,2}a_{2,3}+a_{2,1}a_{1,2}a_{3,3}+a_{3,1}a_{2,2}a_{1,3}}$
%  \end{eqnarray*}
%\end{remarques}

\begin{proposition}
Soit $T$ une matrice triangulaire supérieure
\[\xymatrix @-0.85cm
  {& &\lambda_1\ar@{.}[drdrdrdr] & \star\ar@{.}[drdrdr]\ar@{.}[rrr]& & &
     \star\ar@{.}[ddd] \\
   & &  & & & &  \\
   T\defeq&
     \setbox0=\hbox{$\left(\rule[0pt]{0pt}{1.45cm}\right.$}
     \ht0=0pt\dp0=0pt\wd0=0pt\kern-5pt\lower5pt\box0
     &  & &\setbox0=\hbox{}\ht0=20pt\wd0=20pt\box0 & & &
     \setbox0=\hbox{$\left.\rule[0pt]{0pt}{1.45cm}\right)$}
     \ht0=0pt\dp0=0pt\wd0=0pt\kern-15pt\lower5pt\box0\\
   & &  &(0) & & & \star \\
   & &  & & & & \lambda_n}\]
alors
\[\det T=\prod_{k=1}^n \lambda_k.\]
\end{proposition}

\begin{preuve}
On a~:
\[\det T=\sum_{\sigma\in\gsym{n}}\signat{\sigma} t_{\sigma\p{1},1}\cdots
  t_{\sigma(n),n}\]  
On remarque que si $\sigma\neq\id$, le produit
$t_{\sigma\p{1},1}\cdots t_{\sigma(n),n}$ est nul. En effet, soit $\sigma\in\gsym{n}$
une permutation différente de l'identité. On considère le plus petit entier
$k_0\in\intere{1}{n}$ tel que $\sigma\p{k_0}=k_0$. Par minimalité de $k_0$, on a~:
\[\forall k\in\intere{1}{k_0-1} \quad \sigma(k)=k\]
Puisque $\sigma$ est injective, on en déduit que $\sigma\p{k_0}\geq k_0$.
Comme $\sigma\p{k_0}\neq k_0$, on en déduit que $\sigma\p{k_0}>k_0$ donc
$t_{\sigma\p{k_0},k_0}=0$ car $T$ est triangulaire supérieure. Le produit
$t_{\sigma\p{1},1}\cdots t_{\sigma(n),n}$ est donc nul. On a donc~:
\[\det T=\signat{\id} t_{1,1}\cdots t_{n,n}\]
Puisque $\signat{\id}=1$ et que $t_{k,k} =\lambda_k$ pour tout $k\in\intere{1}{n}$,
on en déduit que~:
\[\det T=\prod_{k=1}^n \lambda_k\]
\end{preuve}

\begin{sol}
Mettons ici la remarque du calcul du déterminant pour une matrice diagonale.
\end{sol}

\begin{remarqueUnique}
\remarque En particulier, le déterminant d'une matrice diagonale est le produit de
  ses coefficients diagonaux.
\end{remarqueUnique}

\begin{proposition}
Soit $A\in\mat{n}{\K}$ une matrice triangulaire supérieure par blocs
\[A=\begin{pmatrix}A_{1}&\star\\0&A_{2}\end{pmatrix}\]
où $n_1+n_2=n$, $A_{1}\in\mat{n_1}{\K}$ et $A_{2}\in\mat{n_2}{\K}$. Alors
\[\det A=\det A_{1} \det A_{2}.\]
\end{proposition}

\begin{remarques}
\remarque En particulier, si $A\in\mat{n+1}{\K}$ est de la forme
  % \[A=\begin{pmatrix}\lambda&B\\0&C\end{pmatrix}\]
\[\xymatrix @-0.85cm
  {& &\lambda & \star\ar@{.}[rrrr]& & & &
     \star \\
   & & 0 \ar@{.}[ddddd]
    & 
    % \ar@{-}[rrrr]\ar@{-}[ddddd]
     & & & & 
    %  \ar@{-}[ddddd]
     \\
   & &  & & & & & \\
   A=&
     \setbox0=\hbox{$\left(\rule[0pt]{0pt}{1.45cm}\right.$}
     \ht0=0pt\dp0=0pt\wd0=0pt\kern-5pt\lower-4pt\box0
     & & & \setbox0=\hbox{}\ht0=15pt\wd0=10pt\box0 
       & B & \setbox0=\hbox{}\ht0=15pt\wd0=10pt\box0 & &
     \setbox0=\hbox{$\left.\rule[0pt]{0pt}{1.45cm}\right)$}
     \ht0=0pt\dp0=0pt\wd0=0pt\kern-15pt\lower-4pt\box0\\
   & &  & & & & & \\
   & &  & & & & & \\
   & & 0 & 
  %  \ar@{-}[rrrr]
    & & & & }\]
  où $\lambda\in\K$ et $B\in\mat{n}{\K}$, alors
  \[\det A=\lambda \det B.\]
% Alors
% \[\det A=\lambda\det B.\]
\remarque De même, si $A\in\mat{n}{\K}$ est une matrice triangulaire inférieure par blocs
  \[A=\begin{pmatrix}A_{1}&0\\ \star&A_{2}\end{pmatrix}\]
  alors $\det A=\det A_{1} \det A_{2}$.
\end{remarques}

% \begin{proposition}

% \end{proposition}

% \begin{preuve}
% Soit $\sigma=\p{n+1 \quad n \quad\cdots\quad 1}$. En appliquant cette permutation
% aux colonnes de $A$, puis aux lignes de $A$, on obtient la matrice~:
% \[\xymatrix @-0.85cm
%   {& & \ar@{-}[rrrr]\ar@{-}[ddddd]& & & & \ar@{-}[ddddd] & 0\ar@{.}[ddddd] \\
%    & & & & & & & \\
%    C=&
%      \setbox0=\hbox{$\left(\rule[0pt]{0pt}{1.6cm}\right.$}
%      \ht0=0pt\dp0=0pt\wd0=0pt\kern-5pt\lower25pt\box0
%      & & \setbox0=\hbox{}\ht0=15pt\wd0=10pt\box0 
%        & B & \setbox0=\hbox{}\ht0=15pt\wd0=10pt\box0 & & &
%      \setbox0=\hbox{$\left.\rule[0pt]{0pt}{1.6cm}\right)$}
%      \ht0=0pt\dp0=0pt\wd0=0pt\kern-15pt\lower25pt\box0\\
%    & &  & & & & & \\
%    & &  & & & & & \\
%    & & \ar@{-}[rrrr] & & & & & 0 \\
%    & & \star\ar@{.}[rrrr] & & & & \star & \lambda}\]
% Comme $\sigma$ est un $\p{n+1}$-cycle, on en déduit que
% $\signat{\sigma}=\p{-1}^{n+2}$. Donc~:
% \[\det C=\p{-1}^{n+2} \p{-1}^{n+2}\det A=\det A\]
% Montrons que $\det C=\lambda \det B$. On a~:
% \[\det C=\sum_{\sigma\in\gsym{n+1}} \signat{\sigma} c_{\sigma\p{1},1}\cdots c_{\sigma{n},n}
%          c_{\sigma{n+1},n+1}\]
% Or, si $\sigma\p{n+1}\neq n+1$, $c_{\sigma{n+1},n+1}=0$, donc on peut restreindre la
% somme aux permutations qui laissent fixe $n+1$~:
% \begin{eqnarray*}
% \det C
% &=& \sum_{\substack{\sigma\in\gsym{n+1}\\ \sigma(n+1)=n+1}}
%     \signat{\sigma} c_{\sigma\p{1},1}\cdots c_{\sigma{n},n} c_{\sigma{n+1},n+1}\\
% &=& \sum_{\substack{\sigma\in\gsym{n+1}\\ \sigma(n+1)=n+1}}
%     \signat{\sigma} c_{\sigma\p{1},1}\cdots c_{\sigma{n},n}
%     \underbrace{c_{n+1,n+1}}_{=\lambda}\\
% &=& \lambda \sum_{\substack{\sigma\in\gsym{n+1}\\ \sigma(n+1)=n+1}}
%     \signat{\sigma} c_{\sigma\p{1},1}\cdots c_{\sigma{n},n}
% \end{eqnarray*}
% Or, si on note $X$ l'ensemble des permutations de $\intere{1}{n+1}$ qui laissent
% fixe $n+1$, l'application~:
% \[\dspappli{f}{X}{\gsym{n}}{\sigma}{
%     \dspappli{f\p{\sigma}}{\intere{1}{n}}{\intere{1}{n}}{k}{\sigma(k)}
%   }\]
% est une bijection. En effet~:
% \begin{itemize}
% \item $f$ est bien définie car si $\sigma$ est une permutation de $\intere{1}{n+1}$
%   qui laisse fixe $n+1$, $\sigma$ étant injective, et $\sigma\p{n+1}=n+1$, on a~:
%   \[\forall k\in\intere{1}{n} \quad \sigma(k)\in\intere{1}{n}\]
%   Donc $f\p{\sigma}$ est bien une application de $\intere{1}{n}$ dans lui-même.
%   De plus, puisque $\sigma$ est injective, $f\p{\sigma}$ est injective donc
%   $f\p{\sigma}\in\gsym{n}$.
% \item $f$ est injective. En effet soit $\sigma_1,\sigma_2$ deux permutations
%   de $\intere{1}{n+1}$ laissant fixe $n+1$ telles que $f\p{\sigma_1}=f\p{\sigma_2}$.
%   Alors $\sigma_1$ et $\sigma_2$ coïncident sur $\intere{1}{n}$. Comme de plus
%   $\sigma_1\p{n+1}=\sigma_2\p{n+1}=n+1$, on en déduit que $\sigma_1=\sigma_2$.
% \item Enfin $f$ est surjective. En effet, soit $g$ une permutation de
%   $\intere{1}{n}$. On définit alors l'application $f$ de $\intere{1}{n+1}$ dans
%   lui-même par~:
%   \[\forall k\in\intere{1}{n+1}\]
% \end{itemize}
% \end{preuve}

% \begin{remarqueUnique}
% \remarque Plus généralement, si une matrice est triangulaire supérieure par
%   blocs, on montre que son déterminant est égal au produit des
%   déterminants des matrices blocs présentes sur la diagonale.
% \end{remarqueUnique}

\begin{proposition}
Soit $A\in\mat{n}{\K}$, $i,j\in\intere{1}{n}$ tels que $i\neq j$ et $\lambda\in\K$.
\begin{itemize}
\item Les opérations élémentaires $L_i \gets L_i + \lambda L_j$ et $C_i\gets C_i+\lambda C_j$
  ne modifient pas les déterminants.
\item Les opérations élémentaires $L_i \gets \lambda L_i$ et $C_i\gets \lambda C_i$ multiplient
  les déterminants par $\lambda$.
\item Les opérations élémentaires $L_i \leftrightarrow L_j$ et $C_i \leftrightarrow C_j$ multiplient
  les déterminants par $-1$.
\end{itemize}
% \begin{itemize}
% \item On multiplie le déterminant de $A$ par $\lambda$ lorsqu'on multiplie une
%   de ses colonnes (resp. lignes) par $\lambda$.
% \item On ne change pas le déterminant de $A$ lorsqu'à une colonne (resp. ligne)
%   de $A$ on ajoute une combinaison linéaire des ses autres colonnes (resp.
%   lignes).
% \item On change le signe du déterminant de $A$ lorsqu'on échange deux de ses
%   colonnes (resp. lignes). Plus généralement, une permutation paire des colonnes
%   (resp. lignes) de $A$ ne change pas le signe de son déterminant, tandis qu'une
%   permutation impaire de ses colonnes (resp. lignes) change son signe.
% \end{itemize}
\end{proposition}


\begin{remarques}
\remarque Plus généralement, si à une colonne, on ajoute une combinaison linéaire des autres
  colonnes, on ne change pas le déterminant. De même, si à une ligne, on ajoute une
  combinaison linéaire des autres lignes, on ne change pas non plus le déterminant.
\remarque Toute permutation des colonnes multiplie le déterminant par la signature
  de cette permutation. De même, toute permutation des lignes multiplie le déterminant
  par la signature de cette permutation.
\end{remarques}

\begin{preuve}
C'est le caractère $n$-linéaire alterné du déterminant.
\end{preuve}

% \begin{exos}
% \exo Montrer que le déterminant
%   \[\begin{vmatrix}
%     3 & 4 & 2\\
%     5 & -3 & -2\\
%     2 & 1 & 6
%     \end{vmatrix}\]
%   est divisible par 9.
% \end{exos}
%
\vspace{2ex}
\begin{exoUnique}
\exo Soit $a,b,c\in\R$. Calculer les déterminants
  \[\begin{vmatrix}
    1 & a & a^2\\
    a & 1 & a\\
    a^2 & a & 1
    \end{vmatrix} \qquad
  \begin{vmatrix}
  1 & \sin a & \cos a\\
  1 & \sin b & \cos b\\
  1 & \sin c & \cos c
  \end{vmatrix}\et
  \begin{vmatrix}
    a+b & 2a & 2a\\
    2b & a+b & 2a\\
    2b & 2b & a+b
    \end{vmatrix}\]
  \begin{sol}
  On trouve $\p{a-1}^2\p{a+1}^2$, $4\sin((b-a)/2)\sin((c-a)/2)\sin((c-b)/2)$ et
  $\p{a+b}\p{a-b}^2$.    
  \end{sol}
\end{exoUnique}

\subsection{Développement par rapport à une colonne}

\begin{definition}
Soit $A\in\mat{n}{\K}$.
\begin{itemize}
\item On appelle \emph{mineur} d'indice $\p{i,j}\in\intere{1}{n}^2$ le déterminant $\Delta_{i,j}$ de la
  matrice obtenue en supprimant la $i$-ème ligne et la $j$-ème colonne de la
  matrice $A$.
\item On appelle \emph{cofacteur} d'indice $\p{i,j}\in\intere{1}{n}^2$, et on note $A_{i,j}$, le scalaire
  $A_{i,j}\defeq\p{-1}^{i+j}\Delta_{i,j}$.
\end{itemize}
\end{definition}

\begin{remarqueUnique}
\remarque On retiendra que la matrice des $(-1)^{i+j}$ comporte des $1$ sur la 
  diagonale et qu'on passe de $\pm 1$ à son opposé lorsqu'on change de ligne ou de colonne.
% \remarque Si $\rg A\leq n-2$, tous ses mineurs, et donc ses cofacteurs, sont
%   nuls.
\end{remarqueUnique}

\begin{proposition}[nom={Développement par rapport à une ligne ou une colonne}]
Soit $A\in\mat{n}{\K}$.
\begin{itemize}
\item Soit $j\in\intere{1}{n}$. Alors
  \begin{eqnarray*}
  \det A
  &=& \sum_{i=1}^n a_{i,j}A_{i,j}\\
  &=& \sum_{i=1}^n \p{-1}^{i+j}a_{i,j}\Delta_{i,j}.
  \end{eqnarray*}
\item Soit $i\in\intere{1}{n}$. Alors
  \begin{eqnarray*}
  \det A
  &=& \sum_{j=1}^n a_{i,j}A_{i,j}\\
  &=& \sum_{j=1}^n \p{-1}^{i+j}a_{i,j}\Delta_{i,j}.
  \end{eqnarray*}  
\end{itemize}
\end{proposition}

\begin{exos}
\exo Soit $u,v\in\R$. Calculer le déterminant de la matrice
  \[\begin{pmatrix}
    -u & v & 0\\
    -2 & 0 & 2v\\\
    0 & -1 & u
    \end{pmatrix}\]
  \begin{sol}
  On trouve 0. On peut remarquer que $vC_1+uC_2+C_3=0$.    
  \end{sol}
\exo Déterminer les $n\in\N$ pour lesquels la matrice suivante est inversible.
  \[\begin{pmatrix}
    1 & 1 &  &  & (0)\\
    1 & 1 & 1 &  & \\
     & \ddots & \ddots & \ddots & \\
     &  & 1 & 1 & 1\\
    (0) &  &   & 1 & 1
    \end{pmatrix}\in\mat{n}{\R}.\]
  \begin{sol}%
  On trouve $u_{n+2}=u_{n+1}-u_n$ puis
  $u_n=\frac{2}{\sqrt{3}}\sin\p{\p{n+1}\frac{\pi}{3}}$. On retrouve le fait
  que la matrice est inversible sauf lorsque $n\equiv 2\ [3]$.     
  \end{sol}
\end{exos}


\begin{proposition}[nom={Déterminant de \nom{Vandermonde}}]
On appelle \nom{Vandermonde} de la famille $(x_0,\ldots,x_n)\in\K^{n+1}$ le déterminant
  \[V(x_0,\ldots,x_n)\defeq \begin{vmatrix}
    1 & x_0 & x_0^2 & \cdots & x_0^n\\
    1 & x_1 & x_1^2 & \cdots & x_1^n\\
    \vdots & & & & \vdots\\
    1 & x_n & x_n^2 & \cdots & x_n^n
    \end{vmatrix}\]
On a
\[V(x_0,\ldots,x_n)=\prod_{0\leq i<j\leq n} (x_j-x_i).\]
\end{proposition}

\begin{remarqueUnique}
\remarque Soit $x_0,\ldots,x_n\in\K$, $n+1$ éléments deux à deux distincts et
  $y_0,\ldots,y_n\in\K$. Soit $P\defeq a_0 +a_1 X+\cdots+a_n X^n\in\polyK[n]$. Alors
  \[\forall k\in\intere{0}{n}\qsep P(x_k)=y_k\]
  si et seulement si
  \[\syslin{
    a_0 &+ a_1 x_0&+\cdots &+a_n x_0^n&=&y_0\hfill\cr
    a_0 &+ a_1 x_1&+\cdots &+a_n x_1^n&=&y_1\hfill\cr
    \vdots\ & & &\vdots\ \ & &\ \vdots\hfill\cr
    a_0 &+ a_1 x_n&+\cdots &+a_n x_n^n&=&y_n.\hfill}\]
  Le déterminant de la matrice de ce système est un \nom{Vandermonde} qui
  est non nul car les $x_k$ sont deux à deux distincts. On retrouve
  le fait que ce système est de \nom{Cramer} et donc qu'il admet une unique
  solution, résultat au cœur de la définition des polynômes interpolateurs de
  \nom{Lagrange}.
\end{remarqueUnique}

\subsection{Comatrice}
\begin{definition}
Soit $A\in\mat{n}{\K}$. On appelle \emph{comatrice} de $A$ et on note $\Com A$ la
matrice des cofacteurs de $A$
\[\forall i,j\in\intere{1}{n} \qsep \cro{\Com A}_{i,j}\defeq
  \p{-1}^{i+j}\Delta_{i,j}.\]
\end{definition}

\begin{proposition}
Soit $A\in\mat{n}{\K}$. Alors
\[A\trans{\p{\Com A}}=\trans{\p{\Com A}}A=\p{\det A}I_n.\]
En particulier, si $A\in\gl{n}{\K}$
\[A^{-1}=\frac{1}{\det A}\trans{\p{\Com A}}.\]
\end{proposition}


\begin{remarqueUnique}
\remarque La matrice
  \[A\defeq\begin{pmatrix}a&b\\c&d\end{pmatrix}\in\mat{2}{\K}\]
  est inversible si et seulement si $\det A=ad-cb$ est non nul. De plus, si tel est le cas
  \[A^{-1}=\frac{1}{ad-cb}\begin{pmatrix}d&-b\\-c&a\end{pmatrix}\]
  Cependant, on évitera d'utiliser une telle formule en pratique.
\end{remarqueUnique}

\begin{exoUnique}
\exo Soit $A\in\mat{n}{\Z}$ une matrice inversible dans $\mat{n}{\R}$.
  Montrer que son inverse est à coefficients entiers si et seulement si
  $\det A\in\ens{-1,1}$.
\end{exoUnique}

% \subsection{Orientation d'un espace vectoriel réel de de dimension finie}

% \begin{definition}
% Soit $E$ un \Rev de dimension finie. On dit qu'un automorphisme $f$ de $E$ est
% \begin{itemize}
% \item \emph{direct} lorsque $\det f>0$.
% \item \emph{indirect} lorsque $\det f<0$.
% \end{itemize}
% On note ${\rm GL}^+(E)$ l'ensemble des automorphismes directs de $E$. On
% définit de même les notions de matrice directes et indirectes ainsi que l'ensemble
% ${\rm GL}_n^+\p{\R}$.
% \end{definition}

% \begin{exoUnique}
% \item Si $E$ est un \Rev de dimension $n$, $-\id$ est direct si $n$ est
%   pair et indirect si $n$ est impair.
% \end{exoUnique}

% \begin{proposition}
% ${\rm GL}^+(E)$ (resp. ${\rm GL}_n^+\p{\R}$)  est un sous-groupe de $\gl{}{E}$
% (resp. $\gl{n}{\R}$).
% \end{proposition}

% \begin{definition}
% Soit $\mathcal{B}$ et $\mathcal{B}'$ deux bases de $E$. Alors les
% deux assertions suivantes sont équivalentes.
% \begin{itemize}
% \item $\det\p{P\p{\mathcal{B},\mathcal{B}'}}>0$.
% \item L'unique automorphisme $f$ qui transforme $\mathcal{B}$ en $\mathcal{B}'$
%   est direct.
% \end{itemize}
% Si tel est le cas, on dit que $\mathcal{B}$ a \emph{même orientation} que
% $\mathcal{B}'$.
% \end{definition}

% \begin{proposition}
% La relation \flqq\ a même orientation que \frqq\ est une relation d'équivalence
% sur l'ensemble des bases de $E$ et possède exactement deux classes
% d'équivalence.
% \end{proposition}

% \begin{definition}
% Choisir une orientation de $E$, c'est choisir une base $\mathcal{B}$ de $E$ que
% l'on définit comme directe. Les bases ayant même orientation que $\mathcal{B}$
% sont dites \emph{directes}, les autres sont dites \emph{indirectes}.
% \end{definition}
%END_BOOK

\end{document}
\section{Exercices}
\setcounter{numeroexercice}{1}
\documentclass{magnolia}

\magtex{tex_driver={pdftex}}
\magfiche{document_nom={Exercices sur les déterminants},
          auteur_nom={François Fayard},
          auteur_mail={fayard.prof@gmail.com}}
\magexos{exos_matiere={maths},
         exos_niveau={mpsi},
         exos_chapitre_numero={25},
         exos_theme={Déterminants}}
\magmisenpage{}
\maglieudiff{}
\magprocess

\begin{document}
%BEGIN_BOOK


\magsection{Déterminant}
\magsubsection{Forme $n$-linéaire alternée}


\exercice{nom={Exercice \textsc{ENS}}}
Soit $E$ un \Kev de dimension $d\in\N$. Pour tout $n\in\N$, calculer la dimension
de l'espace vectoriel $\mathcal{A}_n(E)$ des formes $n$-linéaires alternées sur $E$.

\magsubsection{Déterminant d'une famille de $n$ vecteurs}

\exercice{nom={Exercice}}
Soit $E$ un \Kev de dimension $n$, $f\in\mathcal{L}(E)$ et $\mathcal{B}$ une base
de $E$. Montrer que pour tous $x_1,\ldots,x_n\in E$
\[\sum_{k=1}^n \det{}_{\mathcal{B}}(x_1,\ldots,x_{k-1},f(x_k),x_{k+1},\ldots,x_n)=\tr(f)
  \det{}_{\mathcal{B}}(x_1,\ldots,x_n).\]

\magsubsection{Déterminant d'un endomorphisme}

\exercice{nom={Déterminant de la transposition}}
Soit $\phi$ l'application de $\mat{n}{\K}$ dans lui-même qui à la matrice $M$
associe sa transposée. Calculer le déterminant de $\phi$.
\begin{sol}
$\quad$
\begin{questions}
\question La transposition est une symétrie. La dimension de l'espace des
  matrices antisymétriques est $n\p{n-1}/2$. Donc le déterminant de la
  transposition est
  \[\p{-1}^{\frac{n\p{n-1}}{2}}\]
\end{questions}
\end{sol}

\magsubsection{Déterminant d'une matrice carrée}
\magsection{Calcul de déterminant}
\magsubsection{Méthode du pivot}

\exercice{nom={Calcul de déterminant}}
Soit $a_1,\ldots,a_n\in\R$. Calculer le déterminant de la matrice $A\in\mat{n}{\R}$
donnée par les coefficients
\[\forall i,j\in\intere{1}{n}\qsep a_{i,j}\defeq\sin(a_i+a_j).\]
% \[\left|
%   \begin{array}{cccc}
%   \sin(a_1+a_1) & \sin(a_1+a_2) & \dots & \sin(a_1+a_n)\\
%   \sin(a_2+a_1) &               &       & \vdots\\
%   \vdots        &               &       & \vdots\\
%   \sin(a_n+a_1) & \sin(a_n+a_2) & \dots & \sin(a_n+a_n)\\
%   \end{array}
%   \right|.\]
\begin{sol}
Si on note $A=\p{\sin a_1,\ldots,\sin a_n}$ et
$B=\p{\cos a_1,\ldots,\cos a_n}$, alors $C_j=\p{\cos a_j}A+\p{\sin a_j}B$.
Donc le rang de la matrice est inférieur ou égal à 2.
\begin{itemize}
\item Si $n\geq 3$, le déterminant est nul.
\item Si $n=1$, le déterminant est $\sin\p{2a_1}$.
\item Si $n=2$, le déterminant est $-\sin^2\p{a_1-a_2}$.
\end{itemize}
\end{sol}

\exercice{nom={Calculs de déterminants}}
Soit $a,b,c\in\K$. Calculer et factoriser les déterminants suivants
\[\begin{vmatrix}
  1 & 1 & 1\\
  a+b & c+a & b+c\\
  ab & ca & bc
  \end{vmatrix}\qquad
  \begin{vmatrix}
  a+b & b+c & c+a\\
  a^2+b^2 & b^2+c^2 & c^2+a^2\\
  a^3+b^3 & b^3+c^3 & c^3+a^3
  \end{vmatrix}\]
\[\begin{vmatrix}
  (b+c)^2 & b^2 & c^2\\
  a^2 & (c+a)^2 & c^2\\
  a^2 & b^2 & (a+b)^2
  \end{vmatrix}\]
\begin{sol}
$\quad$
\begin{questions}
\question On trouve $\p{c-b}\p{c-a}\p{b-a}$.
\question On trouve $2abc\p{a+b+c}^3$. On fait apparaître $a+b+c$ par
  l'opération $L_1\gets L_1-L_3$.
\end{questions}
\end{sol}

\exercice{nom={Calculs de déterminants}}
Soit $a,b,c,d\in\K$. Calculer et factoriser les déterminants suivants
\[\begin{vmatrix}
  1 & 1 & 1 & 1\\
  a & b & c & d\\
  a^2 & b^2 & c^2 & d^2\\
  bcd & acd & abd & abc
  \end{vmatrix}\qquad
  \begin{vmatrix}
  1 & a & b & ac\\
  1 & b & c & bd\\
  1 & c & d & ac\\
  1 & d & a & bd
  \end{vmatrix}
  \qquad
  \begin{vmatrix}
  0 & 1 & 1 & 1\\
  1 & 0 & a^2 & b^2\\
  1 & a^2 & 0 & c^2\\
  1 & b^2 & c^2 & 0
  \end{vmatrix}
  \]

\[\begin{vmatrix}
  1+a & b & a & b\\
  b & 1+a & b & a\\
  a & b & 1+a & b\\
  b & a & b & 1+a
  \end{vmatrix}\qquad
  \begin{vmatrix}
  a & c & c & b\\
  c & a & b & c\\
  c & b & a & c\\
  b & c & c & a
  \end{vmatrix}\]

\exercice{nom={Calcul de déterminant}}
Soit $a_1,\ldots,a_n\in\R$. Calculer
\[\left|
  \begin{array}{cccc}
  1        & 1 & \cdots & 1\\
  \cos(a_1) & \cos(a_2) & \cdots & \cos(a_n)\\
  \vdots   &   \vdots   &         & \vdots\\
  \cos ((n-1)a_1) & \cos((n-1)a_2) & \cdots & \cos((n-1)a_n)
  \end{array}
  \right|\]
\begin{sol}
Comme $\cos\p{k a_j}=P_k\p{\cos a_j}$ et que le coefficient dominant
de $P_k$ est $2^{k-1}$, le déterminant est égal à 
\begin{eqnarray*}
& & 2^{\frac{\p{n-2}\p{n-1}}{2}}V\p{\cos a_1,\ldots,\cos a_n}\\
&=& 2^{\frac{\p{n-2}\p{n-1}}{2}}\prod_{1\leq i<j\leq n} \p{\cos a_j-\cos a_i}\\
&=& \p{-1}^{\frac{n\p{n-1}}{2}}2^{\p{n-1}^2}\prod_{1\leq i<j\leq n}
    \sin\p{\frac{a_j+a_i}{2}}\sin\p{\frac{a_j-a_i}{2}}
\end{eqnarray*}
\end{sol}

\exercice{nom={Matrice circulante}}
\begin{questions}
\question Soit $a,b,c\in\C$. On pose
  \[M\defeq\begin{pmatrix}
    a&b&c\\
    c&a&b\\
    b&c&a
  \end{pmatrix} \quad\et\quad J\defeq\begin{pmatrix}
  1 & 1 & 1\\
  1 & \jj & \jj^2\\
  1 & \jj^2 & \jj
  \end{pmatrix}\]
  \begin{questions}
  \question Montrer que $\det J\neq 0$.
  \question Calculer $MJ$ et en déduire $\det M$.
  \end{questions}
\question Mêmes questions avec la \emph{matrice circulante}
  \[M\defeq\begin{pmatrix}
    a_0 & a_1 & \cdots & a_{n-1}\\
    a_{n-1} & a_0 & \cdots & a_{n-2}\\
    \vdots & \vdots & \ddots & \vdots\\
    a_1 & a_2 & \cdots & a_0
  \end{pmatrix} \quad\text{et}\quad 
  J\defeq \p{\omega^{(i-1)(j-1)}}_{1\leq i,j\leq n}\]
  pour $a_0,\ldots,a_{n-1}\in\C$ en posant $\omega\defeq \e^{\ii\frac{2\pi}{n}}$.
\end{questions}

\exercice{nom={Calculs de déterminants}}
Soit $a_1,\ldots,a_n,b_1,\ldots,b_n\in\K$. Calculer le déterminant
\[\left|
\begin{array}{ccccc}
a_1+b_1 & a_1 & \cdots & \cdots & a_1\\
a_2 & a_2+b_2 & a_2 & \cdots & a_2\\
a_3 & \ddots & \ddots & \ddots & \vdots\\
\vdots & & \ddots & \ddots & a_{n-1}\\
a_n & \cdots & \cdots & a_n & a_n+b_n
\end{array}
\right|\]

\begin{sol}
On trouve
  \[\prod_{k=1}^n b_k+\sum_{k=1}^n
    \cro{a_k \prod_{\substack{i=1\\i\neq k}}^n b_i}\]
  Pour cela, il y a plusieurs méthodes~:
  \begin{itemize}
  \item Le plus naturelle (et la plus longue). On enlève la dernière colonne
    à toutes les autres. On développe ensuite par rapport à la dernière ligne.
  \item Méthode plus fine. On utilise la linéarité du déterminant par rapport
    à la dernière colonne pour faire apparaître $\p{a_1,\ldots,a_n}$ dans la
    dernière colonne. Dans cette matrice, on enlève la dernière colonne à
    toutes les autres. Dans l'autre dont la dernière colonne est
    $\p{0,\ldots,0,b_n}$, on développe par rapport à cette dernière colonne.
  \item Méthode encore plus fine. On pose $u=\p{a_1,\ldots,a_n}$ et
    $\mathcal{B}=e_1,\ldots,e_n$ la base canonique de $\K^n$. Alors
    \begin{eqnarray*}
    \Delta=\det{}_{\mathcal{B}}\p{u+b_1 e_1,\ldots,u+b_n e_n}
    \end{eqnarray*}
    On utilise ensuite la linéarité du déterminant par rapport à chacune des
    variables.
  \end{itemize}
\end{sol}

\exercice{nom={Calcul de déterminant}}
Soit $n\in\Ns$ et $p\in\intere{1}{n}$. On définit la matrice $A_{n,p}$
de $\mat{p+1}{\R}$ par
\[A_{n,p}\defeq
  \begin{pmatrix}
  1 & \binom{n}{1} & \binom{n}{2} & \cdots & \binom{n}{p}\\ 
  1 & \binom{n+1}{1} & \binom{n+1}{2} & \cdots & \binom{n+1}{p}\\ 
  \vdots & \vdots &  &  & \vdots\\ 
  1 & \binom{n+p}{1} & \binom{n+p}{2} & \cdots & \binom{n+p}{p}
  \end{pmatrix}\]
Calculer $\det A_{n,p}$.


\exercice{nom={Opérations par blocs}}
\begin{questions}
\question Soit $A,B,C,D\in\mat{n}{\K}$. Quel produit matriciel permet de transformer
  \[\begin{pmatrix}A & B\\ C&D\end{pmatrix} \quad\text{en}\quad
    \begin{pmatrix}A+2C & B+2D\\ C&D\end{pmatrix}\]
\question Montrer que pour tous $A,B\in\mat{n}{\K}$
  \[\begin{vmatrix}I_n & B\\A & I_n\end{vmatrix} = \det(I_n - AB)=\det(I_n-BA).\]
\end{questions}


\magsubsection{Développement par rapport à une colonne}

\exercice{nom={Matrice tridiagonale}}
Soit $x\in\R$. Calculer le déterminant
\[\left|
\begin{array}{ccccc}
1+x^2 & x &  & & (0)\\
x & 1+x^2 & x & & \\
  & \ddots & \ddots & \ddots & \\
  &   & x & 1+x^2 & x\\
(0)  & &  & x & 1+x^2
\end{array}
\right|\]
\begin{sol}
On trouve la relation de récurrence $u_{n+2}=\p{1+x^2}u_{n+1}-x^2 u_n$.\begin{itemize}
\item Si $x^2=1$, on a $u_n=n+1$.
\item Sinon, $u_n=\frac{x^{2n+2}-1}{x^2-1}$.
\end{itemize}
\end{sol}



\exercice{nom={Calcul de rang}}
Résoudre le système linéaire dont la matrice $A$ est définie par
\[\forall i,j\in\intere{1}{n}\qsep \quad a_{i,j}\defeq
  \begin{cases}
  1 & \text{si} \abs{i-j}\leq 1,\\
  0 & \text{sinon.}
  \end{cases}\]
Quel est son rang ?



\magsubsection{Comatrice}


\exercice{nom={Comatrice du produit}}
Soit $\K$ un corps et $A,B\in\mat{n}{\K}$.
\begin{questions}
\question Montrer que les matrices $A-X I_n$ et $B-XI_n$ sont inversibles dans
  $\mat{n}{\fracK}$.
\question Montrer que $\Com(AB)=\Com(A)\Com(B)$.
\end{questions}










% \exercice{nom={Le déterminant de Vandermonde}}
% Soit $n\in\Ns$ et $x_0,\ldots,x_n$ $n+1$ scalaires. On définit le
% déterminant de Vandermonde par :
% \[V(x_0,x_1,\ldots,x_n)=
%   \begin{vmatrix}
%   1 & x_0 & x_0^2 & \dots & x_0^n\\
%   1 & x_1 & x_1^2 & \dots & x_1^n\\
%   \vdots & \vdots & \vdots & & \vdots\\
%   1 & x_n & x_n^2 & \dots & x_n^n
%   \end{vmatrix}\]
% Le but de cet exercice est de montrer par récurrence sur $n$ que :
% $$V(x_0,\ldots,x_n)=\prod_{0\leq i<j \leq n} (x_j-x_i)$$
% de deux manières distinctes.\\
% \begin{questions}
% \question Établir ce résultat en effectuant un pivot de Gauss.
% \question Établir ce résultat en étudiant le polynôme
%  \[V(x_0,\ldots,x_n,X)\]
% \end{questions}



% \exercice{nom={$\gl{n}{\Z}$}}
% Soit $M\in\mat{n}{\Z}$. Donner une condition nécessaire et
% suffisante pour que $M$ soit inversible et que
% $M^{-1}\in\mat{n}{\Z}$.
% \begin{sol}
% $\quad$
% \begin{questions}
% \question Il faut et il suffit que $\det A=\pm 1$.
% \end{questions}  
% \end{sol}

% \magsection{Déterminants en géométrie affine}
% \exercice{nom={Coordonnées barycentriques}}
% Soit $A,B,C$ trois points du plan non alignés.
% \begin{questions}
% \question On se donne $M_1,M_2,M_3$ trois points du plan dont les coordonnées
%   barycentriques dans le système $(A,B,C)$ sont respectivement
%   $(\alpha_1,\beta_1,\gamma_1)$, $(\alpha_2,\beta_2,\gamma_2)$ et
%   $(\alpha_3,\beta_3,\gamma_3)$. Montrer que $M_1,M_2$ et $M_3$ sont alignés si
%   et seulement si :
%   $$
%   \left|
%   \begin{array}{ccc}
%   \alpha_1 & \alpha_2 & \alpha_3\\
%   \beta_1 & \beta_2 & \beta_3\\
%   \gamma_1 & \gamma_2 & \gamma_3\\
%   \end{array}
%   \right|=0
%   $$
% \question Soit $A'$, $B'$, $C'$ trois points situés respectivement sur les
%   droites $\p{BC}$, $\p{CA}$ et $\p{AB}$ et distincts de $A,B,C$. Montrer que
%   $A'$, $B'$, $C'$ sont alignés si et seulement si~:
%   \[\frac{\overline{A'B}}{\overline{A'C}} \cdot
%     \frac{\overline{B'C}}{\overline{B'A}} \cdot
%     \frac{\overline{C'A}}{\overline{C'B}}=1\]
% \end{questions}


% \begin{sol}
% $\quad$
% \begin{questions}
% \question On peut se placer dans le repère $(A,\ve{AB},\ve{AC})$. Les points
%   $M_1,M_2,M_3$ sont alignés si et seulement si il existe $a,b,c\in\R$ tels
%   que $\p{a,b}\neq\p{0,0}$ et $M_1,M_2,M_3$ sont sur la droite d'équation
%   $ax+by+c=0$.
% \question Si on note $\alpha=\overline{A'B}/\overline{A'C}$, alors $A'$ est
%   barycentre de $B$ et $C$ avec les coefficients $1$ et $-\alpha$. Il suffit
%   ensuite d'appliquer la question précédente.
% \end{questions}
% \end{sol}


% \exercice{nom={Droites concourantes}}
% Dans le plan rapporté à un repère affine, on se donne $3$ droites
% $\mathcal{D}_1,\mathcal{D}_2$ et $\mathcal{D}_3$ d'équations respectives~:
% $$(\mathcal{D}_i) \quad \alpha_i x+\beta_i y+\gamma_i=0$$
% On suppose qu'au moins deux des droites ne sont pas parrallèles. Montrer
% que les trois droites sont concourantes si et seulement si :

% $$
% \left|
% \begin{array}{ccc}
% \alpha_1 & \alpha_2 & \alpha_3\\
% \beta_1 & \beta_2 & \beta_3\\
% \gamma_1 & \gamma_2 & \gamma_3\\
% \end{array}
% \right|=0
% $$




%END_BOOK

\end{document}




\chapter{Fonctions de deux variables}
\setcounter{numeroexercicecours}{1}
\documentclass{magnolia}


\magtex{tex_driver={pdftex},
        tex_packages={xypic}}
\magfiche{document_nom={Cours sur les fonctions de plusieurs variables},
          auteur_nom={François Fayard},
          auteur_mail={fayard.prof@gmail.com}}
\magcours{cours_matiere={maths},
          cours_niveau={mpsi},
          cours_chapitre_numero={24},
          cours_chapitre={Fonctions de deux variables}}
\magmisenpage{}
\maglieudiff{}
\magprocess

\begin{document}

%BEGIN_BOOK
\magtoc

\section{Limite, continuité}



\subsection{Notion d'ouvert}

\begin{proposition}
On note $\norme{.}$ la norme dérivant du produit scalaire usuel sur $\R^2$.
On rappelle que
\begin{eqnarray*}
\forall x\in\R^2, & & \norme{x}\geq 0\\
\forall x\in\R^2, & & \norme{x}=0 \ssi x=\p{0,0}\\
\forall x\in\R^2 \qsep \forall \lambda\in\R, & &
        \norme{\lambda x}=\abs{\lambda}\norme{x}\\
\forall x,y\in\R^2, & & \norme{x+y}\leq\norme{x}+\norme{y}\\
\forall x,y\in\R^2, & & \abs{\norme{x}-\norme{y}}\leq\norme{x-y}
\end{eqnarray*}
\end{proposition}

\begin{proposition}
\begin{eqnarray*}
\forall \p{x_1,x_2}\in\R^2, & & \abs{x_1}\leq\norme{\p{x_1,x_2}} \et 
       \abs{x_2}\leq\norme{\p{x_1,x_2}}\\
\forall \p{x_1,x_2}\in\R^2, & & \norme{\p{x_1,x_2}}\leq\abs{x_1}+\abs{x_2}
\end{eqnarray*}
\end{proposition}

\begin{definition}
Soit $x_0\in\R^2$ et $r>0$.
\begin{itemize}
\item On appelle \emph{boule ouverte} de centre $x_0$ et de rayon $r$ l'ensemble
  \[B_O\p{x_0,r}\defeq\enstq{x\in\R^2}{\norme{x-x_0}<r}.\]
\item On appelle \emph{boule fermée} de centre $x_0$ et de rayon $r$ l'ensemble
  \[B_F\p{x_0,r}\defeq\enstq{x\in\R^2}{\norme{x-x_0}\leq r}.\]
\end{itemize}
\end{definition}

\begin{definition}
On dit qu'une partie $\mathcal{U}$ de $\R^2$ est un \emph{ouvert} lorsque
\[\forall x\in\mathcal{U} \qsep \exists \eta>0 \qsep
  B_F\p{x,\eta} \subset \mathcal{U}.\]
\end{definition}

\begin{proposition}
$\quad$
\begin{itemize}
\item $\emptyset$ et $\R^2$ sont des ouverts.
\item Une union d'ouverts est un ouvert.
\item Une intersection finie d'ouverts est un ouvert.
\end{itemize}
\end{proposition}

\begin{proposition}
$\quad$
\begin{itemize}
\item Les boules ouvertes sont des ouverts.
\item Si $ax+by+c=0$ est l'équation d'une droite, le demi-plan d'équation
  \[ax+by+c>0\]
  est un ouvert.
\item Si on enlève un nombre fini de points à un ouvert, il reste ouvert.
\end{itemize}
\end{proposition}

\begin{exoUnique}
\exo Soit $x_0\in\R^2$. Montrer que
  \[\bigcap_{n\in\Ns} B_O\p{x_0,\frac{1}{n}}\]
  n'est pas un ouvert.
\end{exoUnique}

\begin{definition}
Soit $\mathcal{U}$ une partie de $\R^2$ et $a\in\R^2$. On dit que $a$ est
\emph{adhérent} à $\mathcal{U}$ lorsque
\[\forall\epsilon>0 \qsep B_F\p{a,\epsilon}\cap\mathcal{U}\neq\emptyset.\]
\end{definition}

\begin{exoUnique}
\exemple Soit $x_0\in\R^2$ et $r>0$. Montrer que l'ensemble des points adhérents à
  la boule ouverte de centre $x_0$ et de rayon $r$ est la boule fermée de
  même centre et de même rayon.
\end{exoUnique}


\begin{definition}
On appelle \emph{fonction réelle de deux variables} toute fonction définie sur une
partie ouverte $\mathcal{U}$ de $\R^2$ à valeurs dans $\R$.
\end{definition}

\subsection{Limite}



\begin{definition}
Soit $f:\mathcal{U}\subset\R^2\to\R$ et $a\in\R^2$ un point adhérent à $\mathcal{U}$.
\begin{itemize}
\item Soit $l\in\R$. On dit que $f(x)$ \emph{tend} vers $l$ lorsque $x$ tend
  vers $a$ lorsque
  \[\forall \epsilon>0 \qsep \exists \eta>0 \qsep \forall x\in\mathcal{U}
    \qsep \norme{x-a}\leq\eta \implique \abs{f(x)-l}\leq\epsilon.\]
\item On dit que $f(x)$ \emph{tend} vers $+\infty$ lorsque $x$ tend vers $a$
  lorsque
  \[\forall m\in\R \qsep \exists \eta>0 \qsep \forall x\in\mathcal{U}
    \qsep \norme{x-a}\leq\eta \implique f(x)\geq m.\]
\item On dit que $f(x)$ \emph{tend} vers $-\infty$ lorsque $x$ tend vers $a$
  lorsque
  \[\forall M\in\R \qsep \exists \eta>0 \qsep \forall x\in\mathcal{U}
    \qsep \norme{x-a}\leq\eta \implique f(x)\leq M.\]
\end{itemize}
\end{definition}

\begin{remarqueUnique}
\remarque Soit $\p{x_0,y_0}\in\R^2$. Alors
  \[x\tendvers{\p{x,y}}{\p{x_0,y_0}} x_0, \qquad
    y\tendvers{\p{x,y}}{\p{x_0,y_0}} y_0 \et
    \norme{\p{x,y}} \tendvers{\p{x,y}}{\p{x_0,y_0}} \norme{\p{x_0,y_0}}.\]
\end{remarqueUnique}

\begin{proposition}
Soit $f$ et $g:\mathcal{U}\subset\R^2\to\R$ 
 et $a\in\R^2$ un point adhérent à $\mathcal{U}$. On suppose que $f(x)$
tend vers $l_1\in\R$ et $g(x)$ tend vers $l_2\in\R$ lorsque $x$ tend vers $a$. Alors
\begin{itemize}
\item Si $\lambda,\mu\in\R$
  \[\lambda f(x)+\mu g(x)\tendvers{x}{a} \lambda l_1+\mu l_2.\]
\item On a
  \[f(x)g(x)\tendvers{x}{a} l_1 l_2.\]
\item Si $l_2\neq 0$, alors il existe une boule de centre $a$ sur
  laquelle $g$ ne s'annule pas. De plus
  \[\frac{f(x)}{g(x)}\tendvers{x}{a}\frac{l_1}{l_2}.\]
\end{itemize}
\end{proposition}

\begin{proposition}
Soit $f:\mathcal{U}\subset\R^2\to\R$, $a\in\R^2$ un
point adhérent à $\mathcal{U}$ et $g:\mathcal{D}_g\subset\R\to\R$ une fonction telle que
$f\p{\mathcal{U}}\subset\mathcal{D}_g$. On
suppose que
\[f(x)\tendvers{x}{a} l_f\in\Rbar \et
  g(y)\tendvers{y}{l_f} l_g\in\Rbar.\]
Alors
\[g\p{f(x)}\tendvers{x}{a} l_g.\]
\end{proposition}

\begin{proposition}
Soit $f$, $g$ et $h:\mathcal{U}\subset\R^2\to\R$ et $a\in\R^2$ un point adhérent à $\mathcal{U}$.
\begin{itemize}
\item On suppose que
  \[\forall x\in\mathcal{U} \qsep f(x)\leq g(x)\leq h(x).\]
  et que $f(x)$ et $h(x)$ admettent la même limite finie $l\in\R$ lorsque
  $x$ tend vers $a$. Alors
  \[f(x)\tendvers{x}{a} l.\]
\item On suppose qu'il existe $l\in\R$ tel que
  \[\forall x\in\mathcal{U} \qsep \abs{f(x)-l} \leq g(x).\]
  et que $g(x)$ tend vers 0 lorsque $x$ tend vers $a$. Alors
  \[f(x)\tendvers{x}{a} l.\]
\item On suppose que
  \[\forall x\in\mathcal{U} \qsep f(x)\leq g(x).\]
  \begin{itemize}
  \item Si $f(x)\tendvers{x}{a} +\infty$, alors $g(x)\tendvers{x}{a} +\infty$.
  \item Si $g(x)\tendvers{x}{a} -\infty$, alors $f(x)\tendvers{x}{a} -\infty$.
  \end{itemize}
\end{itemize}
\end{proposition}

\begin{remarqueUnique}
\remarque Si il existe une fonction $\phi:\RP\to\RP$ telle que
  \[\phi(r)\tendvers{r}{0}0 \et
    \forall r\geq 0 \qsep \forall \theta\in\R \qsep 
    \abs{f\p{a_1+r\cos\theta,a_2+r\sin\theta}-l}\leq \phi\p{r}\]
  alors $f(x)\tendvers{x}{a}l$.
\end{remarqueUnique}

\begin{exoUnique}
\exo Soit $f$ la fonction définie par
  \[\forall\p{x,y}\in\R^2\setminus\ens{\p{0,0}} \qsep
    f\p{x,y}\defeq\frac{x^2 y}{x^2+y^2}.\]
  Montrer que $f$ admet une limite en $\p{0,0}$.
\end{exoUnique}

\subsection{Continuité}

\begin{definition}
On dit qu'une fonction $f:\mathcal{U}\subset\R^2\to\R$ est \emph{continue} en $x_0\in\mathcal{U}$ lorsque
\[f(x)\tendvers{x}{x_0} f\p{x_0}.\]
\end{definition}

\begin{remarques}
\remarque On dit qu'une fonction $f:\mathcal{U}\subset\R^2\to\R$ est continue lorsqu'elle est continue en tout
  point de $\mathcal{U}$.
% \remarque Les fonctions $(x,y)\mapsto x$ et $(x,y)\mapsto y$ sont continues.
\remarque Les fonctions
  \[\dspappli{f}{\R^2}{\R}{(x,y)}{x} \quad\et\quad
    \dspappli{g}{\R^2}{\R}{(x,y)}{y}\]
  sont continues.
\end{remarques}

\begin{proposition}
Soit $f$ et $g:\mathcal{U}\subset\R^2\to\R$ deux fonctions continues en $x_0\in\mathcal{U}$.  Alors
\begin{itemize}
\item Si $\lambda,\mu\in\R$, $\lambda f+\mu g$ est continue en $x_0$.
\item $fg$ est continue en $x_0$.
\item Si $g\p{x_0}\neq 0$, il existe une boule de centre $x_0$ sur
  laquelle $g$ ne s'annule pas et $f/g$ est continue en $x_0$.
\end{itemize}
\end{proposition}

% \begin{exempleUnique}
% \remarque D'après les théorèmes usuels, la fonction $f$ définie sur $\R^2$ par
%   \[\forall \p{x,y}\in\R^2 \qsep f\p{x,y}\defeq\frac{xy}{1+x^2+y^2}\]
%   est continue sur $\R^2$.
% \end{exempleUnique}

\begin{proposition}
Soit $f:\mathcal{U}\subset\R^2\to\R$ et $g:\mathcal{V}\subset\R\to\R$ deux fonctions
telles que $f\p{\mathcal{U}}\subset\mathcal{V}$. Si
$f$ est continue en $x_0\in\mathcal{U}$ et $g$ est continue en $f\p{x_0}$, alors $g\circ f$
est continue en $x_0$.
\end{proposition}

\subsection{Application partielle}

\begin{definition}
Soit $f:\mathcal{U}\subset\R^2\to\R$ et $x\defeq\p{x_1,x_2}\in\mathcal{U}$. On
définit les fonctions réelles d'une variable réelle $f_{p_1}$ et $f_{p_2}$ par
\[f_{p_1}(t)\defeq f\p{t,x_2} \et f_{p_2}(t)\defeq f\p{x_1,t}.\]
$f_{p_1}$ et $f_{p_2}$ sont appelées \emph{applications partielles} de $f$ au point
$x=\p{x_1,x_2}$.
\end{definition}

\begin{exoUnique}
\exemple Soit $f$ la fonction définie sur $\R^2$ par
  \[\forall\p{x,y}\in\R^2 \qsep f\p{x,y}\defeq\frac{\p{2-y}\cos\p{xy}}{1+x^2}.\]
  Déterminer les applications partielles de $f$ en $(0,0)$.
  \begin{sol}
  les applications partielles en $\p{0,0}$ sont respectivement
  $f_1(t)=2/\p{1+x^2}$ et $f_2(t)=2-y$.    
  \end{sol}
\end{exoUnique}

\begin{proposition}
\begin{itemize}
\item Si $f(x)$ tend vers $l\in\Rbar$ lorsque $x$ tend vers $a\defeq\p{a_1,a_2}$,
  $f_{p_1}(t)$ tend vers $l$ lorsque $t$ tend vers $a_1$ et $f_{p_2}(t)$
  tend vers $l$ lorsque $t$ tend vers $a_2$.
\item Si $f$ est continue en $x\defeq\p{x_1,x_2}$, $f_{p_1}$ est continue en $x_1$ et
  $f_{p_2}$ est continue en $x_2$.
\end{itemize}
\end{proposition}

\begin{remarqueUnique}
\remarque Nous verrons que les réciproques de ces théorèmes sont fausses. Par
  exemple, $f_{p_1}$ peut être continue en $x_1$ et $f_{p_2}$ peut être continue
  en $x_2$ sans que $f$ soit continue en $\p{x_1,x_2}$.
\end{remarqueUnique}

\begin{exoUnique}
\exo Montrer que la fonction $f$ définie sur $\R^2\setminus\ens{\p{0,0}}$
  par
  \[\forall\p{x,y}\in\R^2\setminus\ens{\p{0,0}} \qsep
    f\p{x,y}\defeq\frac{x^2-y^2}{x^2+y^2}\]
  n'a pas de limite en $\p{0,0}$.
  \begin{sol}
  La limite selon $y=0$ est 1. Celle selon $x=0$ est $-1$.
  \end{sol}
\end{exoUnique}

\subsection{Extension aux fonctions à valeurs dans $\R^2$}

\begin{definition}
Soit $f:\mathcal{U}\subset\R^2\to\R^2$. Les fonctions $f_{1}$ et $f_{2}:\mathcal{U}\subset\R^2\to\R$ définies par
\[\forall x\in\mathcal{U} \qsep f(x)=\p{f_{1}(x),f_{2}(x)}.\]
sont appelées \emph{fonctions coordonnées} de $f$.
\end{definition}

\begin{remarqueUnique}
\remarque On étend les notions de limite et de continuité aux fonctions à valeurs dans $\R^2$ en remplaçant
  les valeurs absolues par des normes dans les définitions données plus haut.
  Par exemple, si $f:\mathcal{U}\subset\R^2\to\R^2$, $a\in\R^2$ est un
  point adhérent à $\mathcal{U}$ et $l\in\R^2$, on dit que $f(x)$ tend vers $l$ lorsque $x$ tend vers $a$
  lorsque
  \[\forall \epsilon>0\qsep \exists \eta>0\qsep \forall x\in\mathcal{U}\qsep
    \norme{x-a}\leq\eta \implique \norme{f(x)-l}\leq\epsilon.\]
\end{remarqueUnique}

\begin{proposition}
Soit $f:\mathcal{U}\subset\R^2\to\R^2$.
\begin{itemize}
\item Si $a$ est un point adhérent à $\mathcal{U}$ et $l\defeq\p{l_1,l_2}\in\R^2$
  alors $f(x)$ tend vers $l$ lorsque $x$ tend vers $a$ si et seulement si
  \[f_{1}(x) \tendvers{x}{a} l_1 \et f_{2}(x) \tendvers{x}{a} l_2.\]
\item Si $x_0\in\mathcal{U}$, $f$ est continue en $x_0$ si et seulement si
  $f_{1}$ et $f_{2}$ le sont.
\end{itemize}
\end{proposition}

\begin{proposition}
Soit $f:\mathcal{U}\subset\R^p\to\R^q$ et $g:\mathcal{V}\subset\R^q\to\R^r$
avec $p,q,r\in\intere{1}{2}$. On suppose que
$f\p{\mathcal{U}}\subset\mathcal{V}$.
\begin{itemize}
\item Soit $a\in\R^p$ un point adhérent à $\mathcal{U}$. Si $f(x)$ tend
  vers $b\in\R^q$ lorsque $x$ tend vers $a$ et $g(y)$ tend vers $l$ lorsque
  $y$ tend vers $b$, alors
  \[g\p{f(x)}\tendvers{x}{a} l.\]
\item Soit $x_0\in\mathcal{U}$. Si $f$ est continue en $x_0$ et si $g$ est
  continue en $f\p{x_0}$, alors $g\circ f$ est continue en $x_0$.
\end{itemize}
\end{proposition}

\begin{remarqueUnique}
\remarque Soit $f:\mathcal{U}\subset\R^2\to\R$ et $a$ un point adhérent à
  $\mathcal{U}$ en lequel $f$ admet une limite $l\in\Rbar$. Si
  $g:I\to\mathcal{U}\subset\R^2$ où $I$ est un intervalle de $\R$, $b$ est une
  borne de $I$ et $g(t)$ tend vers $a$ lorsque $t$ tend vers $b$, alors
  \[f\p{g(t)}\tendvers{t}{b}  l.\]
  On peut ainsi, en choisissant différentes fonctions $g$, se faire une idée de
  la limite éventuelle de $f$ en $a$, ou prouver que $f$ n'admet pas de limite
  en $a$.
\end{remarqueUnique}


\begin{exos}
\exo Soit $f$ la fonction définie par
  \[\forall\p{x,y}\in\R^2\setminus\ens{\p{0,0}} \qsep f\p{x,y}\defeq
    \frac{xy}{x^2+y^2}.\]
  Étudier la limite éventuelle de $f$ en $\p{0,0}$.
  \begin{sol}
  Passer en coordonnées polaires.
  \end{sol}
\exo Soit $f$ la fonction définie sur ${\RPs}^2$ par
  \[\forall (x,y)\in{\RPs}^2 \qsep f\p{x,y}\defeq x^y.\]
  Étudier la limite éventuelle de $f$ en $\p{0,0}$.
  \begin{sol}
  Si on arrive en diagonale, la limite est $1$. Cependant, si $x=t$ et
  $y=-1/\p{\ln t}$, la limite est $e^{-1}$.
  \end{sol}
\end{exos}

\section{Dérivation}



\subsection{Dérivée partielle}

\begin{definition}
Soit $f:\mathcal{U}\subset\R^2\to\R$ et $x\defeq\p{x_1,x_2}\in\mathcal{U}$.
\begin{itemize}
\item On dit que $f$ admet \emph{une dérivée partielle par rapport à la première variable} en
  $x$ lorsque
  \[\frac{f\p{x_1+t,x_2}-f\p{x_1,x_2}}{t}\]
  admet une limite finie lorsque $t$ tend vers 0. Si tel est le cas, cette
  limite est notée
  \[\parfrac{f}{x_1}(x).\]
\item On dit que $f$ admet \emph{une dérivée partielle par rapport à la seconde variable} en
  $x$ lorsque
  \[\frac{f\p{x_1,x_2+t}-f\p{x_1,x_2}}{t}\]
  admet une limite finie lorsque $t$ tend vers 0. Si tel est le cas, cette
  limite est notée
  \[\parfrac{f}{x_2}(x).\]
\end{itemize}
\end{definition}


\begin{remarqueUnique}
\remarque Soit $f:\R^2\to\R$ une fonction admettant des dérivées partielles en tout point
  par rapport à la première et à la seconde variable, $x\defeq (x_1,x_2)$
  et $f_{p_1}$, $f_{p_2}:\R\to\R$ les fonctions partielles définies par
  \[\forall t\in\R\qsep f_{p_1}(t)\defeq f(t, x_2)
    \quad\et\quad
    f_{p_2}(t)\defeq (x_1,t).\]
  Alors $f_{p_1}$ et $f_{p_2}$ sont dérivables sur $\R$ et
  \[\forall t\in\R\qsep f_{p_1}'(t)=\parfrac{f}{x_1}(t,x_2) \quad\et\quad
    f_{p_2}'(t)=\parfrac{f}{x_2}(x_1,t).\]
\end{remarqueUnique}

\begin{definition}
Soit $f:\mathcal{U}\subset\R^2\to\R$, $x\in\mathcal{U}$ et $h\in\R^2$. On dit
que $f$ admet une \emph{dérivée en $x$ selon le vecteur $h$} lorsque
\[\frac{f\p{x+th}-f(x)}{t}\]
admet une limite finie lorsque $t$ tend vers 0. Si tel est le cas, cette limite
est notée
\[{\rm D}f_x(h).\]
\end{definition}

\begin{remarques}
\remarque Si $e_1\defeq(1,0)$, $f$ admet une dérivée en $x$
  selon le vecteur $e_1$ si et seulement si elle admet une dérivée partielle par rapport à la
  première variable en $x$. De plus, si tel est le cas
  \[{\rm D}f_{x}(e_1)=\parfrac{f}{x_1}(x).\]
  De même, $f$ admet une dérivée en $x$ selon le vecteur $e_2\defeq(0,1)$ si et seulement
  si elle est admet une dérivée partielle par rapport à la seconde variable en $x$.
\remarque Soit $f:\mathcal{U}\subset\R^2\to\R$. Alors $f$ admet une dérivée en tout
  $x\in\mathcal{U}$ selon le vecteur nul et ${\rm D}f_x(0)=0$.
  De plus, si $f$ admet une dérivée en $x$ selon le vecteur $h\in\R^2$ alors, pour tout
  $\lambda\in\R$, elle admet une dérivée en $x$ selon le vecteur $\lambda h$
  et
  \[{\rm D}f_x(\lambda h)=\lambda {\rm D}f_x(h).\]
\end{remarques}

\begin{exoUnique}
\exo Soit $f$ la fonction définie sur $\R^2$ par
  \[\forall\p{x,y}\in\R^2 \qsep f\p{x,y}\defeq
    \begin{cases}
    \frac{x^2 y}{x^4+y^2} & \text{si $\p{x,y}\neq\p{0,0}$}\\
    0 & \text{si $\p{x,y}=0$.}
    \end{cases}\]
  Montrer que $f$ est dérivable en $\p{0,0}$ selon tout vecteur mais n'est pas
  continue en $\p{0,0}$.
  \begin{sol}
  Calculer la dérivée selon le vecteur $\p{\cos\theta,\sin\theta}$. On trouve
  0 si $\sin\theta=0$ et $\cos^2 \theta/\sin\theta$ sinon. Cependant
  $f\p{x,x^2}=1/2$ qui ne tend pas vers 0 en 0, donc $f$ n'est pas continue
  en 0.
  \end{sol}
\end{exoUnique}

\begin{definition}
On dit qu'une fonction $f:\mathcal{U}\subset\R^2\to\R$ est de \emph{classe
$\classec{1}$} lorsque
\[\parfrac{f}{x_1} \et \parfrac{f}{x_2}\]
sont définies et continues sur $\mathcal{U}$.
\end{definition}

\begin{remarques}
\remarque Si $f$ est de classe $\classec{1}$, les applications partielles
  sont de classe $\classec{1}$.
\remarque Les fonctions
  \[\dspappli{f}{\R^2}{\R}{(x,y)}{x} \quad\et\quad
    \dspappli{g}{\R^2}{\R}{(x,y)}{y}\]
  sont de classe $\classec{1}$.
\end{remarques}

\begin{proposition}
Soit $f$ et $g:\mathcal{U}\subset\R^2\to\R$ deux fonctions de classe $\classec{1}$.
Alors
\begin{itemize}
\item Si $\lambda,\mu\in\R$, $\lambda f+\mu g$ est de classe $\classec{1}$ et
  \[\forall i\in\intere{1}{2} \qsep \forall x\in\mathcal{U} \qsep
    \parfrac{\p{\lambda f+\mu g}}{x_i}(x)=\lambda\parfrac{f}{x_i}(x)+\mu
    \parfrac{g}{x_i}(x).\]
\item $fg$ est de classe $\classec{1}$ et
  \[\forall i\in\intere{1}{2} \qsep \forall x\in\mathcal{U} \qsep
    \parfrac{\p{fg}}{x_i}(x)=\parfrac{f}{x_i}(x)g(x)+
    f(x)\parfrac{g}{x_i}(x).\]
\item Si $g$ ne s'annule pas sur $\mathcal{U}$, $f/g$ est de classe
  $\classec{1}$ et
  \[\forall i\in\intere{1}{2} \qsep \forall x\in\mathcal{U} \qsep
    \parfrac{\p{\dsp\frac{f}{g}}}{x_i}(x)=
    \frac{\dsp\parfrac{f}{x_i}(x)g(x)-
    f(x)\dsp\parfrac{g}{x_i}(x)}{g(x)^2}.\]
\end{itemize}
\end{proposition}

\begin{proposition}
Soit $f:\mathcal{U}\subset\R^2\to\R$ et $g:\mathcal{V}\subset\R\to\R$ deux fonctions de
classe $\classec{1}$ telles que $f(\mathcal{U})\subset\mathcal{V}$. Alors
$g\circ f$ est de classe $\classec{1}$ et
\[\forall i\in\intere{1}{2}\qsep \forall x\in\mathcal{U}\qsep
  \parfrac{(g\circ f)}{x_i}(x)=g'(f(x))\parfrac{f}{x_i}(x).\]
\end{proposition}

\begin{exoUnique}
\exo Trouver l'ensemble des fonctions $f:\R^2\to\R$ de classe $\classec{1}$ telles que
  \[\forall\p{x,y}\in\R^2 \qsep \parfrac{f}{x} \p{x,y}=0.\]
\exo Déterminer l'ensemble des fonctions $f:\R^2\to\R$ de classe $\classec{1}$ telles que
  \[\forall\p{x,y}\in\R^2 \qsep \parfrac{f}{y} \p{x,y}= x+y.\]
  \begin{sol}
  On trouve $f\p{x,y}=y^2/2+xy+\phi(x)$ où $\phi$ est une fonction de classe
  $\classec{1}$ quelconque sur $\R$.
  \end{sol}
\end{exoUnique}

\subsection{Développement limité, gradient}

\begin{definition}
Soit $f:\mathcal{U}\subset\R^2\to\R$ une fonction définie en $(0,0)$. On dit que $f(h)$ est \emph{négligeable} devant $\norme{h}$ en $\p{0,0}$ lorsque
\[\forall\epsilon>0 \qsep \exists \eta>0 \qsep \forall h\in\mathcal{U}
  \qsep \norme{h}\leq\eta \implique \abs{f(h)}\leq\epsilon\norme{h}.\]
Si tel est le cas, on note
\[f(h)=\petito{h}{\p{0,0}}{\norme{h}}.\]
\end{definition}

% \begin{remarqueUnique}
% \remarque Si $f(h)$ est négligeable devant $\norme{h}$ en $(0,0)$, alors $f(h)$ tend vers 0 lorsque $h$ tend vers
%   $(0,0)$ et $f$ admet une dérivée nulle selon tout vecteur en $(0,0)$.
% \end{remarqueUnique}

\begin{proposition}
Soit $f:\mathcal{U}\subset\R^2\to\R$ une fonction de classe $\classec{1}$ et
$x\in\mathcal{U}$. Alors, en notant $h=\p{h_1,h_2}$
\[f\p{x+h}=f(x)+\parfrac{f}{x_1}(x)h_1+\parfrac{f}{x_2}(x)h_2+
  \petito{h}{\p{0,0}}{\norme{h}}.\]
\end{proposition}

\begin{remarqueUnique}
\remarque Si $f:\mathcal{U}\subset\R^2\to\R$ est une fonction de classe $\classec{1}$ et $(x_0,y_0)\in\mathcal{U}$ alors,
  le plan d'équation
  \[z=f(x_0,y_0)+\parfrac{f}{x}(x_0,y_0)(x-x_0)+\parfrac{f}{y}(x_0,y_0)(y-y_0)\]
  est appelé \emph{plan tangent} en $(x_0,y_0)$ à la surface d'équation $z=f(x,y)$.
\end{remarqueUnique}

\begin{proposition}
Soit $f:\mathcal{U}\subset\R^2\to\R$ une fonction de classe $\classec{1}$.
Alors $f$ est continue.
\end{proposition}

\begin{proposition}
Soit $f:\mathcal{U}\subset\R^2\to\R$ une fonction de classe $\classec{1}$ et
$x\in\mathcal{U}$. Alors $f$ admet une dérivée en $x$ selon tout vecteur
$h\defeq\p{h_1,h_2}$ et
\[{\rm D}f_x(h)=\parfrac{f}{x_1}(x)h_1 
  +\parfrac{f}{x_2}(x)h_2.\]
Autrement dit, ${\rm D}f_x$ est une forme linéaire sur $\R^2$.
\end{proposition}



\begin{exoUnique}
\exo Soit $f$ la fonction définie sur $\R^2$ par
  \[\forall\p{x,y}\in\R^2 \qsep f\p{x,y}\defeq\frac{xy}{1+x^2}.\]
  Montrer que $f$ est de classe $\classec{1}$ et calculer
  $\text{D}f_{\p{x,y}}\p{h_1,h_2}$.
  \begin{sol}
  On trouve
  \[\text{d}f_{\p{x,y}}\p{h_1,h_2}=y\cdot\frac{1-x^2}{\p{1+x^2}^2}h_1+
    \frac{x}{1+x^2}h_2\]
  ce que l'on écrit aussi
  \[\text{d}f=y\cdot\frac{1-x^2}{\p{1+x^2}^2}\text{d}x+
    \frac{x}{1+x^2}\text{d}y\]
  \end{sol}
\end{exoUnique}

\begin{definition}
Soit $f:\mathcal{U}\subset\R^2\to\R$ une fonction de classe $\classec{1}$ et
$x\in\mathcal{U}$. On appelle \emph{gradient} de $f$ en $x$ l'unique vecteur
$\nabla f(x)\in\R^2$ tel que
\[\forall h\in\R^2 \qsep {\rm D}f_{x}(h)=\ps{\nabla f(x)}{h}.\]
Autrement dit
\[\nabla f(x)=\p{\parfrac{f}{x_1}(x),\parfrac{f}{x_2}(x)}.\]
\end{definition}

\begin{remarques}
\remarque Le symbole $\nabla$ est un delta majuscule inversé. Il se prononce \og nabla \fg en référence au nom grec désignant une harpe phénicienne.
\remarque Pour tout $\theta\in\R$, on pose $u_\theta\defeq(\cos\theta,\sin\theta)$. Alors,
d'après Cauchy-Schwarz
\[\forall \theta\in\R\qsep {\rm D}f_x(u_\theta)=\ps{\nabla f(x)}{u_\theta} \leq \norme{\nabla f(x)}\norme{u_\theta}=\norme{\nabla f(x)}\]
cette inégalité étant une égalité si et seulement si $\nabla f(x)$ et $u_\theta$ sont
positivement liés. On en déduit que $\nabla f(x)$ donne la direction dans laquelle
$f$ croît le plus vite.
\end{remarques}


\subsection{Dérivation des fonctions composées}

\begin{proposition}
Soit $f:\mathcal{U}\subset\R^2\to\R$ et $g:\mathcal{V}\subset\R\to\R$ deux
  fonctions de classe $\classec{1}$ telles que
  $f\p{\mathcal{U}}\subset\mathcal{V}$. Alors $g\circ f$ est de classe
  $\classec{1}$ et
  \[\forall x\in\mathcal{U} \qsep
    \cro{\nabla (g\circ f)} (x)=g'\p{f(x)}\nabla f(x).\]
  On écrit aussi, de manière abusive
  \[\nabla\p{g\circ f}=\frac{{\rm d}g}{{\rm d}f}\nabla f.\]
\end{proposition}


\begin{exoUnique}
\exo Calculer $\nabla f$, où $f$ est définie sur
  $\R^2\setminus\ens{\p{0,0}}$ par
  \[f\p{x,y}\defeq\frac{1}{\sqrt{x^2+y^2}}.\]
  \begin{sol}
  On trouve
  \[\nabla\p{\frac{1}{r}}=-\frac{1}{r^2}\ve{u_\theta}\]
  \end{sol}
\end{exoUnique}

% \begin{definition}
% Soit $p,q\in\intere{1}{2}$ et $f:\mathcal{U}\subset\R^p\to\R^q$. On dit que $f$
% est de classe $\classec{1}$ sur $\mathcal{U}$ lorsque quels que soient
% $j\in\intere{1}{p}$ et $i\in\intere{1}{q}$, l'application
% \[\parfrac{f_i}{x_j}\]
% est définie et continue sur $\mathcal{U}$. Si tel est la cas, on définit pour
% tout $x\in\mathcal{U}$ la \emph{matrice jacobienne} de $f$ en $x$ par
% \[J_f(x)\defeq \p{\parfrac{f_i}{x_j}}_{\substack{1\leq i\leq q\\1\leq j\leq p}}
%   \in\mat{q,p}{\R}\]
% \end{definition}

% \begin{exoUnique}
% \exo Soit $f$ la fonction définie sur $\R^2$ par
%   \[\forall\p{r,\theta}\in\R^2 \qsep f\p{r,\theta}\defeq \p{r\cos\theta,r\sin\theta}.\]
%   Calculer la jacobienne de $f$ en $\p{r,\theta}$.
% \end{exoUnique}

% \begin{proposition}
% Soit $p,q,r\in\intere{1}{2}$, $f:\mathcal{U}\subset\R^p\to\R^q$ et
% $g:\mathcal{V}\subset\R^q\to\R^r$ tels que $f\p{\mathcal{U}}\subset\mathcal{V}$.
% Si $f$ et $g$ sont $\classec{1}$, alors $g\circ f$ est $\classec{1}$ et
% \[J_{g\circ f}(x)=J_g\p{f(x)}J_f(x).\]
% \end{proposition}

\begin{proposition}
Soit $f:\mathcal{U}\subset\R\to\R^2$ et $g:\mathcal{V}\subset\R^2\to\R$ deux
fonctions de classe $\classec{1}$ telles que
$f\p{\mathcal{U}}\subset\mathcal{V}$. Alors $g\circ f$ est de classe
$\classec{1}$ et
\begin{eqnarray*}
\forall t\in\mathcal{U} \qsep \p{g\circ f}'(t)
&=&\parfrac{g}{x_1}\p{f(t)}f_1'(t)+\parfrac{g}{x_2}\p{f(t)}f_2'(t)\\
&=&\ps{\nabla g \p{f(t)}}{f'(t)}
\end{eqnarray*}
On écrit aussi, de manière abusive
\begin{eqnarray*}
\frac{{\rm d}\p{g\p{f_1,f_2}}}{{\rm d} t}
&=& \parfrac{g}{f_1}\frac{{\rm d}f_1}{{\rm d} t}+
    \parfrac{g}{f_2}\frac{{\rm d}f_2}{{\rm d} t}\\
&=& \ps{\nabla g}{\frac{{\rm d}f}{{\rm d}t}}
\end{eqnarray*}
\end{proposition}

\begin{exoUnique}
\exo Soit $f:\R^2\to\R$ une fonction de classe $\classec{1}$ et
  $\phi$ la fonction définie sur $\R$ par
  \[\forall t\in\R \qsep \phi(t)\defeq f\p{t^2,t^3}.\]
  Montrer que $\phi$ est de classe $\classec{1}$ et calculer $\phi'(t)$.
\end{exoUnique}

\begin{proposition}
Soit $f:\mathcal{U}\subset\R^2\to\R$ une fonction de classe $\classec{1}$ et
$M:I\to\R^2$ une fonction de classe $\classec{1}$ à valeurs dans 
$\mathcal{U}$. On suppose qu'il existe $c\in\R$ tel que
\[\forall t\in I \qsep f\p{M(t)}=c.\]
Alors, pour tout $t\in I$, $M'(t)$ est orthogonal au gradient
de $f$ en $M(t)$.
\end{proposition}

\begin{remarqueUnique}
\remarque Si $f:\mathcal{U}\subset\R^2\to\R$, on appelle \emph{ligne de niveau} de hauteur $c\in\R$
  l'ensemble
  \[\mathcal{L}_c\defeq\enstq{(x,y)\in\mathcal{U}}{f(x,y)=c}.\]
  La proposition précédente nous assure que le gradient de $f$ est orthogonal à ses lignes de niveau.
\end{remarqueUnique}

\begin{proposition}
Soit $f:\mathcal{U}\subset\R^2\to\R^2$ et $g:\mathcal{V}\subset\R^2\to\R$ deux
fonctions de classe $\classec{1}$ telles que
$f\p{\mathcal{U}}\subset\mathcal{V}$. Alors $g\circ f$ est de classe
$\classec{1}$ et
\[\forall x\in\mathcal{U} \qsep
  \begin{cases}
  \dsp\parfrac{\p{g\circ f}}{x_1}(x) = \dsp\parfrac{g}{y_1}\p{f(x)}
  \dsp\parfrac{f_1}{x_1}(x)+\dsp\parfrac{g}{y_2}\p{f(x)}
  \dsp\parfrac{f_2}{x_1}(x) &\\
  \dsp\parfrac{\p{g\circ f}}{x_2}(x) = \dsp\parfrac{g}{y_1}\p{f(x)}
  \dsp\parfrac{f_1}{x_2}(x)+\dsp\parfrac{g}{y_2}\p{f(x)}
  \dsp\parfrac{f_2}{x_2}(x) &
  \end{cases}\]
  On écrit aussi, de manière abusive
  \[\dsp\parfrac{\p{g\p{f_1,f_2}}}{x_1}=\parfrac{g}{f_1}\parfrac{f_1}{x_1}
    +\parfrac{g}{f_2}\parfrac{f_2}{x_1} \et
    \dsp\parfrac{\p{g\p{f_1,f_2}}}{x_2}=\parfrac{g}{f_1}\parfrac{f_1}{x_2}
    +\parfrac{g}{f_2}\parfrac{f_2}{x_2}.\]
\end{proposition}

\begin{exos}
\exo Soit $f:\R^2\to\R$ une fonction de classe $\classec{1}$. Montrer que la fonction
  $h:\R^2\to\R$ définie par
  \[\forall (x,y)\in\R^2\qsep h\p{x,y}\defeq f\p{2xy,x}\]
  est de classe $\classec{1}$ et calculer ses dérivées partielles.
  \begin{sol}
  On trouve
   \[\parfrac{h}{x}\p{x,y}=2y\parfrac{f}{x}\p{2xy,x}+\parfrac{f}{y}\p{2xy,x}
     \et \parfrac{h}{y}\p{x,y}=2x\parfrac{f}{x}\p{2xy,x}\]
  \end{sol}
% \exo Calcul en coordonnées polaires. On trouve
%   \[\syslin{\parfrac{}{r}&=&\cos\theta\parfrac{}{x}+&\sin\theta\parfrac{}{y}\cr
%     \frac{1}{r}\parfrac{}{\theta}&=&-\sin\theta\parfrac{}{x}+&
%     \cos\theta\parfrac{}{y}}\]
%   Donc
%   \[\syslin{\parfrac{}{x}&=&\cos\theta\parfrac{}{r}-&\sin\theta\frac{1}{r}
%     \parfrac{}{\theta}\cr
%     \parfrac{}{y}&=&\sin\theta\parfrac{}{r}+&
%     \cos\theta\frac{1}{r}\parfrac{}{\theta}}\]
%   Donc
%   \[\nabla f=\parfrac{f}{r}\ve{u_\theta}+\frac{1}{r}\parfrac{f}{\theta}
%     \ve{v_\theta}\]
%   En particulier
%   \[\nabla\p{\frac{1}{r}}=-\frac{1}{r^2}\ve{u_\theta}\]
\exo Soit $\p{a,b}\in\R^2\setminus\ens{\p{0,0}}$. On souhaite déterminer les fonctions
  $f:\R^2\to\R$ de classe $\classec{1}$ telles que
  \[(E) \quad \forall\p{x,y}\in\R^2 \qsep
    a\parfrac{f}{x}\p{x,y}+b\parfrac{f}{y}\p{x,y}=0.\]
  \begin{questions}
  \question Montrer que l'application
    \[\dspappli{\phi}{\R^2}{\R^2}{(u,v)}{(au-bv,bu+av)}\]
    est bijective et calculer $\phi^{-1}$.
  \question Soit $f:\R^2\to\R$ une fonction de classe $\classec{1}$. On définit
    la fonction $g:\R^2\to\R$ par
    \[\forall(u,v)\in\R^2\qsep g(u,v)\defeq f(au-bv,bu+av).\]
    Montrer que $f$ est solution de $(E)$ si et seulement si
    \[\forall(u,v)\in\R^2\qsep \parfrac{g}{u}(u,v)=0.\]
  \question Conclure.
  \end{questions}
  % On pourra introduire la fonction $g:\R^2\to\R$ définie par
  \begin{sol}%
  On fait le changement de variable
  $g\p{u,v}=f\p{au-bv,bu+av}$. On trouve $\parfrac{g}{u}=0$, donc
  $g=\phi\p{v}$ donc $f=\psi\p{bx-ay}$.    
  \end{sol}
\end{exos}



% \begin{exemples}
% \exemple Soit $\mathcal{E}$ l'ellipse de foyer $F_1$ et $F_2$ et de demi-grand
%   axe $a$. Si $M$ est un point de $\mathcal{E}$ la tangente en $M$ est
%   la bissectrice extérieure des demi-droites $[F_1M)$ et $[F_2M)$.
% \end{exemples}

\subsection{Extrémum d'une fonction de deux variables}

\begin{definition}
Soit $f:\mathcal{U}\subset\R^2\to\R$ et $x_0\in\mathcal{U}$. On dit que
\begin{itemize}
\item $f$ présente un \emph{maximum global} en $x_0$ lorsque
  \[\forall x\in\mathcal{U} \qsep f(x)\leq f\p{x_0}.\]
\item $f$ présente un \emph{maximum local} en $x_0$ lorsque
  \[\exists r>0\qsep \forall x\in\mathcal{U} \qsep \norme{x-x_0}\leq r \implique
    f(x)\leq f\p{x_0}.\]
\item $f$ présente un \emph{minimum global} en $x_0$ lorsque
  \[\forall x\in\mathcal{U} \qsep f(x)\geq f\p{x_0}.\]
\item $f$ présente un \emph{minimum local} en $x_0$ lorsque
  \[\exists r>0\qsep \forall x\in\mathcal{U} \qsep \norme{x-x_0}\leq r \implique
    f(x)\geq f\p{x_0}.\]
\end{itemize}
\end{definition}

\begin{remarques}
\remarque Un extrémum global est un extrémum local, la réciproque étant fausse en général.
\remarque Si $f:\R^2\to\R$ admet un minimum local en
  $(x_1,x_2)\in\R^2$, alors les fonctions partielles
  \[\dspappli{f_{p_1}}{\R}{\R}{t}{f(t,x_2)} \quad\et\quad
    \dspappli{f_{p_2}}{\R}{\R}{t}{f(x_1,t)}\]
  admettent un minimum local, respectivement en $x_1$ et $x_2$.
  Cependant, la réciproque est fausse comme le montre l'exemple de la fonction
  $f:\R^2\to\R$ définie par
  \[\forall (x,y)\in\R^2\qsep f(x,y)=x^2-3xy+y^2.\]
\remarque Plus généralement, si $f:\R^2\to\R$ admet un minimul local en $x_0\in\R^2$ alors,
  quel que soit $h\in\R^2$, la fonction $\phi:\R\to\R$ définie par
  \[\forall t\in\R\qsep \phi(t)\defeq f(x_0+th)\]
  admet un minimum local en $0$. 
\end{remarques}

\begin{definition}
Soit $f:\mathcal{U}\subset\R^2\to\R$ une fonction de classe $\classec{1}$. On dit que $x\in\mathcal{U}$
est un point \emph{critique} de $f$ lorsque
\[\parfrac{f}{x_1}\p{x}=0 \et \parfrac{f}{x_2}\p{x}=0.\]
\end{definition}

\begin{proposition}
Soit $f:\mathcal{U}\subset\R^2\to\R$ une fonction de classe $\classec{1}$ définie sur un ouvert $\mathcal{U}$.
Si $f$ présente un extrémum local en $x_0$ alors $x_0$ est un point critique pour $f$.
\end{proposition}

\begin{remarques}
\remarque Les extrémums locaux d'une fonction $f:\mathcal{U}\subset\R^2\to\R$ de classe $\classec{1}$ définie
  sur un ouvert $\mathcal{U}$ sont donc à chercher parmi les points critiques de $f$.
\remarque La réciproque est fausse. Par exemple la fonction $f:\R^2\to\R$ définie par
  \[\forall (x,y)\in\R^2\qsep f(x,y)\defeq x^2-y^2\]
  admet $(0,0)$ pour point critique alors qu'elle ne présente pas d'extrémum local en $(0,0)$.
\end{remarques}

\begin{exoUnique}
\exo Soit $f$ la fonction définie sur $\R^2$ par
  \[\forall\p{x,y}\in\R^2 \qsep f\p{x,y}\defeq -x^2-6xy+9y^2+18x-18y+1.\]
  Déterminer les extrémums éventuels de $f$.
  \begin{sol}
  On trouve un unique point critique $x=3$ et $y=2$. On pose $x=3+u$ et
  $y=2+v$. On a $f\p{3+u,2+v}=10+9v^2-6uv-u^2=\p{3v-u}^2-2u^2$. On est en
  un point selle.
  \end{sol}
\end{exoUnique}



%END_BOOK

\end{document}
\section{Exercices}
\setcounter{numeroexercice}{1}
\documentclass{magnolia}

\magtex{tex_driver={pdftex}}
\magfiche{document_nom={Exercices sur les fonctions de deux variables},
          auteur_nom={François Fayard},
          auteur_mail={fayard.prof@gmail.com}}
\magexos{exos_matiere={maths},
         exos_niveau={mpsi},
         exos_chapitre_numero={3},
         exos_theme={Fonctions de deux variables}}
\magmisenpage{misenpage_presentation={tikzvelvia},
              misenpage_format={a4},
              misenpage_nbcolonnes={1},
              misenpage_sol={non}}
\maglieudiff{lieu_lycee={Aux Lazaristes},
             lieu_classe={MPSI 1},
             lieu_annee={2019--2020}}
\magprocess

\begin{document}

%BEGIN_BOOK
\magsection{Limite, continuité}
\magsubsection{Notion d'ouvert}
\magsubsection{Limite}
\exercice{nom={Calcul de limites}}
Déterminer les limites, si elles existent, des fonctions suivantes en $\p{0,0}$.
\[f\p{x,y}\defeq\p{x+y}\sin\p{\frac{1}{x^2+y^2}}, \qquad
  f\p{x,y}\defeq\frac{x^2-y^2}{x^2+y^2}, \qquad
  f\p{x,y}\defeq\frac{\abs{x+y}}{x^2+y^2},\]
\[f\p{x,y}\defeq\frac{x^3+y^3}{x^2+y^2}, \qquad
  f\p{x,y}\defeq\frac{x^2+y^2-1}{x}\sin x, \qquad
  f\p{x,y}\defeq x^y.\]
\magsubsection{Continuité}
\exercice{nom={Fonction définie par morceaux}}
Soit $f$ la fonction définie sur $\R^2$ par
\[\forall (x,y)\in\R^2 \qsep f(x,y)\defeq
  \begin{cases}
  \frac{1}{2}x^2+y^2-1 & \text{si $x^2+y^2>1$,}\\
  -\frac{1}{2}x^2 & \text{si $x^2+y^2\leq 1$.}
  \end{cases}
\]
Montrer que $f$ est continue en tout point de $\R^2$.

\exercice{nom={L'image d'un convexe}}
Soit $f$ une fonction continue sur un convexe $P$ de $\R^2$. Montrer que $f(P)$
est un intervalle.

\magsubsection{Application partielle}
\magsubsection{Extension aux fonctions à valeurs dans $\R^2$}





\magsection{Dérivation}
\magsubsection{Dérivée partielle}

\exercice{nom={Dérivées partielles et continuité}}
Soit $f$ la fonction définie sur $\R^2$ par
\[\forall (x,y)\in\R^2 \qsep f(x,y)\defeq
  \begin{cases}
  \frac{x^2y}{x^4+y^2} & \text{si $(x,y)\not=(0,0)$,}\\
  0                    & \text{si $(x,y)=(0,0)$.}  
  \end{cases}\]
Montrer que $f$ admet une dérivée selon tout vecteur en $(0,0)$ mais n'est pas
continue en $0$.

\exercice{nom={Régularité d'une fonction}}
Soit $f$ la fonction définie sur $\R^2$ par
\[\forall\p{x,y}\in\R^2 \qsep
  f\p{x,y}\defeq
  \begin{cases}
  \frac{\sin\p{x^2}+\sin\p{y^2}}{\sqrt{x^2+y^2}} &
    \text{si $\p{x,y}\neq\p{0,0}$,}\\
  0 &  \text{si $\p{x,y}=\p{0,0}$.}
  \end{cases}\]
\begin{questions}
\question $f$ admet-elle des dérivées partielles au point $\p{0,0}$~?
\question $f$ est-elle continue en $\p{0,0}$~?
\question $f$ est-elle de classe $\classec{1}$~?   
\end{questions}

\magsubsection{Développement limité, gradient}

\exercice{nom={Calcul}}
Soit $f:\p{x,y}\mapsto xy^2$.
\begin{questions}
\question Déterminer la dérivée de $f$ au point $a\defeq \p{-1,3}$ selon le vecteur
  $h\defeq \p{2,-3}$.
\question Écrire le développement limité de $f$ à l'ordre 1 en $a$.
\end{questions}
% \begin{solution}
% $\quad$
% \begin{questions}
% \question On trouve 36.
% \question C'est $-9+9h_1-6h_2$.
% \end{questions}
% \end{solution}
\magsubsection{Dérivation des fonctions composées}
\exercice{nom={Calcul de dérivées}}
Soit $f$ une fonction de classe $\classec{1}$ sur $\R^2$. Calculer les dérivées
ou dérivées partielles des fonctions suivantes
\[g_1\p{x,y}\defeq f(y,x), \qquad g_2(x)\defeq f(x,x),\]
\[g_3\p{x,y}\defeq f(y,f(x,x)), \qquad g_4(x)\defeq f(x,f(x,x)).\]

\exercice{nom={Équation aux dérivées partielles}}
On pose $\mathcal{U}\defeq\R^2\setminus\enstq{(x,y)\in\R^2}{y=0 \et x\leq 0}$. On souhaite
déterminer l'ensemble des fonctions $f:\mathcal{U}\mapsto\R$ de classe $\classec{1}$ telles que
\[(E) \qquad \forall (x,y)\in\mathcal{U} \qsep
  x\frac{\partial f}{\partial x}(x,y)+y\frac{\partial f}{\partial y}(x,y)=
  \sqrt{x^2+y^2}.\]
\begin{questions}
\question On pose $\mathcal{V}\defeq \RPs\times\intero{-\pi}{\pi}$ et on définit
  \[\dspappli{\phi}{\mathcal{V}}{\mathcal{U}}{(r,\theta)}{(r\cos \theta,r \sin \theta)}\]
  Montrer que $\phi$ est une bijection.
\question Soit $f:\mathcal{U}\to\R$ une fonction de classe $\classec{1}$. On
  définit la fonction $g:\mathcal{V}\to\R$ par
  \[\forall (r,\theta)\in\mathcal{V}\qsep g(r,\theta)\defeq f(r\cos \theta,r\sin\theta).\]
  Montrer que $g$ est de classe $\classec{1}$ et déterminer les dérivées partielles de $g$ en fonction de celles de $f$.
\question En déduire l'ensemble des solutions de $(E)$.
\end{questions}

\magsubsection{Extrémum d'une fonction de deux variables}

% \exercice{nom={Équation aux dérivées partielles}}
% Soit $h$ une fonction de la variable $x$.
% \begin{questions}
% \question Montrer que, pour qu'il existe une fonction
%   $f\in\classec{2}\p{\R^2,\R}$ telle que~:
%   \[\parfrac{f}{x}\p{x,y}=-x\p{y^2+1}h(x) \qquad
%     \parfrac{f}{y}\p{x,y}=2x^2yh(x) \quad \p{1}\]
%   il faut que $h$ vérifie~:
%   \[xh'(x)+3h(x)=0 \quad \p{2}\]
% \question Déterminer les solutions de $\p{2}$, puis trouver $f$ de classe
%   $\classec{2}$ satisfaisant $\p{1}$.
% \end{questions}

% \exercice{nom={Laplacien}}
% Soit $u$ une fonction réelle des variables réelles $x$ et $y$ définie par
% $u\p{x,y}=\p{F\circ r}\p{x,y}$ où $r\p{x,y}=\sqrt{x^2+y^2}$ et $F$ est une
% fonction réelle d'une variable réelle. On pose~:
% \[\Delta u=\frac{\partial^2 u}{\partial x^2}
%           +\frac{\partial^2 u}{\partial y^2}\]
% \begin{questions}
% \question Calculer~:
%   \[\parfrac{r}{x} \qquad \parfrac{r}{y} \qquad
%     \frac{\partial^2 r}{\partial x^2} \qquad \frac{\partial^2 r}{\partial y^2}\]
% \question Prouver que~:
%   \[\Delta u=F''(r)+\frac{F'(r)}{r}\]
% \question En déduire $\Delta u$ lorsque $u\p{x,y}=\ln\p{x^2+y^2}$.
% \end{questions}

% \exercice{nom={Équations aux dérivées partielles}}
% Résoudre les équations aux dérivées partielles suivantes~:
% \[x^2\frac{\partial^2 f}{\partial x\partial y}=1 \qquad
%   \frac{\partial^2 f}{\partial x\partial y}=\frac{x}{y}+a\]
% \[\frac{\partial^2 f}{\partial x^2}=xy \qquad
%   \frac{\partial f}{\partial x}=af\]


% \begin{solution}
% On trouve~:
% \[\forall (x,y)\in\mathcal{U} \quad f(x,y)=
%   \sqrt{x^2+y^2}+\phi\p{2\arctan\frac{y}{\sqrt{x^2+y^2}+x}}\]
% où $\phi$ est une fonction de classe $\classec{1}$ de $\intero{-\pi}{\pi}$
% dans $\R$.    
% \end{solution}

%END_BOOK

\end{document}

\chapter{Familles sommables}
\setcounter{numeroexercicecours}{1}
\documentclass{magnolia}

\magtex{tex_driver={pdftex},
        tex_packages={epigraph,xypic}}
\magfiche{document_nom={Cours sur les séries},
          auteur_nom={François Fayard},
          auteur_mail={fayard.prof@gmail.com}}
\magcours{cours_matiere={maths},
          cours_niveau={mpsi},
          cours_chapitre_numero={21},
          cours_chapitre={Famille sommable}}
\magmisenpage{misenpage_presentation={tikzvelvia},
          misenpage_format={a4},
          misenpage_nbcolonnes={1},
          misenpage_preuve={non},
          misenpage_sol={non}}
\magmisenpage{}
\maglieudiff{}
\magprocess

\begin{document}

%BEGIN_BOOK
\setlength\epigraphwidth{.6\textwidth}
\epigraph{\og Divergent series are the invention of the devil, and it is shameful to base on them any demonstration whatsoever. \fg}{--- {\sc Niels Abel (1802--1829)}}

\setlength\epigraphwidth{.4\textwidth}
\epigraph{$\dsp \sum_{k=1}^{+\infty} k=-\frac{1}{12}$}{--- {\sc Srinivasa Ramanujan (1887--1920)}}

\magtoc
\vspace{2ex}
Dans ce chapitre, $\K$ désignera l'un des corps $\R$ ou $\C$.

\section{Famille sommable}


Les séries nous ont permis de donner un sens, lorsque c'est possible, à
\[\sum_{n=0}^{+\infty} u_n.\]
Cependant, de nombreux problèmes arrivent lorsque l'on souhaite sommer des familles $(u_{i,j})_{(i,j)\in\N^2}$
indexées par deux entiers. Il serait naturel de définir une telle somme, lorsque les séries en jeu
sont convergentes, par
\[\sum_{i=0}^{+\infty} \sum_{j=0}^{+\infty} u_{i,j}.\]
Mais on trouve rapidement des exemples pour lesquels
\[\sum_{i=0}^{+\infty} \sum_{j=0}^{+\infty} u_{i,j} \quad\et\quad
  \sum_{j=0}^{+\infty} \sum_{i=0}^{+\infty} u_{i,j}\]
ont toutes deux un sens et des valeurs différentes. Par exemple, si on définit la famille $(u_{i,j})_{(i,j)\in\N^2}$
par
\[\forall i,j\in\N\qsep u_{i,j}\defeq\begin{cases}
  1 & \text{si $i=j$,}\\
  -1 & \text{si $i=j+1$,}\\
  0 & \text{sinon,}\end{cases}\]
les séries doubles définies plus haut ont toutes deux un sens, mais la première est égale à 1 tandis que la seconde
vaut 0.\\

Contrairement à ce qui se passe dans le cas des sommes finies, il arrive donc que la \og somme \fg des éléments
d'une famille infinie dépende de l'ordre de sommation. L'objet de la théorie des \emph{familles sommables} est d'avoir
un cadre dans lequel la somme de ces familles ne dépend pas de cet ordre.


\subsection{Famille sommable de réels positifs}

\begin{definition}
On pose $[0,+\infty]\defeq [0,+\infty[\cup\ens{+\infty}$.
\begin{itemize}
\item On étend la définition de $+$ sur $[0,+\infty]$ en posant
  \begin{eqnarray*}
  \forall x\in[0,+\infty],& &x + (+\infty) \defeq +\infty\\ 
                          & &(+\infty) + x \defeq +\infty
  \end{eqnarray*}
  On vérifie que $+$ est associative et commutative sur $[0,+\infty]$ et que 0 est élément
  neutre. 
\item On étend la définition de $\times$ en posant
  \begin{eqnarray*}
  \forall x\in]0,+\infty],& &x \times (+\infty) \defeq +\infty\\ 
                          & &(+\infty) \times x \defeq +\infty
  \end{eqnarray*}
  On pose enfin $0\times(+\infty)\defeq 0$ et $(+\infty)\times 0\defeq 0$. On vérifie que $\times$
  est associative et commutative sur $[0,+\infty]$ et que 1 est élément neutre. Enfin,
  $\times$ est distributive par rapport à $+$ sur $[0,+\infty]$.
\item On étend la définition de $\leq$ sur $[0,+\infty]$ en posant
  \[\forall x\in[0,+\infty]\qsep x\leq +\infty.\]
  Muni de $\leq$, $[0,+\infty]$ est un ensemble totalement ordonné.
\end{itemize}
\end{definition}

\begin{remarques}
\remarque Excepté $+\infty$, tous les éléments de $[0,+\infty]$ sont réguliers pour $+$.
  Excepté 0 et $+\infty$, tous les éléments de $[0,+\infty]$ sont réguliers pour $\times$.
\remarque La relation d'ordre $\leq$ reste compatible avec l'addition et la multiplication~:
  il est toujours possible d'ajouter et de multiplier entre elles des inégalités puisque
  ces dernières sont positives.
\end{remarques}

\begin{definition}
Soit $A$ une partie de $[0,+\infty]$. On dit que $A$ admet une \emph{borne supérieure dans
$[0,+\infty]$} lorsque l'ensemble des majorants de $A$ dans $[0,+\infty]$ admet un plus
petit élément. Si tel est le cas, on le note $\supb A$.
\end{definition}

\begin{proposition}
Toute partie de $[0,+\infty]$ admet une borne supérieure dans $[0,+\infty]$.
\end{proposition}

\begin{remarques}
\remarque Soit $A$ une partie de $[0,+\infty]$. Si $+\infty\in A$, alors $\supb A=+\infty$. Sinon,
  $A$ est une partie de $[0,+\infty[$ et
  \begin{itemize}
  \item Si $A$ est vide, alors $\supb A=0$.
  \item Si $A$ n'est pas majorée, alors $\supb A=+\infty$. 
  \item Sinon, $A$ est non vide majorée. Elle admet donc une borne supérieure $\sup A$
    dans $\R$ et $\supb A=\sup A$.
  \end{itemize}
\remarque Soit $A$ une partie de $[0,+\infty]$.
  \begin{itemize}
  \item Si $B$ est une partie de $A$, alors $\supb B\leq\supb A$.
  \item Si $x\in [0,+\infty]$, on définit $x+A$ par
    \[x+A\defeq \ensim{x+a}{a\in A}.\]
    Alors, si $A$ est non vide, $\supb(x+A)=x+\supb A$.
  \item Si $\lambda\in[0,+\infty]$, on définit $\lambda A$ par
    \[\lambda A\defeq \ensim{\lambda a}{a\in A}.\]
    Alors, $\supb(\lambda A)=\lambda\supb A$.
  \end{itemize}
% L'ensemble vide admet une borne supérieure dans $[0,+\infty]$ qui est 0. Si $A$
%   est une partie de $[0,+\infty[$ qui n'est pas majorée dans $[0,+\infty[$, alors elle
%   admet une borne supérieure dans $[0,+\infty]$ qui est $+\infty$.
% \remarque Si $A$ est une partie de $[0,+\infty[$ non vide majorée, elle admet une borne
%   supérieure $\alpha\defeq\sup A\in\R$ au sens usuel. Alors $\alpha$ est la borne supérieure de $A$
%   dans $[0,+\infty]$. 
\end{remarques}

\begin{definition}
Soit $(u_i)_{i\in I}$ une famille d'éléments de $[0,+\infty]$. On appelle \emph{somme des $u_i$}
pour $i\in I$ et on note $\sum_{i\in I} u_i$ la borne supérieure de
\[A\defeq\ensim{\sum_{i\in K} u_i}{\text{$K$ est une partie finie de $I$}}\]
dans $[0,+\infty]$.
\end{definition}

\begin{remarqueUnique}
\remarque Si $I$ est fini, la somme que l'on vient de définir n'est autre que la somme
  $\sum_{i\in I} u_i$ définie de manière classique.
% \remarque Soit $(u_i)_{i\in I}$ une famille de réels positifs.
%   \begin{itemize}
%   \item Alors $\sum_{i\in I} u_i=+\infty$ si et seulement si, quel que soit $m\in\RP$, il existe une partie
%     finie $J$ de $I$ tel que
%     \[\sum_{i\in J} u_i\geq m.\]
%   \item De plus, si $l\in\RP$, alors $\sum_{i\in I} u_i=l$ si et seulement si, quel que
%   soit la partie finie $J$ de $I$, on a
%   \[\sum_{i\in J} u_i\leq l\]
%   et quel que soit $\epsilon>0$, il existe une partie finie $J$ de $I$ telle que
%   \[\sum_{i\in J} u_i\geq l-\epsilon.\] 
%   \end{itemize}
\end{remarqueUnique}

\begin{exoUnique}
\exo Montrer que
  \[\sum_{(i,j)\in\Ns^2} \frac{1}{(i+j)^2}=+\infty.\]
\end{exoUnique}

\begin{definition}
On dit qu'une famille $(u_i)_{i\in I}$ d'éléments de $[0,+\infty]$ est \emph{sommable} lorsque
\[\sum_{i\in I} u_i < +\infty.\]
\end{definition}

\begin{remarqueUnique}
\remarque Si l'un des $x_i$ est égal à $+\infty$, alors
  \[\sum_{i\in I} u_i=+\infty.\]
  En particulier, tous les éléments d'une famille sommable sont réels.
\end{remarqueUnique}


\begin{proposition}
Soit $(u_n)_{n\in\N}$ une suite réelle positive. Alors, la famille $(u_n)_{n\in\N}$ est
sommable si et seulement si la série $\sum u_n$ est convergente. De plus, si tel est le
cas
\[\sum_{n\in\N} u_n=\sum_{n=0}^{+\infty} u_n.\]
\end{proposition}

\begin{remarqueUnique}
\remarque Si la série à termes positifs $\sum u_n$ diverge, alors
  \[\sum_{n\in\N} u_n=+\infty.\]
  C'est pourquoi, certains auteurs se permettent d'écrire
  $\sum_{n=0}^{+\infty} u_n=+\infty$.
\end{remarqueUnique}

\begin{proposition}
Soit $(u_i)_{i\in I}$ une famille d'éléments de $[0,+\infty]$ et $J$ une partie de
  $I$. Alors
  \[\sum_{i\in J} u_i \leq \sum_{i\in I} u_i.\]  
\end{proposition}

\begin{remarqueUnique}
\remarque En particulier, si $(u_i)_{i\in I}$ est sommable, alors $(u_i)_{i\in J}$ est sommable.
\end{remarqueUnique}



\begin{proposition}
Soit $(u_i)_{i\in I}$ et $(v_i)_{i\in I}$ deux familles d'éléments de $[0,+\infty]$ et $\lambda,\mu\in[0,+\infty]$.
Alors
\[\sum_{i\in I} \p{\lambda u_i+\mu v_i}=\lambda \sum_{i\in I} u_i + \mu \sum_{i\in I} v_i.\]
\end{proposition}

  

\begin{proposition}
Soit $(u_i)_{i\in I}$ et $(v_i)_{i\in I}$ deux familles d'éléments de $[0,+\infty]$
telles que
\[\forall i\in I\qsep u_i\leq v_i.\]
Alors
\[\sum_{i\in I} u_i \leq \sum_{i\in I} v_i.\]
\end{proposition}

\begin{remarqueUnique}
\remarque En particulier, si $(v_i)_{i\in I}$ est sommable, alors $(u_i)_{i\in I}$ est sommable.
\end{remarqueUnique}


\begin{proposition}
Soit $(u_i)_{i\in I}$ une famille d'éléments de $[0,+\infty]$ et $\sigma:J\to I$ une bijection.
Alors
\[\sum_{i\in I} u_i=\sum_{j\in J} u_{\sigma(j)}.\]
\end{proposition}


\begin{proposition}[nom={Théorème de sommation par paquets}]
Soit $(u_i)_{i\in I}$ une famille d'éléments de $[0,+\infty]$. Si $(I_j)_{j\in J}$ est
une partition de $I$, alors
\[\sum_{j\in J} \p{\sum_{i\in I_j} u_i}=\sum_{i\in I} u_i.\]
\end{proposition}

\begin{remarqueUnique}
\remarque En particulier, si $(u_i)_{i\in I}$ une famille d'éléments de $[0,+\infty]$ et
    $I_1,I_2\in\mathcal{P}(I)$ sont tels que $I=I_1 \sqcup I_2$, alors
    \[\sum_{i\in I} u_i = \sum_{i\in I_1} u_i + \sum_{i\in I_2} u_i.\]
\end{remarqueUnique}

\begin{exoUnique}
\exo Pour quelles valeurs de $\alpha\in\R$ la famille
  \[\p{\frac{pq}{(p+q)^{\alpha}}}_{(p,q)\in\Ns^2}\]
  est-elle sommable~?
\begin{sol}
Si et seulement si $\alpha>4$. 
\end{sol}
\end{exoUnique}

\begin{proposition}[nom={Théorème de \nom{Fubini}}]
Soit $(u_{i,j})_{(i,j)\in I\times J}$ une famille d'éléments de $[0,+\infty]$. Alors
\[\sum_{j\in J} \p{\sum_{i\in I} u_{i,j}}=
  \sum_{i\in I} \p{\sum_{j\in J} u_{i,j}}=
  \sum_{(i,j)\in I\times J} u_{i,j}.\]
\end{proposition}

% \begin{remarqueUnique}
%   \remarque Soit $(u_{i,j})_{(i,j)\in\N^2}$ une famille d'éléments de $[0,+\infty]$. Alors
%   \[
%   \sum_{(i,j)\in\N^2} u_{i,j} =
%     \sum_{i\in\N} \sum_{j\in\N} u_{i,j} =
%     \sum_{j\in\N} \sum_{i\in\N} u_{i,j}\]
%   % \remarque Nous verrons qu'il existe des familles $(u_{i,j})_{(i,j)\in\N^2}$ telles que les sommes
%   %   \[\sum_{i=0}^{+\infty} \sum_{j=0}^{+\infty} u_{i,j} \quad\et\quad 
%   %     \sum_{j=0}^{+\infty} \sum_{i=0}^{+\infty} u_{i,j}\]
%   %   soient bien définies et distinctes. Bien entendu, de telles familles ne sont pas sommables.
%   \end{remarqueUnique}

\begin{proposition}
Soit $(u_i)_{i\in I}$ et $(v_j)_{j\in J}$ deux familles d'éléments de $[0,+\infty]$.
Alors
\[\sum_{(i,j)\in I\times J} u_i v_j = \p{\sum_{i\in I} u_i} \p{\sum_{j\in J} v_j}.\]
\end{proposition}


\begin{exos}
\exo Soit $a,b\in\RPs$. Montrer que la famille
  \[\p{\e^{-(ap+bq)}}_{(p,q)\in\N^2}\]
  est sommable et calculer sa somme.
\exo On considère la fonction $\zeta$ de \nom{Riemann} définie par
  \[\zeta(x)\defeq\sum_{n\in\Ns} \frac{1}{n^x}.\]
  pour tout $x\in\R$ pour lequel la somme ci-dessus est finie.
  \begin{questions}
  \question Montrer que le domaine de définition de $\zeta$ est $]1,+\infty[$.
  \question Montrer que
    \[\sum_{n\geq 2} (\zeta(n)-1)=1.\] 
  \end{questions}
\end{exos}

% \subsection{Familles sommables de réels positifs}




% \begin{proposition}
% Soit $(u_i)_{i\in I}$ et $(v_i)_{i\in I}$ deux familles d'éléments de $\RP$ telles que
% \[\forall i\in I\qsep 0\leq u_i\leq v_i.\]
% Si $(v_i)_{i\in I}$ est sommable, alors $(u_i)_{i\in I}$ est sommable et
% \[\sum_{i\in I} u_i\leq\sum_{i\in I} v_i.\] 
% \end{proposition}


% \begin{proposition}
% Soit $(u_i)_{i\in I}$ et $(v_i)_{i\in I}$ deux familles sommables d'éléments de $\RP$ et
% $\lambda,\mu\in\RP$. Alors $(\lambda u_i+\mu v_i)_{i\in I}$ est sommable et
% \[\sum_{i\in I} \p{\lambda u_i+\mu v_i} = \lambda \sum_{i\in I} u_i + \mu \sum_{i\in I} v_i.\]
% \end{proposition}


% \begin{proposition}[nom={Théorème de \nom{Fubini}}]
% Soit $(u_{i,j})_{(i,j)\in I\times J}$ une famille d'éléments de $\RP$. Alors
% \[\sum_{(i,j)\in I\times J} u_{i,j}=
%   \sum_{j\in J} \p{\sum_{i\in I} u_{i,j}}=
%   \sum_{i\in I} \p{\sum_{j\in J} u_{i,j}}.\]
% \end{proposition}


\subsection{Famille sommable d'éléments de $\K$}

\begin{definition}
On dit qu'une famille $(u_i)_{i\in I}$ d'éléments de $\K$ est \emph{sommable} lorsque
la famille des réels positifs $(\abs{u_i})_{i\in I}$ est sommable.
\end{definition}

\begin{remarques}
\remarque L'ensemble des familles sommables indexées par $I$ est noté $\ell^1(I,\K)$
  ou $\ell^1(I)$.
  C'est un sous-espace vectoriel de $\K^I$.
\remarque Si $(u_i)_{i\in I}$ est une famille sommable et $J$ est une partie de $I$, alors $(u_i)_{i\in J}$ est sommable.
\end{remarques}

\begin{exos}
\exo Montrer que la famille
  \[\p{\frac{\sin(p+q)}{p^2 q^2}}_{(p,q)\in\Ns^2}\]
  est sommable.
\exo Montrer que la famille
  \[\p{z^{ij}}_{(i,j)\in\Ns^2}\]
  est sommable si et seulement si $\abs{z}<1$.
\end{exos}

\begin{definition}
Pour tout $x\in\R$ on définit respectivement la \emph{partie positive} $x^+$ 
et la \emph{partie négative} $x^-$ de $x$ par
\[x^+\defeq
  \begin{cases}
  x & \text{si $x\geq 0$}\\
  0 & \text{si $x<0$}
  \end{cases}
  \qquad \text{et} \qquad
  x^-\defeq
  \begin{cases}
  \abs{x} & \text{si $x\leq 0$}\\
  0 & \text{si $x>0$.}
  \end{cases}\]
\end{definition}

\begin{proposition}
Pour tout $x\in\R$
\[x=x^+ - x^-,\qquad \abs{x}=x^+ + x^-,\qquad
  0\leq x^+\leq \abs{x} \qquad\text{et}\qquad
  0\leq x^-\leq \abs{x}.\]
\end{proposition}

\begin{definition}
\begin{itemize}
\item Soit $(u_i)_{i\in I}$ une famille de réels sommable. Alors, les familles $(u_i^{+})_{i\in I}$ et
  $(u_i^{-})_{i\in I}$ sont sommables et on définit
  \[\sum_{i\in I} u_i\defeq \sum_{i\in I} u_i^{+} - \sum_{i\in I} u_i^{-}.\]
\item Soit $(u_i)_{i\in I}$ une famille de nombres complexes sommable. En décomposant $u_i = a_i + \ii b_i$ en sa
  partie réelle et sa partie imaginaire, les familles $(a_i)_{i\in I}$ et $(b_i)_{i\in I}$ sont sommables et on définit
  \[\sum_{i\in I} u_i\defeq \sum_{i\in I} a_i + \ii \sum_{i\in I} b_i.\]
\end{itemize}
\end{definition}


\begin{proposition}
Soit $(u_n)_{n\in\N}$ une suite d'éléments de $\K$. Alors, la famille $(u_n)_{n\in\N}$ est
sommable si et seulement si la série $\sum u_n$ est absolument convergente. De plus, si tel est le
cas
\[\sum_{n\in\N} u_n=\sum_{n=0}^{+\infty} u_n.\]
\end{proposition}
  

\begin{proposition}
Soit $(u_i)_{i\in I}$ une famille sommable d'éléments de $\K$ et $l\in\K$. Alors
\[\sum_{i\in I} u_i=l\]
si et seulement si, quel que soit $\epsilon>0$, il existe une partie finie $K$ de $I$ telle que
pour toute partie finie $L$ de $I$ telle que $K\subset L$, on a
\[\abs{\sum_{i\in L} u_i - l}\leq\epsilon.\]
\end{proposition}

\begin{remarqueUnique}
\remarque La définition \og historique \fg d'une famille sommable est la suivante~: on dit
  qu'une famille $(u_i)_{i\in I}$ d'éléments de $\K$ est sommable lorsqu'il existe $l\in\K$ tel que quel
  que soit $\epsilon>0$, il existe une partie finie $K$ de $I$ telle que pour toute partie finie $L$ de
  $I$ telle que $K\subset L$, on a
  \[\abs{\sum_{i\in L} u_i - l}\leq\epsilon.\]
  Si c'est le cas, $l$ est unique, et est appelé somme de la famille $(u_i)_{i\in I}$. La proposition
  précédente nous montre donc que si une famille est sommable pour le sens donné dans ce cours,
  alors elle est sommable pour le sens \og historique \fg. Réciproquement, on peut montrer que si une famille
  est sommable pour le sens \og historique \fg, elle est sommable pour le sens donné dans ce cours.
  Cela montre l'équivalence des deux approches.
\end{remarqueUnique}

\begin{proposition}
Soit $(u_i)_{i\in I}$ et $(v_i)_{i\in I}$ deux familles sommables d'éléments de $\K$ et
$\lambda,\mu\in\K$. Alors $(\lambda u_i+\mu v_i)_{i\in I}$ est sommable et
\[\sum_{i\in I} \p{\lambda u_i+\mu v_i} = \lambda \sum_{i\in I} u_i + \mu \sum_{i\in I} v_i.\]
\end{proposition}

\begin{remarqueUnique}
\remarque Attention, il est possible que $(\lambda u_i+\mu v_i)_{i\in I}$ soit sommable sans
  que $(u_i)_{i\in I}$ et $(v_i)_{i\in I}$ le soient. Dans ce cas, il est bien sur
  interdit d'écrire 
  \[\sum_{i\in I} \p{\lambda u_i+\mu v_i} = \lambda \sum_{i\in I} u_i + \mu \sum_{i\in I} v_i\]
  puisque l'expression à droite de l'égalité n'a aucun sens.
  % Par exemple, bien que la famille
  % \[\p{\frac{1}{n(n+1)}}_{n\in\Ns}\] soit sommable, on ne peut pas écrire
  % \[\sum_{n\in\Ns} \frac{1}{n(n+1)}=\sum_{n\in\Ns} \p{\frac{1}{n}-\frac{1}{n+1}}=
  %   \sum_{n\in\Ns} \frac{1}{n} - \sum_{n\in\Ns}\frac{1}{n+1}\]
  % puisque les familles $\p{\frac{1}{n}}_{n\in\Ns}$ et $\p{\frac{1}{n+1}}_{n\in\Ns}$ ne sont pas sommables.
\end{remarqueUnique}

\begin{proposition}
  Soit $(u_i)_{i\in I}$ et $(v_i)_{i\in I}$ deux familles réelles sommables telles que
  \[\forall i\in I\qsep u_i\leq v_i.\]
  Alors
  \[\sum_{i\in I} u_i \leq \sum_{i\in I} v_i.\]
  \end{proposition}

\begin{proposition}
  Soit $(u_i)_{i\in I}$ une famille sommable d'éléments de $\K$. Alors
  \[\abs{\sum_{i\in I} u_i}\leq\sum_{i\in I} \abs{u_i}.\]
  \end{proposition}

\begin{proposition}
Soit $(u_i)_{i\in I}$ une famille sommable d'éléments de $\K$ et $\sigma:J\to I$ une bijection.
Alors $(u_{\sigma(j)})_{j\in J}$ est sommable et
\[\sum_{i\in I} u_i=\sum_{j\in J} u_{\sigma(j)}.\]
\end{proposition}

\begin{remarques}
\remarque Soit $\sum u_n$ une série absolument convergente. Alors, quel que soit
  la bijection $\sigma:\N\to\N$, la série $\sum u_{\sigma(n)}$ est absolument convergente et
  \[\sum_{n=0}^{+\infty} u_n=\sum_{n=0}^{+\infty} u_{\sigma(n)}.\]
\remarque Cette propriété est fausse pour les séries semi-convergentes. En effet, le théorème de réarrangement
  de \nom{Riemann} montre que si $\sum u_n$ est une série réelle semi-convergente, alors
  quel que soit
  $l\in\R$, il existe une bijection $\sigma:\N\to\N$ telle que
  \[\sum_{n=0}^{+\infty} u_{\sigma(n)} = l.\]
\end{remarques}

\begin{proposition}[nom={Théorème de sommation par paquets}]
Soit $(u_i)_{i\in I}$ une famille sommable d'éléments de $\K$. Si $(I_j)_{j\in J}$ est
une partition de $I$, alors pour tout $j\in J$ la famille $(u_i)_{i\in I_j}$ est
sommable. De plus, la famille $(\sum_{i\in I_j} u_i)_{j\in J}$ est sommable et
\[\sum_{j\in J} \p{\sum_{i\in I_j} u_i}=\sum_{i\in I} u_i.\]
\end{proposition}

\begin{proposition}[nom={Théorème de \nom{Fubini}}]
Soit $(u_{i,j})_{(i,j)\in I\times J}$ une famille sommable d'éléments de $\K$. Alors,
pour tout $j\in J$ la famille $(u_{i,j})_{i\in I}$ est sommable et pour tout $i\in I$,
la famille  $(u_{i,j})_{j\in J}$ est sommable. De plus, les
familles $(\sum_{i\in I} u_{i,j})_{j\in J}$ et $(\sum_{j\in J} u_{i,j})_{i\in I}$ sont
sommables et
\[\sum_{j\in J} \p{\sum_{i\in I} u_{i,j}}=
  \sum_{i\in I} \p{\sum_{j\in J} u_{i,j}}=
  \sum_{(i,j)\in I\times J} u_{i,j}.\]
\end{proposition}


\begin{remarqueUnique}
\remarque Soit $(u_{i,j})_{(i,j)\in\N^2}$ une famille sommable d'éléments de $\K$. Alors
\begin{eqnarray*}
\sum_{(i,j)\in\N^2} u_{i,j} &=&
  \sum_{i\in\N} \sum_{j\in\N} u_{i,j} =
  \sum_{j\in\N} \sum_{i\in\N} u_{i,j}\\
  &=& \sum_{i=0}^{+\infty} \sum_{j=0}^{+\infty} u_{i,j} =
  \sum_{j=0}^{+\infty} \sum_{i=0}^{+\infty} u_{i,j}.
\end{eqnarray*}
% \remarque Nous verrons qu'il existe des familles $(u_{i,j})_{(i,j)\in\N^2}$ telles que les sommes
%   \[\sum_{i=0}^{+\infty} \sum_{j=0}^{+\infty} u_{i,j} \quad\et\quad 
%     \sum_{j=0}^{+\infty} \sum_{i=0}^{+\infty} u_{i,j}\]
%   soient bien définies et distinctes. Bien entendu, de telles familles ne sont pas sommables.
\end{remarqueUnique}

\begin{exoUnique}
  \exo Montrer que la famille
    \[\p{\frac{\e^{\frac{2\ii k\pi}{n}}}{2^n}}_{k\in\Ns, n> k}\]
    est sommable et calculer sa somme.
  % \exo On définit la famille $(u_{i,j})_{(i,j)\in\N^2}$ par
  %   \[\forall i,j\in\N\qsep u_{i,j}\defeq
  %     \begin{cases}
  %     1 & \text{si $i=j$,}\\
  %     -\frac{1}{2^{j-i}} & \text{si $i<j$,}\\
  %     0 & \text{sinon.}
  %     \end{cases}\]
  %   Calculer
  %   \[\sum_{i=0}^{+\infty} \sum_{j=0}^{+\infty} u_{i,j} \quad\et\quad
  %     \sum_{j=0}^{+\infty} \sum_{i=0}^{+\infty} u_{i,j}.\]
  \end{exoUnique}

\begin{proposition}
  Soit $(u_i)_{i\in I}$ et $(v_j)_{j\in J}$ deux familles sommables d'éléments de $\K$.
  Alors la famille $(u_i v_j)_{(i,j)\in I\times J}$ est sommable et
  \[\sum_{(i,j)\in I\times J} u_i v_j = \p{\sum_{i\in I} u_i} \p{\sum_{j\in J} v_j}.\]
  \end{proposition}



\begin{definition}
Soit $(u_n)_{n\in\N}$ et $(v_n)_{n\in\N}$ deux suites. On appelle
\emph{produit de \nom{Cauchy}} de ces suites la suite $(w_n)_{n\in\N}$
définie par
\[\forall n\in\N\qsep w_n\defeq\sum_{k=0}^n u_k v_{n-k}.\]
\end{definition}

\begin{proposition}
Soit $\sum u_n$ et $\sum v_n$ deux séries absolument convergentes. Alors, leur produit de
\nom{Cauchy} est absolument convergent et
\[\sum_{n=0}^{+\infty} \p{\sum_{k=0}^n u_k v_{n-k}}=
  \p{\sum_{n=0}^{+\infty} u_n}\p{\sum_{n=0}^{+\infty} v_n}.\]
\end{proposition}

\begin{exoUnique}
\remarque L'objet de cet exercice est de donner une définition \og moderne \fg de l'exponentielle.
  \begin{questions}
  \question Pour tout $z\in\C$, montrer que la série
    \[\sum \frac{z^n}{n!}\]
    est absolument convergente; on note $\exp(z)$ sa somme.
  \question Montrer que
  \[\forall a,b\in\C\qsep \exp(a+b)=\exp(a) \exp(b).\]
  \end{questions}


% \remarque Montrer que
%   \[\forall x\in\intero{-1}{1}\qsep \sum_{n=1}^{+\infty} n x^{n-1}=\frac{1}{(1-x)^2}.\]
\end{exoUnique}


%END_BOOK
\end{document}
\section{Exercices}
\setcounter{numeroexercice}{1}
\documentclass{magnolia}

\magtex{tex_driver={pdftex}}
\magfiche{document_nom={Exercices sur les développements limités},
          auteur_nom={François Fayard},
          auteur_mail={fayard.prof@gmail.com}}
\magexos{exos_matiere={maths},
         exos_niveau={mpsi},
         exos_chapitre_numero={22},
         exos_theme={Famille sommable}}
\magmisenpage{}
\maglieudiff{}
\magprocess

\begin{document}
%BEGIN_BOOK

\magsection{Famille sommable}
\magsubsection{Famille sommable de réels positifs}

\exercice{nom={Exercice}}
Soit $\alpha\in\R$. Calculer
\[\sum_{n\in\Ns} \sum_{k\geq n} \frac{1}{k^\alpha}, \qquad
  \sum_{m\in\Ns}\sum_{n\in\Ns} \frac{1}{(m+n)^{\alpha}}.\]

\magsubsection{Famille sommable d'éléments de $\K$}

\exercice{nom={Exercice}}
Pour tout $k\in\N$, on pose
\[L_k\defeq \prod_{i=0}^{k-1} (X-i).\]
\begin{questions}
\question Montrer que pour tout $P\in\polyC$, la famille
  \[\p{\frac{P(n)}{n!}}_{n\in\N}\]
  est sommable.
\question Calculer
  \[\sum_{n\in\N} \frac{L_k(n)}{n!}\]
  pour tout $k\in\N$.
\question En déduire
  \[\sum_{n\in\N} \frac{n^4}{n!}.\]
\end{questions}

\exercice{nom={Exercice}}
\begin{questions}
\question Prouver la sommabilité et calculer
\[\sum_{n\in\N} \sum_{k\geq n} \frac{1}{k!}, \qquad
  \sum_{n\in\Ns} \sum_{k\geq n} \frac{(-1)^k}{k^3}, \qquad
  \sum_{n\in\Ns} (-1)^n \sum_{k\geq n} \frac{1}{k^3}.\]
\question Soit $z\in\C$ tel que $\abs{z}<1$. Prouver la sommabilité et calculer
  \[\sum_{n\in\N} \frac{z^{2^n}}{1-z^{2^{n+1}}}.\]
\end{questions}

\exercice{nom={Exercice}}
À quelles conditions nécessaires et suffisantes sur $a,b,z\in\C$ les familles suivantes
sont-elles sommables~?
\[\p{\frac{z^p}{q!}}_{p,q\in\N}, \qquad
  \p{\frac{a^p b^q}{p!q!}}_{p,q\in\N}, \qquad
  \p{\frac{q^p z^p}{p!q!}}_{p,q\in\N}, \qquad
  \p{\binom{p+q}{p}z^{p+q}}_{p,q\in\N}.\]

\exercice{nom={Exercice}}
\begin{questions}
\question Soit $n\in\N$. Décomposer en éléments simples la fraction rationnelle
  \[F_n\defeq \frac{1}{X(X+1)\cdots(X+n)}.\]
\question En déduire que pour tout $z\in\C\setminus(-\N)$
  \[\sum_{n\in\N} \frac{1}{z(z+1)\cdots(z+n)}=\e\sum_{n\in\N} \frac{(-1)^n}{n!(z+n)}.\]
\end{questions}










%[$\star$]
%Quels sont les réels dont le développement décimal est périodique à partir d'un certain rang ?
%



%END_BOOK

\end{document}


\end{document}
